% Options for packages loaded elsewhere
\PassOptionsToPackage{unicode}{hyperref}
\PassOptionsToPackage{hyphens}{url}
\PassOptionsToPackage{dvipsnames,svgnames,x11names}{xcolor}
%
\documentclass[
  11pt,
]{book}
\usepackage{amsmath,amssymb}
\usepackage{iftex}
\ifPDFTeX
  \usepackage[T1]{fontenc}
  \usepackage[utf8]{inputenc}
  \usepackage{textcomp} % provide euro and other symbols
  \DeclareUnicodeCharacter{2212}{-}
  \DeclareUnicodeCharacter{2192}{\ensuremath{\rightarrow}}
  \DeclareUnicodeCharacter{2248}{\ensuremath{\approx}}
  \DeclareUnicodeCharacter{2260}{\ensuremath{\neq}}
\else % if luatex or xetex
  \usepackage{unicode-math} % this also loads fontspec
  \defaultfontfeatures{Scale=MatchLowercase}
  \defaultfontfeatures[\rmfamily]{Ligatures=TeX,Scale=1}
\fi
\usepackage{lmodern}
\ifPDFTeX\else
  % xetex/luatex font selection
\fi
% Use upquote if available, for straight quotes in verbatim environments
\IfFileExists{upquote.sty}{\usepackage{upquote}}{}
\IfFileExists{microtype.sty}{% use microtype if available
  \usepackage[]{microtype}
  \UseMicrotypeSet[protrusion]{basicmath} % disable protrusion for tt fonts
}{}
\makeatletter
\@ifundefined{KOMAClassName}{% if non-KOMA class
  \IfFileExists{parskip.sty}{%
    \usepackage{parskip}
  }{% else
    \setlength{\parindent}{0pt}
    \setlength{\parskip}{6pt plus 2pt minus 1pt}}
}{% if KOMA class
  \KOMAoptions{parskip=half}}
\makeatother
\usepackage{xcolor}
\usepackage[margin=1in]{geometry}
\usepackage{color}
\usepackage{fancyvrb}
\newcommand{\VerbBar}{|}
\newcommand{\VERB}{\Verb[commandchars=\\\{\}]}
\DefineVerbatimEnvironment{Highlighting}{Verbatim}{commandchars=\\\{\}}
% Add ',fontsize=\small' for more characters per line
\newenvironment{Shaded}{}{}
\newcommand{\AlertTok}[1]{\textcolor[rgb]{1.00,0.00,0.00}{\textbf{#1}}}
\newcommand{\AnnotationTok}[1]{\textcolor[rgb]{0.38,0.63,0.69}{\textbf{\textit{#1}}}}
\newcommand{\AttributeTok}[1]{\textcolor[rgb]{0.49,0.56,0.16}{#1}}
\newcommand{\BaseNTok}[1]{\textcolor[rgb]{0.25,0.63,0.44}{#1}}
\newcommand{\BuiltInTok}[1]{\textcolor[rgb]{0.00,0.50,0.00}{#1}}
\newcommand{\CharTok}[1]{\textcolor[rgb]{0.25,0.44,0.63}{#1}}
\newcommand{\CommentTok}[1]{\textcolor[rgb]{0.38,0.63,0.69}{\textit{#1}}}
\newcommand{\CommentVarTok}[1]{\textcolor[rgb]{0.38,0.63,0.69}{\textbf{\textit{#1}}}}
\newcommand{\ConstantTok}[1]{\textcolor[rgb]{0.53,0.00,0.00}{#1}}
\newcommand{\ControlFlowTok}[1]{\textcolor[rgb]{0.00,0.44,0.13}{\textbf{#1}}}
\newcommand{\DataTypeTok}[1]{\textcolor[rgb]{0.56,0.13,0.00}{#1}}
\newcommand{\DecValTok}[1]{\textcolor[rgb]{0.25,0.63,0.44}{#1}}
\newcommand{\DocumentationTok}[1]{\textcolor[rgb]{0.73,0.13,0.13}{\textit{#1}}}
\newcommand{\ErrorTok}[1]{\textcolor[rgb]{1.00,0.00,0.00}{\textbf{#1}}}
\newcommand{\ExtensionTok}[1]{#1}
\newcommand{\FloatTok}[1]{\textcolor[rgb]{0.25,0.63,0.44}{#1}}
\newcommand{\FunctionTok}[1]{\textcolor[rgb]{0.02,0.16,0.49}{#1}}
\newcommand{\ImportTok}[1]{\textcolor[rgb]{0.00,0.50,0.00}{\textbf{#1}}}
\newcommand{\InformationTok}[1]{\textcolor[rgb]{0.38,0.63,0.69}{\textbf{\textit{#1}}}}
\newcommand{\KeywordTok}[1]{\textcolor[rgb]{0.00,0.44,0.13}{\textbf{#1}}}
\newcommand{\NormalTok}[1]{#1}
\newcommand{\OperatorTok}[1]{\textcolor[rgb]{0.40,0.40,0.40}{#1}}
\newcommand{\OtherTok}[1]{\textcolor[rgb]{0.00,0.44,0.13}{#1}}
\newcommand{\PreprocessorTok}[1]{\textcolor[rgb]{0.74,0.48,0.00}{#1}}
\newcommand{\RegionMarkerTok}[1]{#1}
\newcommand{\SpecialCharTok}[1]{\textcolor[rgb]{0.25,0.44,0.63}{#1}}
\newcommand{\SpecialStringTok}[1]{\textcolor[rgb]{0.73,0.40,0.53}{#1}}
\newcommand{\StringTok}[1]{\textcolor[rgb]{0.25,0.44,0.63}{#1}}
\newcommand{\VariableTok}[1]{\textcolor[rgb]{0.10,0.09,0.49}{#1}}
\newcommand{\VerbatimStringTok}[1]{\textcolor[rgb]{0.25,0.44,0.63}{#1}}
\newcommand{\WarningTok}[1]{\textcolor[rgb]{0.38,0.63,0.69}{\textbf{\textit{#1}}}}
\usepackage{longtable,booktabs,array}
\usepackage{calc} % for calculating minipage widths
% Correct order of tables after \paragraph or \subparagraph
\usepackage{etoolbox}
\makeatletter
\patchcmd\longtable{\par}{\if@noskipsec\mbox{}\fi\par}{}{}
\makeatother
% Allow footnotes in longtable head/foot
\IfFileExists{footnotehyper.sty}{\usepackage{footnotehyper}}{\usepackage{footnote}}
\makesavenoteenv{longtable}
\setlength{\emergencystretch}{3em} % prevent overfull lines
\providecommand{\tightlist}{%
  \setlength{\itemsep}{0pt}\setlength{\parskip}{0pt}}
\setcounter{secnumdepth}{-\maxdimen} % remove section numbering
\usepackage{microtype}
\usepackage{bookmark}
\usepackage{fvextra}
\DefineVerbatimEnvironment{Highlighting}{Verbatim}{breaklines,breakanywhere,commandchars=\\\{\}}
\RecustomVerbatimEnvironment{Verbatim}{Verbatim}{breaklines,breakanywhere}
\setlength{\emergencystretch}{3em}
\ifLuaTeX
  \usepackage{selnolig}  % disable illegal ligatures
\fi
\usepackage{bookmark}
\IfFileExists{xurl.sty}{\usepackage{xurl}}{} % add URL line breaks if available
\urlstyle{same}
\hypersetup{
  pdftitle={PAM Reference Book},
  pdfauthor={David Joyner; source documents co-authored with Claude},
  colorlinks=true,
  linkcolor={blue},
  filecolor={Maroon},
  citecolor={Blue},
  urlcolor={blue},
  pdfcreator={LaTeX via pandoc}}

\title{PAM Reference Book}
\author{David Joyner; source documents co-authored with Claude}
\date{May 27, 2026}

\begin{document}
\frontmatter
\maketitle

{
\hypersetup{linkcolor=}
\setcounter{tocdepth}{3}
\tableofcontents
}
\mainmatter
\chapter{PAM --- Pose And Motion}\label{pam-pose-and-motion}

\textbf{Stick-figure animation library for Manim, plus a Fountain
screenplay converter --- built for the \emph{Too Nice to Die} animated
sci-fi production pipeline.}

\emph{PAM v0.9.18 · fountain2pam v0.9.13} \emph{Co-authored by David
Joyner and Claude (Anthropic).}

\begin{center}\rule{0.5\linewidth}{0.5pt}\end{center}

\section{What PAM is}\label{what-pam-is}

PAM (Pose And Motion) is a Python library that animates screenplay
characters as stick figures, built on a graph-theoretic skeleton: joints
are nodes, bones are edges, and every pose is a dict mapping joint names
to (x, y) coordinates. Animation is keyframe-based --- each action
interpolates between named poses or programmatic targets.

PAM ships with:

\begin{itemize}
\tightlist
\item
  \textbf{Three figure types} --- humanoid (with \texttt{male} /
  \texttt{female} / \texttt{child} gender presets), alien (split-torso
  skeleton, \texttt{male} / \texttt{female} variants), and dog
  (quadruped, with \texttt{facing} parameter).
\item
  \textbf{Two prop-character types} --- \texttt{DogGraph} (e.g.~Chekov
  the robot dog) and \texttt{GovernorGraph} (autonomous dodecahedral
  background character).
\item
  \textbf{A prop system} with \textasciitilde30 registered builders,
  organized across five sub-modules: furniture, carried, flora,
  environment, accessories. Includes stateful props
  (\texttt{avatar\_pod}, \texttt{elevator}, \texttt{pocket\_door},
  \texttt{tv\_monitor}) with bound animation methods.
\item
  \textbf{A keyframe pose library} (\textasciitilde50 named poses,
  \textasciitilde11 cycles) covering standing, walking, running,
  jumping, sitting, reaching, lying flat, and emotional reaction beats.
\item
  \textbf{A camera system} built on Manim's \texttt{MovingCameraScene}
  with a framing / MOVE vocabulary (\texttt{wide}, \texttt{medium},
  \texttt{close}, \texttt{insert}; \texttt{pan-up}, \texttt{pan-down},
  \texttt{tilt-up}, \texttt{tilt-down}, \texttt{dolly}).
\item
  \textbf{A Fountain+ converter} (\texttt{fountain2pam.py}) that parses
  Fountain-format screenplays plus PAM-specific annotation keys and
  emits JSON action files plus per-subscene visual prompts.
\item
  \textbf{A canonical character registry}
  (\texttt{tntd\_characters.json}) carrying full color palettes, figure
  types, gender presets, and voice metadata for the TNTD cast.
\end{itemize}

PAM's design priority is \textbf{screenplay-driven blocking} --- get the
action timing, character positions, and dialogue beats correct quickly,
so the output can be used as a reference for higher-fidelity downstream
rendering (Blender, AI video) without re-deriving the choreography each
time.

\begin{center}\rule{0.5\linewidth}{0.5pt}\end{center}

\section{The production pipeline}\label{the-production-pipeline}

\begin{verbatim}
   Fountain screenplay (.fountain)
            |
            v
   +---------------------+
   |   fountain2pam.py   |   parses Fountain+ annotations
   +---------------------+
            |
            +---> screenplay.json     (PAM action sequence)
            +---> prompts.json        (per-subscene visual prompts)
            +---> characters.txt      (synced from CHARACTER blocks)
            |
            v
   +---------------------+
   |    pam_player.py    |   Manim blocking animation
   +---------------------+
            |
            v
       blocking MP4         -- used as timing/blocking reference
            |
            v
   +---------------------+
   |   pam2blender.py    |   exports scene layout to Blender
   +---------------------+
            |
            v
   Blender scene (camera, lights, props, character positions)
            |
            v
   AI video generation (frame-by-frame, externally)
            |
            v
   Final Cut Pro X (assembly, sound, final cut)
\end{verbatim}

The blocking MP4 is the artifact most often regenerated during authoring
--- it answers ``does the scene read?'' before any expensive downstream
rendering.

\begin{center}\rule{0.5\linewidth}{0.5pt}\end{center}

\section{Installation}\label{installation}

Requirements:

\begin{itemize}
\tightlist
\item
  Python 3.10+
\item
  \href{https://www.manim.community/}{Manim Community} (tested with
  0.20.1 on the Cairo backend)
\item
  \href{https://pypi.org/project/screenplain/}{screenplain} --- Fountain
  parser used by \texttt{fountain2pam.py}
\item
  numpy, scipy (transitive via Manim)
\end{itemize}

\begin{Shaded}
\begin{Highlighting}[]
\ExtensionTok{pip}\NormalTok{ install manim screenplain}
\CommentTok{\# clone PAM into your project, then:}
\BuiltInTok{cd}\NormalTok{ pam}
\ExtensionTok{pip}\NormalTok{ install }\AttributeTok{{-}e}\NormalTok{ .}
\end{Highlighting}
\end{Shaded}

Optional:

\begin{itemize}
\tightlist
\item
  LaTeX (TeX Live or MacTeX) --- needed only if you author scenes that
  use LaTeX math labels on characters (the dodecahedron \texttt{D}-style
  labels work without it).
\item
  Blender 4.x --- required only to run the \texttt{pam2blender.py}
  exporter.
\end{itemize}

\begin{center}\rule{0.5\linewidth}{0.5pt}\end{center}

\section{Quick start}\label{quick-start}

A two-character ``hello'' scene in JSON:

\begin{Shaded}
\begin{Highlighting}[]
\OtherTok{[}
  \FunctionTok{\{}\DataTypeTok{"action"}\FunctionTok{:} \StringTok{"title"}\FunctionTok{,} \DataTypeTok{"text"}\FunctionTok{:} \StringTok{"Quick Start"}\FunctionTok{,} \DataTypeTok{"subtitle"}\FunctionTok{:} \StringTok{"PAM v0.9.13"}\FunctionTok{\}}\OtherTok{,}

  \FunctionTok{\{}\DataTypeTok{"action"}\FunctionTok{:} \StringTok{"cast"}\FunctionTok{,} \DataTypeTok{"characters"}\FunctionTok{:} \FunctionTok{\{}
    \DataTypeTok{"alice"}\FunctionTok{:} \FunctionTok{\{}\DataTypeTok{"figure\_type"}\FunctionTok{:} \StringTok{"human"}\FunctionTok{,} \DataTypeTok{"gender"}\FunctionTok{:} \StringTok{"female"}\FunctionTok{,}
              \DataTypeTok{"offset"}\FunctionTok{:} \OtherTok{[}\FloatTok{{-}2.0}\OtherTok{,} \DecValTok{0}\OtherTok{,} \DecValTok{0}\OtherTok{]}\FunctionTok{,}
              \DataTypeTok{"style"}\FunctionTok{:} \FunctionTok{\{}\DataTypeTok{"head\_label"}\FunctionTok{:} \StringTok{"A"}\FunctionTok{,} \DataTypeTok{"edge\_color"}\FunctionTok{:} \StringTok{"\#cc4477"}\FunctionTok{\}\},}
    \DataTypeTok{"bob"}\FunctionTok{:}   \FunctionTok{\{}\DataTypeTok{"figure\_type"}\FunctionTok{:} \StringTok{"human"}\FunctionTok{,} \DataTypeTok{"gender"}\FunctionTok{:} \StringTok{"male"}\FunctionTok{,}
              \DataTypeTok{"offset"}\FunctionTok{:} \OtherTok{[} \FloatTok{2.0}\OtherTok{,} \DecValTok{0}\OtherTok{,} \DecValTok{0}\OtherTok{]}\FunctionTok{,}
              \DataTypeTok{"style"}\FunctionTok{:} \FunctionTok{\{}\DataTypeTok{"head\_label"}\FunctionTok{:} \StringTok{"B"}\FunctionTok{,} \DataTypeTok{"edge\_color"}\FunctionTok{:} \StringTok{"\#4477cc"}\FunctionTok{\}\}}
  \FunctionTok{\}\}}\OtherTok{,}

  \FunctionTok{\{}\DataTypeTok{"action"}\FunctionTok{:} \StringTok{"props"}\FunctionTok{,} \DataTypeTok{"items"}\FunctionTok{:} \FunctionTok{\{}
    \DataTypeTok{"table"}\FunctionTok{:} \FunctionTok{\{}\DataTypeTok{"type"}\FunctionTok{:} \StringTok{"desk"}\FunctionTok{,} \DataTypeTok{"x"}\FunctionTok{:} \DecValTok{0}\ErrorTok{.}\DecValTok{0}\FunctionTok{,} \DataTypeTok{"label"}\FunctionTok{:} \StringTok{"T"}\FunctionTok{\}}
  \FunctionTok{\}\}}\OtherTok{,}

  \FunctionTok{\{}\DataTypeTok{"action"}\FunctionTok{:} \StringTok{"fade\_in"}\FunctionTok{,} \DataTypeTok{"who"}\FunctionTok{:} \StringTok{"alice"}\FunctionTok{\}}\OtherTok{,}
  \FunctionTok{\{}\DataTypeTok{"action"}\FunctionTok{:} \StringTok{"fade\_in"}\FunctionTok{,} \DataTypeTok{"who"}\FunctionTok{:} \StringTok{"bob"}\FunctionTok{\}}\OtherTok{,}

  \FunctionTok{\{}\DataTypeTok{"action"}\FunctionTok{:} \StringTok{"say"}\FunctionTok{,}  \DataTypeTok{"who"}\FunctionTok{:} \StringTok{"alice"}\FunctionTok{,} \DataTypeTok{"text"}\FunctionTok{:} \StringTok{"Hi, Bob!"}\FunctionTok{\}}\OtherTok{,}
  \FunctionTok{\{}\DataTypeTok{"action"}\FunctionTok{:} \StringTok{"wave"}\FunctionTok{,} \DataTypeTok{"who"}\FunctionTok{:} \StringTok{"bob"}\FunctionTok{,} \DataTypeTok{"cycles"}\FunctionTok{:} \DecValTok{2}\FunctionTok{\}}\OtherTok{,}
  \FunctionTok{\{}\DataTypeTok{"action"}\FunctionTok{:} \StringTok{"say"}\FunctionTok{,}  \DataTypeTok{"who"}\FunctionTok{:} \StringTok{"bob"}\FunctionTok{,}   \DataTypeTok{"text"}\FunctionTok{:} \StringTok{"Hi, Alice!"}\FunctionTok{\}}\OtherTok{,}

  \FunctionTok{\{}\DataTypeTok{"action"}\FunctionTok{:} \StringTok{"walk\_to"}\FunctionTok{,} \DataTypeTok{"who"}\FunctionTok{:} \StringTok{"alice"}\FunctionTok{,} \DataTypeTok{"to\_x"}\FunctionTok{:} \DecValTok{{-}0}\ErrorTok{.}\DecValTok{5}\FunctionTok{\}}\OtherTok{,}
  \FunctionTok{\{}\DataTypeTok{"action"}\FunctionTok{:} \StringTok{"walk\_to"}\FunctionTok{,} \DataTypeTok{"who"}\FunctionTok{:} \StringTok{"bob"}\FunctionTok{,}   \DataTypeTok{"to\_x"}\FunctionTok{:}  \DecValTok{0}\ErrorTok{.}\DecValTok{5}\FunctionTok{\}}\OtherTok{,}

  \FunctionTok{\{}\DataTypeTok{"action"}\FunctionTok{:} \StringTok{"say"}\FunctionTok{,} \DataTypeTok{"who"}\FunctionTok{:} \StringTok{"alice"}\FunctionTok{,} \DataTypeTok{"text"}\FunctionTok{:} \StringTok{"Nice desk."}\FunctionTok{\}}
\OtherTok{]}
\end{Highlighting}
\end{Shaded}

Render it:

\begin{Shaded}
\begin{Highlighting}[]
\ExtensionTok{manim} \AttributeTok{{-}pql}\NormalTok{ pam\_player.py PAMScene }\AttributeTok{{-}{-}json}\NormalTok{ scene\_quick\_start.json}
\end{Highlighting}
\end{Shaded}

A working sample scene with the full feature set lives in
\texttt{examples/scene\_quick\_start.json}.

\begin{center}\rule{0.5\linewidth}{0.5pt}\end{center}

\section{Package layout}\label{package-layout}

\begin{verbatim}
pam/
+-- __init__.py             public API exports
+-- pam_player.py           main Manim scene + action dispatcher
+-- characters.py           HumanGraph, AlienGraph, DogGraph,
|                            GovernorGraph
+-- poses.py                ~50 named poses + ~11 cycles +
|                            EXPRESSION_GLYPHS dict
+-- paired_poses.py         multi-character interaction templates
+-- actions.py              gesture and locomotion handlers
|                            (nod, shake_head, shrug, wave, …)
+-- builds.py               figure proportions and joint scaling
|                            (default, narrow, broad, alien, …)
+-- props.py                thin dispatcher — re-exports from
|                            sub-modules
+-- props_core.py           shared infrastructure
+-- props_furniture.py      chair, desk, pocket_door, elevator,
|                            avatar_pod, tv_monitor, desk_lamp
+-- props_carried.py        hat, briefcase, folder, phone (3 styles),
|                            backpack, laptop, bug
+-- props_flora.py          flower, floral_arrangement
+-- props_environment.py    building, sun, moon, dodecahedron,
|                            backdrop, solar_panel
+-- props_accessories.py    silver_hair builder
|                            (name_tag etc. are CAPTION-only)
+-- character_gallery.py    visual regression / documentation
+-- fountain2pam.py         Fountain -> PAM JSON converter
+-- pam2blender.py          PAM JSON -> Blender scene exporter
+-- tntd_characters.json    canonical character registry
\end{verbatim}

\texttt{tntd\_characters.json} is the \textbf{canonical} character
registry --- \texttt{characters.json} (older, present in some clones) is
superseded and known to have at least one error: Factor's
\texttt{figure\_type} is \texttt{alien}, not human.

\begin{center}\rule{0.5\linewidth}{0.5pt}\end{center}

\section{Core concepts}\label{core-concepts}

\subsection{The skeleton graph}\label{the-skeleton-graph}

Each character is built from a graph: nodes are body joints
(\texttt{head}, \texttt{neck}, \texttt{lshoulder}, \texttt{lelbow},
\texttt{lhand}, etc.) and edges are bones. Three skeleton variants ship
with PAM:

\begin{longtable}[]{@{}llll@{}}
\toprule\noalign{}
Skeleton & Joints & Edges & Use \\
\midrule\noalign{}
\endhead
\bottomrule\noalign{}
\endlastfoot
Human & 12 & 11 & \texttt{figure\_type:\ "human"} \\
Alien & 14 & 14 & \texttt{figure\_type:\ "alien"} (split torso) \\
Dog & 11 & 11 & \texttt{figure\_type:\ "dog"} (quadruped) \\
\end{longtable}

A \textbf{pose} is a Python dict mapping each joint name to an (x, y)
tuple. PAM's animation engine interpolates between poses --- actions
like \texttt{walk\_to} are implemented as keyframe-by-keyframe morphs
over a walk cycle, with the figure translated in world space at each
step.

\subsection{Character classes}\label{character-classes}

\begin{verbatim}
HumanGraph        humanoid skeleton, gender presets
AlienGraph        split-torso alien, gender variants
DogGraph          quadruped, with `facing` parameter
GovernorGraph     dodecahedral autonomous prop-character
\end{verbatim}

Each accepts a \texttt{style} dict in the cast block carrying
head\_label, edge\_color, node\_color, node\_stroke, head\_color,
head\_stroke, and highlight\_color.

\subsection{Props}\label{props}

Props are Manim \texttt{VGroup}s with three layers of metadata:

\begin{itemize}
\tightlist
\item
  \textbf{Flat attrs} (backward-compat): \texttt{.pam\_name},
  \texttt{.pam\_type}, \texttt{.pam\_x}, \texttt{.pam\_y},
  \texttt{.pam\_surface\_y}
\item
  \textbf{Scene-graph node} (\texttt{.pam\_node} dict): \texttt{name},
  \texttt{kind}, \texttt{parent}, \texttt{attach}, \texttt{attrs},
  \texttt{x}, \texttt{y}
\item
  \textbf{Attachment points} (\texttt{.pam\_attachments} dict): named
  3-D positions on the geometry where child props or characters can
  connect
\end{itemize}

See \texttt{PROP\_PROPERTIES.md} for the complete prop reference.

\subsection{Actions}\label{actions}

A PAM screenplay is a flat list of action dicts. Every action has an
\texttt{"action"} key naming the handler. Handlers fall into five
categories --- see \hyperref[categorized-action-reference]{Categorized
action reference} below.

\begin{center}\rule{0.5\linewidth}{0.5pt}\end{center}

\section{The PAM JSON screenplay
format}\label{the-pam-json-screenplay-format}

A scene file is a JSON array of action dicts. The conventional order is:

\begin{enumerate}
\def\labelenumi{\arabic{enumi}.}
\tightlist
\item
  \texttt{title} (optional) --- opening title card
\item
  \texttt{cast} --- declare all characters
\item
  \texttt{props} --- declare all initial props
\item
  \texttt{\_manifest\_captions} (optional, documentation-only) ---
  at-a-glance list of caption text and timing
\item
  Scene action sequence (\texttt{fade\_in}, \texttt{say},
  \texttt{walk\_to}, \ldots)
\end{enumerate}

Comment keys are silently skipped by \texttt{pam\_player.py}:

\begin{longtable}[]{@{}ll@{}}
\toprule\noalign{}
Key prefix & Behavior \\
\midrule\noalign{}
\endhead
\bottomrule\noalign{}
\endlastfoot
\texttt{\_comment} & Skipped at any depth \\
\texttt{\_hint} & Skipped; surfaces in debug logs \\
\texttt{\_subscene\_marker} & Used by \texttt{fountain2pam.py} to chunk
prompts.json \\
\texttt{\_shot\_meta} & Camera framing hints for the subscene \\
\end{longtable}

\textbf{Avoid} the unprefixed keys \texttt{\_manifest\_captions},
\texttt{\_note}, etc., at top level inside action dicts ---
\texttt{pam\_player.py} does not skip those in all versions and may try
to read \texttt{step{[}"action"{]}} and fail. Either wrap them in
\texttt{"\_comment"} or place them at the array boundary between
actions.

\subsection{A minimal cast block}\label{a-minimal-cast-block}

\begin{Shaded}
\begin{Highlighting}[]
\FunctionTok{\{}\DataTypeTok{"action"}\FunctionTok{:} \StringTok{"cast"}\FunctionTok{,} \DataTypeTok{"characters"}\FunctionTok{:} \FunctionTok{\{}
  \DataTypeTok{"freydoon"}\FunctionTok{:} \FunctionTok{\{}
    \DataTypeTok{"figure\_type"}\FunctionTok{:} \StringTok{"human"}\FunctionTok{,}
    \DataTypeTok{"gender"}\FunctionTok{:}      \StringTok{"male"}\FunctionTok{,}
    \DataTypeTok{"build"}\FunctionTok{:}       \StringTok{"default"}\FunctionTok{,}
    \DataTypeTok{"pose"}\FunctionTok{:}        \StringTok{"standing\_front"}\FunctionTok{,}
    \DataTypeTok{"scale"}\FunctionTok{:}       \FunctionTok{\{}\DataTypeTok{"sy"}\FunctionTok{:} \DecValTok{0}\ErrorTok{.}\DecValTok{7}\FunctionTok{,} \DataTypeTok{"sx"}\FunctionTok{:} \DecValTok{0}\ErrorTok{.}\DecValTok{7}\FunctionTok{,} \DataTypeTok{"anchor"}\FunctionTok{:} \StringTok{"lankle"}\FunctionTok{\},}
    \DataTypeTok{"offset"}\FunctionTok{:}      \OtherTok{[}\FloatTok{{-}4.5}\OtherTok{,} \DecValTok{0}\OtherTok{,} \DecValTok{0}\OtherTok{]}\FunctionTok{,}
    \DataTypeTok{"style"}\FunctionTok{:} \FunctionTok{\{}
      \DataTypeTok{"head\_label"}\FunctionTok{:}   \StringTok{"Freydoon"}\FunctionTok{,}
      \DataTypeTok{"edge\_color"}\FunctionTok{:}   \StringTok{"\#e8a87c"}\FunctionTok{,}
      \DataTypeTok{"node\_color"}\FunctionTok{:}   \StringTok{"\#7a4a2a"}\FunctionTok{,}
      \DataTypeTok{"node\_stroke"}\FunctionTok{:}  \StringTok{"\#f0c090"}\FunctionTok{,}
      \DataTypeTok{"head\_color"}\FunctionTok{:}   \StringTok{"\#5a3010"}\FunctionTok{,}
      \DataTypeTok{"head\_stroke"}\FunctionTok{:}  \StringTok{"\#f5d0a0"}\FunctionTok{,}
      \DataTypeTok{"highlight\_color"}\FunctionTok{:} \StringTok{"\#fde8c0"}
    \FunctionTok{\}}
  \FunctionTok{\}}
\FunctionTok{\}\}}
\end{Highlighting}
\end{Shaded}

Full field reference: see \texttt{CHARACTER\_PROPERTIES.md}.

\subsection{A minimal props block}\label{a-minimal-props-block}

\begin{Shaded}
\begin{Highlighting}[]
\FunctionTok{\{}\DataTypeTok{"action"}\FunctionTok{:} \StringTok{"props"}\FunctionTok{,} \DataTypeTok{"items"}\FunctionTok{:} \FunctionTok{\{}
  \DataTypeTok{"main\_desk"}\FunctionTok{:} \FunctionTok{\{}\DataTypeTok{"type"}\FunctionTok{:} \StringTok{"desk"}\FunctionTok{,}     \DataTypeTok{"x"}\FunctionTok{:}  \DecValTok{0}\ErrorTok{.}\DecValTok{0}\FunctionTok{,} \DataTypeTok{"label"}\FunctionTok{:} \StringTok{"D"}\FunctionTok{\},}
  \DataTypeTok{"exit\_door"}\FunctionTok{:} \FunctionTok{\{}\DataTypeTok{"type"}\FunctionTok{:} \StringTok{"pocket\_door"}\FunctionTok{,} \DataTypeTok{"x"}\FunctionTok{:} \FloatTok{5.5}\FunctionTok{\}}
\FunctionTok{\}\}}
\end{Highlighting}
\end{Shaded}

Full field reference: see \texttt{PROP\_PROPERTIES.md}.

\begin{center}\rule{0.5\linewidth}{0.5pt}\end{center}

\section{Categorized action
reference}\label{categorized-action-reference}

This section organizes PAM actions by \emph{what they operate on}, so
the right action is easy to find when authoring a scene. Detailed
sub-key tables for each action live in \texttt{CHARACTER\_PROPERTIES.md}
(character-target actions) and \texttt{PROP\_PROPERTIES.md} (prop-target
actions).

\subsection{Quick-lookup table}\label{quick-lookup-table}

\begin{longtable}[]{@{}llll@{}}
\toprule\noalign{}
Action & Target(s) & Since & Parallel-safe? \\
\midrule\noalign{}
\endhead
\bottomrule\noalign{}
\endlastfoot
\texttt{walk\_to} & character & 0.3 & yes (different characters) \\
\texttt{run\_to} & character & 0.3 & yes \\
\texttt{rush\_to} / \texttt{rush\_out} & character & 0.9.5 & yes \\
\texttt{squeeze\_through} & character & 0.9.6 & no \\
\texttt{trot\_to} & dog & 0.9.7 & yes \\
\texttt{turn} & character & 0.3 & yes \\
\texttt{morph} & character & 0.3 & yes \\
\texttt{nod} / \texttt{shake\_head} / \texttt{shrug} & character &
0.9.11 & no \\
\texttt{wave} & character & 0.3 & yes \\
\texttt{jump\_up} & character & 0.9.5 & yes \\
\texttt{dodge} & character & 0.9.5 & yes \\
\texttt{pat} & character & 0.9.5 & no \\
\texttt{react} / \texttt{express} & character & 0.9.5 & yes \\
\texttt{say} & character & 0.3 & no (blocks for read time) \\
\texttt{prop\_say} & prop & 0.7 & no \\
\texttt{clear\_bubble} & character or prop & 0.9.10 & yes \\
\texttt{clear\_all\_bubbles} & scene-level & 0.9.10 & yes \\
\texttt{change\_uniform} & character (torso color) & 0.9.9 & yes \\
\texttt{focus} / \texttt{focus\_reset} & scene-level & 0.9.7 & yes \\
\texttt{fade\_in} / \texttt{fade\_out} & character & 0.3 & yes \\
\texttt{flash} & character or prop & 0.9.7 & yes \\
\texttt{rotate} & character or prop & 0.9.13 & no \\
\texttt{group\_translate} & multi-target & 0.9.6 & WARNING unsafe
today \\
\texttt{reach\_for} & character + prop & 0.9.5 & no \\
\texttt{grab} & character + prop & 0.9.5 & no \\
\texttt{punch\_button} & character + prop & 0.9.5 & no \\
\texttt{place\_on} & character + prop & 0.9.5 & no \\
\texttt{move\_aside} & character + prop & 0.9.5 & no \\
\texttt{stick\_to} & character + prop & 0.9.5 & no \\
\texttt{peel\_from\_hand} & character + prop & 0.9.5 & no \\
\texttt{snap\_photo} & character + prop & 0.9.5 & no \\
\texttt{exit\_through\_doors} & character + prop & 0.9.5 & no \\
\texttt{open\_doors} / \texttt{close\_doors} & prop & 0.9.5 & yes \\
\texttt{open\_lid} / \texttt{close\_lid} & prop & 0.9.5 & yes \\
\texttt{spawn\_prop} / \texttt{remove\_prop} & prop & 0.7 & yes \\
\texttt{caption} / \texttt{persistent\_caption} & scene-level & 0.9.5 &
yes \\
\texttt{sound\_cue} & scene-level & 0.9.5 & yes \\
\texttt{on\_screen\_text} & scene-level & 0.5 & yes \\
\texttt{title} & scene-level & 0.3 & n/a (always first) \\
\texttt{wait} & scene-level & 0.3 & n/a \\
\texttt{parallel} & meta & 0.7 & n/a (it \emph{is} the meta) \\
\texttt{zone\_shift} & scene-level & 0.9.5 & yes \\
\end{longtable}

WARNING \texttt{group\_translate} is the only known parallel-unsafe
action that may plausibly appear inside a \texttt{parallel} block
(walk-and-talk). v1.0.0 will either fix the safety issue or have the
player reject it with a clear error.

\subsection{A. Character-only actions}\label{a.-character-only-actions}

Locomotion, gestures, poses, dialogue --- anything that operates on a
single character without a prop target.

\begin{itemize}
\tightlist
\item
  \textbf{Locomotion:} \texttt{walk\_to}, \texttt{run\_to},
  \texttt{rush\_to}, \texttt{rush\_out}, \texttt{squeeze\_through},
  \texttt{trot\_to} (dogs)
\item
  \textbf{Pose / direction:} \texttt{turn}, \texttt{morph},
  \texttt{set\_pose}
\item
  \textbf{Gestures:} \texttt{nod}, \texttt{shake\_head}, \texttt{shrug},
  \texttt{wave}, \texttt{jump\_up}, \texttt{dodge}, \texttt{pat}
\item
  \textbf{Reactions:} \texttt{react}, \texttt{express} (smirk /
  roll\_eyes glyphs)
\item
  \textbf{Dialogue:} \texttt{say} (with \texttt{persist},
  \texttt{duration}, \texttt{os} bubble styles)
\item
  \textbf{Visibility:} \texttt{fade\_in}, \texttt{fade\_out}
\item
  \textbf{Costume:} \texttt{change\_uniform}
\end{itemize}

\subsection{B. Prop-only actions}\label{b.-prop-only-actions}

Operate on a single prop or change a prop's state.

\begin{itemize}
\tightlist
\item
  \textbf{Spawning / despawning:} \texttt{spawn\_prop},
  \texttt{remove\_prop}
\item
  \textbf{Stateful methods:} \texttt{open\_doors},
  \texttt{close\_doors}, \texttt{open\_lid}, \texttt{close\_lid}
  (dispatched by \texttt{getattr} on the prop VGroup)
\item
  \textbf{Visual cues:} \texttt{flash} (also accepts a character)
\item
  \textbf{Dialogue:} \texttt{prop\_say}
\end{itemize}

\subsection{C. Character + prop
interactions}\label{c.-character-prop-interactions}

Actions that require both a character (\texttt{who}) and a prop
(\texttt{prop} or \texttt{target}). Most are sequential-only (not
parallel-safe) because they mutate the character's hand/arm state.

\begin{itemize}
\tightlist
\item
  \textbf{Pointing:} \texttt{reach\_for}, \texttt{punch\_button}
\item
  \textbf{Carrying:} \texttt{grab}, \texttt{place\_on},
  \texttt{move\_aside}, \texttt{stick\_to}, \texttt{peel\_from\_hand},
  \texttt{snap\_photo}
\item
  \textbf{Exit:} \texttt{exit\_through\_doors} (walks to a
  door/elevator, pauses, fades out)
\end{itemize}

\subsection{D. Scene-level actions}\label{d.-scene-level-actions}

Operate on the scene, camera, or overlay text.

\begin{itemize}
\tightlist
\item
  \textbf{Captions:} \texttt{caption}, \texttt{persistent\_caption},
  \texttt{title}, \texttt{on\_screen\_text}
\item
  \textbf{Diegetic cues:} \texttt{sound\_cue}
\item
  \textbf{Time control:} \texttt{wait}, \texttt{parallel}
\item
  \textbf{Focus:} \texttt{focus}, \texttt{focus\_reset}
\item
  \textbf{Cleanup:} \texttt{clear\_bubble}, \texttt{clear\_prop\_bubble}
  (alias), \texttt{clear\_all\_bubbles}
\item
  \textbf{Zone:} \texttt{zone\_shift} (sub-location marker for the
  prompts file)
\end{itemize}

\subsection{E. Multi-target actions}\label{e.-multi-target-actions}

\begin{itemize}
\tightlist
\item
  \textbf{\texttt{group\_translate}} --- synchronized translation of
  multiple characters and/or props. Used for walk-and-talk treadmill
  resets. Parallel-unsafe today (see warning above).
\end{itemize}

\begin{center}\rule{0.5\linewidth}{0.5pt}\end{center}

\section{Fountain+ annotations}\label{fountain-annotations}

Fountain+ extends standard Fountain with bracketed annotation lines that
drive PAM-specific behavior. All annotations live in
\texttt{{[}{[}\ …\ {]}{]}} note blocks scoped to the preceding beat or
scene heading.

The full reference lives in \textbf{\texttt{fountain2pam-manual.md}} ---
this is a summary table of the keys.

\begin{longtable}[]{@{}
  >{\raggedright\arraybackslash}p{(\columnwidth - 4\tabcolsep) * \real{0.3333}}
  >{\raggedright\arraybackslash}p{(\columnwidth - 4\tabcolsep) * \real{0.3333}}
  >{\raggedright\arraybackslash}p{(\columnwidth - 4\tabcolsep) * \real{0.3333}}@{}}
\toprule\noalign{}
\begin{minipage}[b]{\linewidth}\raggedright
Key
\end{minipage} & \begin{minipage}[b]{\linewidth}\raggedright
Purpose
\end{minipage} & \begin{minipage}[b]{\linewidth}\raggedright
Emits
\end{minipage} \\
\midrule\noalign{}
\endhead
\bottomrule\noalign{}
\endlastfoot
\texttt{CAMERA} & Framing / move / subject / transition &
\texttt{\_shot\_meta} keys on the next subscene marker \\
\texttt{MOOD} & Lighting / atmosphere & prompts.json hint \\
\texttt{KIND} & Scene type (interior / exterior / fantasy / phone-call)
& prompts.json hint \\
\texttt{SCENE\ POPULATION} & Background extras & prompts.json hint \\
\texttt{NEGATIVE} & ``Avoid this'' terms & prompts.json hint \\
\texttt{CAPTION} & On-screen text card & \texttt{caption} action \\
\texttt{SOUND} & Diegetic sound flash & \texttt{sound\_cue} action \\
\texttt{PHONE} & Intercut phone-call mode & \texttt{os\_bubble} on
\texttt{say} \\
\texttt{CHARACTER} & Per-scene character override & cast-block entry \\
\texttt{PRODUCTION\ NOTE} & Metadata only & (nothing --- for humans) \\
\end{longtable}

Example:

\begin{Shaded}
\begin{Highlighting}[]
\NormalTok{INT. POLYMERCORP {-} CEO\textquotesingle{}S OFFICE {-} DAY}

\NormalTok{[[ CAMERA: FRAMING=medium | SUBJECT=ensemble | MOVE=static | TRANSITION=cut ]]}
\NormalTok{[[ MOOD: corporate, harsh fluorescent, late afternoon ]]}
\NormalTok{[[ CAPTION: Office of Mrs Tatiana Bosch, CEO of PolymerCorp. | lower{-}third | 5s | italic ]]}

\NormalTok{BRAD}
\NormalTok{(into phone)}
\NormalTok{We need that shipment Tuesday.}

\NormalTok{[[ PHONE: Brad on landline; intercut with Chava ]]}
\end{Highlighting}
\end{Shaded}

\begin{center}\rule{0.5\linewidth}{0.5pt}\end{center}

\section{Reference documents}\label{reference-documents}

The README is intentionally a front door --- it points to deeper docs
for the details.

\begin{longtable}[]{@{}
  >{\raggedright\arraybackslash}p{(\columnwidth - 2\tabcolsep) * \real{0.5000}}
  >{\raggedright\arraybackslash}p{(\columnwidth - 2\tabcolsep) * \real{0.5000}}@{}}
\toprule\noalign{}
\begin{minipage}[b]{\linewidth}\raggedright
Document
\end{minipage} & \begin{minipage}[b]{\linewidth}\raggedright
Scope
\end{minipage} \\
\midrule\noalign{}
\endhead
\bottomrule\noalign{}
\endlastfoot
\textbf{\texttt{CHARACTER\_PROPERTIES.md}} & Complete character
reference: every field a cast entry accepts, every action that targets a
character, color authority chain, registry behavior, name resolution,
the \texttt{hidden} flag, z-order, and \textasciitilde22 documented
footguns. \\
\textbf{\texttt{PROP\_PROPERTIES.md}} & Complete prop reference:
creation-time fields, prop metadata (flat attrs, \texttt{.pam\_node},
\texttt{.pam\_attachments}), registry methods, all prop-target actions,
stateful-prop dispatch, walk-following, and prop-side gotchas. \\
\textbf{\texttt{fountain2pam-manual.md}} & Full Fountain+ reference:
every annotation key with examples, CAMERA / LIGHTING vocabulary tables,
action-interpreter pattern reference, output file structures, complete
worked screenplay example. \\
\textbf{\texttt{pam\_player-manual.md}} & (planned) Player-side
reference: action handler internals, the camera system
(\texttt{MovingCameraScene}, \texttt{\_apply\_camera}, framing / MOVE
vocabulary), parallel-block semantics, the bubble registry. \\
\textbf{\texttt{BACK\_BURNER.md}} & Items deferred to v1.0.0 or later:
API normalization (\texttt{to\_x}/\texttt{x}), the \texttt{StatefulProp}
mixin refactor, walk-and-talk attached-bubble pattern, etc. \\
\end{longtable}

\begin{center}\rule{0.5\linewidth}{0.5pt}\end{center}

\section{Extending PAM}\label{extending-pam}

\subsection{Adding a new prop type}\label{adding-a-new-prop-type}

\begin{enumerate}
\def\labelenumi{\arabic{enumi}.}
\tightlist
\item
  Implement
  \texttt{build\_\textless{}yourtype\textgreater{}(name,\ x,\ y=...,\ **kwargs)}
  in the relevant \texttt{props\_*.py} sub-module, returning a
  \texttt{VGroup} with \texttt{\_attach\_pam\_attrs} metadata.
\item
  Add an entry to \texttt{PROP\_TYPES} in \texttt{props.py} mapping the
  type name (and any aliases) to the builder.
\item
  If the prop is stateful, attach \texttt{open\_*}, \texttt{close\_*},
  or method-named handlers to the returned VGroup. See
  \texttt{build\_pocket\_door} or \texttt{build\_avatar\_pod} as
  templates.
\item
  Document in \texttt{PROP\_PROPERTIES.md} §6.
\end{enumerate}

\subsection{Adding a new character
build}\label{adding-a-new-character-build}

\begin{enumerate}
\def\labelenumi{\arabic{enumi}.}
\tightlist
\item
  Add a token -\textgreater{} build-name entry to
  \texttt{CHARACTER\_BUILD\_MAP} in \texttt{fountain2pam.py}.
\item
  If the build needs new proportions, add a row to the proportion table
  in \texttt{builds.py}.
\item
  If the build introduces a new figure type, see the \texttt{AlienGraph}
  / \texttt{DogGraph} implementations in \texttt{characters.py} as
  templates.
\end{enumerate}

\subsection{Adding a new action
handler}\label{adding-a-new-action-handler}

\begin{enumerate}
\def\labelenumi{\arabic{enumi}.}
\tightlist
\item
  Add the handler to \texttt{actions.py} (gesture/locomotion) or
  \texttt{pam\_player.py} (scene-level or prop-state actions).
\item
  Add entries to \texttt{\_BEAT\_DURATION\_MAP},
  \texttt{\_summarise\_beats}, and the drama-scoring function in
  \texttt{fountain2pam.py} if the action should be inferred from prose.
\item
  Document the action in \texttt{CHARACTER\_PROPERTIES.md} or
  \texttt{PROP\_PROPERTIES.md} (whichever is the natural home) and add a
  row to the \hyperref[quick-lookup-table]{Quick-lookup table}.
\end{enumerate}

\subsection{Adding a Fountain+ annotation
key}\label{adding-a-fountain-annotation-key}

\begin{enumerate}
\def\labelenumi{\arabic{enumi}.}
\tightlist
\item
  Add the key's parser to \texttt{fountain2pam.py} alongside the
  existing \texttt{MOOD} / \texttt{KIND} / \texttt{CAMERA} parsers.
\item
  Decide whether it emits a PAM action or contributes to
  \texttt{prompts.json} metadata only.
\item
  Document in \texttt{fountain2pam-manual.md}.
\end{enumerate}

\begin{center}\rule{0.5\linewidth}{0.5pt}\end{center}

\section{Version history}\label{version-history}

A summary of headline changes per minor version. Full per-version notes
live in \texttt{CHANGELOG.md}.

\begin{longtable}[]{@{}
  >{\raggedright\arraybackslash}p{(\columnwidth - 2\tabcolsep) * \real{0.5000}}
  >{\raggedright\arraybackslash}p{(\columnwidth - 2\tabcolsep) * \real{0.5000}}@{}}
\toprule\noalign{}
\begin{minipage}[b]{\linewidth}\raggedright
Version
\end{minipage} & \begin{minipage}[b]{\linewidth}\raggedright
Headline changes
\end{minipage} \\
\midrule\noalign{}
\endhead
\bottomrule\noalign{}
\endlastfoot
0.9.18 & \texttt{face\_builder.py} per-character geometry overrides for
hair and triangle hats. Hair: \texttt{hair\_width}, \texttt{hair\_height},
\texttt{hair\_offset\_y} tune the primary cap shape of any
\texttt{hair\_style}. Hats: \texttt{hat\_width} and \texttt{hat\_height}
now also honor the \texttt{"triangle"} and \texttt{"triangle\_inv"}
styles (previously only \texttt{"square"} read them). All overrides are
optional and default to the existing hardcoded values, so existing
characters render unchanged. Alien-ear anchoring fix: \texttt{\_build\_ears}
and \texttt{\_build\_earrings} accept a \texttt{head\_rx} argument so
ears sit at the species-correct head edge. \texttt{pam/tics.py}
fragment registry expanded with \texttt{hand\_twitch\_l} (mirror of
\texttt{hand\_twitch\_r}), \texttt{foot\_tap\_r}, and
\texttt{foot\_tap\_l}, all applicable to \texttt{human} and
\texttt{alien} figures.
\texttt{GovernorGraph.say} bubble placement is now geometry-aware:
bubble bottom sits a fixed gap above
\texttt{self.\_y\ +\ self.\_radius} regardless of prop size, replacing
the v0.9.x hardcoded \texttt{self.\_y\ +\ 0.3} that put the tail tip
inside the dodecahedron's lower half. \texttt{prop\_say} accepts a new
\texttt{y\_offset} key (\texttt{GovernorGraph} only) for per-call
vertical nudging.
Focus/focus\_reset head-dot exclusion upgrade: the v0.9.17 end-state
repair for face-attached figures is superseded by removing the head
dot from the animation target entirely (rebuilt
\texttt{VGroup(edge\_group,\ *non\_head\_dots)}), eliminating the
transient flicker visible at \texttt{rt < 0.05}.
See \hyperref[hat-geometry-reference]{§5.5} for the hair/hat default
tables and syntax examples, \hyperref[tic_profile]{§1.12} for the
tic fragment registry, \hyperref[prop_say]{§3.1.2} for the
\texttt{y\_offset} parameter, and
\hyperref[the-focus_reset-overlay-gotcha]{§8.4} for the focus
head-dot exclusion. \\
0.9.13 & \texttt{rotate} action for characters and props (uses joint
coords + offset for character pivot; native Manim getters for props);
\texttt{wave} direction-key bug fixed (now reads \texttt{direction} from
JSON); \texttt{lying\_flat\_right} / \texttt{lying\_flat\_left} poses
for horizontal speech-bubble anchoring; \texttt{closed} param on
\texttt{build\_laptop}; \texttt{paired\_poses.py} for multi-character
interaction templates. \\
0.9.11 & \texttt{nod}, \texttt{shake\_head}, \texttt{shrug} gesture
actions. \\
0.9.10 & Persistent-bubble cleanup: \texttt{clear\_bubble},
\texttt{clear\_prop\_bubble}, \texttt{clear\_all\_bubbles}. Documented
the \texttt{focus\_reset} overlay gotcha. \\
0.9.9 & \texttt{change\_uniform} action for mid-scene torso recolor.
Cast-block \texttt{uniforms} dict for named variants. \\
0.9.8 & Color authority chain: three-place merge (registry / cast block
/ Fountain+ annotation). \texttt{fountain2pam.py\ -\/-characters}
correctly applies palettes. \\
0.9.7 & \texttt{props.py} split into six sub-modules.
\texttt{build\_backpack} (Fano-plane graph), \texttt{build\_laptop}.
\texttt{\_drag\_attached\_props} for held-prop walk-following.
\texttt{focus} / \texttt{focus\_reset} scene-dimming. \\
0.9.6 & \texttt{group\_translate} for multi-target synchronized motion.
\texttt{squeeze\_through} for narrow-stance walks. \\
0.9.5 & Major action expansion: \texttt{grab}, \texttt{place\_on},
\texttt{move\_aside}, \texttt{reach\_for}, \texttt{punch\_button},
\texttt{stick\_to}, \texttt{snap\_photo}, \texttt{peel\_from\_hand},
\texttt{jump\_up}, \texttt{dodge}, \texttt{react}, \texttt{express},
\texttt{pat}, \texttt{search\_drawers}, \texttt{caption},
\texttt{persistent\_caption}, \texttt{sound\_cue}, \texttt{zone\_shift},
\texttt{rush\_to}, \texttt{rush\_out}, \texttt{exit\_through\_doors}.
Fountain+ keys: \texttt{CAPTION}, \texttt{SOUND}, \texttt{PHONE},
\texttt{PRODUCTION\ NOTE}. \texttt{flash} action for temporary
recolor. \\
0.9.4 & \texttt{pam2blender.py} exporter. Blender workflow integration.
Stateful props: \texttt{build\_pocket\_door}, \texttt{build\_elevator},
\texttt{build\_avatar\_pod} with
\texttt{open\_doors}/\texttt{close\_doors}/\texttt{open\_lid}/\texttt{close\_lid}
methods. \\
0.9.3 & Gender presets for humanoid and alien builds. Character registry
(\texttt{characters.txt} / \texttt{tntd\_characters.json}). Fountain+
\texttt{CHARACTER} annotation. \texttt{character\_gallery.py}. \\
0.9.0 & First public-facing release. Camera system on
\texttt{MovingCameraScene}. Full Fountain+ CAMERA vocabulary. Manim
-\textgreater{} Blender -\textgreater{} AI video pipeline. \\
\end{longtable}

\begin{center}\rule{0.5\linewidth}{0.5pt}\end{center}

\section{License}\label{license}

Copyright © David Joyner. See \texttt{LICENSE} for terms.

\begin{center}\rule{0.5\linewidth}{0.5pt}\end{center}

\emph{PAM is a working artifact of the }Too Nice to Die* production
pipeline. Bug reports, suggestions, and pull requests welcome on GitHub:
\href{https://github.com/wdjoyner/pam}{wdjoyner/pam}.*

\chapter{\texorpdfstring{\texttt{pam\_player.py} Reference
Manual}{pam\_player.py Reference Manual}}\label{pam_player.py-reference-manual}

\textbf{PAM v0.9.13 --- \texttt{pam\_player.py} reference}

Companion document to \texttt{CHARACTER\_PROPERTIES.md} and
\texttt{PROP\_PROPERTIES.md}. This manual covers the \emph{player} side:
how \texttt{pam\_player.py} loads a screenplay, dispatches actions,
manages the camera, tracks persistent bubbles, and what the
architectural split between \texttt{pam\_player.py} and
\texttt{actions.py} actually means in practice.

Where the property docs answer ``what does this action do?'', this
manual answers ``how does the player know what to do, and what happens
when it does it?''

\begin{center}\rule{0.5\linewidth}{0.5pt}\end{center}

\section{1. Overview}\label{overview}

\subsection{\texorpdfstring{1.1 What \texttt{pam\_player.py}
is}{1.1 What pam\_player.py is}}\label{what-pam_player.py-is}

\texttt{pam\_player.py} is a Manim \texttt{MovingCameraScene} subclass
(\texttt{PAMPlayer}) that loads a JSON screenplay from disk, iterates
its action list, and emits a rendered video. It is the bridge between
the authoring layer (Fountain -\textgreater{} JSON via
\texttt{fountain2pam.py}, or hand-edited JSON) and Manim's rendering
pipeline.

The class exposes a single public entry point --- Manim's
\texttt{construct()} method --- which:

\begin{enumerate}
\def\labelenumi{\arabic{enumi}.}
\tightlist
\item
  Reads \texttt{PAM\_SCRIPT} from the environment.
\item
  Loads the JSON action list.
\item
  Sets up the camera (if \texttt{PAM\_CAMERA\_MODE=1}).
\item
  Iterates the action list, dispatching each step.
\item
  Lets Manim's renderer produce the MP4.
\end{enumerate}

There is no command-line argument parsing of PAM's own --- Manim owns
the CLI, so PAM-specific configuration is passed through environment
variables.

\subsection{\texorpdfstring{1.2 The split: \texttt{pam\_player.py} vs
\texttt{actions.py}}{1.2 The split: pam\_player.py vs actions.py}}\label{the-split-pam_player.py-vs-actions.py}

Before v0.9.5, \texttt{\_dispatch\_one} in \texttt{pam\_player.py} was a
\textasciitilde350-line if/elif chain handling every action type. The
v0.9.5 refactor moved character-targeting action handlers into a
separate \texttt{actions.py} module with a registry-based dispatch:

\begin{Shaded}
\begin{Highlighting}[]
\ImportTok{from}\NormalTok{ pam.actions }\ImportTok{import}\NormalTok{ ACTION\_REGISTRY}

\KeywordTok{def}\NormalTok{ \_dispatch\_one(step, name, }\OperatorTok{*}\NormalTok{, collect\_anims}\OperatorTok{=}\VariableTok{False}\NormalTok{):}
\NormalTok{    act }\OperatorTok{=}\NormalTok{ step[}\StringTok{"action"}\NormalTok{]}
\NormalTok{    fig }\OperatorTok{=}\NormalTok{ \_get\_fig(name)}
    \ControlFlowTok{if}\NormalTok{ fig }\KeywordTok{is} \VariableTok{None}\NormalTok{:}
        \ControlFlowTok{return} \VariableTok{None}

    \CommentTok{\# Special cases handled inline below: fade\_in, say, trot\_to}
    \CommentTok{\# Everything else: dispatch via registry.}
\NormalTok{    handler }\OperatorTok{=}\NormalTok{ ACTION\_REGISTRY.get(act)}
    \ControlFlowTok{if}\NormalTok{ handler:}
\NormalTok{        step\_aug }\OperatorTok{=} \BuiltInTok{dict}\NormalTok{(step, \_collect\_anims}\OperatorTok{=}\NormalTok{collect\_anims) }\OperatorTok{\textbackslash{}}
                   \ControlFlowTok{if}\NormalTok{ collect\_anims }\ControlFlowTok{else}\NormalTok{ step}
        \ControlFlowTok{return}\NormalTok{ handler(fig, step\_aug, }\VariableTok{self}\NormalTok{, name,}
\NormalTok{                       props}\OperatorTok{=}\NormalTok{prop\_registry, cast}\OperatorTok{=}\NormalTok{cast)}
\end{Highlighting}
\end{Shaded}

After the refactor, \texttt{\_dispatch\_one} shrank to
\textasciitilde100 lines and only three character-targeting actions
remain inline: \texttt{fade\_in} (because it \emph{creates} the figure),
\texttt{say} (because it needs the \texttt{\_pending\_camera} mechanism,
§7.5), and \texttt{trot\_to} (because dogs are keyed in the prop
registry, not the cast).

Everything else with a character target --- \texttt{walk\_to},
\texttt{run\_to}, \texttt{morph}, \texttt{turn}, \texttt{wave},
\texttt{nod}, \texttt{grab}, \texttt{punch\_button}, etc. --- lives in
\texttt{actions.py}.

\textbf{What stays in \texttt{pam\_player.py}:}

\begin{itemize}
\tightlist
\item
  The main action loop (the iteration itself)
\item
  The camera system (\texttt{\_apply\_camera}, \texttt{\_camera\_anim},
  \texttt{\_FRAMING\_CAMERA})
\item
  The bubble registry (\texttt{\_persistent\_bubbles},
  \texttt{\_sweep\_expired\_bubbles})
\item
  Cast registration (\texttt{cast} action)
\item
  Prop registration (\texttt{props}, \texttt{spawn\_prop},
  \texttt{remove\_prop})
\item
  The \texttt{parallel} meta-dispatcher
\item
  The three inline special cases above
\item
  Scene overlays (\texttt{title}, \texttt{caption}, \texttt{sound\_cue},
  \texttt{wait})
\item
  \texttt{\_subscene\_marker} and \texttt{zone\_shift} (camera triggers)
\end{itemize}

\textbf{What lives in \texttt{actions.py}:}

\begin{itemize}
\tightlist
\item
  All character-targeting handlers that don't need player-internal state
  (\textasciitilde30+ actions)
\item
  The \texttt{ACTION\_REGISTRY} dict mapping action names to handlers
\item
  The \texttt{KNOWN\_ACTIONS} frozenset (used by
  \texttt{fountain2pam.py} for validation)
\item
  The \texttt{\_CANNOT\_PARALLEL} frozenset (used by
  \texttt{pam\_player.py}'s parallel handler)
\end{itemize}

\subsection{1.3 File responsibilities}\label{file-responsibilities}

\begin{longtable}[]{@{}
  >{\raggedright\arraybackslash}p{(\columnwidth - 2\tabcolsep) * \real{0.5000}}
  >{\raggedright\arraybackslash}p{(\columnwidth - 2\tabcolsep) * \real{0.5000}}@{}}
\toprule\noalign{}
\begin{minipage}[b]{\linewidth}\raggedright
File
\end{minipage} & \begin{minipage}[b]{\linewidth}\raggedright
Owns
\end{minipage} \\
\midrule\noalign{}
\endhead
\bottomrule\noalign{}
\endlastfoot
\texttt{pam\_player.py} & Main loop, camera, bubbles, cast/prop
registries, parallel dispatch, the three inline handlers \\
\texttt{actions.py} & Character-targeting action handlers,
\texttt{ACTION\_REGISTRY}, \texttt{\_CANNOT\_PARALLEL} \\
\texttt{characters.py} & \texttt{HumanGraph}, \texttt{AlienGraph},
\texttt{DogGraph}, \texttt{GovernorGraph} classes --- including
\texttt{say()} and \texttt{morph\_to()} methods \\
\texttt{figure.py} & Shared figure infrastructure (bubble layout, scale
anchors) --- referenced by character classes \\
\texttt{poses.py} & Named pose dicts and walk/run/jump cycles \\
\texttt{paired\_poses.py} & Multi-character interaction pose
templates \\
\texttt{props.py} + \texttt{props\_*.py} & Prop builders and registry \\
\end{longtable}

When debugging an action, the question ``where does this code live?''
follows the rule: if the action targets a character only and doesn't
need camera or bubble state, it's in \texttt{actions.py}; if it touches
the camera, the bubble registry, the cast/prop registries, or the
\texttt{parallel} block, it's in \texttt{pam\_player.py}.

\begin{center}\rule{0.5\linewidth}{0.5pt}\end{center}

\section{2. Running pam\_player}\label{running-pam_player}

\subsection{\texorpdfstring{2.1 Required:
\texttt{PAM\_SCRIPT}}{2.1 Required: PAM\_SCRIPT}}\label{required-pam_script}

The path to the JSON screenplay file. Required ---
\texttt{pam\_player.py} raises a clear error if it's unset or points to
a missing file.

\begin{Shaded}
\begin{Highlighting}[]
\VariableTok{PAM\_SCRIPT}\OperatorTok{=}\NormalTok{scene\_25\_hospital.json }\ExtensionTok{manim} \AttributeTok{{-}pql}\NormalTok{ pam\_player.py PAMPlayer}
\end{Highlighting}
\end{Shaded}

The path is resolved relative to the working directory at run time.

\subsection{2.2 Optional environment
variables}\label{optional-environment-variables}

\textbf{\texttt{PAM\_CAMERA\_MODE}} --- gates the entire camera system.

\begin{longtable}[]{@{}
  >{\raggedright\arraybackslash}p{(\columnwidth - 2\tabcolsep) * \real{0.5000}}
  >{\raggedright\arraybackslash}p{(\columnwidth - 2\tabcolsep) * \real{0.5000}}@{}}
\toprule\noalign{}
\begin{minipage}[b]{\linewidth}\raggedright
Value
\end{minipage} & \begin{minipage}[b]{\linewidth}\raggedright
Effect
\end{minipage} \\
\midrule\noalign{}
\endhead
\bottomrule\noalign{}
\endlastfoot
unset / \texttt{0} & No camera moves. \texttt{\_subscene\_marker} steps
are silently skipped. Frame stays at Manim's default. \\
\texttt{1} & Camera system active. \texttt{\_subscene\_marker} triggers
reframes per \texttt{\_shot\_meta}. Initial wide/ensemble reset fires
before the action loop. \\
\end{longtable}

\textbf{Crucial gotcha:} without \texttt{PAM\_CAMERA\_MODE=1}, every
\texttt{\_subscene\_marker} in the JSON is a no-op. If a scene has an
\texttt{insert} framing on a phone that doesn't seem to fire, this is
the first thing to check.

\textbf{\texttt{PAM\_PROMPTS}} --- path to the prompts JSON written by
\texttt{fountain2pam.py} alongside the action JSON. Only consulted when
\texttt{PAM\_CAMERA\_MODE=1}. Provides per-subscene metadata (lighting
hints, freeform notes) that complement \texttt{\_shot\_meta}. Default:
\texttt{\textless{}PAM\_SCRIPT\ stem\textgreater{}\_prompts.json}.

\textbf{\texttt{PAM\_SHOW\_CLOCK}} --- set to \texttt{1} to render a
live elapsed-time display in the upper-right corner. Format:
\texttt{M:SS.ss} (minutes, seconds, hundredths). Intended for
low-quality preview renders only; the clock is a Manim mobject and
consumes real render time.

\begin{Shaded}
\begin{Highlighting}[]
\VariableTok{PAM\_SCRIPT}\OperatorTok{=}\NormalTok{scene.json }\VariableTok{PAM\_SHOW\_CLOCK}\OperatorTok{=}\NormalTok{1 }\ExtensionTok{manim} \AttributeTok{{-}pql}\NormalTok{ pam\_player.py PAMPlayer}
\end{Highlighting}
\end{Shaded}

\subsection{\texorpdfstring{2.3 The \texttt{pam-render}
wrapper}{2.3 The pam-render wrapper}}\label{the-pam-render-wrapper}

A shell script (typically in the repo root or
\texttt{\textasciitilde{}/bin}) wraps the environment-variable plumbing:

\begin{Shaded}
\begin{Highlighting}[]
\CommentTok{\#!/bin/bash}
\CommentTok{\# pam{-}render — render a PAM screenplay}
\VariableTok{SCRIPT}\OperatorTok{=}\StringTok{"}\VariableTok{$\{1}\OperatorTok{:{-}}\NormalTok{screenplay.json}\VariableTok{\}}\StringTok{"}
\VariableTok{OUTPUT}\OperatorTok{=}\StringTok{"}\VariableTok{$\{2}\OperatorTok{:{-}}\NormalTok{PAMPlayer}\VariableTok{\}}\StringTok{"}
\VariableTok{QUALITY}\OperatorTok{=}\StringTok{"}\VariableTok{$\{3}\OperatorTok{:{-}}\NormalTok{l}\VariableTok{\}}\StringTok{"}

\VariableTok{PAM\_SCRIPT}\OperatorTok{=}\StringTok{"}\VariableTok{$SCRIPT}\StringTok{"} \ExtensionTok{manim} \AttributeTok{{-}pq}\StringTok{"}\VariableTok{$QUALITY}\StringTok{"} \DataTypeTok{\textbackslash{}}
  \AttributeTok{{-}o} \StringTok{"}\VariableTok{$\{OUTPUT\}}\StringTok{.mp4"}\NormalTok{ pam\_player.py PAMPlayer}
\end{Highlighting}
\end{Shaded}

Usage:

\begin{Shaded}
\begin{Highlighting}[]
\ExtensionTok{./pam{-}render}\NormalTok{ scene\_25.json scene\_25\_preview        }\CommentTok{\# 480p}
\ExtensionTok{./pam{-}render}\NormalTok{ scene\_25.json scene\_25\_final     h    }\CommentTok{\# 1080p}
\ExtensionTok{./pam{-}render}\NormalTok{ scene\_25.json scene\_25\_master    k    }\CommentTok{\# 4K}
\end{Highlighting}
\end{Shaded}

Extended versions of \texttt{pam-render} may accept \texttt{-\/-camera},
\texttt{-\/-show-clock}, and \texttt{-\/-prompts} flags that set the
corresponding environment variables before invoking Manim.

\subsection{2.4 Manim flags and quality
presets}\label{manim-flags-and-quality-presets}

\begin{longtable}[]{@{}llll@{}}
\toprule\noalign{}
Flag & Quality & Resolution & FPS \\
\midrule\noalign{}
\endhead
\bottomrule\noalign{}
\endlastfoot
\texttt{-ql} & low & 480p & 15 \\
\texttt{-qm} & medium & 720p & 30 \\
\texttt{-qh} & high & 1080p & 60 \\
\texttt{-qk} & 4K & 2160p & 60 \\
\texttt{-p} & preview & (opens player when done) & --- \\
\end{longtable}

Authoring iteration usually runs at \texttt{-pql} (preview + low) for
fast feedback; final renders use \texttt{-qh} or \texttt{-qk}.

\begin{center}\rule{0.5\linewidth}{0.5pt}\end{center}

\section{3. World coordinates}\label{world-coordinates}

\subsection{3.1 The reference frame}\label{the-reference-frame}

PAM characters at the default \texttt{scale=0.7} sit in a coordinate
frame where:

\begin{longtable}[]{@{}ll@{}}
\toprule\noalign{}
Landmark & Approximate y \\
\midrule\noalign{}
\endhead
\bottomrule\noalign{}
\endlastfoot
Floor & -2.6 \\
Ankle / feet & -2.0 \\
Waist & -0.5 \\
Shoulders & 0.5 \\
Head & 1.2 \\
Governor (hovering) & 1.5 \\
\end{longtable}

X coordinates run left (negative) to right (positive). The
\textbf{stage} --- the area characters should occupy --- spans roughly
\texttt{x\ in\ {[}-6,\ 6{]}}.

\subsection{3.2 Stage extents and safe
zones}\label{stage-extents-and-safe-zones}

The default Manim frame width is 14.2 units; the visible screen extends
to \texttt{x\ in\ {[}-7.1,\ 7.1{]}} and \texttt{y\ in\ {[}-4,\ 4{]}}.

Speech bubbles extend \textasciitilde2.5 units from the speaker's head.
Safe x ranges for speakers:

\begin{itemize}
\tightlist
\item
  \texttt{side:\ "right"} bubble --- speaker must be at
  \texttt{x\ \textless{}=\ 4.0}
\item
  \texttt{side:\ "left"} bubble --- speaker must be at
  \texttt{x\ \textgreater{}=\ -4.0}
\end{itemize}

Characters parked off-camera can sit at \texttt{x\ approx.\ ±6.5} (just
inside the stage edge).

\subsection{3.3 Default frame extents}\label{default-frame-extents}

The initial camera state (set both by Manim's defaults and by the
explicit initial reset in \texttt{construct()} when camera mode is
active):

\begin{longtable}[]{@{}ll@{}}
\toprule\noalign{}
Property & Value \\
\midrule\noalign{}
\endhead
\bottomrule\noalign{}
\endlastfoot
Frame width & 14.2 \\
Frame center & \texttt{(0.0,\ -0.5,\ 0.0)} \\
\end{longtable}

This corresponds to \texttt{framing:\ "wide"},
\texttt{subject:\ "ensemble"}.

\begin{center}\rule{0.5\linewidth}{0.5pt}\end{center}

\section{4. The action dispatcher}\label{the-action-dispatcher}

\subsection{4.1 The main loop}\label{the-main-loop}

The heart of \texttt{pam\_player.py} is the
\texttt{for\ step\ in\ actions:} loop in \texttt{construct()}. Each
iteration does roughly:

\begin{Shaded}
\begin{Highlighting}[]
\ControlFlowTok{for}\NormalTok{ step }\KeywordTok{in}\NormalTok{ actions:}
    \ControlFlowTok{if} \StringTok{"\_comment"} \KeywordTok{in}\NormalTok{ step }\KeywordTok{or} \StringTok{"\_hint"} \KeywordTok{in}\NormalTok{ step:}
        \ControlFlowTok{continue}

    \ControlFlowTok{if}\NormalTok{ \_persistent\_bubbles:}
\NormalTok{        \_sweep\_expired\_bubbles()}

\NormalTok{    act }\OperatorTok{=}\NormalTok{ step.get(}\StringTok{"action"}\NormalTok{)}

    \CommentTok{\# Main{-}loop handlers (cast, props, spawn\_prop, parallel,}
    \CommentTok{\# \_subscene\_marker, fade\_in, say, prop\_say, trot\_to, …)}
    \ControlFlowTok{if}\NormalTok{ act }\OperatorTok{==} \StringTok{"cast"}\NormalTok{: …}
    \ControlFlowTok{elif}\NormalTok{ act }\OperatorTok{==} \StringTok{"props"}\NormalTok{: …}
    \ControlFlowTok{elif}\NormalTok{ act }\OperatorTok{==} \StringTok{"parallel"}\NormalTok{: …}
    \ControlFlowTok{elif}\NormalTok{ ...}
    \ControlFlowTok{elif} \StringTok{"\_subscene\_marker"} \KeywordTok{in}\NormalTok{ step:}
\NormalTok{        ...}
    \ControlFlowTok{else}\NormalTok{:}
        \CommentTok{\# Resolve targets and dispatch via the registry}
        \ControlFlowTok{for}\NormalTok{ name }\KeywordTok{in}\NormalTok{ \_targets(step):}
\NormalTok{            \_dispatch\_one(step, name)}
\end{Highlighting}
\end{Shaded}

The main loop's structure mixes:

\begin{enumerate}
\def\labelenumi{\arabic{enumi}.}
\tightlist
\item
  \textbf{Annotation skip} --- \texttt{\_comment} / \texttt{\_hint} keys
  cause \texttt{continue}.
\item
  \textbf{Bubble bookkeeping} --- \texttt{\_sweep\_expired\_bubbles()}
  at the top.
\item
  \textbf{Inline special cases} --- handlers that need player-level
  state.
\item
  \textbf{Registry dispatch} --- for everything else, via
  \texttt{\_dispatch\_one}.
\end{enumerate}

\subsection{\texorpdfstring{4.2 Skipped keys (\texttt{\_comment},
\texttt{\_hint})}{4.2 Skipped keys (\_comment, \_hint)}}\label{skipped-keys-_comment-_hint}

Any step containing \texttt{\_comment} or \texttt{\_hint} is skipped
entirely. This allows scene authors to annotate the JSON:

\begin{Shaded}
\begin{Highlighting}[]
\FunctionTok{\{}\DataTypeTok{"\_comment"}\FunctionTok{:} \StringTok{"REVIEW: this beat may need a beat of silence first"}\FunctionTok{\}}\ErrorTok{,}
\FunctionTok{\{}\DataTypeTok{"action"}\FunctionTok{:} \StringTok{"say"}\FunctionTok{,} \DataTypeTok{"who"}\FunctionTok{:} \StringTok{"freydoon"}\FunctionTok{,} \DataTypeTok{"text"}\FunctionTok{:} \StringTok{"…wait."}\FunctionTok{\}}
\end{Highlighting}
\end{Shaded}

\textbf{Gotcha:} because \texttt{\_comment} is checked at the top of the
loop, \emph{any} step containing \texttt{\_comment} is skipped --- even
if it also has an \texttt{"action"} key. An action with a
\texttt{\_comment} key inside it is a no-op. If you need both, place the
comment in a separate dict \emph{before} the action.

Other underscore-prefixed top-level keys are \textbf{not} skipped:

\begin{itemize}
\tightlist
\item
  \texttt{\_subscene\_marker} is handled specially (§5.8)
\item
  \texttt{\_shot\_meta} is consumed by the subscene marker handler
\item
  \texttt{\_manifest\_captions}, \texttt{\_note}, etc. --- \textbf{not}
  skipped, and the player will try to read \texttt{step{[}"action"{]}}
  and fail. Wrap these in \texttt{\_comment} if you need documentation
  blocks.
\end{itemize}

\subsection{\texorpdfstring{4.3 Target resolution: \texttt{\_targets},
\texttt{\_get\_fig},
\texttt{\_get\_prop}}{4.3 Target resolution: \_targets, \_get\_fig, \_get\_prop}}\label{target-resolution-_targets-_get_fig-_get_prop}

Target resolution happens in three helpers:

\textbf{\texttt{\_targets(step)}} --- returns a list of character names
from a step's \texttt{who} key. Accepts:

\begin{longtable}[]{@{}ll@{}}
\toprule\noalign{}
\texttt{who} value & Returns \\
\midrule\noalign{}
\endhead
\bottomrule\noalign{}
\endlastfoot
\texttt{"freydoon"} & \texttt{{[}"freydoon"{]}} \\
\texttt{{[}"freydoon",\ "bevers"{]}} &
\texttt{{[}"freydoon",\ "bevers"{]}} \\
\texttt{"all"} & every key in \texttt{cast} \\
(missing) & \texttt{{[}{]}} \\
\end{longtable}

\textbf{\texttt{\_get\_fig(name)}} --- returns the live figure object
for a character, or \texttt{None} if the character isn't faded in. Reads
\texttt{cast{[}name{]}{[}"fig"{]}}.

\textbf{\texttt{\_get\_prop(pname)}} --- returns the live prop VGroup,
or \texttt{None}. Reads from the prop registry.

\subsection{\texorpdfstring{4.4 \texttt{\_dispatch\_one} and the
ACTION\_REGISTRY}{4.4 \_dispatch\_one and the ACTION\_REGISTRY}}\label{dispatch_one-and-the-action_registry}

For any step that isn't an inline special case, the main loop calls
\texttt{\_dispatch\_one(step,\ name)} per resolved character name.
Internally:

\begin{Shaded}
\begin{Highlighting}[]
\KeywordTok{def}\NormalTok{ \_dispatch\_one(step, name, }\OperatorTok{*}\NormalTok{, collect\_anims}\OperatorTok{=}\VariableTok{False}\NormalTok{):}
\NormalTok{    act }\OperatorTok{=}\NormalTok{ step[}\StringTok{"action"}\NormalTok{]}
\NormalTok{    fig }\OperatorTok{=}\NormalTok{ \_get\_fig(name)}
    \ControlFlowTok{if}\NormalTok{ fig }\KeywordTok{is} \VariableTok{None}\NormalTok{:}
        \ControlFlowTok{return} \VariableTok{None}

    \CommentTok{\# Inline cases (fade\_in, say, trot\_to) handled here…}

    \CommentTok{\# Otherwise: ACTION\_REGISTRY dispatch}
\NormalTok{    handler }\OperatorTok{=}\NormalTok{ ACTION\_REGISTRY.get(act)}
    \ControlFlowTok{if}\NormalTok{ handler }\KeywordTok{is} \VariableTok{None}\NormalTok{:}
        \ControlFlowTok{return} \VariableTok{None}    \CommentTok{\# silent no{-}op for unknown actions}
    \ControlFlowTok{return}\NormalTok{ handler(fig, step, }\VariableTok{self}\NormalTok{, name,}
\NormalTok{                   props}\OperatorTok{=}\NormalTok{prop\_registry, cast}\OperatorTok{=}\NormalTok{cast)}
\end{Highlighting}
\end{Shaded}

The registry is built in \texttt{actions.py} as a dict mapping action
names to handler functions:

\begin{Shaded}
\begin{Highlighting}[]
\NormalTok{ACTION\_REGISTRY: }\BuiltInTok{dict}\NormalTok{[}\BuiltInTok{str}\NormalTok{, }\BuiltInTok{callable}\NormalTok{] }\OperatorTok{=}\NormalTok{ \{}
    \StringTok{"walk\_to"}\NormalTok{:   act\_walk\_to,}
    \StringTok{"run\_to"}\NormalTok{:    act\_run\_to,}
    \StringTok{"turn"}\NormalTok{:      act\_turn,}
    \StringTok{"morph"}\NormalTok{:     act\_morph,}
    \StringTok{"wave"}\NormalTok{:      act\_wave,}
    \StringTok{"nod"}\NormalTok{:       act\_nod,}
    \StringTok{"shake\_head"}\NormalTok{: act\_shake\_head,}
    \StringTok{"shrug"}\NormalTok{:     act\_shrug,}
    \StringTok{"grab"}\NormalTok{:      act\_grab,}
    \StringTok{"punch\_button"}\NormalTok{: act\_punch\_button,}
    \StringTok{"reach\_for"}\NormalTok{: act\_reach\_for,}
    \CommentTok{\# … \textasciitilde{}30+ entries …}
\NormalTok{\}}
\end{Highlighting}
\end{Shaded}

\textbf{Unknown-action behavior.} If an action name doesn't appear in
the registry and isn't a main-loop special case,
\texttt{\_dispatch\_one} returns \texttt{None} silently. This is the
failure mode when a typo in the JSON produces nothing on screen.
Validate authored JSON against \texttt{KNOWN\_ACTIONS} (also exported by
\texttt{actions.py}) to catch these at parse time.

\subsection{4.5 Handler signature
convention}\label{handler-signature-convention}

Every handler in \texttt{actions.py} follows the same signature:

\begin{Shaded}
\begin{Highlighting}[]
\KeywordTok{def}\NormalTok{ act\_walk\_to(fig, step, scene, name}\OperatorTok{=}\StringTok{"I.G. NoreMe"}\NormalTok{,}
                \OperatorTok{*}\NormalTok{, props}\OperatorTok{=}\VariableTok{None}\NormalTok{, cast}\OperatorTok{=}\VariableTok{None}\NormalTok{):}
    \CommentTok{"""Walk a character to a target x position."""}
\NormalTok{    to\_x }\OperatorTok{=}\NormalTok{ step.get(}\StringTok{"to\_x"}\NormalTok{, step.get(}\StringTok{"x"}\NormalTok{))  }\CommentTok{\# accept either}
\NormalTok{    rt   }\OperatorTok{=}\NormalTok{ step.get(}\StringTok{"rt"}\NormalTok{, }\FloatTok{0.6}\NormalTok{)}
\NormalTok{    …}
    \ControlFlowTok{if}\NormalTok{ step.get(}\StringTok{"\_collect\_anims"}\NormalTok{):}
        \ControlFlowTok{return}\NormalTok{ fig.\_walk\_anims(to\_x, rt}\OperatorTok{=}\NormalTok{rt)}
\NormalTok{    fig.walk\_to(scene, to\_x, rt}\OperatorTok{=}\NormalTok{rt)}
    \ControlFlowTok{return} \VariableTok{None}
\end{Highlighting}
\end{Shaded}

The signature elements:

\begin{longtable}[]{@{}ll@{}}
\toprule\noalign{}
Arg & Meaning \\
\midrule\noalign{}
\endhead
\bottomrule\noalign{}
\endlastfoot
\texttt{fig} & The live figure object (already resolved) \\
\texttt{step} & The action dict \\
\texttt{scene} & The PAMPlayer instance (Manim scene) \\
\texttt{name} & Character key (for state updates) \\
\texttt{props=} & The prop registry (for actions that target props) \\
\texttt{cast=} & The cast registry (for actions touching others) \\
\end{longtable}

The \texttt{name="I.G.\ NoreMe"} default is a sentinel --- handlers that
require a real name check for it explicitly.

\subsection{\texorpdfstring{4.6 The \texttt{\_collect\_anims}
thread}{4.6 The \_collect\_anims thread}}\label{the-_collect_anims-thread}

When a sub-action runs inside a \texttt{parallel} block, the dispatcher
passes \texttt{collect\_anims=True}. This is threaded into the step dict
as \texttt{step{[}"\_collect\_anims"{]}\ =\ True}, so handlers see it
without a signature change.

A handler that supports parallel collection returns a list of Manim
animations instead of playing them itself:

\begin{Shaded}
\begin{Highlighting}[]
\ControlFlowTok{if}\NormalTok{ step.get(}\StringTok{"\_collect\_anims"}\NormalTok{):}
    \ControlFlowTok{return}\NormalTok{ fig.\_pose\_anims(target\_pose, new\_offset)}
\NormalTok{fig.morph\_to(scene, target\_pose, new\_offset}\OperatorTok{=}\NormalTok{new\_offset, rt}\OperatorTok{=}\NormalTok{rt)}
\ControlFlowTok{return} \VariableTok{None}
\end{Highlighting}
\end{Shaded}

The \texttt{parallel} handler then plays all collected animations
concurrently in a single \texttt{scene.play(*all\_anims,\ run\_time=rt)}
call.

Handlers that don't support \texttt{\_collect\_anims} (the
parallel-unsafe set) ignore the flag and just run. The parallel handler
routes them to a sequential fallback path --- see §6.

\begin{center}\rule{0.5\linewidth}{0.5pt}\end{center}

\section{5. Special main-loop
handlers}\label{special-main-loop-handlers}

These handlers live inline in \texttt{construct()}'s main loop, not in
\texttt{actions.py}, because they touch player-level state (registries,
camera, bubbles, scene clock).

\subsection{\texorpdfstring{5.1 \texttt{cast} --- character
registration}{5.1 cast --- character registration}}\label{cast-character-registration}

Registers all characters at scene start. Reads the \texttt{characters}
sub-dict and populates the \texttt{cast} dict with one entry per
character:

\begin{Shaded}
\begin{Highlighting}[]
\NormalTok{cast[cname] }\OperatorTok{=}\NormalTok{ \{}
    \StringTok{"fig"}\NormalTok{:         }\VariableTok{None}\NormalTok{,          }\CommentTok{\# populated on fade\_in}
    \StringTok{"figure\_type"}\NormalTok{: }\StringTok{"human"}\NormalTok{,       }\CommentTok{\# or "alien", "dog", "governor"}
    \StringTok{"pose"}\NormalTok{:        }\StringTok{"standing\_front"}\NormalTok{,}
    \StringTok{"offset"}\NormalTok{:      [}\OperatorTok{{-}}\FloatTok{2.0}\NormalTok{, }\DecValTok{0}\NormalTok{, }\DecValTok{0}\NormalTok{],}
    \StringTok{"style"}\NormalTok{:       \{}\StringTok{"head\_label"}\NormalTok{: }\StringTok{"F"}\NormalTok{, }\StringTok{"edge\_color"}\NormalTok{: }\StringTok{"\#e8a87c"}\NormalTok{, …\},}
    \StringTok{"build"}\NormalTok{:       }\StringTok{"default"}\NormalTok{,}
    \StringTok{"scale"}\NormalTok{:       \{}\StringTok{"sy"}\NormalTok{: }\FloatTok{0.7}\NormalTok{, }\StringTok{"sx"}\NormalTok{: }\FloatTok{0.7}\NormalTok{, }\StringTok{"anchor"}\NormalTok{: }\StringTok{"lankle"}\NormalTok{\},}
    \StringTok{"spawn"}\NormalTok{:       \{\},            }\CommentTok{\# for prop{-}characters}
\NormalTok{\}}
\end{Highlighting}
\end{Shaded}

The \texttt{fig} slot stays \texttt{None} until \texttt{fade\_in}
actually builds the figure (§5.4). This separation lets \texttt{cast}
declare characters without committing to their appearance --- they may
be revealed only later in the scene.

\subsection{\texorpdfstring{5.2 \texttt{props} --- initial prop
registration}{5.2 props --- initial prop registration}}\label{props-initial-prop-registration}

Same pattern for props:

\begin{Shaded}
\begin{Highlighting}[]
\NormalTok{\{}\StringTok{"action"}\NormalTok{: }\StringTok{"props"}\NormalTok{, }\StringTok{"items"}\NormalTok{: \{}
  \StringTok{"main\_desk"}\NormalTok{: \{}\StringTok{"type"}\NormalTok{: }\StringTok{"desk"}\NormalTok{, }\StringTok{"x"}\NormalTok{: }\FloatTok{0.0}\NormalTok{, }\StringTok{"label"}\NormalTok{: }\StringTok{"D"}\NormalTok{\},}
  \StringTok{"exit\_door"}\NormalTok{: \{}\StringTok{"type"}\NormalTok{: }\StringTok{"pocket\_door"}\NormalTok{, }\StringTok{"x"}\NormalTok{: }\FloatTok{5.5}\NormalTok{\}}
\NormalTok{\}\}}
\end{Highlighting}
\end{Shaded}

Each entry is built immediately via \texttt{build\_prop(name,\ **spec)}
and added to the prop registry. Props declared in \texttt{props} appear
at opacity 1 from scene start unless they carry
\texttt{"hidden":\ true}.

\subsection{\texorpdfstring{5.3 \texttt{spawn\_prop} /
\texttt{remove\_prop}}{5.3 spawn\_prop / remove\_prop}}\label{spawn_prop-remove_prop}

Mid-scene appearance and despawn. \texttt{spawn\_prop} builds the prop
and fades it in:

\begin{Shaded}
\begin{Highlighting}[]
\NormalTok{prop }\OperatorTok{=}\NormalTok{ build\_prop(name, }\OperatorTok{**}\NormalTok{spawn\_spec)}
\NormalTok{prop\_registry.add(name, prop)}
\NormalTok{scene.play(FadeIn(prop), run\_time}\OperatorTok{=}\NormalTok{rt)}
\end{Highlighting}
\end{Shaded}

\texttt{remove\_prop} does the inverse:

\begin{Shaded}
\begin{Highlighting}[]
\NormalTok{prop }\OperatorTok{=}\NormalTok{ prop\_registry.get(name)}
\NormalTok{scene.play(FadeOut(prop), run\_time}\OperatorTok{=}\NormalTok{rt)}
\NormalTok{prop\_registry.remove(name)}
\end{Highlighting}
\end{Shaded}

See \texttt{PROP\_PROPERTIES.md} §3 for the full registry behavior,
including the symmetric-clear gotcha when a prop has been
parent-attached.

\subsection{\texorpdfstring{5.4 \texttt{fade\_in} --- figure
construction}{5.4 fade\_in --- figure construction}}\label{fade_in-figure-construction}

\texttt{fade\_in} stays inline because it \emph{creates} the figure
object:

\begin{Shaded}
\begin{Highlighting}[]
\ControlFlowTok{if}\NormalTok{ act }\OperatorTok{==} \StringTok{"fade\_in"}\NormalTok{:}
\NormalTok{    entry }\OperatorTok{=}\NormalTok{ cast[name]}
    \ControlFlowTok{if}\NormalTok{ entry[}\StringTok{"figure\_type"}\NormalTok{] }\OperatorTok{==} \StringTok{"human"}\NormalTok{:}
\NormalTok{        fig }\OperatorTok{=}\NormalTok{ HumanGraph(}\OperatorTok{**}\NormalTok{entry[}\StringTok{"style"}\NormalTok{], build}\OperatorTok{=}\NormalTok{entry[}\StringTok{"build"}\NormalTok{])}
    \ControlFlowTok{elif}\NormalTok{ entry[}\StringTok{"figure\_type"}\NormalTok{] }\OperatorTok{==} \StringTok{"alien"}\NormalTok{:}
\NormalTok{        fig }\OperatorTok{=}\NormalTok{ AlienGraph(}\OperatorTok{**}\NormalTok{entry[}\StringTok{"style"}\NormalTok{], gender}\OperatorTok{=}\NormalTok{entry.get(}\StringTok{"gender"}\NormalTok{))}
    \ControlFlowTok{elif}\NormalTok{ entry[}\StringTok{"figure\_type"}\NormalTok{] }\OperatorTok{==} \StringTok{"dog"}\NormalTok{:}
\NormalTok{        fig }\OperatorTok{=}\NormalTok{ DogGraph(facing}\OperatorTok{=}\NormalTok{entry.get(}\StringTok{"facing"}\NormalTok{, }\StringTok{"right"}\NormalTok{))}
    \ControlFlowTok{elif}\NormalTok{ entry[}\StringTok{"figure\_type"}\NormalTok{] }\OperatorTok{==} \StringTok{"governor"}\NormalTok{:}
\NormalTok{        fig }\OperatorTok{=}\NormalTok{ GovernorGraph(}\OperatorTok{**}\NormalTok{entry[}\StringTok{"style"}\NormalTok{])}

\NormalTok{    fig.pose }\OperatorTok{=}\NormalTok{ POSES[entry[}\StringTok{"pose"}\NormalTok{]]}
\NormalTok{    fig.offset }\OperatorTok{=}\NormalTok{ np.array(entry[}\StringTok{"offset"}\NormalTok{])}
    \ControlFlowTok{if}\NormalTok{ entry.get(}\StringTok{"scale"}\NormalTok{):}
\NormalTok{        fig.apply\_scale(}\OperatorTok{**}\NormalTok{entry[}\StringTok{"scale"}\NormalTok{])}

\NormalTok{    fig.fade\_in(}\VariableTok{self}\NormalTok{, rt}\OperatorTok{=}\NormalTok{step.get(}\StringTok{"rt"}\NormalTok{, }\FloatTok{0.5}\NormalTok{))}
\NormalTok{    entry[}\StringTok{"fig"}\NormalTok{] }\OperatorTok{=}\NormalTok{ fig}
    \ControlFlowTok{continue}
\end{Highlighting}
\end{Shaded}

After \texttt{fade\_in}, the figure is in the cast registry and
discoverable by \texttt{\_get\_fig(name)}. All subsequent actions can
target it.

\subsection{\texorpdfstring{5.5 \texttt{say} / \texttt{prop\_say} ---
bubble plus concurrent
camera}{5.5 say / prop\_say --- bubble plus concurrent camera}}\label{say-prop_say-bubble-plus-concurrent-camera}

The \texttt{say} handler stays inline because it implements the
\textbf{\texttt{\_pending\_camera} mechanism} (§7.5): if a
\texttt{\_subscene\_marker} just fired with an animated camera move, the
move is held in \texttt{self.\_pending\_camera} and included in the
\texttt{say}'s \texttt{scene.play(…)} call so the camera glides during
the bubble fade-in:

\begin{Shaded}
\begin{Highlighting}[]
\ControlFlowTok{if}\NormalTok{ act }\OperatorTok{==} \StringTok{"say"}\NormalTok{:}
\NormalTok{    fig }\OperatorTok{=}\NormalTok{ \_get\_fig(name)}
\NormalTok{    text }\OperatorTok{=}\NormalTok{ step[}\StringTok{"text"}\NormalTok{]}
\NormalTok{    extra\_anims }\OperatorTok{=}\NormalTok{ []}
    \ControlFlowTok{if} \VariableTok{self}\NormalTok{.\_pending\_camera }\KeywordTok{is} \KeywordTok{not} \VariableTok{None}\NormalTok{:}
\NormalTok{        extra\_anims.append(}\VariableTok{self}\NormalTok{.\_pending\_camera)}
        \VariableTok{self}\NormalTok{.\_pending\_camera }\OperatorTok{=} \VariableTok{None}

\NormalTok{    fig.say(text, scene}\OperatorTok{=}\VariableTok{self}\NormalTok{,}
\NormalTok{            extra\_anims}\OperatorTok{=}\NormalTok{extra\_anims,}
\NormalTok{            …)}
\end{Highlighting}
\end{Shaded}

\texttt{prop\_say} follows the same pattern but builds a bubble anchored
above the prop's \texttt{pam\_x} / \texttt{pam\_y} rather than the
character's head.

\textbf{Persistent bubbles.} Either handler, when given
\texttt{persist:\ true} or \texttt{duration:\ t}, registers the bubble
in \texttt{\_persistent\_bubbles}:

\begin{Shaded}
\begin{Highlighting}[]
\ControlFlowTok{if}\NormalTok{ step.get(}\StringTok{"persist"}\NormalTok{) }\KeywordTok{or}\NormalTok{ step.get(}\StringTok{"duration"}\NormalTok{):}
\NormalTok{    expires\_at }\OperatorTok{=}\NormalTok{ (}\VariableTok{self}\NormalTok{.\_scene\_clock[}\DecValTok{0}\NormalTok{] }\OperatorTok{+}\NormalTok{ step[}\StringTok{"duration"}\NormalTok{]}
                  \ControlFlowTok{if}\NormalTok{ step.get(}\StringTok{"duration"}\NormalTok{) }\ControlFlowTok{else} \VariableTok{None}\NormalTok{)}
    \VariableTok{self}\NormalTok{.\_persistent\_bubbles[name\_or\_pname] }\OperatorTok{=}\NormalTok{ \{}
        \StringTok{"bubble"}\NormalTok{: bubble,}
        \StringTok{"expires\_at"}\NormalTok{: expires\_at,}
\NormalTok{    \}}
\end{Highlighting}
\end{Shaded}

See §8 for the full bubble lifecycle.

\subsection{\texorpdfstring{5.6 \texttt{parallel} ---
meta-dispatch}{5.6 parallel --- meta-dispatch}}\label{parallel-meta-dispatch}

The \texttt{parallel} handler is structurally a sub-loop:

\begin{Shaded}
\begin{Highlighting}[]
\ControlFlowTok{if}\NormalTok{ act }\OperatorTok{==} \StringTok{"parallel"}\NormalTok{:}
\NormalTok{    sub\_actions }\OperatorTok{=}\NormalTok{ step.get(}\StringTok{"do"}\NormalTok{, [])}
\NormalTok{    rt }\OperatorTok{=}\NormalTok{ step.get(}\StringTok{"rt"}\NormalTok{, }\FloatTok{0.4}\NormalTok{)}

    \CommentTok{\# Partition sub{-}actions into:}
    \CommentTok{\#  {-} locomotion (walk\_to, run\_to, …) — special keyframe path}
    \CommentTok{\#  {-} simple (everything parallel{-}safe) — collect anims}
\NormalTok{    locomotion }\OperatorTok{=}\NormalTok{ [s }\ControlFlowTok{for}\NormalTok{ s }\KeywordTok{in}\NormalTok{ sub\_actions }\ControlFlowTok{if}\NormalTok{ s[}\StringTok{"action"}\NormalTok{] }\KeywordTok{in}\NormalTok{ \_LOCO]}
\NormalTok{    simple     }\OperatorTok{=}\NormalTok{ [s }\ControlFlowTok{for}\NormalTok{ s }\KeywordTok{in}\NormalTok{ sub\_actions }\ControlFlowTok{if}\NormalTok{ s[}\StringTok{"action"}\NormalTok{] }\KeywordTok{not} \KeywordTok{in}\NormalTok{ \_LOCO]}

    \CommentTok{\# Pass 1: locomotion — interpolated keyframes across all movers}
    \ControlFlowTok{if}\NormalTok{ locomotion:}
\NormalTok{        … \_walk\_anims }\ControlFlowTok{for}\NormalTok{ each character, played together …}

    \CommentTok{\# Pass 2: simple parallel{-}safe sub{-}actions}
    \ControlFlowTok{if}\NormalTok{ simple:}
\NormalTok{        all\_anims }\OperatorTok{=}\NormalTok{ []}
        \ControlFlowTok{for}\NormalTok{ sub }\KeywordTok{in}\NormalTok{ simple:}
            \ControlFlowTok{for}\NormalTok{ tname }\KeywordTok{in}\NormalTok{ \_targets(sub):}
\NormalTok{                anims }\OperatorTok{=}\NormalTok{ \_dispatch\_one(sub, tname, collect\_anims}\OperatorTok{=}\VariableTok{True}\NormalTok{)}
                \ControlFlowTok{if}\NormalTok{ anims:}
\NormalTok{                    all\_anims.extend(anims)}
        \ControlFlowTok{if}\NormalTok{ all\_anims:}
\NormalTok{            scene.play(}\OperatorTok{*}\NormalTok{all\_anims, run\_time}\OperatorTok{=}\NormalTok{rt)}

    \CommentTok{\# Post{-}pass: update fig.pose for turn, clear cast for fade\_out}
\NormalTok{    …}
    \ControlFlowTok{continue}
\end{Highlighting}
\end{Shaded}

The two-pass structure exists because locomotion sub-actions need
keyframe-by-keyframe coordination (so two walkers stay in step), but
single-pose sub-actions like \texttt{wave} or \texttt{nod} can just have
their animation lists concatenated.

\subsection{\texorpdfstring{5.7 \texttt{trot\_to} --- DogGraph
routing}{5.7 trot\_to --- DogGraph routing}}\label{trot_to-doggraph-routing}

\texttt{trot\_to} is special because \texttt{DogGraph} instances live in
the \textbf{prop registry}, not the cast. When a step has
\texttt{"action":\ "trot\_to"} and a \texttt{"prop"} key (not
\texttt{"who"}), the main loop routes it directly to a DogGraph-aware
path:

\begin{Shaded}
\begin{Highlighting}[]
\ControlFlowTok{if}\NormalTok{ act }\OperatorTok{==} \StringTok{"trot\_to"} \KeywordTok{and} \StringTok{"prop"} \KeywordTok{in}\NormalTok{ step }\KeywordTok{and} \StringTok{"who"} \KeywordTok{not} \KeywordTok{in}\NormalTok{ step:}
\NormalTok{    \_dispatch\_one(step, step[}\StringTok{"prop"}\NormalTok{])}
    \ControlFlowTok{continue}
\end{Highlighting}
\end{Shaded}

\texttt{\_dispatch\_one} then resolves the dog via the prop registry
instead of the cast.

\textbf{Known inconsistency:} \texttt{trot\_to} reads
\texttt{step{[}"x"{]}} while \texttt{walk\_to} reads
\texttt{step{[}"to\_x"{]}}. Flagged for v1.0.0 normalization.

\subsection{\texorpdfstring{5.8 \texttt{\_subscene\_marker} --- camera
reframe
trigger}{5.8 \_subscene\_marker --- camera reframe trigger}}\label{subscene_marker-camera-reframe-trigger}

A \texttt{\_subscene\_marker} step has the form:

\begin{Shaded}
\begin{Highlighting}[]
\FunctionTok{\{}
  \DataTypeTok{"\_subscene\_marker"}\FunctionTok{:} \StringTok{"scene\_25\_kitchen\_ss03"}\FunctionTok{,}
  \DataTypeTok{"\_shot\_meta"}\FunctionTok{:} \FunctionTok{\{}
    \DataTypeTok{"framing"}\FunctionTok{:}    \StringTok{"medium"}\FunctionTok{,}
    \DataTypeTok{"subject"}\FunctionTok{:}    \StringTok{"freydoon"}\FunctionTok{,}
    \DataTypeTok{"move"}\FunctionTok{:}       \StringTok{"push"}\FunctionTok{,}
    \DataTypeTok{"transition"}\FunctionTok{:} \StringTok{"cut"}\FunctionTok{,}
    \DataTypeTok{"lighting"}\FunctionTok{:}   \KeywordTok{null}\FunctionTok{,}
    \DataTypeTok{"freeform"}\FunctionTok{:}   \KeywordTok{false}
  \FunctionTok{\}}
\FunctionTok{\}}
\end{Highlighting}
\end{Shaded}

When \texttt{PAM\_CAMERA\_MODE=1}, the main loop:

\begin{enumerate}
\def\labelenumi{\arabic{enumi}.}
\tightlist
\item
  Looks up the marker ID in \texttt{\_subscene\_index} (built from
  \texttt{PAM\_PROMPTS}), falling back to
  \texttt{step{[}"\_shot\_meta"{]}} if not found.
\item
  Calls \texttt{\_apply\_camera(meta,\ scene,\ char\_x\_positions)} for
  instant reframes, or stores \texttt{\_camera\_anim(...)} in
  \texttt{\_pending\_camera} for animated moves to fire with the next
  \texttt{say}.
\end{enumerate}

When \texttt{PAM\_CAMERA\_MODE} is unset, the entire subscene-marker
block is skipped silently.

\subsection{\texorpdfstring{5.9 \texttt{zone\_shift} --- sub-location
marker}{5.9 zone\_shift --- sub-location marker}}\label{zone_shift-sub-location-marker}

A \texttt{zone\_shift} step marks a sub-location within a scene:

\begin{Shaded}
\begin{Highlighting}[]
\FunctionTok{\{}\DataTypeTok{"action"}\FunctionTok{:} \StringTok{"zone\_shift"}\FunctionTok{,} \DataTypeTok{"zone"}\FunctionTok{:} \StringTok{"secretary\_s\_pod"}\FunctionTok{,}
 \DataTypeTok{"label"}\FunctionTok{:} \StringTok{"Secretary\textquotesingle{}s pod"}\FunctionTok{\}}
\end{Highlighting}
\end{Shaded}

Used by \texttt{fountain2pam.py} to chunk per-subscene prompts; in
\texttt{pam\_player.py} it's an instantaneous camera-region shift (zero
animation cost) and does not close the current subscene.

\subsection{5.10 Scene overlays}\label{scene-overlays}

\begin{longtable}[]{@{}
  >{\raggedright\arraybackslash}p{(\columnwidth - 2\tabcolsep) * \real{0.5000}}
  >{\raggedright\arraybackslash}p{(\columnwidth - 2\tabcolsep) * \real{0.5000}}@{}}
\toprule\noalign{}
\begin{minipage}[b]{\linewidth}\raggedright
Action
\end{minipage} & \begin{minipage}[b]{\linewidth}\raggedright
Effect
\end{minipage} \\
\midrule\noalign{}
\endhead
\bottomrule\noalign{}
\endlastfoot
\texttt{title} & Opening title card. Should be the first action. \\
\texttt{caption} & On-screen text with \texttt{position}
(\texttt{bottom}/\texttt{top}/\texttt{lower-third}), \texttt{duration},
\texttt{style} (\texttt{normal}/\texttt{italic}/\texttt{bold}). Fades
in, holds, fades out. \\
\texttt{persistent\_caption} & Same as \texttt{caption} but doesn't
auto-dismiss. \\
\texttt{overlay\_caption} & v0.9.11 --- like \texttt{caption} but with a
\texttt{position:\ "center"} that places the bar at
\texttt{\_frame\_cy\ +\ 0.5} (above standing characters). \\
\texttt{sound\_cue} & Flashes a diegetic sound label on screen. \\
\texttt{wait} & Pauses the scene for \texttt{t} seconds. \\
\texttt{on\_screen\_text} & Multi-line text card with \texttt{hold}
parameter. \\
\texttt{clear\_bubble} / \texttt{clear\_all\_bubbles} & See §8.3. \\
\end{longtable}

\begin{center}\rule{0.5\linewidth}{0.5pt}\end{center}

\section{6. The parallel block}\label{the-parallel-block}

\subsection{\texorpdfstring{6.1 The \texttt{\_CANNOT\_PARALLEL}
set}{6.1 The \_CANNOT\_PARALLEL set}}\label{the-_cannot_parallel-set}

\texttt{actions.py} exports a frozenset of action names that cannot run
inside a \texttt{parallel} block:

\begin{Shaded}
\begin{Highlighting}[]
\NormalTok{\_CANNOT\_PARALLEL }\OperatorTok{=} \BuiltInTok{frozenset}\NormalTok{(\{}
    \StringTok{"walk\_to"}\NormalTok{, }\StringTok{"run\_to"}\NormalTok{, }\StringTok{"sit\_down"}\NormalTok{, }\StringTok{"stand\_up"}\NormalTok{, }\StringTok{"wave"}\NormalTok{,}
    \StringTok{"wave\_left"}\NormalTok{, }\StringTok{"wave\_right"}\NormalTok{, }\StringTok{"carry"}\NormalTok{, }\StringTok{"exit\_through"}\NormalTok{,}
    \StringTok{"exit\_through\_doors"}\NormalTok{, }\StringTok{"rush\_to"}\NormalTok{, }\StringTok{"rush\_out"}\NormalTok{, }\StringTok{"squeeze\_through"}\NormalTok{,}
    \StringTok{"jump\_up"}\NormalTok{, }\StringTok{"pat"}\NormalTok{, }\StringTok{"search\_drawers"}\NormalTok{, }\StringTok{"pick\_up\_phone"}\NormalTok{, }\StringTok{"hang\_up"}\NormalTok{,}
    \StringTok{"grab"}\NormalTok{, }\StringTok{"punch\_button"}\NormalTok{, }\StringTok{"reach\_character"}\NormalTok{, }\StringTok{"peel\_from\_hand"}\NormalTok{,}
    \StringTok{"group\_translate"}\NormalTok{, }\StringTok{"nod"}\NormalTok{, }\StringTok{"shake\_head"}\NormalTok{, }\StringTok{"shrug"}\NormalTok{,}
    \StringTok{"grab\_arm"}\NormalTok{, }\StringTok{"twist\_arm\_behind"}\NormalTok{, }\StringTok{"release\_arm"}\NormalTok{, }\StringTok{"rotate"}\NormalTok{,}
\NormalTok{\})}
\end{Highlighting}
\end{Shaded}

When the parallel handler encounters one of these inside a
\texttt{parallel} block, it prints a console warning and runs the
sub-action sequentially (outside the \texttt{scene.play} batch).

\textbf{Locomotion is the exception.} \texttt{walk\_to} and
\texttt{run\_to} are technically in \texttt{\_CANNOT\_PARALLEL} because
they can't be \texttt{\_collect\_anims}'d generically, but the parallel
handler has a \emph{dedicated path} (§6.3) that interpolates their
keyframes in lockstep.

\subsection{6.2 Two-pass collection}\label{two-pass-collection}

The parallel handler partitions sub-actions into two buckets:

\begin{enumerate}
\def\labelenumi{\arabic{enumi}.}
\tightlist
\item
  \textbf{Locomotion} --- handled via the keyframe-by-keyframe special
  path
\item
  \textbf{Simple parallel-safe} --- collected as Manim animation lists
  and played in a single \texttt{scene.play(*anims,\ run\_time=rt)} call
\end{enumerate}

Each bucket is played in its own \texttt{scene.play} invocation; they
don't interleave.

\subsection{6.3 Locomotion sub-actions}\label{locomotion-sub-actions}

The locomotion path zips per-character walk keyframes:

\begin{Shaded}
\begin{Highlighting}[]
\NormalTok{loco\_steps }\OperatorTok{=} \BuiltInTok{max}\NormalTok{(fig.\_walk\_steps(to\_x) }\ControlFlowTok{for}\NormalTok{ to\_x }\KeywordTok{in}\NormalTok{ destinations)}

\ControlFlowTok{for}\NormalTok{ step\_i }\KeywordTok{in} \BuiltInTok{range}\NormalTok{(loco\_steps):}
\NormalTok{    frame\_anims }\OperatorTok{=}\NormalTok{ []}
    \ControlFlowTok{for}\NormalTok{ mover, dest }\KeywordTok{in} \BuiltInTok{zip}\NormalTok{(movers, destinations):}
\NormalTok{        pose, new\_off }\OperatorTok{=}\NormalTok{ mover.\_walk\_keyframe(step\_i, dest)}
\NormalTok{        frame\_anims.extend(mover.\_pose\_anims(pose, new\_off))}
\NormalTok{        mover.pose, mover.offset }\OperatorTok{=}\NormalTok{ pose, new\_off}
\NormalTok{    scene.play(}\OperatorTok{*}\NormalTok{frame\_anims, run\_time}\OperatorTok{=}\NormalTok{loco\_rt, rate\_func}\OperatorTok{=}\NormalTok{smooth)}
\end{Highlighting}
\end{Shaded}

This is the mechanism that lets Lucy and Lenny run together while Chekov
trots behind --- every character's footfall lands on the same beat. The
\texttt{rt\_per\_kf} parameter on \texttt{parallel} controls the
per-keyframe duration (default 0.12s).

\subsection{6.4 Post-pass state updates}\label{post-pass-state-updates}

After the \texttt{scene.play(...)} call completes, the parallel handler
runs a post-pass to update figure state for actions whose state changes
happen \emph{after} the animation rather than within it:

\begin{Shaded}
\begin{Highlighting}[]
\ControlFlowTok{for}\NormalTok{ sub }\KeywordTok{in}\NormalTok{ sub\_actions:}
    \ControlFlowTok{if}\NormalTok{ sub[}\StringTok{"action"}\NormalTok{] }\OperatorTok{==} \StringTok{"turn"}\NormalTok{:}
        \ControlFlowTok{for}\NormalTok{ tname }\KeywordTok{in}\NormalTok{ \_targets(sub):}
\NormalTok{            fig }\OperatorTok{=}\NormalTok{ \_get\_fig(tname)}
            \ControlFlowTok{if}\NormalTok{ fig:}
\NormalTok{                fig.pose }\OperatorTok{=}\NormalTok{ \_resolve\_pose(sub.get(}\StringTok{"pose"}\NormalTok{),}
\NormalTok{                                         fig.\_bp[}\StringTok{"standing\_side"}\NormalTok{])}
    \ControlFlowTok{elif}\NormalTok{ sub[}\StringTok{"action"}\NormalTok{] }\OperatorTok{==} \StringTok{"fade\_out"}\NormalTok{:}
        \ControlFlowTok{for}\NormalTok{ tname }\KeywordTok{in}\NormalTok{ \_targets(sub):}
\NormalTok{            cast[tname][}\StringTok{"fig"}\NormalTok{] }\OperatorTok{=} \VariableTok{None}
\end{Highlighting}
\end{Shaded}

Without this post-pass, a \texttt{turn} inside a \texttt{parallel} block
would animate the figure to the side pose but leave \texttt{fig.pose}
set to \texttt{standing\_front}, breaking the next \texttt{walk\_to}.

\begin{center}\rule{0.5\linewidth}{0.5pt}\end{center}

\section{7. The camera system}\label{the-camera-system}

\subsection{\texorpdfstring{7.1 \texttt{MovingCameraScene}
foundations}{7.1 MovingCameraScene foundations}}\label{movingcamerascene-foundations}

\texttt{PAMPlayer} inherits from Manim's \texttt{MovingCameraScene},
which exposes \texttt{self.camera.frame} --- a \texttt{ScreenRectangle}
whose \texttt{width} and \texttt{center} control the visible viewport.
PAM's camera system manipulates this frame:

\begin{itemize}
\tightlist
\item
  \textbf{Width} controls zoom (smaller width = more zoomed in)
\item
  \textbf{Center} controls pan (x for horizontal, y for vertical)
\end{itemize}

The camera is purely 2D --- there's no perspective, no depth-of-field,
no real ``tilt.'' The \texttt{pan-up} action approximates a tilt by
panning the frame upward while applying a shear to tall props (§7.8).

\subsection{\texorpdfstring{7.2 The \texttt{\_FRAMING\_CAMERA}
dict}{7.2 The \_FRAMING\_CAMERA dict}}\label{the-_framing_camera-dict}

Module-level constant mapping framing names to
\texttt{(frame\_width,\ frame\_center\_y)}:

\begin{Shaded}
\begin{Highlighting}[]
\NormalTok{\_FRAMING\_CAMERA: }\BuiltInTok{dict}\NormalTok{[}\BuiltInTok{str}\NormalTok{, }\BuiltInTok{tuple}\NormalTok{] }\OperatorTok{=}\NormalTok{ \{}
    \StringTok{"wide"}\NormalTok{:         (}\FloatTok{14.2}\NormalTok{, }\OperatorTok{{-}}\FloatTok{0.5}\NormalTok{),   }\CommentTok{\# full stage}
    \StringTok{"medium"}\NormalTok{:       (}\FloatTok{8.0}\NormalTok{,  }\OperatorTok{{-}}\FloatTok{0.2}\NormalTok{),   }\CommentTok{\# waist{-}up}
    \StringTok{"medium{-}close"}\NormalTok{: (}\FloatTok{5.5}\NormalTok{,   }\FloatTok{0.3}\NormalTok{),   }\CommentTok{\# chest{-}up}
    \StringTok{"close"}\NormalTok{:        (}\FloatTok{3.5}\NormalTok{,   }\FloatTok{0.9}\NormalTok{),   }\CommentTok{\# face and shoulders}
    \StringTok{"ots{-}left"}\NormalTok{:     (}\FloatTok{7.0}\NormalTok{,   }\FloatTok{0.0}\NormalTok{),}
    \StringTok{"ots{-}right"}\NormalTok{:    (}\FloatTok{7.0}\NormalTok{,   }\FloatTok{0.0}\NormalTok{),}
    \StringTok{"oneshot"}\NormalTok{:      (}\FloatTok{5.0}\NormalTok{,   }\FloatTok{0.3}\NormalTok{),   }\CommentTok{\# single character}
    \StringTok{"insert"}\NormalTok{:       (}\FloatTok{3.0}\NormalTok{,   }\FloatTok{1.5}\NormalTok{),   }\CommentTok{\# prop / detail level}
\NormalTok{\}}
\end{Highlighting}
\end{Shaded}

Smaller width = tighter framing. \texttt{center\_y} should track the
visual center of the framed subject --- head-level for \texttt{close},
mid-body for \texttt{medium}.

\subsection{\texorpdfstring{7.3 The \texttt{\_MOVE\_RT}
dict}{7.3 The \_MOVE\_RT dict}}\label{the-_move_rt-dict}

Maps move names to animation durations:

\begin{Shaded}
\begin{Highlighting}[]
\NormalTok{\_MOVE\_RT: }\BuiltInTok{dict}\NormalTok{[}\BuiltInTok{str}\NormalTok{, }\BuiltInTok{float}\NormalTok{] }\OperatorTok{=}\NormalTok{ \{}
    \StringTok{"static"}\NormalTok{:     }\FloatTok{0.0}\NormalTok{,    }\CommentTok{\# instant cut}
    \StringTok{"push"}\NormalTok{:       }\FloatTok{0.8}\NormalTok{,}
    \StringTok{"pull"}\NormalTok{:       }\FloatTok{0.8}\NormalTok{,}
    \StringTok{"pan{-}follow"}\NormalTok{: }\FloatTok{0.5}\NormalTok{,}
    \StringTok{"drift"}\NormalTok{:      }\FloatTok{1.5}\NormalTok{,}
\NormalTok{\}}
\end{Highlighting}
\end{Shaded}

\texttt{static} moves bypass \texttt{scene.play} entirely --- the frame
is repositioned instantly between actions. All other moves animate over
their listed duration.

\subsection{\texorpdfstring{7.4 \texttt{\_apply\_camera} --- instant and
animated
reframes}{7.4 \_apply\_camera --- instant and animated reframes}}\label{apply_camera-instant-and-animated-reframes}

The authoritative camera reposition function:

\begin{Shaded}
\begin{Highlighting}[]
\KeywordTok{def}\NormalTok{ \_apply\_camera(meta, scene, char\_x\_positions):}
\NormalTok{    framing }\OperatorTok{=}\NormalTok{ meta.get(}\StringTok{"framing"}\NormalTok{, }\StringTok{"wide"}\NormalTok{).lower()}
\NormalTok{    move    }\OperatorTok{=}\NormalTok{ meta.get(}\StringTok{"move"}\NormalTok{,    }\StringTok{"static"}\NormalTok{).lower()}
\NormalTok{    subject }\OperatorTok{=}\NormalTok{ meta.get(}\StringTok{"subject"}\NormalTok{, }\StringTok{"ensemble"}\NormalTok{).lower()}

\NormalTok{    target\_w, target\_y }\OperatorTok{=}\NormalTok{ \_FRAMING\_CAMERA[framing]}
\NormalTok{    rt                 }\OperatorTok{=}\NormalTok{ \_MOVE\_RT.get(move, }\FloatTok{0.0}\NormalTok{)}

    \CommentTok{\# Resolve subject to a target x}
    \ControlFlowTok{if}\NormalTok{ subject }\KeywordTok{in}\NormalTok{ (}\StringTok{"ensemble"}\NormalTok{, }\StringTok{"none"}\NormalTok{, }\StringTok{""}\NormalTok{):}
\NormalTok{        target\_x }\OperatorTok{=} \FloatTok{0.0}
    \ControlFlowTok{else}\NormalTok{:}
\NormalTok{        raw\_x  }\OperatorTok{=}\NormalTok{ char\_x\_positions.get(subject, }\FloatTok{0.0}\NormalTok{)}
\NormalTok{        half\_w }\OperatorTok{=}\NormalTok{ target\_w }\OperatorTok{/} \DecValTok{2}
\NormalTok{        target\_x }\OperatorTok{=} \BuiltInTok{float}\NormalTok{(np.clip(raw\_x,}
                                 \OperatorTok{{-}}\FloatTok{7.1} \OperatorTok{+}\NormalTok{ half\_w, }\FloatTok{7.1} \OperatorTok{{-}}\NormalTok{ half\_w))}

\NormalTok{    frame }\OperatorTok{=}\NormalTok{ scene.camera.frame}
    \ControlFlowTok{if}\NormalTok{ rt }\OperatorTok{\textgreater{}} \DecValTok{0}\NormalTok{:}
\NormalTok{        scene.play(}
\NormalTok{            frame.animate.set\_width(target\_w)}
\NormalTok{                  .move\_to([target\_x, target\_y, }\DecValTok{0}\NormalTok{]),}
\NormalTok{            run\_time}\OperatorTok{=}\NormalTok{rt, rate\_func}\OperatorTok{=}\NormalTok{smooth)}
    \ControlFlowTok{else}\NormalTok{:}
\NormalTok{        frame.set\_width(target\_w)}
\NormalTok{        frame.move\_to([target\_x, target\_y, }\DecValTok{0}\NormalTok{])}
\end{Highlighting}
\end{Shaded}

The clamp on \texttt{target\_x} keeps the frame inside the visible stage
so that following a character near the edge doesn't reveal off-stage
space.

\subsection{\texorpdfstring{7.5 \texttt{\_camera\_anim} and
\texttt{\_pending\_camera}}{7.5 \_camera\_anim and \_pending\_camera}}\label{camera_anim-and-_pending_camera}

A variant that \textbf{returns} an animation instead of playing it. Used
when a camera move should run \emph{concurrently} with a speech bubble:

\begin{Shaded}
\begin{Highlighting}[]
\KeywordTok{def}\NormalTok{ \_camera\_anim(meta, scene, char\_x\_positions):}
    \CommentTok{"""Return a camera animation (or None) — does NOT play it."""}
\NormalTok{    …}
    \ControlFlowTok{if}\NormalTok{ rt }\OperatorTok{==} \FloatTok{0.0}\NormalTok{:}
        \ControlFlowTok{return} \VariableTok{None}    \CommentTok{\# static — handled by \_apply\_camera elsewhere}
    \ControlFlowTok{return}\NormalTok{ frame.animate.set\_width(target\_w).move\_to(}
\NormalTok{        np.array([target\_x, target\_y, }\DecValTok{0}\NormalTok{]))}
\end{Highlighting}
\end{Shaded}

The \texttt{\_pending\_camera} mechanism in \texttt{construct()}:

\begin{Shaded}
\begin{Highlighting}[]
\VariableTok{self}\NormalTok{.\_pending\_camera }\OperatorTok{=} \VariableTok{None}

\CommentTok{\# When a \_subscene\_marker fires with an animated move:}
\NormalTok{anim }\OperatorTok{=}\NormalTok{ \_camera\_anim(meta, }\VariableTok{self}\NormalTok{, char\_x\_positions)}
\VariableTok{self}\NormalTok{.\_pending\_camera }\OperatorTok{=}\NormalTok{ anim}

\CommentTok{\# When the next \textasciigrave{}say\textasciigrave{} fires:}
\NormalTok{extra\_anims }\OperatorTok{=}\NormalTok{ []}
\ControlFlowTok{if} \VariableTok{self}\NormalTok{.\_pending\_camera }\KeywordTok{is} \KeywordTok{not} \VariableTok{None}\NormalTok{:}
\NormalTok{    extra\_anims.append(}\VariableTok{self}\NormalTok{.\_pending\_camera)}
    \VariableTok{self}\NormalTok{.\_pending\_camera }\OperatorTok{=} \VariableTok{None}
\NormalTok{fig.say(text, scene}\OperatorTok{=}\VariableTok{self}\NormalTok{, extra\_anims}\OperatorTok{=}\NormalTok{extra\_anims, …)}
\end{Highlighting}
\end{Shaded}

This makes camera moves invisible \emph{as moves} --- they feel like the
camera is gliding toward the speaker during their first beat of
dialogue. Without \texttt{\_pending\_camera}, the camera would jump,
pause, then the bubble would appear separately.

If no \texttt{say} follows the marker before the next marker, the
pending camera is consumed by the next marker's \texttt{\_apply\_camera}
call directly.

\subsection{\texorpdfstring{7.6 The \texttt{\_shot\_meta}
dict}{7.6 The \_shot\_meta dict}}\label{the-_shot_meta-dict}

Every \texttt{\_subscene\_marker} carries an inline
\texttt{\_shot\_meta} dict with six keys:

\begin{longtable}[]{@{}
  >{\raggedright\arraybackslash}p{(\columnwidth - 4\tabcolsep) * \real{0.3333}}
  >{\raggedright\arraybackslash}p{(\columnwidth - 4\tabcolsep) * \real{0.3333}}
  >{\raggedright\arraybackslash}p{(\columnwidth - 4\tabcolsep) * \real{0.3333}}@{}}
\toprule\noalign{}
\begin{minipage}[b]{\linewidth}\raggedright
Key
\end{minipage} & \begin{minipage}[b]{\linewidth}\raggedright
Values
\end{minipage} & \begin{minipage}[b]{\linewidth}\raggedright
Default
\end{minipage} \\
\midrule\noalign{}
\endhead
\bottomrule\noalign{}
\endlastfoot
\texttt{framing} & \texttt{wide} / \texttt{medium} /
\texttt{medium-close} / \texttt{close} / \texttt{ots-left} /
\texttt{ots-right} / \texttt{oneshot} / \texttt{insert} &
\texttt{wide} \\
\texttt{subject} & character key / prop key / \texttt{ensemble} /
\texttt{none} & \texttt{ensemble} \\
\texttt{move} & \texttt{static} / \texttt{push} / \texttt{pull} /
\texttt{pan-follow} / \texttt{pan-up} / \texttt{pan-down} /
\texttt{drift} & \texttt{static} \\
\texttt{transition} & \texttt{cut} / \texttt{hold} / \texttt{hold-empty}
/ \texttt{smash} & \texttt{cut} \\
\texttt{lighting} & (lighting vocab) & \texttt{null} \\
\texttt{freeform} & bool & \texttt{false} \\
\end{longtable}

\texttt{null} values mean ``inherit current state'' --- they're not the
same as explicit defaults. A subscene with \texttt{framing:\ null} after
a \texttt{pan-follow} will inherit the panned frame. A subscene with
\texttt{framing:\ "wide"} will reset to wide. This distinction matters
when a previous shot moved the camera.

\subsection{7.7 Initial camera reset}\label{initial-camera-reset}

When \texttt{PAM\_CAMERA\_MODE=1}, \texttt{construct()} forces the
camera to wide/ensemble \emph{before} the action loop:

\begin{Shaded}
\begin{Highlighting}[]
\ControlFlowTok{if}\NormalTok{ \_camera\_mode:}
\NormalTok{    init\_frame }\OperatorTok{=} \VariableTok{self}\NormalTok{.camera.frame}
\NormalTok{    init\_frame.width }\OperatorTok{=} \FloatTok{14.2}
\NormalTok{    init\_frame.move\_to(np.array([}\FloatTok{0.0}\NormalTok{, }\OperatorTok{{-}}\FloatTok{0.5}\NormalTok{, }\DecValTok{0}\NormalTok{]))}
\end{Highlighting}
\end{Shaded}

Without this reset, the camera could inherit a zoomed or off-center
state from a previous render in the same Manim process.

\subsection{\texorpdfstring{7.8 \texttt{pan-up} / \texttt{pan-down} and
the tilt
implementation}{7.8 pan-up / pan-down and the tilt implementation}}\label{pan-up-pan-down-and-the-tilt-implementation}

\texttt{pan-up} is implemented by \texttt{\_execute\_pan\_up}, which:

\begin{enumerate}
\def\labelenumi{\arabic{enumi}.}
\tightlist
\item
  Zooms in to \texttt{medium-close} on the named subject.
\item
  Pans the frame upward at the configured tilt duration (2.5s).
\item
  Optionally applies a vertical shear to the target building prop to
  fake keystone convergence.
\item
  Holds briefly at the top.
\end{enumerate}

The target prop must have a \texttt{pam\_height} attribute ---
\texttt{build\_building} sets this, but other tall props may not. If the
subject doesn't have \texttt{pam\_height}, the camera tilts to a default
upper-frame position.

\texttt{pan-down} mirrors the implementation, starting high and panning
to ground level.

\subsection{\texorpdfstring{7.9 The \texttt{insert} framing for prop
close-ups}{7.9 The insert framing for prop close-ups}}\label{the-insert-framing-for-prop-close-ups}

\texttt{insert} is the tightest framing (frame width 3.0, center\_y 1.5)
intended for prop detail shots:

\begin{Shaded}
\begin{Highlighting}[]
\NormalTok{[[ CAMERA: FRAMING=insert | SUBJECT=smartphone | MOVE=null | TRANSITION=cut ]]}
\end{Highlighting}
\end{Shaded}

When \texttt{subject} resolves to a prop name (rather than a character),
the camera centers on the prop's \texttt{pam\_x} / \texttt{pam\_y}
instead of char\_x\_positions. Used for phone screens, monitor displays,
audio\_video\_bug shots, and similar tiny-prop moments.

\begin{center}\rule{0.5\linewidth}{0.5pt}\end{center}

\section{8. The bubble registry}\label{the-bubble-registry}

\subsection{\texorpdfstring{8.1 \texttt{\_persistent\_bubbles}
structure}{8.1 \_persistent\_bubbles structure}}\label{persistent_bubbles-structure}

A class-level dict on \texttt{PAMPlayer}:

\begin{Shaded}
\begin{Highlighting}[]
\NormalTok{\_persistent\_bubbles: }\BuiltInTok{dict}\NormalTok{[}\BuiltInTok{str}\NormalTok{, }\BuiltInTok{dict}\NormalTok{] }\OperatorTok{=}\NormalTok{ \{\}}
\CommentTok{\# Structure:}
\CommentTok{\#   key   = character name OR prop registry name}
\CommentTok{\#   value = \{"bubble": VGroup, "expires\_at": float | None\}}
\end{Highlighting}
\end{Shaded}

An entry is added whenever a \texttt{say} or \texttt{prop\_say} includes
either \texttt{persist:\ true} (no expiry --- clear explicitly) or
\texttt{duration:\ t} (expires at scene-clock + t seconds).

\subsection{\texorpdfstring{8.2 \texttt{\_scene\_clock} and
\texttt{\_sweep\_expired\_bubbles}}{8.2 \_scene\_clock and \_sweep\_expired\_bubbles}}\label{scene_clock-and-_sweep_expired_bubbles}

The scene clock is a mutable holder:

\begin{Shaded}
\begin{Highlighting}[]
\NormalTok{\_scene\_clock: }\BuiltInTok{list} \OperatorTok{=}\NormalTok{ [}\FloatTok{0.0}\NormalTok{]}
\end{Highlighting}
\end{Shaded}

It's incremented after every \texttt{scene.play(...)} and
\texttt{scene.wait(...)} call by the run\_time, tracking wall-clock
seconds of scene time.

At the top of each main-loop iteration:

\begin{Shaded}
\begin{Highlighting}[]
\KeywordTok{def}\NormalTok{ \_sweep\_expired\_bubbles(}\VariableTok{self}\NormalTok{):}
\NormalTok{    now }\OperatorTok{=} \VariableTok{self}\NormalTok{.\_scene\_clock[}\DecValTok{0}\NormalTok{]}
\NormalTok{    expired }\OperatorTok{=}\NormalTok{ [k }\ControlFlowTok{for}\NormalTok{ k, e }\KeywordTok{in} \VariableTok{self}\NormalTok{.\_persistent\_bubbles.items()}
               \ControlFlowTok{if}\NormalTok{ e[}\StringTok{"expires\_at"}\NormalTok{] }\KeywordTok{is} \KeywordTok{not} \VariableTok{None} \KeywordTok{and}\NormalTok{ e[}\StringTok{"expires\_at"}\NormalTok{] }\OperatorTok{\textless{}=}\NormalTok{ now]}
    \ControlFlowTok{for}\NormalTok{ k }\KeywordTok{in}\NormalTok{ expired:}
\NormalTok{        bubble }\OperatorTok{=} \VariableTok{self}\NormalTok{.\_persistent\_bubbles.pop(k)[}\StringTok{"bubble"}\NormalTok{]}
        \VariableTok{self}\NormalTok{.play(FadeOut(bubble), run\_time}\OperatorTok{=}\FloatTok{0.25}\NormalTok{)}
\end{Highlighting}
\end{Shaded}

This is what makes \texttt{duration:\ 2.5} work --- the bubble persists
until its expiry passes, at which point the next main-loop iteration
fades it out automatically.

\subsection{\texorpdfstring{8.3 \texttt{clear\_bubble} /
\texttt{clear\_prop\_bubble} /
\texttt{clear\_all\_bubbles}}{8.3 clear\_bubble / clear\_prop\_bubble / clear\_all\_bubbles}}\label{clear_bubble-clear_prop_bubble-clear_all_bubbles}

Explicit-dismissal actions for \texttt{persist:\ true} bubbles (which
never auto-expire).

\begin{Shaded}
\begin{Highlighting}[]
\FunctionTok{\{}\DataTypeTok{"action"}\FunctionTok{:} \StringTok{"clear\_bubble"}\FunctionTok{,} \DataTypeTok{"who"}\FunctionTok{:} \StringTok{"thalia"}\FunctionTok{,} \DataTypeTok{"rt"}\FunctionTok{:} \DecValTok{0}\ErrorTok{.}\DecValTok{25}\FunctionTok{\}}
\FunctionTok{\{}\DataTypeTok{"action"}\FunctionTok{:} \StringTok{"clear\_bubble"}\FunctionTok{,} \DataTypeTok{"prop"}\FunctionTok{:} \StringTok{"phone1"}\FunctionTok{\}}
\FunctionTok{\{}\DataTypeTok{"action"}\FunctionTok{:} \StringTok{"clear\_all\_bubbles"}\FunctionTok{,} \DataTypeTok{"rt"}\FunctionTok{:} \DecValTok{0}\ErrorTok{.}\DecValTok{25}\FunctionTok{\}}
\end{Highlighting}
\end{Shaded}

Each looks up the entry in \texttt{\_persistent\_bubbles} by
\texttt{who} or \texttt{prop}, plays a \texttt{FadeOut}, and removes the
entry. \texttt{clear\_all\_bubbles} iterates the dict and fades
everything concurrently in a single \texttt{scene.play} call.

\texttt{clear\_prop\_bubble} is a v0.9.9 alias for
\texttt{clear\_bubble} retained for scenes authored before the rename.

% =====================================================================
% Replaces lines 2328--2359 of pam-reference-book_v0916-expanded.tex.
% Drop in between the §8.3 closing line ("...rename.") and the existing
% §8.5 ("Bubble world-space anchoring") subsection.  No other edits to
% the surrounding file are required: §8.5 keeps its number, §8.6 still
% follows, etc.  The \label is preserved so existing cross-references
% (search for \ref{the-focus_reset-overlay-gotcha}) keep resolving.
% =====================================================================

\subsection{\texorpdfstring{8.4 The \texttt{focus} / \texttt{focus\_reset}
re-paint gotcha}{8.4 The focus / focus\_reset re-paint
gotcha}}\label{the-focus_reset-overlay-gotcha}

\textbf{The general rule.} \texttt{focus} and \texttt{focus\_reset}
animate \texttt{fig.group.animate.set\_opacity(x)} on each affected
figure. Because \texttt{fig.group} is a property returning
\texttt{VGroup(edge\_group,\ dot\_group)}, the animation walks every
edge and every joint dot --- including any dot whose opacity was
previously set to \texttt{0} by an attach-style action. Two consequences
fire together at the end of the \texttt{play()} call:

\begin{enumerate}
\tightlist
\item
  Anything you had hidden by setting its opacity to \texttt{0} inside
  \texttt{fig.group} is re-revealed at the new target opacity.
\item
  Manim's painter's algorithm re-asserts the animated mobjects above
  anything that was painted on top of them between when they were
  originally added and when \texttt{focus} fired.
\end{enumerate}

So any decoration that lives \emph{above} an opacity-hidden part of
\texttt{fig.group} is at risk: the part underneath comes back into view
\emph{and} jumps in front of the decoration. The rule has two known
instances.

\medskip

\textbf{Example A: persistent speech bubbles (v0.9.10, unfixed).}

\textbf{Symptom.} A persistent bubble (from \texttt{say} or
\texttt{prop\_say} with \texttt{persist:\ true} or \texttt{duration:\ t})
shows on screen. A \texttt{focus\_reset} fires, restoring all characters
to full opacity. The bubble is now visually overlaid by the character
mobjects.

\textbf{Cause.} The bubble was added to the scene after the characters.
\texttt{focus\_reset} re-plays the characters, and the painter's
algorithm puts them on top of the bubble.

\textbf{Workaround.} Issue a \texttt{clear\_bubble} or
\texttt{clear\_all\_bubbles} before the \texttt{focus\_reset}:

\begin{Shaded}
\begin{Highlighting}[]
\FunctionTok{\{}\DataTypeTok{"action"}\FunctionTok{:} \StringTok{"clear\_all\_bubbles"}\FunctionTok{\}}\ErrorTok{,}
\FunctionTok{\{}\DataTypeTok{"action"}\FunctionTok{:} \StringTok{"focus\_reset"}\FunctionTok{\}}
\end{Highlighting}
\end{Shaded}

Or tune the bubble's \texttt{duration} so it expires before the reset
fires. This case is not patched in-handler because the right semantic
answer (``should bubbles always render above characters?'') is part of
the v1.0.0 bubble-lifecycle migration --- see
\texttt{BACK\_BURNER.md} and \S 8.5 (world-space anchoring), which is
the larger refactor this gotcha is waiting on.

\medskip

\textbf{Example B: attached faces (v0.9.15, fixed in v0.9.17).}

\textbf{Symptom.} A character has had \texttt{attach\_face} applied,
so a face PNG is showing on the head. A \texttt{focus} fires with
\texttt{bright:\ 1.0} (or a \texttt{focus\_reset}) targeting that
character. The face PNG is replaced by the underlying head dot --- a
dark oval with the character's label, e.g.\ ``Bevers''.

\textbf{Cause.} \texttt{attach\_face} hides the head dot by setting
\texttt{fig.dots["head"].opacity\ =\ 0} and adds a \texttt{head\_face}
\texttt{ImageMobject} above it. The head dot is a child of
\texttt{fig.dots}, hence of \texttt{fig.group}, so the focus
animation drives its opacity back to \texttt{1.0} and the play also
re-orders it above \texttt{head\_face}. The labeled dot reappears,
clobbering the face. (At \texttt{dim} opacity the same mechanism
applies but the dot only animates to \texttt{0.30}, so visually the
dimmed face still dominates --- which is why the bug shows
asymmetrically on the brightened characters only.)

\textbf{Fix (in-handler, v0.9.17, end-state repair).} After each
focus/reset \texttt{play()}, the handler walks every figure it touched
(both brightened and dimmed) and, for any figure with a live
\texttt{head\_face}, reasserts
\texttt{fig.dots["head"].set\_opacity(0)} synchronously and calls
\texttt{self.bring\_to\_front(fig.head\_face)}. This corrected the
final frame but left intermediate animation frames showing the dot
peek through --- visible flicker (5--10 frames at 24/30fps) on focus
calls with \texttt{rt < 0.05}.

\textbf{Upgrade (v0.9.18, transient elimination).} The handler now
selects its animation target conditionally for each figure. For
figures with no \texttt{head\_face}, it animates \texttt{fig.group}
exactly as before. For figures with a live \texttt{head\_face}, it
animates a rebuilt \texttt{VGroup(fig.edge\_group,\ *non\_head\_dots)}
that EXCLUDES \texttt{fig.dots["head"]} entirely. The head dot is
never in the animation, so no interpolated frame can show it above
the face. The v0.9.17 post-play repair is retained as defensive
belt-and-suspenders --- its \texttt{set\_opacity(0)} call is now a
no-op (the dot was never touched), but its
\texttt{bring\_to\_front(head\_face)} is still meaningful for z-order
reassertion if any unrelated later animation re-paints over the face.
See the \texttt{\_animation\_group} helper in
\texttt{pam\_player.py}'s focus/focus\_reset handler. No screenplay
changes are required.

\textbf{Refinement (v0.9.18.1, residual neck/shoulder flicker).}
Excluding the head dot left a subtler artifact: the animated group
still contains \texttt{fig.edge\_group} and the non-head dots, which
include the neck edge and the shoulder dots/struts that a
head-and-shoulders face bust covers in the static frame. Manim's
\texttt{play()} re-adds the animated group to the front of the scene
each frame, so those covered body elements paint ABOVE the face for
the duration of the animation, then snap back behind it when the
post-play \texttt{bring\_to\_front} runs --- a brief flash of
neck/shoulder lines over the face. Rather than enumerate exactly
which edges and dots the bust covers (build-specific, and the set
would drift as the art evolves), the handler installs a per-frame
updater on each touched figure's \texttt{head\_face} for the span of
the \texttt{play()}. The updater re-asserts
\texttt{bring\_to\_front(head\_face)} every frame, holding the face
on top regardless of what is re-added behind it. The updaters are
removed immediately after the \texttt{play()}; steady-state z-order
is then handled by \texttt{\_repair\_face\_overlays} as before.
Because nothing persistent changes, there is no interaction with
bubble or caption z-order outside the focus window. See the
\texttt{\_install\_face\_front\_keepers} /
\texttt{\_remove\_face\_front\_keepers} helpers in the focus handler.

\medskip

\textbf{Why the two examples are treated differently.} The face case
has one unambiguous correct rendering: the face PNG was put there
deliberately to obscure the dot, so any reveal of the dot is a bug.
The bubble case is genuinely ambiguous: a persistent bubble may or
may not be intended to outlast a focus reset, and the bubble's
anchoring story (\S 8.5) is itself in flux. The face fix can be a
local handler patch; the bubble fix needs the larger refactor.

\medskip

\textbf{Generalization to future attach methods.} The mechanism is
not face-specific. Any future \texttt{attach\_*} method that hides
part of \texttt{fig.group} by setting a child's opacity to \texttt{0}
and painting something on top will collide with this rule. Two
options for new code:

\begin{enumerate}
\tightlist
\item
  Teach the v0.9.18 \texttt{\_animation\_group} helper to also
  exclude the hidden child for any figure that has the new attached
  attribute. Same exclusion pattern, parameterized by the attached
  attribute name.
\item
  Avoid the hide-with-opacity pattern entirely: \texttt{scene.remove}
  the hidden child instead, and re-add it on \texttt{detach}.
\end{enumerate}


Option (2) is the cleaner long-term design but breaks any code that
expects the hidden joint dot to remain in \texttt{fig.dots} for
geometry queries (\texttt{lshoulder}-style lookups, the prop
attachment machinery in \texttt{\_drag\_attached\_props}, etc.).
Option (1) is what v0.9.17 chose.

\subsection{8.5 Bubble world-space
anchoring}\label{bubble-world-space-anchoring}

Bubbles attach to a fixed world-space position at creation, then
\textbf{do not track} the speaker or prop afterward. Implications:

\begin{itemize}
\tightlist
\item
  \texttt{say} with \texttt{persist:\ true} followed by
  \texttt{walk\_to} leaves the bubble behind at the character's starting
  position.
\item
  \texttt{prop\_say} followed by moving the prop leaves the bubble
  behind.
\end{itemize}

This is why walk-and-talk dialogue cannot use \texttt{persist} --- the
bubble would visibly detach. The v1.0.0 bubble-lifecycle migration
(figures own their bubbles, see \texttt{BACK\_BURNER.md}) is intended to
fix this by making the bubble follow the head joint per-frame.

For walk-and-talk today, break dialogue into short non-persisted
\texttt{say} beats per walk step, or design the dialogue around standing
rest points (treadmill pattern: walk-and-talk -\textgreater{} bench rest
-\textgreater{} group-translate reset).

\begin{center}\rule{0.5\linewidth}{0.5pt}\end{center}

\section{9. Internals quick reference}\label{internals-quick-reference}

\subsection{9.1 Module-level constants}\label{module-level-constants}

\begin{longtable}[]{@{}ll@{}}
\toprule\noalign{}
Constant & Purpose \\
\midrule\noalign{}
\endhead
\bottomrule\noalign{}
\endlastfoot
\texttt{\_FRAMING\_CAMERA} & framing -\textgreater{} (frame\_width,
frame\_center\_y) \\
\texttt{\_MOVE\_RT} & move -\textgreater{} animation duration in
seconds \\
\texttt{\_WIDE\_W} & 14.2 --- full-stage frame width \\
\texttt{\_WIDE\_Y} & -0.5 --- default frame center y \\
\texttt{\_DEFAULT} & sentinel for single-character mode \\
\texttt{BG\_COLOR} & scene background \\
\texttt{LABEL\_COLOR} & default character label color \\
\texttt{PADDING\_WAIT} & post-bubble pause (configurable) \\
\texttt{\_LOCO} & set of locomotion action names \\
\end{longtable}

\subsection{9.2 Module-level helpers}\label{module-level-helpers}

\begin{longtable}[]{@{}
  >{\raggedright\arraybackslash}p{(\columnwidth - 2\tabcolsep) * \real{0.5000}}
  >{\raggedright\arraybackslash}p{(\columnwidth - 2\tabcolsep) * \real{0.5000}}@{}}
\toprule\noalign{}
\begin{minipage}[b]{\linewidth}\raggedright
Helper
\end{minipage} & \begin{minipage}[b]{\linewidth}\raggedright
Purpose
\end{minipage} \\
\midrule\noalign{}
\endhead
\bottomrule\noalign{}
\endlastfoot
\texttt{\_resolve\_pose(name,\ default,\ fig=None)} & Look up a pose by
name in \texttt{POSES} \\
\texttt{\_targets(step)} & Resolve \texttt{who} to a list of character
names \\
\texttt{\_get\_fig(name)} & Look up a live figure in the cast \\
\texttt{\_get\_prop(pname)} & Look up a prop in the registry \\
\texttt{\_apply\_camera(meta,\ scene,\ char\_x\_positions)} & Instant or
animated reframe \\
\texttt{\_camera\_anim(meta,\ scene,\ char\_x\_positions)} & Returns an
animation (does not play) \\
\texttt{\_execute\_pan\_up(meta,\ scene,\ props,\ char\_x\_positions)} &
Tilt-up implementation \\
\texttt{\_sweep\_expired\_bubbles()} & Bubble bookkeeping (instance
method) \\
\end{longtable}

\subsection{9.3 Environment-variable
summary}\label{environment-variable-summary}

\begin{longtable}[]{@{}
  >{\raggedright\arraybackslash}p{(\columnwidth - 4\tabcolsep) * \real{0.3333}}
  >{\raggedright\arraybackslash}p{(\columnwidth - 4\tabcolsep) * \real{0.3333}}
  >{\raggedright\arraybackslash}p{(\columnwidth - 4\tabcolsep) * \real{0.3333}}@{}}
\toprule\noalign{}
\begin{minipage}[b]{\linewidth}\raggedright
Variable
\end{minipage} & \begin{minipage}[b]{\linewidth}\raggedright
Required?
\end{minipage} & \begin{minipage}[b]{\linewidth}\raggedright
Effect
\end{minipage} \\
\midrule\noalign{}
\endhead
\bottomrule\noalign{}
\endlastfoot
\texttt{PAM\_SCRIPT} & yes & Path to JSON screenplay \\
\texttt{PAM\_CAMERA\_MODE} & no & \texttt{1} enables camera system;
otherwise subscene markers are skipped \\
\texttt{PAM\_PROMPTS} & no & Path to prompts JSON; default
\texttt{\textless{}PAM\_SCRIPT\ stem\textgreater{}\_prompts.json} \\
\texttt{PAM\_SHOW\_CLOCK} & no & \texttt{1} displays elapsed-time
overlay \\
\end{longtable}

\subsection{9.4 Output paths}\label{output-paths}

Manim writes by default to:

\begin{verbatim}
./media/videos/pam_player/<quality>/<SceneName>.mp4
\end{verbatim}

\ldots where \texttt{\textless{}quality\textgreater{}} is
\texttt{480p15}, \texttt{720p30}, \texttt{1080p60}, or \texttt{2160p60}
depending on the \texttt{-q} flag, and
\texttt{\textless{}SceneName\textgreater{}} is \texttt{PAMPlayer} unless
overridden with \texttt{-o}.

For convenience, the \texttt{pam-render} wrapper accepts an output name
as its second argument and stages the final MP4 alongside the script.

\begin{center}\rule{0.5\linewidth}{0.5pt}\end{center}

\section{10. Runtime additions in the attached v0.9.14--v0.9.X
sources}\label{runtime-additions-in-the-attached-v0.9.14v0.9.x-sources}

This section records runtime behavior present in the attached source
files but not fully covered by the earlier \texttt{pam\_player.py}
reference. It is intentionally practical: each entry gives the
screenplay syntax, the registry affected, and the failure mode to expect
while debugging.

\subsection{\texorpdfstring{10.1 Scene-object teardown:
\texttt{remove\_scene\_object} and
\texttt{clear\_all\_scene\_objects}}{10.1 Scene-object teardown: remove\_scene\_object and clear\_all\_scene\_objects}}\label{scene-object-teardown-remove_scene_object-and-clear_all_scene_objects}

\texttt{scene\_objects} are large, camera-addressable background objects
stored in the player's private \texttt{\_scene\_objects} registry. They
are deliberately separate from interactive props in the
\texttt{PropRegistry}; this avoids name collisions and keeps stage
dressing from being treated as grabbable, parentable, or speech-capable
props.

Two teardown verbs complete that separation:

\begin{longtable}[]{@{}
  >{\raggedright\arraybackslash}p{(\columnwidth - 4\tabcolsep) * \real{0.3333}}
  >{\raggedright\arraybackslash}p{(\columnwidth - 4\tabcolsep) * \real{0.3333}}
  >{\raggedright\arraybackslash}p{(\columnwidth - 4\tabcolsep) * \real{0.3333}}@{}}
\toprule\noalign{}
\begin{minipage}[b]{\linewidth}\raggedright
Action
\end{minipage} & \begin{minipage}[b]{\linewidth}\raggedright
Target registry
\end{minipage} & \begin{minipage}[b]{\linewidth}\raggedright
Purpose
\end{minipage} \\
\midrule\noalign{}
\endhead
\bottomrule\noalign{}
\endlastfoot
\texttt{remove\_scene\_object} & \texttt{\_scene\_objects} only & Fade
out one named scene object. \\
\texttt{clear\_all\_scene\_objects} & \texttt{\_scene\_objects} only &
Fade out every registered scene object concurrently. \\
\end{longtable}

\texttt{remove\_scene\_object} uses the sub-key \texttt{prop} for
symmetry with \texttt{remove\_prop}, but it does \textbf{not} search the
prop registry. If the named object is not in \texttt{\_scene\_objects},
the action is a silent no-op. It does not remove characters, dogs,
Governor figures, or interactive props.

\begin{Shaded}
\begin{Highlighting}[]
\OtherTok{[}
  \FunctionTok{\{}\DataTypeTok{"action"}\FunctionTok{:} \StringTok{"scene\_objects"}\FunctionTok{,} \DataTypeTok{"items"}\FunctionTok{:} \FunctionTok{\{}
    \DataTypeTok{"warehouse{-}wall"}\FunctionTok{:} \FunctionTok{\{}\DataTypeTok{"type"}\FunctionTok{:} \StringTok{"wall"}\FunctionTok{,} \DataTypeTok{"x"}\FunctionTok{:} \DecValTok{0}\FunctionTok{,} \DataTypeTok{"y"}\FunctionTok{:} \DecValTok{0}\FunctionTok{\},}
    \DataTypeTok{"skyline"}\FunctionTok{:}        \FunctionTok{\{}\DataTypeTok{"type"}\FunctionTok{:} \StringTok{"backdrop"}\FunctionTok{,} \DataTypeTok{"x"}\FunctionTok{:} \DecValTok{0}\FunctionTok{,} \DataTypeTok{"y"}\FunctionTok{:} \FloatTok{1.5}\FunctionTok{\}}
  \FunctionTok{\}\}}\OtherTok{,}

  \FunctionTok{\{}\DataTypeTok{"action"}\FunctionTok{:} \StringTok{"remove\_scene\_object"}\FunctionTok{,} \DataTypeTok{"prop"}\FunctionTok{:} \StringTok{"warehouse{-}wall"}\FunctionTok{,} \DataTypeTok{"rt"}\FunctionTok{:} \DecValTok{0}\ErrorTok{.}\DecValTok{4}\FunctionTok{\}}\OtherTok{,}
  \FunctionTok{\{}\DataTypeTok{"action"}\FunctionTok{:} \StringTok{"clear\_all\_scene\_objects"}\FunctionTok{,} \DataTypeTok{"rt"}\FunctionTok{:} \DecValTok{0}\ErrorTok{.}\DecValTok{5}\FunctionTok{\}}
\OtherTok{]}
\end{Highlighting}
\end{Shaded}

Use \texttt{remove\_prop} for interactive props and
\texttt{remove\_scene\_object} for non-interactive scene dressing. A
common authoring error is to declare a background item under
\texttt{scene\_objects} and later try to delete it with
\texttt{remove\_prop}; the result is that nothing happens, because the
item is not in the prop registry.

\subsection{10.2 Render-time preview
clock}\label{render-time-preview-clock}

The player supports a preview clock overlay controlled by
\texttt{PAM\_SHOW\_CLOCK=1}, or by the
\texttt{pam-render\ -\/-show-clock} wrapper. The overlay appears near
the title bar in the upper-right and shows elapsed scene time as
\texttt{M:SS.ss}. It is meant for low-quality preview renders,
especially when tuning dialogue holds, camera drift, pauses, or
synchronized entrances.

\begin{Shaded}
\begin{Highlighting}[]
\VariableTok{PAM\_SCRIPT}\OperatorTok{=}\NormalTok{my\_scene.json }\VariableTok{PAM\_SHOW\_CLOCK}\OperatorTok{=}\NormalTok{1 }\ExtensionTok{manim} \AttributeTok{{-}pql}\NormalTok{ pam\_player.py PAMPlayer}
\ExtensionTok{./pam{-}render} \AttributeTok{{-}{-}script}\NormalTok{ my\_scene.json }\AttributeTok{{-}{-}show{-}clock}
\end{Highlighting}
\end{Shaded}

Do not use the clock for final renders; it is a timing diagnostic.

\subsection{10.3 Path C speaker resolution and dual
registration}\label{path-c-speaker-resolution-and-dual-registration}

The attached \texttt{actions.py} adds
\texttt{\_resolve\_speaker(step,\ cast,\ props)}, a small but important
target-resolution helper. It reads \texttt{who} first, then
\texttt{prop}, and looks the resulting name up in both registries. The
purpose is to make dog-like figures work consistently whether an action
reaches them through the cast side or the prop side.

The practical rule is:

\begin{longtable}[]{@{}
  >{\raggedright\arraybackslash}p{(\columnwidth - 6\tabcolsep) * \real{0.2143}}
  >{\raggedleft\arraybackslash}p{(\columnwidth - 6\tabcolsep) * \real{0.2857}}
  >{\raggedleft\arraybackslash}p{(\columnwidth - 6\tabcolsep) * \real{0.2857}}
  >{\raggedright\arraybackslash}p{(\columnwidth - 6\tabcolsep) * \real{0.2143}}@{}}
\toprule\noalign{}
\begin{minipage}[b]{\linewidth}\raggedright
Object kind
\end{minipage} & \begin{minipage}[b]{\linewidth}\raggedleft
Cast registry
\end{minipage} & \begin{minipage}[b]{\linewidth}\raggedleft
Prop registry
\end{minipage} & \begin{minipage}[b]{\linewidth}\raggedright
Notes
\end{minipage} \\
\midrule\noalign{}
\endhead
\bottomrule\noalign{}
\endlastfoot
Human / alien character & yes & no & Normal biped figure. \\
Plain prop & no & yes & Desk, phone, door, button panel, etc. \\
DogGraph / quadruped figure & yes & yes & Dual-registered under the same
key. \\
\end{longtable}

For authoring, this means a dog can be addressed consistently by
\texttt{who} in character-like actions and by \texttt{prop} in prop-like
or speech actions:

\begin{Shaded}
\begin{Highlighting}[]
\OtherTok{[}
  \FunctionTok{\{}\DataTypeTok{"action"}\FunctionTok{:} \StringTok{"fade\_in"}\FunctionTok{,}  \DataTypeTok{"who"}\FunctionTok{:} \StringTok{"chekov"}\FunctionTok{\}}\OtherTok{,}
  \FunctionTok{\{}\DataTypeTok{"action"}\FunctionTok{:} \StringTok{"trot\_to"}\FunctionTok{,}  \DataTypeTok{"who"}\FunctionTok{:} \StringTok{"chekov"}\FunctionTok{,} \DataTypeTok{"x"}\FunctionTok{:} \FloatTok{2.0}\FunctionTok{\}}\OtherTok{,}
  \FunctionTok{\{}\DataTypeTok{"action"}\FunctionTok{:} \StringTok{"prop\_say"}\FunctionTok{,} \DataTypeTok{"prop"}\FunctionTok{:} \StringTok{"chekov"}\FunctionTok{,} \DataTypeTok{"text"}\FunctionTok{:} \StringTok{"Woof."}\FunctionTok{\}}\OtherTok{,}
  \FunctionTok{\{}\DataTypeTok{"action"}\FunctionTok{:} \StringTok{"react"}\FunctionTok{,}    \DataTypeTok{"who"}\FunctionTok{:} \StringTok{"chekov"}\FunctionTok{\}}
\OtherTok{]}
\end{Highlighting}
\end{Shaded}

The helper is deliberately quiet on misses. Individual action handlers
decide whether a missing target is an error, a warning, or a no-op.

\subsection{\texorpdfstring{10.4 Explicit bubble colors on
\texttt{prop\_say} and figure
speech}{10.4 Explicit bubble colors on prop\_say and figure speech}}\label{explicit-bubble-colors-on-prop_say-and-figure-speech}

The attached player and figure code allow speech bubbles to be colored
explicitly. For ordinary character \texttt{say}, the figure style may
define \texttt{bubble\_color} and \texttt{bubble\_text\_color}. For
prop-style speech, the step can supply explicit colors.

\begin{Shaded}
\begin{Highlighting}[]
\FunctionTok{\{}
  \DataTypeTok{"action"}\FunctionTok{:} \StringTok{"cast"}\FunctionTok{,}
  \DataTypeTok{"characters"}\FunctionTok{:} \FunctionTok{\{}
    \DataTypeTok{"nona"}\FunctionTok{:} \FunctionTok{\{}
      \DataTypeTok{"figure\_type"}\FunctionTok{:} \StringTok{"alien"}\FunctionTok{,}
      \DataTypeTok{"style"}\FunctionTok{:} \FunctionTok{\{}
        \DataTypeTok{"edge\_color"}\FunctionTok{:} \StringTok{"\#8a6030"}\FunctionTok{,}
        \DataTypeTok{"bubble\_color"}\FunctionTok{:} \StringTok{"\#8a6030"}\FunctionTok{,}
        \DataTypeTok{"bubble\_text\_color"}\FunctionTok{:} \StringTok{"\#201408"}
      \FunctionTok{\}}
    \FunctionTok{\}}
  \FunctionTok{\}}
\FunctionTok{\}}
\end{Highlighting}
\end{Shaded}

\begin{Shaded}
\begin{Highlighting}[]
\FunctionTok{\{}\DataTypeTok{"action"}\FunctionTok{:} \StringTok{"prop\_say"}\FunctionTok{,} \DataTypeTok{"prop"}\FunctionTok{:} \StringTok{"governor"}\FunctionTok{,}
 \DataTypeTok{"text"}\FunctionTok{:} \StringTok{"Processing."}\FunctionTok{,}
 \DataTypeTok{"bubble\_color"}\FunctionTok{:} \StringTok{"\#382050"}\FunctionTok{,}
 \DataTypeTok{"text\_color"}\FunctionTok{:} \StringTok{"\#ffffff"}\FunctionTok{,}
 \DataTypeTok{"border\_color"}\FunctionTok{:} \StringTok{"\#c8a020"}\FunctionTok{\}}
\end{Highlighting}
\end{Shaded}

If no explicit speech-bubble color is provided, character speech falls
back to the figure's \texttt{edge\_color}. The human/alien
\texttt{say()} implementation uses the bubble color as a saturated
border and derives a pale tinted-white fill so that dark text remains
legible.


\emph{End of pam\_player-manual.md v0.9.13 draft. Companion to
\texttt{CHARACTER\_PROPERTIES.md} and \texttt{PROP\_PROPERTIES.md}.
Revision pass expected after the v1.0.0 cleanup of
\texttt{to\_x}/\texttt{x} normalization, \texttt{group\_translate}
parallel-safety, and the bubble-lifecycle migration.}

\chapter{Character Properties
Reference}\label{character-properties-reference}

\section{Overview}\label{overview-1}

\textbf{Version basis:} PAM v0.9.15.

A complete reference for everything that can be attached to or done with
a character in PAM. Use this document as the single source of truth for
character-side authoring; the main README links here rather than
duplicating these details.

The reference is organized by \textbf{lifecycle}:

\begin{enumerate}
\def\labelenumi{\arabic{enumi}.}
\tightlist
\item
  \textbf{Creation-time properties} --- fields on a cast-block entry.
  What a character \emph{is}.
\item
  \textbf{Modifier actions} --- change character state without
  locomotion or dialogue.
\item
  \textbf{Action targets} --- verbs that take a character as primary
  target.
\item
  \textbf{Meta} --- name resolution, the \texttt{hidden} flag, z-order,
  the color authority chain, registry behavior.
\item
  \textbf{Face \& Appearance} --- face data, expression variants,
  \texttt{attach\_face}, hats, layering, and face debugging.
\end{enumerate}

This document is intended to evolve into the v1.0.0 frozen reference.
Where current behavior is provisional or under active design, the
relevant subsection is flagged.

\begin{center}\rule{0.5\linewidth}{0.5pt}\end{center}

\section{Character Properties
Sections}\label{character-properties-sections}

\subsection{1. Creation-time properties}\label{creation-time-properties}

Characters are declared in a \texttt{cast} action block, typically at
the top of a screenplay:

\begin{Shaded}
\begin{Highlighting}[]
\FunctionTok{\{}
  \DataTypeTok{"action"}\FunctionTok{:} \StringTok{"cast"}\FunctionTok{,}
  \DataTypeTok{"characters"}\FunctionTok{:} \FunctionTok{\{}
    \DataTypeTok{"bevers"}\FunctionTok{:} \FunctionTok{\{}
      \DataTypeTok{"figure\_type"}\FunctionTok{:} \StringTok{"alien"}\FunctionTok{,}
      \DataTypeTok{"build"}\FunctionTok{:} \StringTok{"alien"}\FunctionTok{,}
      \DataTypeTok{"gender"}\FunctionTok{:} \StringTok{"male"}\FunctionTok{,}
      \DataTypeTok{"pose"}\FunctionTok{:} \StringTok{"standing\_front"}\FunctionTok{,}
      \DataTypeTok{"offset"}\FunctionTok{:} \OtherTok{[}\DecValTok{3}\OtherTok{,} \DecValTok{0}\OtherTok{,} \DecValTok{0}\OtherTok{]}\FunctionTok{,}
      \DataTypeTok{"scale"}\FunctionTok{:} \FunctionTok{\{}\DataTypeTok{"sy"}\FunctionTok{:} \DecValTok{0}\ErrorTok{.}\DecValTok{7}\FunctionTok{,} \DataTypeTok{"sx"}\FunctionTok{:} \DecValTok{0}\ErrorTok{.}\DecValTok{7}\FunctionTok{,} \DataTypeTok{"anchor"}\FunctionTok{:} \StringTok{"lankle"}\FunctionTok{\},}
      \DataTypeTok{"color"}\FunctionTok{:} \StringTok{"\#44bb66"}\FunctionTok{,}
      \DataTypeTok{"torso\_color"}\FunctionTok{:} \StringTok{"\#2a6090"}\FunctionTok{,}
      \DataTypeTok{"style"}\FunctionTok{:} \FunctionTok{\{}\DataTypeTok{"head\_label"}\FunctionTok{:} \StringTok{"Bevers"}\FunctionTok{\},}
      \DataTypeTok{"tic\_profile"}\FunctionTok{:} \OtherTok{[}\FunctionTok{\{}\DataTypeTok{"fragment"}\FunctionTok{:} \StringTok{"hand\_twitch\_r"}\FunctionTok{,} \DataTypeTok{"triggers"}\FunctionTok{:} \OtherTok{[}\StringTok{"react"}\OtherTok{]}\FunctionTok{\}}\OtherTok{]}
    \FunctionTok{\}}
  \FunctionTok{\}}
\FunctionTok{\}}
\end{Highlighting}
\end{Shaded}

Every key inside a character entry except the cast-key itself
(\texttt{"bevers"} above) is optional in principle, but
\texttt{figure\_type} should always be set explicitly to avoid the
default-human fallback misclassifying alien or quadruped characters.

The same fields are also accepted in a single-character
\texttt{fade\_in} step, which spawns a character mid-scene without a
prior \texttt{cast} entry:

\begin{Shaded}
\begin{Highlighting}[]
\FunctionTok{\{}\DataTypeTok{"action"}\FunctionTok{:} \StringTok{"fade\_in"}\FunctionTok{,} \DataTypeTok{"who"}\FunctionTok{:} \StringTok{"athena"}\FunctionTok{,}
 \DataTypeTok{"figure\_type"}\FunctionTok{:} \StringTok{"human"}\FunctionTok{,} \DataTypeTok{"build"}\FunctionTok{:} \StringTok{"narrow"}\FunctionTok{,} \DataTypeTok{"gender"}\FunctionTok{:} \StringTok{"female"}\FunctionTok{,}
 \DataTypeTok{"offset"}\FunctionTok{:} \OtherTok{[}\DecValTok{{-}2}\OtherTok{,} \DecValTok{0}\OtherTok{,} \DecValTok{0}\OtherTok{]}\FunctionTok{\}}
\end{Highlighting}
\end{Shaded}

For multi-scene projects, store canonical character entries once in
\texttt{tntd\_characters.json} (or your project's equivalent) and let
\texttt{fountain2pam.py} merge them into each scene's cast block at
conversion time. See
\S \cite{field-precedence-and-the-canonical-registry}.%{§1.13}.

\begin{center}\rule{0.5\linewidth}{0.5pt}\end{center}

\subsubsection{\texorpdfstring{1.1
\texttt{figure\_type}}{1.1 figure\_type}}\label{figure_type}

\emph{Type:} \texttt{str} · \emph{Default:} \texttt{"human"}

The figure class used to render the character. Determines skeleton
topology, locomotion vocabulary, and which actions are meaningful.

\textbf{Accepted values:}

\begin{itemize}
\tightlist
\item
  \texttt{"human"} --- \texttt{HumanGraph}. 15-vertex biped skeleton.
  Full action vocabulary including walk, run, gesture, restraint, prop
  interaction.
\item
  \texttt{"alien"} --- \texttt{AlienGraph}. Subclass of
  \texttt{HumanGraph} with a split-torso joint structure (extra
  \texttt{torso\_left}/\texttt{torso\_right} joints). All human actions
  work; in addition, alien-specific poses live in the build's pose
  registry. Use \texttt{build="alien"} or \texttt{build="alien\_female"}
  to match.
\item
  \texttt{"dog"} --- \texttt{DogGraph}. 19-joint quadruped, side-view
  default. Locomotion is \texttt{trot\_to} (not \texttt{walk\_to} /
  \texttt{run\_to}); \texttt{say} is a no-op for dogs --- use
  \texttt{prop\_say} keyed on the dog's name. Path C: dual-registered as
  both a cast member and a prop, so \texttt{who:\ chekov} and
  \texttt{prop:\ chekov} both resolve.
\item
  \texttt{"dodecahedron"} --- \texttt{GovernorGraph}.
  Procedurally-animated dodecahedron (Schlegel diagram or 3D spinning,
  depending on \texttt{style}). Used for AI characters like the
  Governor. Rotates autonomously; \texttt{say} is a no-op --- use
  \texttt{prop\_say}.
\end{itemize}

\textbf{Default behavior if omitted:} treated as \texttt{"human"}. This
is a known footgun for alien characters --- always set
\texttt{figure\_type} explicitly when declaring an alien.

\textbf{Example:}

\begin{Shaded}
\begin{Highlighting}[]
\ErrorTok{"chekov":}   \FunctionTok{\{}\DataTypeTok{"figure\_type"}\FunctionTok{:} \StringTok{"dog"}\FunctionTok{\}}\ErrorTok{,}
\ErrorTok{"governor":} \FunctionTok{\{}\DataTypeTok{"figure\_type"}\FunctionTok{:} \StringTok{"dodecahedron"}\FunctionTok{\}}\ErrorTok{,}
\ErrorTok{"thalia":}   \FunctionTok{\{}\DataTypeTok{"figure\_type"}\FunctionTok{:} \StringTok{"alien"}\FunctionTok{,} \DataTypeTok{"build"}\FunctionTok{:} \StringTok{"alien\_female"}\FunctionTok{,} \DataTypeTok{"gender"}\FunctionTok{:} \StringTok{"female"}\FunctionTok{\}}
\end{Highlighting}
\end{Shaded}

\begin{center}\rule{0.5\linewidth}{0.5pt}\end{center}

\subsubsection{\texorpdfstring{1.2
\texttt{build}}{1.2 build}}\label{build}

\emph{Type:} \texttt{str} · \emph{Default:} derived from \texttt{gender}

The proportions preset. Controls skeleton dimensions (shoulder width,
leg length, torso height) and certain build-specific poses. Builds are
defined in \texttt{pam/builds.py}.

\textbf{Accepted values:}

\begin{itemize}
\tightlist
\item
  \texttt{"default"} --- neutral male-coded proportions. Default for
  \texttt{gender="male"}.
\item
  \texttt{"narrow"} --- narrower shoulders, slightly shorter. Default
  for \texttt{gender="female"} and \texttt{gender="child"}.
\item
  \texttt{"broad"} --- wider shoulders, heavier frame.
\item
  \texttt{"alien"} --- alien proportions with split torso. Use with
  \texttt{figure\_type="alien"}, \texttt{gender="male"}.
\item
  \texttt{"alien\_female"} --- alien proportions, narrower variant. Use
  with \texttt{figure\_type="alien"}, \texttt{gender="female"}.
\end{itemize}

\textbf{Interaction with \texttt{gender}:} if \texttt{build} is omitted,
the value is inferred from \texttt{gender} via \texttt{GENDER\_DEFAULTS}
(in \texttt{figure.py}). Setting \texttt{build} explicitly always wins.

\textbf{For aliens:} the human builds (\texttt{default},
\texttt{narrow}, \texttt{broad}) have no
\texttt{torso\_left}/\texttt{torso\_right} joints, so they cannot render
alien poses. An alien character with \texttt{build="default"} will
render with human proportions but break on alien-specific morphs. Always
pair \texttt{figure\_type="alien"} with \texttt{build="alien"} or
\texttt{build="alien\_female"}.

\textbf{Example:}

\begin{Shaded}
\begin{Highlighting}[]
\ErrorTok{"chava":}  \FunctionTok{\{}\DataTypeTok{"figure\_type"}\FunctionTok{:} \StringTok{"human"}\FunctionTok{,} \DataTypeTok{"build"}\FunctionTok{:} \StringTok{"narrow"}\FunctionTok{,} \DataTypeTok{"gender"}\FunctionTok{:} \StringTok{"female"}\FunctionTok{\}}\ErrorTok{,}
\ErrorTok{"factor":} \FunctionTok{\{}\DataTypeTok{"figure\_type"}\FunctionTok{:} \StringTok{"alien"}\FunctionTok{,} \DataTypeTok{"build"}\FunctionTok{:} \StringTok{"alien"}\FunctionTok{,}  \DataTypeTok{"gender"}\FunctionTok{:} \StringTok{"male"}\FunctionTok{\}}\ErrorTok{,}
\ErrorTok{"thalia":} \FunctionTok{\{}\DataTypeTok{"figure\_type"}\FunctionTok{:} \StringTok{"alien"}\FunctionTok{,} \DataTypeTok{"build"}\FunctionTok{:} \StringTok{"alien\_female"}\FunctionTok{,} \DataTypeTok{"gender"}\FunctionTok{:} \StringTok{"female"}\FunctionTok{\}}
\end{Highlighting}
\end{Shaded}

\begin{center}\rule{0.5\linewidth}{0.5pt}\end{center}

\subsubsection{\texorpdfstring{1.3
\texttt{gender}}{1.3 gender}}\label{gender}

\emph{Type:} \texttt{str} · \emph{Default:} \texttt{"male"}

Drives the build inference (see \cite{build})
%{§1.2})
and is
consumed by downstream tooling (e.g.~voice casting in any future
text-to-speech integration). Does not directly affect rendering unless
\texttt{build} is also omitted.

\textbf{Accepted values:} \texttt{"male"}, \texttt{"female"},
\texttt{"child"}.

\textbf{Default warning:} the converter prints a warning if
\texttt{gender} is missing and defaults to \texttt{"male"}. Always set
explicitly.

\textbf{Example:}

\begin{Shaded}
\begin{Highlighting}[]
\ErrorTok{"bart":} \FunctionTok{\{}\DataTypeTok{"figure\_type"}\FunctionTok{:} \StringTok{"human"}\FunctionTok{,} \DataTypeTok{"gender"}\FunctionTok{:} \StringTok{"child"}\FunctionTok{\}}
\end{Highlighting}
\end{Shaded}

\begin{center}\rule{0.5\linewidth}{0.5pt}\end{center}

\subsubsection{\texorpdfstring{1.4
\texttt{color}}{1.4 color}}\label{color}

\emph{Type:} \texttt{str} (hex string) · \emph{Default:} build's
\texttt{edge\_color}

Single-color shorthand for the character's edge/stroke color. Sets
\texttt{edge\_color} on the style palette; the remaining palette keys
(\texttt{node\_color}, \texttt{node\_stroke}, \texttt{head\_color},
\texttt{head\_stroke}, \texttt{highlight\_color}) are derived from
\texttt{color} via \texttt{\_style\_from\_color} in \texttt{figure.py}.

\textbf{Interaction with \texttt{style}:} explicit values inside
\texttt{style} override anything derived from \texttt{color}. Common
pattern: set \texttt{color} and \texttt{torso\_color} at the top level,
then add \texttt{head\_label} (and only \texttt{head\_label}) in
\texttt{style}.

\textbf{Authority:} when \texttt{tntd\_characters.json} (or another
canonical registry) is in play, the registry's \texttt{color} is
authoritative --- \texttt{fountain2pam.py} re-derives the full palette
from the canonical \texttt{color} and overwrites any per-scene values
during conversion. See
\hyperref[113-field-precedence-and-the-canonical-registry]{§1.13}.

\textbf{Example:}

\begin{Shaded}
\begin{Highlighting}[]
\ErrorTok{"bevers":} \FunctionTok{\{}
  \DataTypeTok{"color"}\FunctionTok{:} \StringTok{"\#44bb66"}\FunctionTok{,}
  \DataTypeTok{"torso\_color"}\FunctionTok{:} \StringTok{"\#2a6090"}\FunctionTok{,}
  \DataTypeTok{"style"}\FunctionTok{:} \FunctionTok{\{}\DataTypeTok{"head\_label"}\FunctionTok{:} \StringTok{"Bevers"}\FunctionTok{\}}
\FunctionTok{\}}
\end{Highlighting}
\end{Shaded}

\begin{center}\rule{0.5\linewidth}{0.5pt}\end{center}

\subsubsection{\texorpdfstring{1.5
\texttt{torso\_color}}{1.5 torso\_color}}\label{torso_color}

\emph{Type:} \texttt{str} (hex string) · \emph{Default:} none
(single-zone rendering)

Two-zone color (v0.9.4). When present, the torso edges
(\texttt{lshoulder}--\texttt{torso}, \texttt{rshoulder}--\texttt{torso},
\texttt{lhip}--\texttt{torso}, \texttt{rhip}--\texttt{torso}) and torso
fill are drawn in \texttt{torso\_color} instead of the main
\texttt{edge\_color}. Reads as a uniform, shirt, or sash.

\textbf{Omit for:} single-color rendering --- character drawn entirely
in \texttt{edge\_color} / derived palette.

\textbf{Example:}

\begin{Shaded}
\begin{Highlighting}[]
\ErrorTok{"chava":}  \FunctionTok{\{}\DataTypeTok{"color"}\FunctionTok{:} \StringTok{"\#e8943a"}\FunctionTok{,} \DataTypeTok{"torso\_color"}\FunctionTok{:} \StringTok{"\#7a4010"}\FunctionTok{\}}\ErrorTok{,}
\ErrorTok{"thalia":} \FunctionTok{\{}\DataTypeTok{"color"}\FunctionTok{:} \StringTok{"\#e055aa"}\FunctionTok{,} \DataTypeTok{"torso\_color"}\FunctionTok{:} \StringTok{"\#2a6090"}\FunctionTok{\}}
\end{Highlighting}
\end{Shaded}

Pair with \texttt{uniforms} (\cite{uniforms}) for named
torso variants that swap via the \texttt{change\_uniform} action.

\begin{center}\rule{0.5\linewidth}{0.5pt}\end{center}

\subsubsection{\texorpdfstring{1.6
\texttt{style}}{1.6 style}}\label{style}

\emph{Type:} \texttt{dict} · \emph{Default:} \texttt{DEFAULT\_STYLE} (or
derived from \texttt{color})

Explicit palette and label overrides. Every key inside \texttt{style} is
optional; missing keys fall through to the derived palette.

\textbf{Sub-keys:}

\begin{longtable}[]{@{}
  >{\raggedright\arraybackslash}p{(\columnwidth - 6\tabcolsep) * \real{0.2500}}
  >{\raggedright\arraybackslash}p{(\columnwidth - 6\tabcolsep) * \real{0.2500}}
  >{\raggedright\arraybackslash}p{(\columnwidth - 6\tabcolsep) * \real{0.2500}}
  >{\raggedright\arraybackslash}p{(\columnwidth - 6\tabcolsep) * \real{0.2500}}@{}}
\toprule\noalign{}
\begin{minipage}[b]{\linewidth}\raggedright
Key
\end{minipage} & \begin{minipage}[b]{\linewidth}\raggedright
Type
\end{minipage} & \begin{minipage}[b]{\linewidth}\raggedright
Default
\end{minipage} & \begin{minipage}[b]{\linewidth}\raggedright
Description
\end{minipage} \\
\midrule\noalign{}
\endhead
\bottomrule\noalign{}
\endlastfoot
\texttt{edge\_color} & \texttt{str} (hex) & \texttt{"\#3a7bd5"} &
Bone/edge color. Overridden by top-level \texttt{color} if present, then
by \texttt{style.edge\_color} if present (style wins). \\
\texttt{node\_color} & \texttt{str} (hex) & \texttt{"\#1e3a5f"} & Joint
fill color. \\
\texttt{node\_stroke} & \texttt{str} (hex) & \texttt{"\#5b9cf6"} & Joint
outline color. \\
\texttt{head\_color} & \texttt{str} (hex) & \texttt{"\#0d2340"} & Head
circle fill. \\
\texttt{head\_stroke} & \texttt{str} (hex) & \texttt{"\#7ec8ff"} & Head
circle outline. \\
\texttt{head\_label} & \texttt{str} & \texttt{"v\_0"} & Text inside the
head circle. Multi-character labels work but are visually cramped;
one-to-two characters reads cleanest at typical scales. \\
\texttt{head\_font} & \texttt{str} & \texttt{"Courier\ New"} & Font
family for the head label \emph{and} speech bubbles. \\
\texttt{head\_font\_sz} & \texttt{int} & \texttt{14} & Font size for the
head label. \\
\texttt{head\_radius} & \texttt{float} & \texttt{0.28} & Head circle
radius in world units. \texttt{DogGraph} defaults to \texttt{0.22}. \\
\texttt{node\_radius} & \texttt{float} & \texttt{0.14} & Joint circle
radius. \texttt{DogGraph} defaults to \texttt{0.11}. \\
\texttt{edge\_width} & \texttt{float} & \texttt{2.5} & Stroke width for
edges. \\
\texttt{highlight\_color} & \texttt{str} (hex) & \texttt{"\#7ec8ff"} &
Color used by \texttt{highlight\_edges()} and the v0.9.15
\texttt{attach\_face} halo (not implemented in current MVP). \\
\end{longtable}

\textbf{DogGraph-additional sub-keys:}

\begin{longtable}[]{@{}
  >{\raggedright\arraybackslash}p{(\columnwidth - 6\tabcolsep) * \real{0.2500}}
  >{\raggedright\arraybackslash}p{(\columnwidth - 6\tabcolsep) * \real{0.2500}}
  >{\raggedright\arraybackslash}p{(\columnwidth - 6\tabcolsep) * \real{0.2500}}
  >{\raggedright\arraybackslash}p{(\columnwidth - 6\tabcolsep) * \real{0.2500}}@{}}
\toprule\noalign{}
\begin{minipage}[b]{\linewidth}\raggedright
Key
\end{minipage} & \begin{minipage}[b]{\linewidth}\raggedright
Type
\end{minipage} & \begin{minipage}[b]{\linewidth}\raggedright
Default
\end{minipage} & \begin{minipage}[b]{\linewidth}\raggedright
Description
\end{minipage} \\
\midrule\noalign{}
\endhead
\bottomrule\noalign{}
\endlastfoot
\texttt{far\_edge\_color} & \texttt{str} (hex) & matches
\texttt{edge\_color} & Far-side legs (perspective cue). \\
\texttt{far\_edge\_opacity} & \texttt{float} & \texttt{0.35} & Opacity
for far-side legs. \\
\texttt{far\_node\_opacity} & \texttt{float} & \texttt{0.30} & Opacity
for far-side joints. \\
\texttt{far\_edge\_width} & \texttt{float} & \texttt{1.5} & Stroke width
for far-side legs. \\
\end{longtable}

\textbf{Example --- full palette:}

\begin{Shaded}
\begin{Highlighting}[]
\ErrorTok{"bevers":} \FunctionTok{\{}
  \DataTypeTok{"style"}\FunctionTok{:} \FunctionTok{\{}
    \DataTypeTok{"head\_label"}\FunctionTok{:} \StringTok{"Bevers"}\FunctionTok{,}
    \DataTypeTok{"edge\_color"}\FunctionTok{:} \StringTok{"\#44bb66"}\FunctionTok{,}
    \DataTypeTok{"node\_color"}\FunctionTok{:} \StringTok{"\#2a6090"}\FunctionTok{,}
    \DataTypeTok{"node\_stroke"}\FunctionTok{:} \StringTok{"\#6ce38e"}\FunctionTok{,}
    \DataTypeTok{"head\_color"}\FunctionTok{:} \StringTok{"\#08170c"}\FunctionTok{,}
    \DataTypeTok{"head\_stroke"}\FunctionTok{:} \StringTok{"\#80f7a2"}\FunctionTok{,}
    \DataTypeTok{"highlight\_color"}\FunctionTok{:} \StringTok{"\#94ffb6"}
  \FunctionTok{\}}
\FunctionTok{\}}
\end{Highlighting}
\end{Shaded}

\textbf{Example --- label only (palette derived from \texttt{color}):}

\begin{Shaded}
\begin{Highlighting}[]
\ErrorTok{"chava":} \FunctionTok{\{}
  \DataTypeTok{"color"}\FunctionTok{:} \StringTok{"\#e8943a"}\FunctionTok{,}
  \DataTypeTok{"torso\_color"}\FunctionTok{:} \StringTok{"\#7a4010"}\FunctionTok{,}
  \DataTypeTok{"style"}\FunctionTok{:} \FunctionTok{\{}\DataTypeTok{"head\_label"}\FunctionTok{:} \StringTok{"Chava"}\FunctionTok{\}}
\FunctionTok{\}}
\end{Highlighting}
\end{Shaded}

\begin{center}\rule{0.5\linewidth}{0.5pt}\end{center}

\subsubsection{\texorpdfstring{1.7 \texttt{pose}}{1.7 pose}}\label{pose}

\emph{Type:} \texttt{str} · \emph{Default:} \texttt{"standing\_front"}

The initial pose at fade-in. Must be a key in the \texttt{POSES}
registry (\texttt{pam/poses.py}) or in the character's build-specific
pose dict.

\textbf{Common values:}

\begin{itemize}
\tightlist
\item
  \texttt{"standing\_front"} --- facing camera. Default. Required for
  \texttt{say} to read correctly in scenes where the bubble lands above
  the head.
\item
  \texttt{"standing\_side"} --- side-view. Required for walking,
  running, and any locomotion verb. \texttt{say} works in
  \texttt{standing\_side} without constraint.
\item
  \texttt{"sitting\_mid"}, \texttt{"sitting\_down"} --- mid-sit and
  fully-seated. Used after \texttt{sit\_down}.
\item
  \texttt{"lying\_flat\_right"}, \texttt{"lying\_flat\_left"} ---
  horizontal poses (v0.9.14). Speech bubble anchoring works correctly
  only with these named poses; an ad-hoc rotation will mis-anchor
  bubbles.
\item
  \texttt{"at\_attention"} --- feet together, arms at sides. Formal
  stance.
\item
  \texttt{"look\_up"}, \texttt{"look\_up\_point"} --- head tilted up.
  Used for sky-watching beats.
\item
  Dog: \texttt{"dog\_standing"}, \texttt{"dog\_trot\_a"},
  \texttt{"dog\_trot\_b"}.
\end{itemize}

\textbf{Full pose registry:} see \texttt{POSES} in
\texttt{pam/poses.py}. Categories: standing · grappled · sitting · lying
· wave (right-arm) · lwave (left-arm) · walk · run · carry (chest) ·
side\_carry (briefcase low) · look · attention · reach · rush · squeeze
· fall · dodge · jump · pat · dog.

\textbf{Build awareness:} if the character's build registers its own
poses (e.g.~alien builds with split-torso variants), those override the
global \texttt{POSES} entry of the same name. An alien character in
\texttt{"standing\_front"} gets the alien-skeleton standing pose, not
the human one.

\textbf{Example:}

\begin{Shaded}
\begin{Highlighting}[]
\ErrorTok{"freydoon":} \FunctionTok{\{}\DataTypeTok{"pose"}\FunctionTok{:} \StringTok{"standing\_front"}\FunctionTok{\}}\ErrorTok{,}
\ErrorTok{"sergeant":} \FunctionTok{\{}\DataTypeTok{"pose"}\FunctionTok{:} \StringTok{"at\_attention"}\FunctionTok{\}}\ErrorTok{,}
\ErrorTok{"chekov":}   \FunctionTok{\{}\DataTypeTok{"figure\_type"}\FunctionTok{:} \StringTok{"dog"}\FunctionTok{,} \DataTypeTok{"pose"}\FunctionTok{:} \StringTok{"dog\_standing"}\FunctionTok{\}}
\end{Highlighting}
\end{Shaded}

\begin{center}\rule{0.5\linewidth}{0.5pt}\end{center}

\subsubsection{\texorpdfstring{1.8
\texttt{offset}}{1.8 offset}}\label{offset}

\emph{Type:} \texttt{list} of 2 or 3 floats · \emph{Default:}
\texttt{{[}0,\ 0,\ 0{]}}

World-space position at fade-in, as \texttt{{[}x,\ y,\ z{]}}. The third
component is conventionally \texttt{0} (PAM is 2D); a 2-element list is
also accepted.

\textbf{Coordinate system:}

\begin{itemize}
\tightlist
\item
  \texttt{x} ranges roughly \texttt{-7.1} to \texttt{+7.1} for a
  default-framing camera. Negative is left, positive is right.
\item
  \texttt{y\ =\ 0} is the ground line. Positive \texttt{y} raises the
  character. For a character standing on a raised surface (a step, a
  platform), set \texttt{y} to the surface height.
\item
  \texttt{z} is unused in current rendering but reserved for future
  depth.
\end{itemize}

\textbf{Relation to \texttt{walk\_to} / \texttt{run\_to}:} locomotion
verbs update the figure's \emph{runtime} offset but do not write back to
the cast spec. A character who walks left and then is faded out + faded
back in returns to their original cast-spec \texttt{offset}, not their
walked-to position. Authors who want continuity across fade cycles must
\texttt{walk\_to} the desired x explicitly after fade-in.

\textbf{Example:}

\begin{Shaded}
\begin{Highlighting}[]
\ErrorTok{"bevers":}  \FunctionTok{\{}\DataTypeTok{"offset"}\FunctionTok{:} \OtherTok{[}\DecValTok{3}\OtherTok{,} \DecValTok{0}\OtherTok{,} \DecValTok{0}\OtherTok{]}\FunctionTok{\}}\ErrorTok{,}
\ErrorTok{"freydoon":} \FunctionTok{\{}\DataTypeTok{"offset"}\FunctionTok{:} \OtherTok{[}\FloatTok{{-}5.5}\OtherTok{,} \DecValTok{0}\OtherTok{,} \DecValTok{0}\OtherTok{]}\FunctionTok{\}}\ErrorTok{,}
\ErrorTok{"narrator":} \FunctionTok{\{}\DataTypeTok{"offset"}\FunctionTok{:} \OtherTok{[}\DecValTok{0}\OtherTok{,} \FloatTok{2.5}\OtherTok{,} \DecValTok{0}\OtherTok{]}\FunctionTok{\}}
\end{Highlighting}
\end{Shaded}

\begin{center}\rule{0.5\linewidth}{0.5pt}\end{center}

\subsubsection{\texorpdfstring{1.9
\texttt{scale}}{1.9 scale}}\label{scale}

\emph{Type:} \texttt{dict} with keys \texttt{sy}, \texttt{sx},
\texttt{anchor} · \emph{Default:} \texttt{1.0} uniform

Initial size, expressed as separate y and x scale factors around an
anchor joint. Different from the \texttt{scale} \emph{action} (which
animates a scale change live).

\textbf{Sub-keys:}

\begin{longtable}[]{@{}
  >{\raggedright\arraybackslash}p{(\columnwidth - 4\tabcolsep) * \real{0.3333}}
  >{\raggedright\arraybackslash}p{(\columnwidth - 4\tabcolsep) * \real{0.3333}}
  >{\raggedright\arraybackslash}p{(\columnwidth - 4\tabcolsep) * \real{0.3333}}@{}}
\toprule\noalign{}
\begin{minipage}[b]{\linewidth}\raggedright
Key
\end{minipage} & \begin{minipage}[b]{\linewidth}\raggedright
Type
\end{minipage} & \begin{minipage}[b]{\linewidth}\raggedright
Description
\end{minipage} \\
\midrule\noalign{}
\endhead
\bottomrule\noalign{}
\endlastfoot
\texttt{sy} & \texttt{float} & Vertical scale. \texttt{1.0} = full size;
\texttt{0.7} is canonical for TNTD's reference framing. \\
\texttt{sx} & \texttt{float} & Horizontal scale. Typically equal to
\texttt{sy} for uniform scaling; setting \texttt{sx\ \textless{}\ sy}
produces a thinner silhouette. \\
\texttt{anchor} & \texttt{str} & Joint name that stays fixed during
scaling. \texttt{"lankle"} (left ankle) keeps the character standing on
the ground line as size changes. Other useful anchors: \texttt{"torso"}
(centroid), \texttt{"head"} (for ``looming'' effects). \\
\end{longtable}

\textbf{Shorthand:} a bare float is accepted and expanded to
\texttt{\{"sy":\ v,\ "sx":\ v,\ "anchor":\ "lankle"\}}. This is the form
\texttt{fountain2pam.py} writes when \texttt{scale=0.48} is set in a
Fountain+ \texttt{CHARACTER} annotation.

\textbf{TNTD canonical values:}

\begin{itemize}
\tightlist
\item
  \texttt{0.7} --- adult character, standard framing.
\item
  \texttt{0.65} --- slightly smaller for narrower or background
  characters.
\item
  \texttt{0.48}, \texttt{0.32} --- distant or miniaturized (e.g.~Chava
  observed through a window).
\end{itemize}

\textbf{Node radius and scaling:} prior to v0.9.x patches, joint node
radii were set at construction time and did not scale with
\texttt{sy}/\texttt{sx}, producing oversized nodes on a scaled-down
character. This is fixed in current \texttt{figure.py}; node radius
scales with the character.

\textbf{Example:}

\begin{Shaded}
\begin{Highlighting}[]
\ErrorTok{"athena":}  \FunctionTok{\{}\DataTypeTok{"scale"}\FunctionTok{:} \FunctionTok{\{}\DataTypeTok{"sy"}\FunctionTok{:} \DecValTok{0}\ErrorTok{.}\DecValTok{65}\FunctionTok{,} \DataTypeTok{"sx"}\FunctionTok{:} \DecValTok{0}\ErrorTok{.}\DecValTok{65}\FunctionTok{,} \DataTypeTok{"anchor"}\FunctionTok{:} \StringTok{"lankle"}\FunctionTok{\}\}}\ErrorTok{,}
\ErrorTok{"freydoon":} \FunctionTok{\{}\DataTypeTok{"scale"}\FunctionTok{:} \FunctionTok{\{}\DataTypeTok{"sy"}\FunctionTok{:} \DecValTok{0}\ErrorTok{.}\DecValTok{7}\FunctionTok{,}  \DataTypeTok{"sx"}\FunctionTok{:} \DecValTok{0}\ErrorTok{.}\DecValTok{7}\FunctionTok{,}  \DataTypeTok{"anchor"}\FunctionTok{:} \StringTok{"lankle"}\FunctionTok{\}\}}\ErrorTok{,}
\ErrorTok{"distant\_chava":} \FunctionTok{\{}\DataTypeTok{"scale"}\FunctionTok{:} \FunctionTok{\{}\DataTypeTok{"sy"}\FunctionTok{:} \DecValTok{0}\ErrorTok{.}\DecValTok{32}\FunctionTok{,} \DataTypeTok{"sx"}\FunctionTok{:} \DecValTok{0}\ErrorTok{.}\DecValTok{32}\FunctionTok{,} \DataTypeTok{"anchor"}\FunctionTok{:} \StringTok{"lankle"}\FunctionTok{\}\}}
\end{Highlighting}
\end{Shaded}

\begin{center}\rule{0.5\linewidth}{0.5pt}\end{center}

\subsubsection{\texorpdfstring{1.10 \texttt{facing} (dog
only)}{1.10 facing (dog only)}}\label{facing-dog-only}

\emph{Type:} \texttt{str} · \emph{Default:} \texttt{"right"} ·
\emph{Applies to:} \texttt{figure\_type:\ "dog"}

Direction the dog faces at spawn time. Affects which side of the
skeleton is drawn solid vs.~dashed (near-side legs solid, far-side
dashed).

\textbf{Accepted values:} \texttt{"left"}, \texttt{"right"}.

Has no effect on \texttt{HumanGraph} / \texttt{AlienGraph} /
\texttt{GovernorGraph} --- those use the \texttt{pose} field to express
facing (\texttt{standing\_side} is right-facing by default; use
\texttt{pose="standing\_side"} plus a \texttt{turn} action to face
left).

\textbf{Example:}

\begin{Shaded}
\begin{Highlighting}[]
\ErrorTok{"chekov":} \FunctionTok{\{}\DataTypeTok{"figure\_type"}\FunctionTok{:} \StringTok{"dog"}\FunctionTok{,} \DataTypeTok{"facing"}\FunctionTok{:} \StringTok{"left"}\FunctionTok{,}
           \DataTypeTok{"offset"}\FunctionTok{:} \OtherTok{[}\DecValTok{{-}2}\OtherTok{,} \DecValTok{0}\OtherTok{,} \DecValTok{0}\OtherTok{]}\FunctionTok{\}}
\end{Highlighting}
\end{Shaded}

\begin{center}\rule{0.5\linewidth}{0.5pt}\end{center}

\subsubsection{\texorpdfstring{1.11
\texttt{uniforms}}{1.11 uniforms}}\label{uniforms}

\emph{Type:} \texttt{dict} mapping name -\textgreater{} uniform-spec ·
\emph{Default:} none

Named torso-color variants for the \texttt{change\_uniform} action. Each
uniform may set \texttt{torso\_color} (required), \texttt{color}
(optional, overrides the character's main edge color), and
\texttt{description} (free-form annotation for authoring; ignored at
runtime).

\textbf{Reserved uniform name:} \texttt{"default"} is the uniform
applied at fade-in. If you declare a \texttt{uniforms} dict, also
declare \texttt{"default"} so the character has a baseline to return to.

\textbf{Example:}

\begin{Shaded}
\begin{Highlighting}[]
\ErrorTok{"chava":} \FunctionTok{\{}
  \DataTypeTok{"color"}\FunctionTok{:} \StringTok{"\#e8943a"}\FunctionTok{,}
  \DataTypeTok{"torso\_color"}\FunctionTok{:} \StringTok{"\#7a4010"}\FunctionTok{,}
  \DataTypeTok{"uniforms"}\FunctionTok{:} \FunctionTok{\{}
    \DataTypeTok{"default"}\FunctionTok{:}  \FunctionTok{\{}\DataTypeTok{"torso\_color"}\FunctionTok{:} \StringTok{"\#7a4010"}\FunctionTok{,}
                 \DataTypeTok{"description"}\FunctionTok{:} \StringTok{"Default spy look — dark brown torso"}\FunctionTok{\},}
    \DataTypeTok{"delivery"}\FunctionTok{:} \FunctionTok{\{}\DataTypeTok{"torso\_color"}\FunctionTok{:} \StringTok{"\#6b4226"}\FunctionTok{,} \DataTypeTok{"color"}\FunctionTok{:} \StringTok{"\#c88040"}\FunctionTok{,}
                 \DataTypeTok{"description"}\FunctionTok{:} \StringTok{"Delivery disguise — warm brown uniform"}\FunctionTok{\},}
    \DataTypeTok{"casual"}\FunctionTok{:}   \FunctionTok{\{}\DataTypeTok{"torso\_color"}\FunctionTok{:} \StringTok{"\#e8b830"}\FunctionTok{,}
                 \DataTypeTok{"description"}\FunctionTok{:} \StringTok{"Off{-}duty — bright golden yellow torso"}\FunctionTok{\}}
  \FunctionTok{\}}
\FunctionTok{\}}
\end{Highlighting}
\end{Shaded}

Authors switch uniforms mid-scene with:

\begin{Shaded}
\begin{Highlighting}[]
\FunctionTok{\{}\DataTypeTok{"action"}\FunctionTok{:} \StringTok{"change\_uniform"}\FunctionTok{,} \DataTypeTok{"who"}\FunctionTok{:} \StringTok{"chava"}\FunctionTok{,} \DataTypeTok{"uniform"}\FunctionTok{:} \StringTok{"delivery"}\FunctionTok{\}}
\end{Highlighting}
\end{Shaded}

Full \texttt{change\_uniform} parameters are documented in the section on
modifier actions.

\begin{center}\rule{0.5\linewidth}{0.5pt}\end{center}

\subsubsection{\texorpdfstring{1.12
\texttt{tic\_profile}}{1.12 tic\_profile}}\label{tic_profile}

\emph{Type:} \texttt{list} of fragment-spec dicts · \emph{Default:} none

Speech-tic profile (v0.9.14). A list of small motion fragments that fire
on specified triggers --- a fidget on \texttt{react}, a tail-wag at
\texttt{sentence\_end}, etc.

\textbf{Spec structure:} each entry is a dict with:

\begin{longtable}[]{@{}
  >{\raggedright\arraybackslash}p{(\columnwidth - 4\tabcolsep) * \real{0.3333}}
  >{\raggedright\arraybackslash}p{(\columnwidth - 4\tabcolsep) * \real{0.3333}}
  >{\raggedright\arraybackslash}p{(\columnwidth - 4\tabcolsep) * \real{0.3333}}@{}}
\toprule\noalign{}
\begin{minipage}[b]{\linewidth}\raggedright
Key
\end{minipage} & \begin{minipage}[b]{\linewidth}\raggedright
Type
\end{minipage} & \begin{minipage}[b]{\linewidth}\raggedright
Description
\end{minipage} \\
\midrule\noalign{}
\endhead
\bottomrule\noalign{}
\endlastfoot
\texttt{fragment} & \texttt{str} & Fragment name from
\texttt{TIC\_FRAGMENTS} in \texttt{pam/tics.py}. See the registry table
below for current fragments and their body-type applicability. \\
\texttt{triggers} & \texttt{list{[}str{]}} & Trigger names that fire
this fragment. Built-in triggers: \texttt{"react"} (fired by the
explicit \texttt{react} action), \texttt{"sentence\_end"} (fired
automatically after every \texttt{say} / \texttt{prop\_say}). \\
\end{longtable}

\textbf{TIC\_FRAGMENTS registry:} the full set of fragments currently
defined in \texttt{pam/tics.py}. Each fragment names the joint it
animates, a delta-based cycle, and the figure types it applies to.

\begin{longtable}[]{@{}
  >{\raggedright\arraybackslash}p{(\columnwidth - 6\tabcolsep) * \real{0.2500}}
  >{\raggedright\arraybackslash}p{(\columnwidth - 6\tabcolsep) * \real{0.1700}}
  >{\raggedright\arraybackslash}p{(\columnwidth - 6\tabcolsep) * \real{0.2000}}
  >{\raggedright\arraybackslash}p{(\columnwidth - 6\tabcolsep) * \real{0.3800}}@{}}
\toprule\noalign{}
\begin{minipage}[b]{\linewidth}\raggedright
Fragment
\end{minipage} & \begin{minipage}[b]{\linewidth}\raggedright
Joint
\end{minipage} & \begin{minipage}[b]{\linewidth}\raggedright
Applies to
\end{minipage} & \begin{minipage}[b]{\linewidth}\raggedright
Description
\end{minipage} \\
\midrule\noalign{}
\endhead
\bottomrule\noalign{}
\endlastfoot
\texttt{"hand\_twitch\_r"} & \texttt{rwrist} & \texttt{human},
\texttt{alien} & Right-hand fidget: small outward jiggle with a slight
up-lift, then inward overshoot return. Reads as ``fidget'', not
``flail.'' \\
\texttt{"hand\_twitch\_l"} & \texttt{lwrist} & \texttt{human},
\texttt{alien} & Left-hand fidget. x-mirror of \texttt{hand\_twitch\_r}
(v0.9.18). Use for left-handed characters or staging beats where the
right hand is occupied. \\
\texttt{"foot\_tap\_r"} & \texttt{rankle} & \texttt{human},
\texttt{alien} & Right-foot tap: pure-vertical lift then drop with a
small overshoot below the rest position before auto-return
(v0.9.18). \\
\texttt{"foot\_tap\_l"} & \texttt{lankle} & \texttt{human},
\texttt{alien} & Left-foot tap. Same vertical motion as
\texttt{foot\_tap\_r} on the opposite ankle (v0.9.18). \\
\texttt{"tail\_wag"} & \texttt{tail} & \texttt{dog} & Quadruped tail
back-and-forth swing. Three keyframes (up-forward, down-back,
up-forward) before auto-return for a clearly oscillating motion. \\
\end{longtable}

\textbf{Body-type filtering:} the tic system silently skips fragments
whose \texttt{applies\_to} list doesn't include the character's
\texttt{figure\_type}. A \texttt{tail\_wag} declared on a biped is a
no-op; a \texttt{hand\_twitch\_r} declared on a dog is a no-op.

\textbf{Per-step opt-out:} any individual \texttt{say} /
\texttt{prop\_say} / \texttt{react} step can set
\texttt{"suppress\_tics":\ true} to skip the otherwise-automatic
\texttt{sentence\_end} trigger for that step.

\textbf{Example --- Bevers's anxious-hand-fidget:}

\begin{Shaded}
\begin{Highlighting}[]
\ErrorTok{"bevers":} \FunctionTok{\{}
  \DataTypeTok{"figure\_type"}\FunctionTok{:} \StringTok{"alien"}\FunctionTok{,}
  \DataTypeTok{"tic\_profile"}\FunctionTok{:} \OtherTok{[}
    \FunctionTok{\{}\DataTypeTok{"fragment"}\FunctionTok{:} \StringTok{"hand\_twitch\_r"}\FunctionTok{,} \DataTypeTok{"triggers"}\FunctionTok{:} \OtherTok{[}\StringTok{"react"}\OtherTok{]}\FunctionTok{\}}
  \OtherTok{]}
\FunctionTok{\}}
\end{Highlighting}
\end{Shaded}

\textbf{Example --- left-handed variant for a hand-occupied beat.}
If a scene stages the right hand off camera or holding a prop, swap
to \texttt{hand\_twitch\_l} so the fidget reads:

\begin{Shaded}
\begin{Highlighting}[]
\ErrorTok{"bevers":} \FunctionTok{\{}
  \DataTypeTok{"figure\_type"}\FunctionTok{:} \StringTok{"alien"}\FunctionTok{,}
  \DataTypeTok{"tic\_profile"}\FunctionTok{:} \OtherTok{[}
    \FunctionTok{\{}\DataTypeTok{"fragment"}\FunctionTok{:} \StringTok{"hand\_twitch\_l"}\FunctionTok{,} \DataTypeTok{"triggers"}\FunctionTok{:} \OtherTok{[}\StringTok{"react"}\OtherTok{]}\FunctionTok{\}}
  \OtherTok{]}
\FunctionTok{\}}
\end{Highlighting}
\end{Shaded}

\textbf{Example --- a foot-tap on every sentence end.} The
\texttt{foot\_tap\_*} fragments fire on \texttt{sentence\_end} as well
as \texttt{react}, so a character can develop an impatience read across
a whole dialogue scene without per-line authoring:

\begin{Shaded}
\begin{Highlighting}[]
\ErrorTok{"thalia":} \FunctionTok{\{}
  \DataTypeTok{"figure\_type"}\FunctionTok{:} \StringTok{"alien"}\FunctionTok{,}
  \DataTypeTok{"tic\_profile"}\FunctionTok{:} \OtherTok{[}
    \FunctionTok{\{}\DataTypeTok{"fragment"}\FunctionTok{:} \StringTok{"foot\_tap\_r"}\FunctionTok{,} \DataTypeTok{"triggers"}\FunctionTok{:} \OtherTok{[}\StringTok{"sentence\_end"}\OtherTok{]}\FunctionTok{\}}
  \OtherTok{]}
\FunctionTok{\}}
\end{Highlighting}
\end{Shaded}

\textbf{Example --- combining two fragments on one character.} A
character can have multiple entries in \texttt{tic\_profile}, each
with its own trigger set:

\begin{Shaded}
\begin{Highlighting}[]
\ErrorTok{"freydoon":} \FunctionTok{\{}
  \DataTypeTok{"figure\_type"}\FunctionTok{:} \StringTok{"alien"}\FunctionTok{,}
  \DataTypeTok{"tic\_profile"}\FunctionTok{:} \OtherTok{[}
    \FunctionTok{\{}\DataTypeTok{"fragment"}\FunctionTok{:} \StringTok{"hand\_twitch\_r"}\FunctionTok{,} \DataTypeTok{"triggers"}\FunctionTok{:} \OtherTok{[}\StringTok{"react"}\OtherTok{]}\FunctionTok{\}}\FunctionTok{,}
    \FunctionTok{\{}\DataTypeTok{"fragment"}\FunctionTok{:} \StringTok{"foot\_tap\_l"}\FunctionTok{,} \DataTypeTok{"triggers"}\FunctionTok{:} \OtherTok{[}\StringTok{"sentence\_end"}\OtherTok{]}\FunctionTok{\}}
  \OtherTok{]}
\FunctionTok{\}}
\end{Highlighting}
\end{Shaded}

\textbf{Example --- Chekov's tail-wag on every sentence end:}

\begin{Shaded}
\begin{Highlighting}[]
\ErrorTok{"chekov":} \FunctionTok{\{}
  \DataTypeTok{"figure\_type"}\FunctionTok{:} \StringTok{"dog"}\FunctionTok{,}
  \DataTypeTok{"tic\_profile"}\FunctionTok{:} \OtherTok{[}
    \FunctionTok{\{}\DataTypeTok{"fragment"}\FunctionTok{:} \StringTok{"tail\_wag"}\FunctionTok{,} \DataTypeTok{"triggers"}\FunctionTok{:} \OtherTok{[}\StringTok{"sentence\_end"}\OtherTok{,} \StringTok{"react"}\OtherTok{]}\FunctionTok{\}}
  \OtherTok{]}
\FunctionTok{\}}
\end{Highlighting}
\end{Shaded}

See \texttt{pam/tics.py} for fragment definitions and the design
rationale for delta-based fragments.

\begin{center}\rule{0.5\linewidth}{0.5pt}\end{center}

\subsubsection{1.13 Field precedence and the canonical
registry}\label{field-precedence-and-the-canonical-registry}

In multi-scene projects, characters typically have a canonical
definition in \texttt{tntd\_characters.json} (or your project's
equivalent registry file). \texttt{fountain2pam.py} merges this into
each scene's cast block at conversion time.

\textbf{Canonical (registry-stored) fields:}

\begin{itemize}
\tightlist
\item
  \texttt{figure\_type}, \texttt{build}, \texttt{gender}
\item
  \texttt{color}, \texttt{torso\_color}, \texttt{style}
\item
  \texttt{default\_scale} (the default scale at which the character
  renders; stored in the registry as a convenience)
\end{itemize}

\textbf{Scene-specific (per-scene only) fields:}

\begin{itemize}
\tightlist
\item
  \texttt{offset}, \texttt{scale}, \texttt{pose}
\item
  \texttt{facing} (dog only)
\end{itemize}

The registry is loaded \emph{before} conversion: missing canonical
fields in a scene's cast block are filled in from the registry. The
scene always wins for scene-specific fields --- a character can stand at
different \texttt{offset}s in different scenes.

\textbf{Color authority (v0.9.8 fix):} if the registry has a
\texttt{color} field, the converter re-derives the full six-key palette
from it and stamps it into the character's \texttt{style} dict,
overwriting any per-scene values during conversion. This means: to
change a character's canonical color, edit the registry file directly.
Per-scene \texttt{color} and \texttt{style.edge\_color} overrides are
not durable across registry-driven conversions.

\textbf{Registry path:} \texttt{fountain2pam.py} looks for
\texttt{characters.json} next to the \texttt{.fountain} file by default.
Override with \texttt{-\/-characters\ /path/to/tntd\_characters.json}.

\textbf{Example registry entry:}

\begin{Shaded}
\begin{Highlighting}[]
\FunctionTok{\{}
  \DataTypeTok{"bevers"}\FunctionTok{:} \FunctionTok{\{}
    \DataTypeTok{"figure\_type"}\FunctionTok{:} \StringTok{"alien"}\FunctionTok{,}
    \DataTypeTok{"build"}\FunctionTok{:} \StringTok{"alien"}\FunctionTok{,}
    \DataTypeTok{"gender"}\FunctionTok{:} \StringTok{"male"}\FunctionTok{,}
    \DataTypeTok{"color"}\FunctionTok{:} \StringTok{"\#44bb66"}\FunctionTok{,}
    \DataTypeTok{"torso\_color"}\FunctionTok{:} \StringTok{"\#2a6090"}\FunctionTok{,}
    \DataTypeTok{"style"}\FunctionTok{:} \FunctionTok{\{}\DataTypeTok{"head\_label"}\FunctionTok{:} \StringTok{"Bevers"}\FunctionTok{\},}
    \DataTypeTok{"default\_scale"}\FunctionTok{:} \FunctionTok{\{}\DataTypeTok{"sy"}\FunctionTok{:} \DecValTok{0}\ErrorTok{.}\DecValTok{7}\FunctionTok{,} \DataTypeTok{"sx"}\FunctionTok{:} \DecValTok{0}\ErrorTok{.}\DecValTok{7}\FunctionTok{,} \DataTypeTok{"anchor"}\FunctionTok{:} \StringTok{"lankle"}\FunctionTok{\},}
    \DataTypeTok{"tic\_profile"}\FunctionTok{:} \OtherTok{[}
      \FunctionTok{\{}\DataTypeTok{"fragment"}\FunctionTok{:} \StringTok{"hand\_twitch\_r"}\FunctionTok{,} \DataTypeTok{"triggers"}\FunctionTok{:} \OtherTok{[}\StringTok{"react"}\OtherTok{]}\FunctionTok{\}}
    \OtherTok{]}
  \FunctionTok{\}}
\FunctionTok{\}}
\end{Highlighting}
\end{Shaded}

A per-scene cast block that references
\texttt{"bevers":\ \{"offset":\ {[}3,\ 0,\ 0{]}\}} will receive all
canonical fields above merged in automatically.

\begin{center}\rule{0.5\linewidth}{0.5pt}\end{center}

\emph{Sections 2 (modifier actions), 3 (action targets), and 4 (meta)
are delivered in subsequent substeps. This document will be consolidated
into a single \texttt{CHARACTER\_PROPERTIES.md} when all sections are
complete.}

\subsection{2. Modifier actions}\label{modifier-actions}

A \emph{modifier action} changes character state without locomotion or
dialogue. This section covers visibility (\texttt{fade\_in},
\texttt{fade\_out}), transforms (\texttt{morph}, \texttt{scale},
\texttt{rotate}), decoration (\texttt{change\_uniform},
\texttt{attach\_face}, \texttt{detach\_face}), framing (\texttt{focus},
\texttt{focus\_reset}), and bubble cleanup (\texttt{clear\_bubble},
\texttt{clear\_all\_bubbles}).

Throughout, each action's heading carries three flags:

\begin{itemize}
\tightlist
\item
  \textbf{Parallel-safe} --- whether the action may appear inside a
  \texttt{parallel} block. Unsafe actions fire sequentially even when
  nested inside \texttt{parallel}, with a console warning.
\item
  \textbf{Animation} --- \texttt{none} (instant state change), or
  \texttt{timed} (the action consumes scene time via a \texttt{play()}
  call). Timed actions accept \texttt{rt} / \texttt{duration} to
  override the default.
\item
  \textbf{Required keys} --- minimum JSON shape that does something
  useful.
\end{itemize}

Cross-references to actions documented in §3 (\texttt{walk\_to},
\texttt{say}, etc.) use forward links rather than duplicating
parameters.

\begin{center}\rule{0.5\linewidth}{0.5pt}\end{center}

\subsubsection{\texorpdfstring{2.1
\texttt{fade\_in}}{2.1 fade\_in}}\label{fade_in}

\emph{Parallel-safe:} no · \emph{Animation:} timed · \emph{Required
keys:} \texttt{who}

Bring a character onto the screen with an animated entrance (edges draw
first, then joints grow). Constructs the appropriate figure class based
on \texttt{figure\_type} and writes the live figure into
\texttt{cast{[}name{]}{[}"fig"{]}}.

\textbf{JSON shape (minimum, character already in cast):}

\begin{Shaded}
\begin{Highlighting}[]
\FunctionTok{\{}\DataTypeTok{"action"}\FunctionTok{:} \StringTok{"fade\_in"}\FunctionTok{,} \DataTypeTok{"who"}\FunctionTok{:} \StringTok{"bevers"}\FunctionTok{\}}
\end{Highlighting}
\end{Shaded}

\textbf{JSON shape (mid-scene spawn, character not in cast):} all the
creation-time fields from §1 may be supplied inline.

\begin{Shaded}
\begin{Highlighting}[]
\FunctionTok{\{}\DataTypeTok{"action"}\FunctionTok{:} \StringTok{"fade\_in"}\FunctionTok{,} \DataTypeTok{"who"}\FunctionTok{:} \StringTok{"athena"}\FunctionTok{,}
 \DataTypeTok{"figure\_type"}\FunctionTok{:} \StringTok{"human"}\FunctionTok{,} \DataTypeTok{"build"}\FunctionTok{:} \StringTok{"narrow"}\FunctionTok{,} \DataTypeTok{"gender"}\FunctionTok{:} \StringTok{"female"}\FunctionTok{,}
 \DataTypeTok{"offset"}\FunctionTok{:} \OtherTok{[}\DecValTok{{-}2}\OtherTok{,} \DecValTok{0}\OtherTok{,} \DecValTok{0}\OtherTok{]}\FunctionTok{,} \DataTypeTok{"pose"}\FunctionTok{:} \StringTok{"standing\_front"}\FunctionTok{,}
 \DataTypeTok{"scale"}\FunctionTok{:} \FunctionTok{\{}\DataTypeTok{"sy"}\FunctionTok{:} \DecValTok{0}\ErrorTok{.}\DecValTok{7}\FunctionTok{,} \DataTypeTok{"sx"}\FunctionTok{:} \DecValTok{0}\ErrorTok{.}\DecValTok{7}\FunctionTok{,} \DataTypeTok{"anchor"}\FunctionTok{:} \StringTok{"lankle"}\FunctionTok{\},}
 \DataTypeTok{"color"}\FunctionTok{:} \StringTok{"\#5b9cf6"}\FunctionTok{,}
 \DataTypeTok{"style"}\FunctionTok{:} \FunctionTok{\{}\DataTypeTok{"head\_label"}\FunctionTok{:} \StringTok{"Athena"}\FunctionTok{\}\}}
\end{Highlighting}
\end{Shaded}

\textbf{Parameters:}

\begin{longtable}[]{@{}
  >{\raggedright\arraybackslash}p{(\columnwidth - 8\tabcolsep) * \real{0.2000}}
  >{\raggedright\arraybackslash}p{(\columnwidth - 8\tabcolsep) * \real{0.2000}}
  >{\raggedright\arraybackslash}p{(\columnwidth - 8\tabcolsep) * \real{0.2000}}
  >{\raggedright\arraybackslash}p{(\columnwidth - 8\tabcolsep) * \real{0.2000}}
  >{\raggedright\arraybackslash}p{(\columnwidth - 8\tabcolsep) * \real{0.2000}}@{}}
\toprule\noalign{}
\begin{minipage}[b]{\linewidth}\raggedright
Key
\end{minipage} & \begin{minipage}[b]{\linewidth}\raggedright
Type
\end{minipage} & \begin{minipage}[b]{\linewidth}\raggedright
Required
\end{minipage} & \begin{minipage}[b]{\linewidth}\raggedright
Default
\end{minipage} & \begin{minipage}[b]{\linewidth}\raggedright
Description
\end{minipage} \\
\midrule\noalign{}
\endhead
\bottomrule\noalign{}
\endlastfoot
\texttt{who} & \texttt{str} & yes & --- & Character key from the cast
block, or a new name to spawn fresh. \\
\texttt{figure\_type} & \texttt{str} & no & from cast, else
\texttt{"human"} & See §1.1. \\
\texttt{build} & \texttt{str} & no & from cast, else \texttt{"default"}
& See §1.2. \\
\texttt{gender} & \texttt{str} & no & from cast, else \texttt{None} &
See §1.3. \\
\texttt{offset} & \texttt{list{[}float{]}} & no & from cast, else
\texttt{{[}0,0,0{]}} & World-space spawn position (3-element). \\
\texttt{x} & \texttt{float} & no & --- & Convenience: sets
\texttt{offset\ =\ {[}x,\ 0,\ 0{]}}. Overrides \texttt{offset} if both
are present. \\
\texttt{pose} & \texttt{str} & no & from cast, else
\texttt{"standing\_front"} & Initial pose name. See §1.7. \\
\texttt{scale} & \texttt{dict} & no & from cast &
\texttt{\{sy,\ sx,\ anchor\}} dict. See §1.9. \\
\texttt{style} & \texttt{dict} & no & from cast & Full style palette.
See §1.6. \\
\texttt{color} & \texttt{str} & no & from cast & Single-color shorthand.
See §1.4. \\
\texttt{torso\_color} & \texttt{str} & no & from cast & Two-zone torso
color. See §1.5. \\
\texttt{facing} & \texttt{str} & no & from cast, else \texttt{"right"} &
Dog-only direction. See §1.10. \\
\texttt{rt} & \texttt{float} & no & figure default & Edge fade-in time
in seconds. \\
\texttt{duration} & \texttt{float} & no & --- & Synonym for
\texttt{rt}. \\
\end{longtable}

\textbf{Behavior:}

\begin{itemize}
\tightlist
\item
  All canonical fields (figure\_type, build, gender, color,
  torso\_color, style) fall through from the cast spec if not given on
  the \texttt{fade\_in} step.
\item
  Scene-specific fields (offset, scale, pose, facing) likewise fall
  through from the cast spec.
\item
  The \texttt{tic\_profile} from the cast spec is wired onto the live
  figure as \texttt{fig.tic\_profile} so trigger handlers can read it
  without a cast-side lookup.
\item
  Dog figures (Path C, v0.9.14) are dual-registered: a \texttt{DogGraph}
  faded in via \texttt{fade\_in} is also added to the prop registry
  under the same name, so \texttt{who:\ chekov} and
  \texttt{prop:\ chekov} both resolve. Skipped if a prop with that name
  already exists.
\end{itemize}

\textbf{Defaults derive from the cast spec, not the cast registry file:}
if \texttt{tntd\_characters.json} has Bevers's color but Bevers is
missing from the cast block of this scene, \texttt{fade\_in} won't find
anything --- the registry merge happens at \texttt{fountain2pam.py}
conversion time, not at \texttt{pam\_player.py} runtime.

\textbf{Re-fade-in:} calling \texttt{fade\_in} for a character who has
been faded out is supported. The figure is rebuilt fresh from the cast
spec; any prior runtime state (walked-to position, attached face,
applied uniform) is gone. This is the documented mechanism for
non-persistence of face attachment (see §2.7).

\begin{center}\rule{0.5\linewidth}{0.5pt}\end{center}

\subsubsection{\texorpdfstring{2.2
\texttt{fade\_out}}{2.2 fade\_out}}\label{fade_out}

\emph{Parallel-safe:} yes · \emph{Animation:} timed · \emph{Required
keys:} \texttt{who}

Fade the character off-screen with a synchronized fade of edges and
joints, and clear their cast slot. Any attached face is torn down
concurrently with the body (v0.9.15).

\textbf{JSON shape:}

\begin{Shaded}
\begin{Highlighting}[]
\FunctionTok{\{}\DataTypeTok{"action"}\FunctionTok{:} \StringTok{"fade\_out"}\FunctionTok{,} \DataTypeTok{"who"}\FunctionTok{:} \StringTok{"bevers"}\FunctionTok{,} \DataTypeTok{"rt"}\FunctionTok{:} \DecValTok{0}\ErrorTok{.}\DecValTok{8}\FunctionTok{\}}
\end{Highlighting}
\end{Shaded}

\textbf{Parameters:}

\begin{longtable}[]{@{}
  >{\raggedright\arraybackslash}p{(\columnwidth - 8\tabcolsep) * \real{0.2000}}
  >{\raggedright\arraybackslash}p{(\columnwidth - 8\tabcolsep) * \real{0.2000}}
  >{\raggedright\arraybackslash}p{(\columnwidth - 8\tabcolsep) * \real{0.2000}}
  >{\raggedright\arraybackslash}p{(\columnwidth - 8\tabcolsep) * \real{0.2000}}
  >{\raggedright\arraybackslash}p{(\columnwidth - 8\tabcolsep) * \real{0.2000}}@{}}
\toprule\noalign{}
\begin{minipage}[b]{\linewidth}\raggedright
Key
\end{minipage} & \begin{minipage}[b]{\linewidth}\raggedright
Type
\end{minipage} & \begin{minipage}[b]{\linewidth}\raggedright
Required
\end{minipage} & \begin{minipage}[b]{\linewidth}\raggedright
Default
\end{minipage} & \begin{minipage}[b]{\linewidth}\raggedright
Description
\end{minipage} \\
\midrule\noalign{}
\endhead
\bottomrule\noalign{}
\endlastfoot
\texttt{who} & \texttt{str} & yes & --- & Character key. \texttt{"all"}
is supported as a convenience to fade out every active cast member
sequentially. \\
\texttt{rt} & \texttt{float} & no & \texttt{1.0} & Fade-out duration in
seconds. \\
\texttt{duration} & \texttt{float} & no & --- & Synonym for
\texttt{rt}. \\
\end{longtable}

\textbf{Behavior:}

\begin{itemize}
\tightlist
\item
  \texttt{cast{[}name{]}{[}"fig"{]}} is cleared to \texttt{None} after
  the fade.
\item
  If \texttt{fig.head\_face} is set (v0.9.15 face attachment), the
  updater is removed and \texttt{FadeOut(face)} is added to the same
  \texttt{play()} call as the body fades --- single concurrent
  animation.
\item
  The figure object itself is discarded; opacity state on its head dot
  is not restored. A subsequent \texttt{fade\_in} builds a fresh figure.
\end{itemize}

\textbf{Parallel-safe} because the action collects animations via the
\texttt{\_collect\_anims} path and returns a list of \texttt{FadeOut}
mobjects for the parallel composer.

\begin{center}\rule{0.5\linewidth}{0.5pt}\end{center}

\subsubsection{\texorpdfstring{2.3
\texttt{morph}}{2.3 morph}}\label{morph}

\emph{Parallel-safe:} yes · \emph{Animation:} timed · \emph{Required
keys:} \texttt{who}, \texttt{pose}

Animate the character from its current pose to a new named pose, with an
optional position offset. The character's \texttt{fig.pose} and
\texttt{fig.offset} are updated to match.

\textbf{JSON shape:}

\begin{Shaded}
\begin{Highlighting}[]
\FunctionTok{\{}\DataTypeTok{"action"}\FunctionTok{:} \StringTok{"morph"}\FunctionTok{,} \DataTypeTok{"who"}\FunctionTok{:} \StringTok{"bevers"}\FunctionTok{,} \DataTypeTok{"pose"}\FunctionTok{:} \StringTok{"standing\_side"}\FunctionTok{,} \DataTypeTok{"rt"}\FunctionTok{:} \DecValTok{0}\ErrorTok{.}\DecValTok{4}\FunctionTok{\}}
\end{Highlighting}
\end{Shaded}

\textbf{Parameters:}

\begin{longtable}[]{@{}
  >{\raggedright\arraybackslash}p{(\columnwidth - 8\tabcolsep) * \real{0.2000}}
  >{\raggedright\arraybackslash}p{(\columnwidth - 8\tabcolsep) * \real{0.2000}}
  >{\raggedright\arraybackslash}p{(\columnwidth - 8\tabcolsep) * \real{0.2000}}
  >{\raggedright\arraybackslash}p{(\columnwidth - 8\tabcolsep) * \real{0.2000}}
  >{\raggedright\arraybackslash}p{(\columnwidth - 8\tabcolsep) * \real{0.2000}}@{}}
\toprule\noalign{}
\begin{minipage}[b]{\linewidth}\raggedright
Key
\end{minipage} & \begin{minipage}[b]{\linewidth}\raggedright
Type
\end{minipage} & \begin{minipage}[b]{\linewidth}\raggedright
Required
\end{minipage} & \begin{minipage}[b]{\linewidth}\raggedright
Default
\end{minipage} & \begin{minipage}[b]{\linewidth}\raggedright
Description
\end{minipage} \\
\midrule\noalign{}
\endhead
\bottomrule\noalign{}
\endlastfoot
\texttt{who} & \texttt{str} & yes & --- & Character key. \\
\texttt{pose} & \texttt{str} & yes & --- & Target pose name from
\texttt{POSES} registry (see §1.7). Build-specific poses override global
registry. \\
\texttt{rt} & \texttt{float} & no & \texttt{0.4} & Morph duration in
seconds. \\
\texttt{dx} & \texttt{float} & no & \texttt{0.0} & Add to the figure's
current x-offset. Useful for poses where the visual center shifts. \\
\texttt{dy} & \texttt{float} & no & \texttt{0.0} & Add to the figure's
current y-offset. \\
\end{longtable}

\textbf{Examples:}

\begin{Shaded}
\begin{Highlighting}[]
\FunctionTok{\{}\DataTypeTok{"action"}\FunctionTok{:} \StringTok{"morph"}\FunctionTok{,} \DataTypeTok{"who"}\FunctionTok{:} \StringTok{"freydoon"}\FunctionTok{,} \DataTypeTok{"pose"}\FunctionTok{:} \StringTok{"look\_up"}\FunctionTok{\}}
\FunctionTok{\{}\DataTypeTok{"action"}\FunctionTok{:} \StringTok{"morph"}\FunctionTok{,} \DataTypeTok{"who"}\FunctionTok{:} \StringTok{"thalia"}\FunctionTok{,}   \DataTypeTok{"pose"}\FunctionTok{:} \StringTok{"standing\_side"}\FunctionTok{,} \DataTypeTok{"rt"}\FunctionTok{:} \DecValTok{0}\ErrorTok{.}\DecValTok{3}\FunctionTok{\}}
\FunctionTok{\{}\DataTypeTok{"action"}\FunctionTok{:} \StringTok{"morph"}\FunctionTok{,} \DataTypeTok{"who"}\FunctionTok{:} \StringTok{"brad"}\FunctionTok{,}     \DataTypeTok{"pose"}\FunctionTok{:} \StringTok{"lying\_flat\_right"}\FunctionTok{,} \DataTypeTok{"dy"}\FunctionTok{:} \DecValTok{{-}0}\ErrorTok{.}\DecValTok{3}\FunctionTok{\}}
\end{Highlighting}
\end{Shaded}

\textbf{Notes:}

\begin{itemize}
\tightlist
\item
  For initial poses at fade-in, use the cast spec's \texttt{pose} field
  (§1.7) rather than a separate \texttt{morph} step.
\item
  For posture transitions (\texttt{sit\_down}, \texttt{stand\_up}),
  prefer the dedicated action (§3) which handles the multi-keyframe
  cycle correctly.
\item
  Morphing back to \texttt{standing\_front} undoes any prior
  \texttt{rotate} (since \texttt{rotate} doesn't update internal pose;
  see §2.5).
\end{itemize}

\begin{center}\rule{0.5\linewidth}{0.5pt}\end{center}

\subsubsection{\texorpdfstring{2.4
\texttt{scale}}{2.4 scale}}\label{scale-1}

\emph{Parallel-safe:} yes · \emph{Animation:} timed · \emph{Required
keys:} \texttt{who}

Apply a persistent scale to the figure. Unlike the creation-time
\texttt{scale} dict (§1.9), this is an action that animates the scale
change live. The new scale persists until the next \texttt{scale}
action.

\textbf{JSON shape:}

\begin{Shaded}
\begin{Highlighting}[]
\FunctionTok{\{}\DataTypeTok{"action"}\FunctionTok{:} \StringTok{"scale"}\FunctionTok{,} \DataTypeTok{"who"}\FunctionTok{:} \StringTok{"chava"}\FunctionTok{,} \DataTypeTok{"sy"}\FunctionTok{:} \DecValTok{0}\ErrorTok{.}\DecValTok{48}\FunctionTok{,} \DataTypeTok{"sx"}\FunctionTok{:} \DecValTok{0}\ErrorTok{.}\DecValTok{48}\FunctionTok{,}
 \DataTypeTok{"anchor"}\FunctionTok{:} \StringTok{"lankle"}\FunctionTok{,} \DataTypeTok{"rt"}\FunctionTok{:} \DecValTok{0}\ErrorTok{.}\DecValTok{6}\FunctionTok{\}}
\end{Highlighting}
\end{Shaded}

\textbf{Parameters:}

\begin{longtable}[]{@{}
  >{\raggedright\arraybackslash}p{(\columnwidth - 8\tabcolsep) * \real{0.2000}}
  >{\raggedright\arraybackslash}p{(\columnwidth - 8\tabcolsep) * \real{0.2000}}
  >{\raggedright\arraybackslash}p{(\columnwidth - 8\tabcolsep) * \real{0.2000}}
  >{\raggedright\arraybackslash}p{(\columnwidth - 8\tabcolsep) * \real{0.2000}}
  >{\raggedright\arraybackslash}p{(\columnwidth - 8\tabcolsep) * \real{0.2000}}@{}}
\toprule\noalign{}
\begin{minipage}[b]{\linewidth}\raggedright
Key
\end{minipage} & \begin{minipage}[b]{\linewidth}\raggedright
Type
\end{minipage} & \begin{minipage}[b]{\linewidth}\raggedright
Required
\end{minipage} & \begin{minipage}[b]{\linewidth}\raggedright
Default
\end{minipage} & \begin{minipage}[b]{\linewidth}\raggedright
Description
\end{minipage} \\
\midrule\noalign{}
\endhead
\bottomrule\noalign{}
\endlastfoot
\texttt{who} & \texttt{str} & yes & --- & Character key. \\
\texttt{sy} & \texttt{float} & no & \texttt{1.0} & New vertical scale.
\textbf{Not a multiplier} --- replaces the previous \texttt{sy}
outright. \\
\texttt{sx} & \texttt{float} & no & \texttt{1.0} & New horizontal
scale. \\
\texttt{anchor} & \texttt{str} & no & \texttt{"lankle"} & Joint name to
keep fixed during the rescale. See §1.9. \\
\texttt{rt} & \texttt{float} & no & \texttt{0.8} & Scale duration in
seconds. \\
\end{longtable}

\textbf{Note on the ``not a multiplier'' semantics:} if Chava is at
scale 0.7 and you write
\texttt{\{"action":\ "scale",\ "who":\ "chava",\ "sy":\ 0.5\}}, her new
scale is 0.5 (not 0.35). To halve a character, compute the target value
yourself.

\textbf{Use cases:}

\begin{itemize}
\tightlist
\item
  Distance suggestion --- drop a character to scale 0.32 to read as
  ``across the room'' or ``outside the window.''
\item
  Growth or shrinking effects --- sci-fi avatar transitions.
\item
  Establishing-shot to close-up transitions when paired with a camera
  move.
\end{itemize}

\begin{center}\rule{0.5\linewidth}{0.5pt}\end{center}

\subsubsection{\texorpdfstring{2.5
\texttt{rotate}}{2.5 rotate}}\label{rotate}

\emph{Parallel-safe:} no · \emph{Animation:} timed · \emph{Required
keys:} \texttt{who} or \texttt{prop}, plus an angle

Rotate a character or a prop around a chosen pivot point. Useful for
laying a character flat on a stretcher, knocking a figure off-balance,
falling, opening a pod lid, or swinging a door on its hinge.

\textbf{JSON shape:}

\begin{Shaded}
\begin{Highlighting}[]
\FunctionTok{\{}\DataTypeTok{"action"}\FunctionTok{:} \StringTok{"rotate"}\FunctionTok{,} \DataTypeTok{"who"}\FunctionTok{:} \StringTok{"freydoon"}\FunctionTok{,}
 \DataTypeTok{"angle\_deg"}\FunctionTok{:} \DecValTok{{-}90}\FunctionTok{,} \DataTypeTok{"pivot"}\FunctionTok{:} \StringTok{"bottom"}\FunctionTok{,} \DataTypeTok{"rt"}\FunctionTok{:} \DecValTok{0}\ErrorTok{.}\DecValTok{6}\FunctionTok{\}}
\end{Highlighting}
\end{Shaded}

\textbf{Parameters:}

\begin{longtable}[]{@{}
  >{\raggedright\arraybackslash}p{(\columnwidth - 8\tabcolsep) * \real{0.2000}}
  >{\raggedright\arraybackslash}p{(\columnwidth - 8\tabcolsep) * \real{0.2000}}
  >{\raggedright\arraybackslash}p{(\columnwidth - 8\tabcolsep) * \real{0.2000}}
  >{\raggedright\arraybackslash}p{(\columnwidth - 8\tabcolsep) * \real{0.2000}}
  >{\raggedright\arraybackslash}p{(\columnwidth - 8\tabcolsep) * \real{0.2000}}@{}}
\toprule\noalign{}
\begin{minipage}[b]{\linewidth}\raggedright
Key
\end{minipage} & \begin{minipage}[b]{\linewidth}\raggedright
Type
\end{minipage} & \begin{minipage}[b]{\linewidth}\raggedright
Required
\end{minipage} & \begin{minipage}[b]{\linewidth}\raggedright
Default
\end{minipage} & \begin{minipage}[b]{\linewidth}\raggedright
Description
\end{minipage} \\
\midrule\noalign{}
\endhead
\bottomrule\noalign{}
\endlastfoot
\texttt{who} & \texttt{str} & one of & --- & Character key. Mutually
exclusive with \texttt{prop}; if both are given, \texttt{prop} wins. \\
\texttt{prop} & \texttt{str} & one of & --- & Prop registry key. \\
\texttt{angle\_deg} & \texttt{float} & one of & --- & Rotation angle in
degrees. Preferred. \\
\texttt{angle} & \texttt{float} & one of & --- & Rotation angle in
radians. Used if \texttt{angle\_deg} is absent. \\
\texttt{pivot} & \texttt{str} or \texttt{list} & no & \texttt{"bottom"}
& One of \texttt{"bottom"}, \texttt{"center"}, \texttt{"top"},
\texttt{"left"}, \texttt{"right"}, or an explicit \texttt{{[}x,\ y{]}}
world point. For characters, \texttt{"bottom"} is the feet; for props,
the floor line. \\
\texttt{rt} & \texttt{float} & no & \texttt{0.5} & Animation duration.
\texttt{0} produces an instant non-animated rotation. \\
\end{longtable}

\textbf{Examples:}

\begin{Shaded}
\begin{Highlighting}[]
\FunctionTok{\{}\DataTypeTok{"action"}\FunctionTok{:} \StringTok{"rotate"}\FunctionTok{,} \DataTypeTok{"who"}\FunctionTok{:} \StringTok{"freydoon"}\FunctionTok{,} \DataTypeTok{"angle\_deg"}\FunctionTok{:} \DecValTok{{-}90}\FunctionTok{,} \DataTypeTok{"pivot"}\FunctionTok{:} \StringTok{"bottom"}\FunctionTok{\}}
\FunctionTok{\{}\DataTypeTok{"action"}\FunctionTok{:} \StringTok{"rotate"}\FunctionTok{,} \DataTypeTok{"prop"}\FunctionTok{:} \StringTok{"bevers\_pod"}\FunctionTok{,} \DataTypeTok{"angle\_deg"}\FunctionTok{:} \DecValTok{70}\FunctionTok{,} \DataTypeTok{"pivot"}\FunctionTok{:} \StringTok{"left"}\FunctionTok{\}}
\FunctionTok{\{}\DataTypeTok{"action"}\FunctionTok{:} \StringTok{"rotate"}\FunctionTok{,} \DataTypeTok{"prop"}\FunctionTok{:} \StringTok{"stretcher"}\FunctionTok{,}  \DataTypeTok{"angle\_deg"}\FunctionTok{:} \DecValTok{4}\FunctionTok{,}  \DataTypeTok{"rt"}\FunctionTok{:} \DecValTok{0}\ErrorTok{.}\DecValTok{15}\FunctionTok{\}}
\end{Highlighting}
\end{Shaded}

\textbf{Caveat --- characters only:} rotation is a visual-only
transform. The figure's internal \texttt{pose} and \texttt{offset} are
not updated. If you \texttt{morph}, \texttt{walk\_to}, or otherwise
reanimate a rotated character, the next \texttt{morph\_to} snaps it back
to upright. Issue a counter-rotation (\texttt{angle\_deg:\ +90} after a
\texttt{-90} lay-flat) before further animation.

\textbf{Props do not have this caveat} --- their geometry is rotated
permanently and subsequent \texttt{move\_aside} /
\texttt{group\_translate} / \texttt{remove\_prop} operations compose
correctly with the rotation.

\textbf{Interaction with \texttt{attach\_face} (v0.9.15):} the face
mobject tracks head position but does not rotate with the body. A
character rotated to lie flat will have a horizontal body and an upright
face --- cartoon convention. See §2.7 for the design note.

\begin{center}\rule{0.5\linewidth}{0.5pt}\end{center}

\subsubsection{\texorpdfstring{2.6
\texttt{change\_uniform}}{2.6 change\_uniform}}\label{change_uniform}

\emph{Parallel-safe:} no · \emph{Animation:} timed · \emph{Required
keys:} \texttt{who}

Apply a named uniform variant from the character's \texttt{uniforms}
dict (§1.11). Updates the torso color (and optionally the edge color)
live; the variant persists until the next \texttt{change\_uniform}.

\textbf{JSON shape:}

\begin{Shaded}
\begin{Highlighting}[]
\FunctionTok{\{}\DataTypeTok{"action"}\FunctionTok{:} \StringTok{"change\_uniform"}\FunctionTok{,} \DataTypeTok{"who"}\FunctionTok{:} \StringTok{"chava"}\FunctionTok{,} \DataTypeTok{"uniform"}\FunctionTok{:} \StringTok{"delivery"}\FunctionTok{,} \DataTypeTok{"rt"}\FunctionTok{:} \DecValTok{0}\ErrorTok{.}\DecValTok{3}\FunctionTok{\}}
\end{Highlighting}
\end{Shaded}

\textbf{Parameters:}

\begin{longtable}[]{@{}
  >{\raggedright\arraybackslash}p{(\columnwidth - 8\tabcolsep) * \real{0.2000}}
  >{\raggedright\arraybackslash}p{(\columnwidth - 8\tabcolsep) * \real{0.2000}}
  >{\raggedright\arraybackslash}p{(\columnwidth - 8\tabcolsep) * \real{0.2000}}
  >{\raggedright\arraybackslash}p{(\columnwidth - 8\tabcolsep) * \real{0.2000}}
  >{\raggedright\arraybackslash}p{(\columnwidth - 8\tabcolsep) * \real{0.2000}}@{}}
\toprule\noalign{}
\begin{minipage}[b]{\linewidth}\raggedright
Key
\end{minipage} & \begin{minipage}[b]{\linewidth}\raggedright
Type
\end{minipage} & \begin{minipage}[b]{\linewidth}\raggedright
Required
\end{minipage} & \begin{minipage}[b]{\linewidth}\raggedright
Default
\end{minipage} & \begin{minipage}[b]{\linewidth}\raggedright
Description
\end{minipage} \\
\midrule\noalign{}
\endhead
\bottomrule\noalign{}
\endlastfoot
\texttt{who} & \texttt{str} & yes & --- & Character key. \\
\texttt{uniform} & \texttt{str} & yes (or override) & --- & Variant name
from the character's \texttt{uniforms} dict. \\
\texttt{torso\_color} & \texttt{str} (hex) & no & --- & Inline override.
Used if \texttt{uniform} is absent or the named variant doesn't define
\texttt{torso\_color}. \\
\texttt{color} & \texttt{str} (hex) & no & --- & Inline override for the
edge color. Updates all non-torso edges and joints. \\
\texttt{rt} & \texttt{float} & no & \texttt{0.3} & Hold time after
applying. Not a true animation --- the color change is currently
instant; \texttt{rt} is a \texttt{wait()} so nearby actions have
breathing room. Pass \texttt{0} to skip the wait. \\
\end{longtable}

\textbf{Resolution order for \texttt{torso\_color}:}

\begin{enumerate}
\def\labelenumi{\arabic{enumi}.}
\tightlist
\item
  The named uniform's \texttt{torso\_color}.
\item
  Inline \texttt{torso\_color} on the step.
\item
  Cast spec's top-level \texttt{torso\_color}.
\end{enumerate}

If none of the three resolves, the action prints a warning and skips.

\textbf{Examples:}

\begin{Shaded}
\begin{Highlighting}[]
\FunctionTok{\{}\DataTypeTok{"action"}\FunctionTok{:} \StringTok{"change\_uniform"}\FunctionTok{,} \DataTypeTok{"who"}\FunctionTok{:} \StringTok{"chava"}\FunctionTok{,} \DataTypeTok{"uniform"}\FunctionTok{:} \StringTok{"delivery"}\FunctionTok{\}}
\FunctionTok{\{}\DataTypeTok{"action"}\FunctionTok{:} \StringTok{"change\_uniform"}\FunctionTok{,} \DataTypeTok{"who"}\FunctionTok{:} \StringTok{"nona"}\FunctionTok{,}  \DataTypeTok{"uniform"}\FunctionTok{:} \StringTok{"formal"}\FunctionTok{\}}
\FunctionTok{\{}\DataTypeTok{"action"}\FunctionTok{:} \StringTok{"change\_uniform"}\FunctionTok{,} \DataTypeTok{"who"}\FunctionTok{:} \StringTok{"chava"}\FunctionTok{,} \DataTypeTok{"uniform"}\FunctionTok{:} \StringTok{"default"}\FunctionTok{\}}
\end{Highlighting}
\end{Shaded}

\textbf{Reserved uniform name:} \texttt{"default"} is the baseline ---
declare it in \texttt{uniforms} if you want a stable name to return to.

\begin{center}\rule{0.5\linewidth}{0.5pt}\end{center}

\subsubsection{\texorpdfstring{2.7
\texttt{attach\_face}}{2.7 attach\_face}}\label{attach_face}

\emph{Parallel-safe:} yes (returns empty animation list) ·
\emph{Animation:} none · \emph{Required keys:} \texttt{who},
\texttt{image} · \emph{Version:} v0.9.15

Attach a PNG face image to a character's head dot. The image is loaded
as an \texttt{ImageMobject}, sized via the \texttt{scale} parameter,
positioned at the head dot's current center, and given a Manim updater
so it tracks the head every frame (through walks, morphs, scale changes,
rotations). The head dot is hidden via opacity-0 so the face replaces it
visually.

\textbf{JSON shape:}

\begin{Shaded}
\begin{Highlighting}[]
\FunctionTok{\{}\DataTypeTok{"action"}\FunctionTok{:} \StringTok{"attach\_face"}\FunctionTok{,} \DataTypeTok{"who"}\FunctionTok{:} \StringTok{"freydoon"}\FunctionTok{,}
 \DataTypeTok{"image"}\FunctionTok{:} \StringTok{"bevers\_neutral\_face.png"}\FunctionTok{,} \DataTypeTok{"scale"}\FunctionTok{:} \DecValTok{0}\ErrorTok{.}\DecValTok{4}\FunctionTok{\}}
\end{Highlighting}
\end{Shaded}

\textbf{Parameters:}

\begin{longtable}[]{@{}
  >{\raggedright\arraybackslash}p{(\columnwidth - 8\tabcolsep) * \real{0.2000}}
  >{\raggedright\arraybackslash}p{(\columnwidth - 8\tabcolsep) * \real{0.2000}}
  >{\raggedright\arraybackslash}p{(\columnwidth - 8\tabcolsep) * \real{0.2000}}
  >{\raggedright\arraybackslash}p{(\columnwidth - 8\tabcolsep) * \real{0.2000}}
  >{\raggedright\arraybackslash}p{(\columnwidth - 8\tabcolsep) * \real{0.2000}}@{}}
\toprule\noalign{}
\begin{minipage}[b]{\linewidth}\raggedright
Key
\end{minipage} & \begin{minipage}[b]{\linewidth}\raggedright
Type
\end{minipage} & \begin{minipage}[b]{\linewidth}\raggedright
Required
\end{minipage} & \begin{minipage}[b]{\linewidth}\raggedright
Default
\end{minipage} & \begin{minipage}[b]{\linewidth}\raggedright
Description
\end{minipage} \\
\midrule\noalign{}
\endhead
\bottomrule\noalign{}
\endlastfoot
\texttt{who} & \texttt{str} & yes & --- & Character key. \\
\texttt{image} & \texttt{str} & yes & --- & Filename in
\texttt{pam/assets/}. PNG, JPG, or JPEG. Resolved to
\texttt{\textless{}pam\ package\ dir\textgreater{}/assets/\textless{}filename\textgreater{}}. \\
\texttt{scale} & \texttt{float} & no & \texttt{1.0} & Explicit scale
factor applied to the loaded image. Tune per face graphic. \\
\end{longtable}

\textbf{Behavior:}

\begin{itemize}
\tightlist
\item
  The head dot's VGroup opacity is set to 0 --- both circle and label
  are hidden. The dot itself persists as the position anchor for the
  updater and for \texttt{fig.say()} bubble anchoring.
\item
  The face mobject is stored on \texttt{fig.head\_face}, with the
  updater on \texttt{fig.head\_face\_updater}.
\item
  If a face is already attached, the old one is torn down (updater
  removed, mobject removed from scene, head dot opacity restored) before
  the new one is attached. This gives \texttt{change\_face} semantics
  for free --- re-call \texttt{attach\_face} with a different filename
  to swap.
\item
  No fade-in animation --- the face snaps in instantly. Parallel-safe.
\end{itemize}

\textbf{Silent skips (graceful no-ops with console warning):}

\begin{itemize}
\tightlist
\item
  Character has no \texttt{"head"} entry in \texttt{fig.dots}
  (quadrupeds, future body types).
\item
  \texttt{image} key missing.
\item
  File extension not in PNG/JPG/JPEG.
\item
  File not found at the resolved path.
\end{itemize}

\textbf{Persistence:} faces do \textbf{not} persist across
\texttt{fade\_out} + \texttt{fade\_in}. \texttt{fade\_out} tears the
face down concurrently with the body fade; a subsequent
\texttt{fade\_in} builds a fresh figure with no face. Re-call
\texttt{attach\_face} after \texttt{fade\_in} if you want the face back.

\textbf{Rotation:} the face tracks head \emph{position} but not head
\emph{orientation}. A character rotated via §2.5 keeps an upright face
by design (cartoon convention). Out-of-scope for the v0.9.15 MVP; an
opt-in \texttt{rotate\_with\_body:\ true} parameter is on the back
burner.

\textbf{Examples:}

\begin{Shaded}
\begin{Highlighting}[]
\FunctionTok{\{}\DataTypeTok{"action"}\FunctionTok{:} \StringTok{"attach\_face"}\FunctionTok{,} \DataTypeTok{"who"}\FunctionTok{:} \StringTok{"freydoon"}\FunctionTok{,}
 \DataTypeTok{"image"}\FunctionTok{:} \StringTok{"bevers\_neutral\_face.png"}\FunctionTok{,} \DataTypeTok{"scale"}\FunctionTok{:} \DecValTok{0}\ErrorTok{.}\DecValTok{4}\FunctionTok{\}}

\FunctionTok{\{}\DataTypeTok{"action"}\FunctionTok{:} \StringTok{"attach\_face"}\FunctionTok{,} \DataTypeTok{"who"}\FunctionTok{:} \StringTok{"freydoon"}\FunctionTok{,}
 \DataTypeTok{"image"}\FunctionTok{:} \StringTok{"bevers\_worried\_face.png"}\FunctionTok{,} \DataTypeTok{"scale"}\FunctionTok{:} \DecValTok{0}\ErrorTok{.}\DecValTok{4}\FunctionTok{\}}

\FunctionTok{\{}\DataTypeTok{"action"}\FunctionTok{:} \StringTok{"attach\_face"}\FunctionTok{,} \DataTypeTok{"who"}\FunctionTok{:} \StringTok{"thalia"}\FunctionTok{,}
 \DataTypeTok{"image"}\FunctionTok{:} \StringTok{"thalia\_smirk.png"}\FunctionTok{,} \DataTypeTok{"scale"}\FunctionTok{:} \DecValTok{0}\ErrorTok{.}\DecValTok{5}\FunctionTok{\}}
\end{Highlighting}
\end{Shaded}

\textbf{Asset path convention:} drop face PNGs into
\texttt{pam/assets/}; reference them by bare filename. The handler
resolves to
\texttt{\textless{}pam\ pkg\ dir\textgreater{}/assets/\textless{}filename\textgreater{}}.

\textbf{Known Manim limitation (Manim CE 0.20.1, Cairo backend):} RGBA
PNGs with transparent corners render with a visible alpha-channel
artifact instead of compositing cleanly with the scene background.
Workaround: produce opaque PNGs with the corners filled in the scene's
\texttt{BG\_COLOR} (\texttt{\#0a0e1a} by default). See the demo's
\texttt{make\_demo\_faces.py} for the pattern.

\begin{center}\rule{0.5\linewidth}{0.5pt}\end{center}

\subsubsection{\texorpdfstring{2.8
\texttt{detach\_face}}{2.8 detach\_face}}\label{detach_face}

\emph{Parallel-safe:} yes (returns empty animation list) ·
\emph{Animation:} none · \emph{Required keys:} \texttt{who} ·
\emph{Version:} v0.9.15

Remove an attached face: stop the updater, remove the image mobject from
the scene, restore head dot visibility. No-op if no face is attached.

\textbf{JSON shape:}

\begin{Shaded}
\begin{Highlighting}[]
\FunctionTok{\{}\DataTypeTok{"action"}\FunctionTok{:} \StringTok{"detach\_face"}\FunctionTok{,} \DataTypeTok{"who"}\FunctionTok{:} \StringTok{"freydoon"}\FunctionTok{\}}
\end{Highlighting}
\end{Shaded}

\textbf{Parameters:}

\begin{longtable}[]{@{}lllll@{}}
\toprule\noalign{}
Key & Type & Required & Default & Description \\
\midrule\noalign{}
\endhead
\bottomrule\noalign{}
\endlastfoot
\texttt{who} & \texttt{str} & yes & --- & Character key. \\
\end{longtable}

\textbf{Use cases:}

\begin{itemize}
\tightlist
\item
  Consciousness-departs-but-body-remains beat (the avatar leaves
  Freydoon's body, but Freydoon stays on screen).
\item
  Test teardown --- useful in demos where you want to verify the head
  dot returns to full visibility.
\end{itemize}

Most authoring doesn't need \texttt{detach\_face} --- \texttt{fade\_out}
auto-removes any attached face.

\begin{center}\rule{0.5\linewidth}{0.5pt}\end{center}

\subsubsection{\texorpdfstring{2.9
\texttt{focus}}{2.9 focus}}\label{focus}

\emph{Parallel-safe:} no · \emph{Animation:} timed · \emph{Required
keys:} \texttt{on} or \texttt{dim}

Dim background characters and/or brighten foreground ones to direct
audience attention. Animates each character's group opacity. Parsed from
Fountain+ \texttt{FOCUS:} keys by \texttt{fountain2pam.py}.

\textbf{JSON shape:}

\begin{Shaded}
\begin{Highlighting}[]
\FunctionTok{\{}\DataTypeTok{"action"}\FunctionTok{:} \StringTok{"focus"}\FunctionTok{,} \DataTypeTok{"on"}\FunctionTok{:} \OtherTok{[}\StringTok{"bevers"}\OtherTok{]}\FunctionTok{,} \DataTypeTok{"dim"}\FunctionTok{:} \StringTok{"all\_others"}\FunctionTok{,} \DataTypeTok{"rt"}\FunctionTok{:} \DecValTok{0}\ErrorTok{.}\DecValTok{4}\FunctionTok{\}}
\end{Highlighting}
\end{Shaded}

\textbf{Parameters:}

\begin{longtable}[]{@{}
  >{\raggedright\arraybackslash}p{(\columnwidth - 8\tabcolsep) * \real{0.2000}}
  >{\raggedright\arraybackslash}p{(\columnwidth - 8\tabcolsep) * \real{0.2000}}
  >{\raggedright\arraybackslash}p{(\columnwidth - 8\tabcolsep) * \real{0.2000}}
  >{\raggedright\arraybackslash}p{(\columnwidth - 8\tabcolsep) * \real{0.2000}}
  >{\raggedright\arraybackslash}p{(\columnwidth - 8\tabcolsep) * \real{0.2000}}@{}}
\toprule\noalign{}
\begin{minipage}[b]{\linewidth}\raggedright
Key
\end{minipage} & \begin{minipage}[b]{\linewidth}\raggedright
Type
\end{minipage} & \begin{minipage}[b]{\linewidth}\raggedright
Required
\end{minipage} & \begin{minipage}[b]{\linewidth}\raggedright
Default
\end{minipage} & \begin{minipage}[b]{\linewidth}\raggedright
Description
\end{minipage} \\
\midrule\noalign{}
\endhead
\bottomrule\noalign{}
\endlastfoot
\texttt{on} & \texttt{list{[}str{]}} or \texttt{{[}"all"{]}} & one of &
\texttt{{[}{]}} & Character keys to keep bright. The literal
\texttt{{[}"all"{]}} triggers a reset (equivalent to
\texttt{focus\_reset}). \\
\texttt{dim} & \texttt{list{[}str{]}} or \texttt{"all\_others"} & one of
& \texttt{{[}{]}} & Character keys to dim. The string
\texttt{"all\_others"} expands to every cast member not in
\texttt{on}. \\
\texttt{opacity} & \texttt{float} & no & \texttt{0.30} & Target opacity
for dimmed characters. \\
\texttt{bright} & \texttt{float} & no & \texttt{1.0} & Target opacity
for focused characters. \\
\texttt{rt} & \texttt{float} & no & \texttt{0.4} & Transition
duration. \\
\end{longtable}

\textbf{Behavior:}

\begin{itemize}
\tightlist
\item
  Each character's resulting opacity is stored on
  \texttt{cast{[}name{]}{[}"opacity"{]}} so subsequent \texttt{focus}
  calls can diff against it.
\item
  The animation is a single \texttt{play()} call combining all character
  opacity animations concurrently.
\end{itemize}

\textbf{Examples:}

\begin{Shaded}
\begin{Highlighting}[]
\FunctionTok{\{}\DataTypeTok{"action"}\FunctionTok{:} \StringTok{"focus"}\FunctionTok{,} \DataTypeTok{"on"}\FunctionTok{:} \OtherTok{[}\StringTok{"thalia"}\OtherTok{]}\FunctionTok{,} \DataTypeTok{"dim"}\FunctionTok{:} \StringTok{"all\_others"}\FunctionTok{\}}
\FunctionTok{\{}\DataTypeTok{"action"}\FunctionTok{:} \StringTok{"focus"}\FunctionTok{,} \DataTypeTok{"on"}\FunctionTok{:} \OtherTok{[}\StringTok{"bevers"}\OtherTok{,} \StringTok{"athena"}\OtherTok{]}\FunctionTok{,} \DataTypeTok{"dim"}\FunctionTok{:} \OtherTok{[}\StringTok{"freydoon"}\OtherTok{,} \StringTok{"brad"}\OtherTok{]}\FunctionTok{\}}
\FunctionTok{\{}\DataTypeTok{"action"}\FunctionTok{:} \StringTok{"focus"}\FunctionTok{,} \DataTypeTok{"on"}\FunctionTok{:} \OtherTok{[}\StringTok{"nona"}\OtherTok{]}\FunctionTok{,} \DataTypeTok{"dim"}\FunctionTok{:} \StringTok{"all\_others"}\FunctionTok{,} \DataTypeTok{"opacity"}\FunctionTok{:} \DecValTok{0}\ErrorTok{.}\DecValTok{15}\FunctionTok{,} \DataTypeTok{"rt"}\FunctionTok{:} \DecValTok{0}\ErrorTok{.}\DecValTok{6}\FunctionTok{\}}
\end{Highlighting}
\end{Shaded}

\textbf{Important interaction with speech bubbles:} see §2.13.

\begin{center}\rule{0.5\linewidth}{0.5pt}\end{center}

\subsubsection{\texorpdfstring{2.10
\texttt{focus\_reset}}{2.10 focus\_reset}}\label{focus_reset}

\emph{Parallel-safe:} no · \emph{Animation:} timed · \emph{Required
keys:} none

Restore all characters to full opacity in a single concurrent animation.

\textbf{JSON shape:}

\begin{Shaded}
\begin{Highlighting}[]
\FunctionTok{\{}\DataTypeTok{"action"}\FunctionTok{:} \StringTok{"focus\_reset"}\FunctionTok{,} \DataTypeTok{"rt"}\FunctionTok{:} \DecValTok{0}\ErrorTok{.}\DecValTok{4}\FunctionTok{\}}
\end{Highlighting}
\end{Shaded}

\textbf{Parameters:}

\begin{longtable}[]{@{}lllll@{}}
\toprule\noalign{}
Key & Type & Required & Default & Description \\
\midrule\noalign{}
\endhead
\bottomrule\noalign{}
\endlastfoot
\texttt{rt} & \texttt{float} & no & \texttt{0.4} & Restore duration. \\
\end{longtable}

Equivalent to \texttt{\{"action":\ "focus",\ "on":\ {[}"all"{]}\}}.
\texttt{cast{[}name{]}{[}"opacity"{]}} is reset to \texttt{1.0} for
every cast member.

\textbf{Important interaction with persistent speech bubbles:} see
§2.13.

\begin{center}\rule{0.5\linewidth}{0.5pt}\end{center}

\subsubsection{\texorpdfstring{2.11
\texttt{clear\_bubble}}{2.11 clear\_bubble}}\label{clear_bubble}

\emph{Parallel-safe:} no · \emph{Animation:} timed · \emph{Required
keys:} \texttt{who} or \texttt{prop}

Dismiss a single persistent bubble (one that was created by a
\texttt{say} or \texttt{prop\_say} step with \texttt{"persist":\ true}
or a non-zero \texttt{"duration"}).

\textbf{JSON shape:}

\begin{Shaded}
\begin{Highlighting}[]
\FunctionTok{\{}\DataTypeTok{"action"}\FunctionTok{:} \StringTok{"clear\_bubble"}\FunctionTok{,} \DataTypeTok{"who"}\FunctionTok{:} \StringTok{"bevers"}\FunctionTok{,} \DataTypeTok{"rt"}\FunctionTok{:} \DecValTok{0}\ErrorTok{.}\DecValTok{25}\FunctionTok{\}}
\end{Highlighting}
\end{Shaded}

\textbf{Parameters:}

\begin{longtable}[]{@{}
  >{\raggedright\arraybackslash}p{(\columnwidth - 8\tabcolsep) * \real{0.2000}}
  >{\raggedright\arraybackslash}p{(\columnwidth - 8\tabcolsep) * \real{0.2000}}
  >{\raggedright\arraybackslash}p{(\columnwidth - 8\tabcolsep) * \real{0.2000}}
  >{\raggedright\arraybackslash}p{(\columnwidth - 8\tabcolsep) * \real{0.2000}}
  >{\raggedright\arraybackslash}p{(\columnwidth - 8\tabcolsep) * \real{0.2000}}@{}}
\toprule\noalign{}
\begin{minipage}[b]{\linewidth}\raggedright
Key
\end{minipage} & \begin{minipage}[b]{\linewidth}\raggedright
Type
\end{minipage} & \begin{minipage}[b]{\linewidth}\raggedright
Required
\end{minipage} & \begin{minipage}[b]{\linewidth}\raggedright
Default
\end{minipage} & \begin{minipage}[b]{\linewidth}\raggedright
Description
\end{minipage} \\
\midrule\noalign{}
\endhead
\bottomrule\noalign{}
\endlastfoot
\texttt{who} & \texttt{str} & one of & --- & Character whose persistent
bubble should be cleared. \\
\texttt{prop} & \texttt{str} & one of & --- & Prop whose persistent
bubble should be cleared. An alias \texttt{clear\_prop\_bubble} is also
accepted. \\
\texttt{rt} & \texttt{float} & no & per-bubble \texttt{pam\_rt\_out},
else \texttt{0.25} & Fade-out duration. \\
\texttt{run\_time} & \texttt{float} & no & --- & Synonym for
\texttt{rt}. \\
\end{longtable}

\textbf{Behavior:}

\begin{itemize}
\tightlist
\item
  The named bubble is removed from the internal
  \texttt{\_persistent\_bubbles} dict and faded out.
\item
  Silent no-op if no bubble is registered under the given key --- does
  not crash a scene with a stale or doubled clear.
\item
  The bubble registry is global to the scene; bubbles persist across
  subscene boundaries unless explicitly cleared.
\end{itemize}

\textbf{No \texttt{who}/\texttt{prop}:} the action silently no-ops
rather than crashing.

\begin{center}\rule{0.5\linewidth}{0.5pt}\end{center}

\subsubsection{\texorpdfstring{2.12
\texttt{clear\_all\_bubbles}}{2.12 clear\_all\_bubbles}}\label{clear_all_bubbles}

\emph{Parallel-safe:} no · \emph{Animation:} timed · \emph{Required
keys:} none

Dismiss every persistent bubble currently on screen.

\textbf{JSON shape:}

\begin{Shaded}
\begin{Highlighting}[]
\FunctionTok{\{}\DataTypeTok{"action"}\FunctionTok{:} \StringTok{"clear\_all\_bubbles"}\FunctionTok{,} \DataTypeTok{"rt"}\FunctionTok{:} \DecValTok{0}\ErrorTok{.}\DecValTok{25}\FunctionTok{\}}
\end{Highlighting}
\end{Shaded}

\textbf{Parameters:}

\begin{longtable}[]{@{}
  >{\raggedright\arraybackslash}p{(\columnwidth - 8\tabcolsep) * \real{0.2000}}
  >{\raggedright\arraybackslash}p{(\columnwidth - 8\tabcolsep) * \real{0.2000}}
  >{\raggedright\arraybackslash}p{(\columnwidth - 8\tabcolsep) * \real{0.2000}}
  >{\raggedright\arraybackslash}p{(\columnwidth - 8\tabcolsep) * \real{0.2000}}
  >{\raggedright\arraybackslash}p{(\columnwidth - 8\tabcolsep) * \real{0.2000}}@{}}
\toprule\noalign{}
\begin{minipage}[b]{\linewidth}\raggedright
Key
\end{minipage} & \begin{minipage}[b]{\linewidth}\raggedright
Type
\end{minipage} & \begin{minipage}[b]{\linewidth}\raggedright
Required
\end{minipage} & \begin{minipage}[b]{\linewidth}\raggedright
Default
\end{minipage} & \begin{minipage}[b]{\linewidth}\raggedright
Description
\end{minipage} \\
\midrule\noalign{}
\endhead
\bottomrule\noalign{}
\endlastfoot
\texttt{rt} & \texttt{float} & no & \texttt{0.25} & Fade-out duration.
Applied uniformly --- per-bubble \texttt{pam\_rt\_out} is ignored
because the fades run concurrently. \\
\texttt{run\_time} & \texttt{float} & no & --- & Synonym for
\texttt{rt}. \\
\end{longtable}

\textbf{Use cases:}

\begin{itemize}
\tightlist
\item
  Scene or subscene boundaries. Subscene transitions do not auto-clear
  persistent bubbles, so an explicit \texttt{clear\_all\_bubbles} at the
  end of a beat is the documented cleanup pattern.
\item
  Before any \texttt{focus\_reset} if persistent bubbles are on screen
  --- see §2.13.
\end{itemize}

\begin{center}\rule{0.5\linewidth}{0.5pt}\end{center}

\subsubsection{\texorpdfstring{2.13 The bubble / \texttt{focus\_reset}
ordering
gotcha}{2.13 The bubble / focus\_reset ordering gotcha}}\label{the-bubble-focus_reset-ordering-gotcha}

\emph{Documented in v0.9.10. Resolved-status: footgun pending fix.}

\texttt{focus\_reset} (and \texttt{focus} with
\texttt{on:\ {[}"all"{]}}) restores character opacity and z-order in a
single \texttt{play()} call. When a \emph{persistent} bubble is on
screen at the moment of reset, the restored character can be drawn on
top of the bubble, partially or fully occluding it.

\textbf{The mechanism:} persistent bubbles are added to the scene via
\texttt{scene.add(bubble)} before \texttt{focus\_reset} runs.
\texttt{focus\_reset} animates opacity on figure groups, which Manim
treats as updating the existing mobject in-place. If the figure's
z-order index is higher than the bubble's, the figure draws on top.

\textbf{Authoring workaround:} clear the bubble before the focus reset.

\begin{Shaded}
\begin{Highlighting}[]
\FunctionTok{\{}\DataTypeTok{"action"}\FunctionTok{:} \StringTok{"clear\_bubble"}\FunctionTok{,} \DataTypeTok{"who"}\FunctionTok{:} \StringTok{"bevers"}\FunctionTok{\}}\ErrorTok{,}
\FunctionTok{\{}\DataTypeTok{"action"}\FunctionTok{:} \StringTok{"focus\_reset"}\FunctionTok{\}}
\end{Highlighting}
\end{Shaded}

\ldots or use \texttt{clear\_all\_bubbles} if multiple persistent
bubbles are active:

\begin{Shaded}
\begin{Highlighting}[]
\FunctionTok{\{}\DataTypeTok{"action"}\FunctionTok{:} \StringTok{"clear\_all\_bubbles"}\FunctionTok{,} \DataTypeTok{"rt"}\FunctionTok{:} \DecValTok{0}\ErrorTok{.}\DecValTok{2}\FunctionTok{\}}\ErrorTok{,}
\FunctionTok{\{}\DataTypeTok{"action"}\FunctionTok{:} \StringTok{"focus\_reset"}\FunctionTok{,} \DataTypeTok{"rt"}\FunctionTok{:} \DecValTok{0}\ErrorTok{.}\DecValTok{4}\FunctionTok{\}}
\end{Highlighting}
\end{Shaded}

\textbf{Alternative workaround --- \texttt{say} with \texttt{duration}:}
if the persistent bubble's purpose is ``stay visible during the focus
shot,'' set \texttt{"duration"} on the \texttt{say} step to a value that
expires \emph{before} the \texttt{focus\_reset}. Avoids needing an
explicit \texttt{clear\_bubble}.

\textbf{Why not \texttt{say} with \texttt{persist} in walk-and-talk:}
persistent bubbles anchor to world-space coordinates at the moment
they're created. A character that walks under a persistent bubble leaves
the bubble behind --- it doesn't track the head dot like
\texttt{attach\_face} does. This is a separate limitation; the
head-anchored-bubble approach is on the back burner.

\begin{center}\rule{0.5\linewidth}{0.5pt}\end{center}

\emph{Section 3 (action targets) and Section 4 (meta) are delivered in
subsequent substeps.}

\subsection{3. Action targets}\label{action-targets}

A character can serve as the target of any of the verbs documented in
this section. Coverage is organized functionally:

\begin{itemize}
\tightlist
\item
  \textbf{§3.1} --- Speech: \texttt{say}, \texttt{prop\_say}.
\item
  \textbf{§3.2} --- Locomotion (single character): \texttt{walk\_to},
  \texttt{run\_to}, \texttt{rush\_to}, \texttt{rush\_out},
  \texttt{trot\_to}, \texttt{walk\_to\_prop}, \texttt{run\_to\_prop},
  \texttt{squeeze\_through}, \texttt{exit\_through},
  \texttt{exit\_through\_doors}, \texttt{jump\_up}, \texttt{fall\_down},
  \texttt{dodge}.
\item
  \textbf{§3.3} --- Posture transitions: \texttt{sit\_down},
  \texttt{stand\_up}.
\item
  \textbf{§3.4} --- Orientation: \texttt{turn}, \texttt{face}.
\item
  \textbf{§3.5} --- Expressive gestures: \texttt{wave},
  \texttt{wave\_left}, \texttt{wave\_right}, \texttt{nod},
  \texttt{shake\_head}, \texttt{shrug}, \texttt{point\_at},
  \texttt{express}, \texttt{react}.
\item
  \textbf{§3.6} --- Prop interaction: \texttt{grab}, \texttt{pick\_up},
  \texttt{pick\_up\_phone}, \texttt{hang\_up}, \texttt{place\_on},
  \texttt{put\_down}, \texttt{peel\_from\_hand}, \texttt{punch\_button},
  \texttt{reach\_for}, \texttt{search\_drawers}, \texttt{snap\_photo},
  \texttt{stick\_to}, \texttt{move\_aside}, \texttt{carry}.
\item
  \textbf{§3.7} --- Two-figure interaction: \texttt{kiss},
  \texttt{hold\_hands}, \texttt{pat}, \texttt{pat\_head},
  \texttt{hand\_to}, \texttt{grab\_arm}, \texttt{twist\_arm\_behind},
  \texttt{release\_arm}, \texttt{reach\_character}.
\item
  \textbf{§3.8} --- Multi-target: \texttt{group\_translate}.
\end{itemize}

Section conventions established in §2 carry forward: each action shows
\emph{Parallel-safe}, \emph{Animation}, and \emph{Required keys}. Where
an action wraps an internal animation, the \textbf{default \texttt{rt}}
is documented per action and is usually overridable.

\subsubsection{Shared conventions}\label{shared-conventions}

\begin{itemize}
\tightlist
\item
  \textbf{\texttt{who}} --- always a character key from the cast block.
  When omitted from a step's docs below, it is required.
\item
  \textbf{\texttt{prop}} --- a key from the prop registry. Used when the
  action targets a prop (\texttt{grab}, \texttt{walk\_to\_prop}, etc.).
\item
  \textbf{\texttt{target}} --- either a prop key or a character key,
  depending on the action. Some actions accept both
  (e.g.~\texttt{reach\_for}, \texttt{face}); others restrict to props
  only (\texttt{stick\_to}) or characters only
  (\texttt{reach\_character}).
\item
  \textbf{\texttt{rt}} --- animation run time in seconds. Defaults vary
  per action; documented per action below.
\item
  \textbf{\texttt{hold}} --- pause time in seconds after the main
  animation, before returning to a rest pose.
\end{itemize}

\subsubsection{The parallel-unsafe set}\label{the-parallel-unsafe-set}

The following actions print a console warning and run sequentially when
nested inside a \texttt{parallel} block (from
\texttt{\_CANNOT\_PARALLEL} in \texttt{actions.py}):

\begin{verbatim}
walk_to, run_to, sit_down, stand_up, wave, wave_left, wave_right,
carry, exit_through, exit_through_doors, rush_to, rush_out,
squeeze_through, jump_up, pat, search_drawers, pick_up_phone,
hang_up, grab, punch_button, reach_character, peel_from_hand,
group_translate, nod, shake_head, shrug, grab_arm, twist_arm_behind,
release_arm, rotate
\end{verbatim}

Every other character-targeting action is parallel-safe in principle
(returns animations via \texttt{\_collect\_anims}). The
\texttt{parallel} + \texttt{say} constraint is documented separately in
§3.1.

\begin{center}\rule{0.5\linewidth}{0.5pt}\end{center}

\subsubsection{3.1 Speech}\label{speech}

\paragraph{\texorpdfstring{3.1.1 \texttt{say}}{3.1.1 say}}\label{say}

\emph{Parallel-safe:} \textbf{no} · \emph{Animation:} timed ·
\emph{Required keys:} \texttt{who}, \texttt{text}

Display a speech bubble above a character's head. The bubble's box is
automatically sized to fit the text, clamped to stay within frame
margins, and anchored to the head dot's current position at fade-in.

\textbf{JSON shape:}

\begin{Shaded}
\begin{Highlighting}[]
\FunctionTok{\{}\DataTypeTok{"action"}\FunctionTok{:} \StringTok{"say"}\FunctionTok{,} \DataTypeTok{"who"}\FunctionTok{:} \StringTok{"bevers"}\FunctionTok{,} \DataTypeTok{"text"}\FunctionTok{:} \StringTok{"Hello, world."}\FunctionTok{,}
 \DataTypeTok{"hold"}\FunctionTok{:} \FloatTok{1.5}\FunctionTok{,} \DataTypeTok{"side"}\FunctionTok{:} \StringTok{"right"}\FunctionTok{\}}
\end{Highlighting}
\end{Shaded}

\textbf{Parameters:}

\begin{longtable}[]{@{}
  >{\raggedright\arraybackslash}p{(\columnwidth - 8\tabcolsep) * \real{0.2000}}
  >{\raggedright\arraybackslash}p{(\columnwidth - 8\tabcolsep) * \real{0.2000}}
  >{\raggedright\arraybackslash}p{(\columnwidth - 8\tabcolsep) * \real{0.2000}}
  >{\raggedright\arraybackslash}p{(\columnwidth - 8\tabcolsep) * \real{0.2000}}
  >{\raggedright\arraybackslash}p{(\columnwidth - 8\tabcolsep) * \real{0.2000}}@{}}
\toprule\noalign{}
\begin{minipage}[b]{\linewidth}\raggedright
Key
\end{minipage} & \begin{minipage}[b]{\linewidth}\raggedright
Type
\end{minipage} & \begin{minipage}[b]{\linewidth}\raggedright
Required
\end{minipage} & \begin{minipage}[b]{\linewidth}\raggedright
Default
\end{minipage} & \begin{minipage}[b]{\linewidth}\raggedright
Description
\end{minipage} \\
\midrule\noalign{}
\endhead
\bottomrule\noalign{}
\endlastfoot
\texttt{who} & \texttt{str} & yes & --- & Character key. \\
\texttt{text} & \texttt{str} & yes & --- & Bubble text. Multi-line is
supported with \texttt{\textbackslash{}n}. \\
\texttt{hold} & \texttt{float} & no & \texttt{1.2} & Seconds to display
the bubble before fade-out. \\
\texttt{rt\_in} & \texttt{float} & no & \texttt{0.4} & Fade-in
duration. \\
\texttt{rt\_out} & \texttt{float} & no & \texttt{0.3} & Fade-out
duration. \\
\texttt{side} & \texttt{str} & no & \texttt{"right"} & Bubble placement
relative to head: \texttt{"left"} or \texttt{"right"}. \\
\texttt{font\_size} & \texttt{int} & no & \texttt{20} & Text font
size. \\
\texttt{style} & \texttt{str} & no & \texttt{"normal"} &
\texttt{"normal"} (solid bubble) or \texttt{"os"} / \texttt{"phone"}
(dashed border for off-screen / phone dialogue). \\
\texttt{persist} & \texttt{bool} & no & \texttt{false} & If true, bubble
stays on screen until a \texttt{clear\_bubble} action. See §2.11. \\
\texttt{duration} & \texttt{float} & no & --- & If set, bubble stays for
\texttt{duration} seconds (measured from fade-in start), then
auto-clears at the next action. \\
\end{longtable}

\textbf{Behavior:}

\begin{itemize}
\tightlist
\item
  The bubble anchors to the head dot at the moment of fade-in. It does
  \textbf{not} track a moving figure --- see ``parallel + say'' below.
\item
  If \texttt{persist} or \texttt{duration} is set, the bubble is added
  to the \texttt{\_persistent\_bubbles} registry under
  \texttt{char:\{who\}}. Subsequent \texttt{say} calls on the same
  character fade the prior bubble out first to avoid stacking.
\item
  Tic system integration (v0.9.14): a \texttt{say} step automatically
  fires \texttt{sentence\_end} triggers in the character's
  \texttt{tic\_profile}. Set \texttt{"suppress\_tics":\ true} on the
  step to skip.
\end{itemize}

\textbf{Parallel + say:} \texttt{say} is not in
\texttt{\_CANNOT\_PARALLEL}, but it cannot be nested inside a
\texttt{parallel} block at the screenplay level --- the player's
parallel composer treats \texttt{say} specially and refuses it. This is
tribal knowledge worth a formal docs callout (back-burner).

\textbf{Speech bubble anchoring caveat:} the bubble is positioned in
world space at the moment of creation. A character who walks while a
\texttt{persist} bubble is up will leave the bubble behind. For
walk-and-talk, interleave \texttt{walk\_to} and short non-persist
\texttt{say} steps rather than running a long persistent bubble; the
head-anchored-bubble design for walk-and-talk is on the back burner.

\textbf{O.S. / phone bubbles:}

\begin{Shaded}
\begin{Highlighting}[]
\FunctionTok{\{}\DataTypeTok{"action"}\FunctionTok{:} \StringTok{"say"}\FunctionTok{,} \DataTypeTok{"who"}\FunctionTok{:} \StringTok{"freydoon"}\FunctionTok{,}
 \DataTypeTok{"text"}\FunctionTok{:} \StringTok{"(over the phone) I\textquotesingle{}m on my way."}\FunctionTok{,} \DataTypeTok{"style"}\FunctionTok{:} \StringTok{"os"}\FunctionTok{\}}
\end{Highlighting}
\end{Shaded}

The dashed border reads as off-screen / disembodied dialogue. Use
sparingly --- readable but visually distinct.

\textbf{Bubble for a horizontal character:} the named poses
\texttt{lying\_flat\_right} and \texttt{lying\_flat\_left} (v0.9.14) are
required for correct bubble anchoring when a character is lying down. An
ad-hoc \texttt{rotate} to horizontal will mis-anchor the bubble.

\begin{center}\rule{0.5\linewidth}{0.5pt}\end{center}

\paragraph{\texorpdfstring{3.1.2
\texttt{prop\_say}}{3.1.2 prop\_say}}\label{prop_say}

\emph{Parallel-safe:} no · \emph{Animation:} timed · \emph{Required
keys:} \texttt{prop}, \texttt{text}

Display a speech bubble next to a prop. Used for prop-characters
(\texttt{GovernorGraph}, \texttt{DogGraph}, news anchors on TV monitors)
that have no head joint for \texttt{say} to anchor to.

\textbf{JSON shape:}

\begin{Shaded}
\begin{Highlighting}[]
\FunctionTok{\{}\DataTypeTok{"action"}\FunctionTok{:} \StringTok{"prop\_say"}\FunctionTok{,} \DataTypeTok{"prop"}\FunctionTok{:} \StringTok{"governor"}\FunctionTok{,}
 \DataTypeTok{"text"}\FunctionTok{:} \StringTok{"Approved."}\FunctionTok{,} \DataTypeTok{"hold"}\FunctionTok{:} \FloatTok{1.2}\FunctionTok{\}}
\end{Highlighting}
\end{Shaded}

\textbf{Parameters:}

\begin{longtable}[]{@{}
  >{\raggedright\arraybackslash}p{(\columnwidth - 8\tabcolsep) * \real{0.2000}}
  >{\raggedright\arraybackslash}p{(\columnwidth - 8\tabcolsep) * \real{0.2000}}
  >{\raggedright\arraybackslash}p{(\columnwidth - 8\tabcolsep) * \real{0.2000}}
  >{\raggedright\arraybackslash}p{(\columnwidth - 8\tabcolsep) * \real{0.2000}}
  >{\raggedright\arraybackslash}p{(\columnwidth - 8\tabcolsep) * \real{0.2000}}@{}}
\toprule\noalign{}
\begin{minipage}[b]{\linewidth}\raggedright
Key
\end{minipage} & \begin{minipage}[b]{\linewidth}\raggedright
Type
\end{minipage} & \begin{minipage}[b]{\linewidth}\raggedright
Required
\end{minipage} & \begin{minipage}[b]{\linewidth}\raggedright
Default
\end{minipage} & \begin{minipage}[b]{\linewidth}\raggedright
Description
\end{minipage} \\
\midrule\noalign{}
\endhead
\bottomrule\noalign{}
\endlastfoot
\texttt{prop} & \texttt{str} & yes & --- & Prop registry key. Path C
dogs (\texttt{who:\ chekov} or \texttt{prop:\ chekov}) resolve via
either key. \\
\texttt{text} & \texttt{str} & yes & --- & Bubble text. \\
\texttt{hold} & \texttt{float} & no & \texttt{1.2} & Seconds to
display. \\
\texttt{style} & \texttt{str} & no & \texttt{"normal"} &
\texttt{"normal"} or \texttt{"os"} (dashed border). \\
\texttt{y\_offset} & \texttt{float} & no & \texttt{0.0} & \emph{(v0.9.18)}
Vertical nudge applied on top of the geometry-aware default. Currently
consumed only by \texttt{GovernorGraph}; \texttt{DogGraph} and the
generic-prop branch ignore it. Positive lifts the bubble, negative
lowers it. \\
\texttt{persist}, \texttt{duration}, \texttt{rt\_in}, \texttt{rt\_out},
\texttt{font\_size} & --- & --- & --- & Same as \texttt{say}. \\
\end{longtable}

\textbf{Behavior:}

\begin{itemize}
\tightlist
\item
  For \texttt{GovernorGraph} and \texttt{DogGraph}, \texttt{prop\_say}
  uses the figure's own bubble layout (these classes implement their own
  bubble geometry).
\item
  \emph{(v0.9.18)} \texttt{GovernorGraph} bubble placement is now
  geometry-aware. The bubble's bottom edge sits \texttt{0.10} units
  above the dodecahedron's top
  (\texttt{self.\_y\ +\ self.\_radius}), regardless of prop size, so
  the bubble grows away from the prop instead of overlapping it. The
  tail tip lands inside the top portion of the dodecahedron, which
  is the comic-book convention for ``bubble points at speaker.''
  Previously the bubble center was hardcoded at \texttt{self.\_y\ +\ 0.3},
  which for a standard-sized Governor put the tail tip below the
  prop's centre. Use \texttt{y\_offset} to nudge the result up or
  down if a specific scene still reads off.
\item
  For generic props (tv\_monitor, cellphone, etc.), the bubble is built
  manually with placement near the prop's top surface. Bubble color
  resolves in priority order: explicit \texttt{bubble\_color} /
  \texttt{text\_color} on the step -\textgreater{} parent character's
  palette (if the prop has \texttt{pam\_follows}) -\textgreater{} prop's
  \texttt{screen\_color} -\textgreater{} default tan.
\item
  Word wrapping in \texttt{prop\_say} is currently implemented for
  \texttt{DogGraph} and \texttt{GovernorGraph} only --- generic props
  receive un-wrapped text.
\end{itemize}

\textbf{Examples:}

\begin{Shaded}
\begin{Highlighting}[]
\FunctionTok{\{}\DataTypeTok{"action"}\FunctionTok{:} \StringTok{"prop\_say"}\FunctionTok{,} \DataTypeTok{"prop"}\FunctionTok{:} \StringTok{"chekov"}\FunctionTok{,}  \DataTypeTok{"text"}\FunctionTok{:} \StringTok{"Woof!"}\FunctionTok{\}}
\FunctionTok{\{}\DataTypeTok{"action"}\FunctionTok{:} \StringTok{"prop\_say"}\FunctionTok{,} \DataTypeTok{"prop"}\FunctionTok{:} \StringTok{"athena\_laptop"}\FunctionTok{,}
 \DataTypeTok{"text"}\FunctionTok{:} \StringTok{"WiFi scan complete."}\FunctionTok{,} \DataTypeTok{"hold"}\FunctionTok{:} \FloatTok{1.8}\FunctionTok{\}}
\FunctionTok{\{}\DataTypeTok{"action"}\FunctionTok{:} \StringTok{"prop\_say"}\FunctionTok{,} \DataTypeTok{"prop"}\FunctionTok{:} \StringTok{"news\_anchor"}\FunctionTok{,} \DataTypeTok{"text"}\FunctionTok{:} \StringTok{"Breaking news."}\FunctionTok{,}
 \DataTypeTok{"style"}\FunctionTok{:} \StringTok{"os"}\FunctionTok{\}}
\end{Highlighting}
\end{Shaded}

\textbf{Example --- Governor with bubble nudged up (v0.9.18).} A
camera angle that places the Governor near the bottom of the frame
sometimes leaves the default bubble too close to the lower edge.
\texttt{y\_offset} lifts it without changing the prop's position:

\begin{Shaded}
\begin{Highlighting}[]
\FunctionTok{\{}\DataTypeTok{"action"}\FunctionTok{:} \StringTok{"prop\_say"}\FunctionTok{,} \DataTypeTok{"prop"}\FunctionTok{:} \StringTok{"governor"}\FunctionTok{,}
 \DataTypeTok{"text"}\FunctionTok{:} \StringTok{"Decision rendered."}\FunctionTok{,}
 \DataTypeTok{"y\_offset"}\FunctionTok{:} \FloatTok{0.3}\FunctionTok{\}}
\end{Highlighting}
\end{Shaded}

\begin{center}\rule{0.5\linewidth}{0.5pt}\end{center}

\subsubsection{3.2 Locomotion (single
character)}\label{locomotion-single-character}

\paragraph{\texorpdfstring{3.2.1
\texttt{walk\_to}}{3.2.1 walk\_to}}\label{walk_to}

\emph{Parallel-safe:} no · \emph{Animation:} timed · \emph{Required
keys:} \texttt{who}, \texttt{x}

Walk the figure to an x position. Drives the four-keyframe walk cycle
for each step's worth of horizontal distance; pose must be a side-view
at start (or the character first morphs to \texttt{standing\_side}
automatically --- confirm against \texttt{figure.py}).

\textbf{Parameters:}

\begin{longtable}[]{@{}
  >{\raggedright\arraybackslash}p{(\columnwidth - 8\tabcolsep) * \real{0.2000}}
  >{\raggedright\arraybackslash}p{(\columnwidth - 8\tabcolsep) * \real{0.2000}}
  >{\raggedright\arraybackslash}p{(\columnwidth - 8\tabcolsep) * \real{0.2000}}
  >{\raggedright\arraybackslash}p{(\columnwidth - 8\tabcolsep) * \real{0.2000}}
  >{\raggedright\arraybackslash}p{(\columnwidth - 8\tabcolsep) * \real{0.2000}}@{}}
\toprule\noalign{}
\begin{minipage}[b]{\linewidth}\raggedright
Key
\end{minipage} & \begin{minipage}[b]{\linewidth}\raggedright
Type
\end{minipage} & \begin{minipage}[b]{\linewidth}\raggedright
Required
\end{minipage} & \begin{minipage}[b]{\linewidth}\raggedright
Default
\end{minipage} & \begin{minipage}[b]{\linewidth}\raggedright
Description
\end{minipage} \\
\midrule\noalign{}
\endhead
\bottomrule\noalign{}
\endlastfoot
\texttt{who} & \texttt{str} & yes & --- & Character key. \\
\texttt{x} & \texttt{float} & yes & --- & Target x in world coordinates.
Negative for left, positive for right. \\
\end{longtable}

\textbf{Behavior:}

\begin{itemize}
\tightlist
\item
  Walk speed is derived from the figure's internal stride length;
  \texttt{walk\_to} figures out how many steps to take based on the x
  delta.
\item
  Attached props update their positions via
  \texttt{\_drag\_attached\_props} so a character carrying a laptop
  keeps the laptop with them.
\item
  \texttt{walk\_to} updates \texttt{fig.offset} to the new x but does
  \textbf{not} write back to \texttt{cast{[}name{]}{[}"offset"{]}}.
  After a \texttt{fade\_out} + \texttt{fade\_in}, the character returns
  to the original cast-spec offset.
\end{itemize}

\textbf{Example:}

\begin{Shaded}
\begin{Highlighting}[]
\FunctionTok{\{}\DataTypeTok{"action"}\FunctionTok{:} \StringTok{"walk\_to"}\FunctionTok{,} \DataTypeTok{"who"}\FunctionTok{:} \StringTok{"bevers"}\FunctionTok{,} \DataTypeTok{"x"}\FunctionTok{:} \FloatTok{{-}2.5}\FunctionTok{\}}
\end{Highlighting}
\end{Shaded}

\begin{center}\rule{0.5\linewidth}{0.5pt}\end{center}

\paragraph{\texorpdfstring{3.2.2
\texttt{run\_to}}{3.2.2 run\_to}}\label{run_to}

\emph{Parallel-safe:} no · \emph{Animation:} timed · \emph{Required
keys:} \texttt{who}, \texttt{x}

Same as \texttt{walk\_to}, but uses the run cycle (faster, more dynamic
posture).

\begin{Shaded}
\begin{Highlighting}[]
\FunctionTok{\{}\DataTypeTok{"action"}\FunctionTok{:} \StringTok{"run\_to"}\FunctionTok{,} \DataTypeTok{"who"}\FunctionTok{:} \StringTok{"freydoon"}\FunctionTok{,} \DataTypeTok{"x"}\FunctionTok{:} \FloatTok{4.0}\FunctionTok{\}}
\end{Highlighting}
\end{Shaded}

\begin{center}\rule{0.5\linewidth}{0.5pt}\end{center}

\paragraph{\texorpdfstring{3.2.3
\texttt{rush\_to}}{3.2.3 rush\_to}}\label{rush_to}

\emph{Parallel-safe:} no · \emph{Animation:} timed · \emph{Required
keys:} \texttt{who}, \texttt{x}

Fast walk with a forward-lean pose. Visually reads as ``in a hurry''
without the full sprint of \texttt{run\_to}.

\begin{Shaded}
\begin{Highlighting}[]
\FunctionTok{\{}\DataTypeTok{"action"}\FunctionTok{:} \StringTok{"rush\_to"}\FunctionTok{,} \DataTypeTok{"who"}\FunctionTok{:} \StringTok{"thalia"}\FunctionTok{,} \DataTypeTok{"x"}\FunctionTok{:} \FloatTok{3.0}\FunctionTok{\}}
\end{Highlighting}
\end{Shaded}

\begin{center}\rule{0.5\linewidth}{0.5pt}\end{center}

\paragraph{\texorpdfstring{3.2.4
\texttt{rush\_out}}{3.2.4 rush\_out}}\label{rush_out}

\emph{Parallel-safe:} no · \emph{Animation:} timed · \emph{Required
keys:} \texttt{who}, \texttt{x}

Alias of \texttt{rush\_to} --- identical mechanics, name only differs
for readability when the destination is off-screen.

\begin{Shaded}
\begin{Highlighting}[]
\FunctionTok{\{}\DataTypeTok{"action"}\FunctionTok{:} \StringTok{"rush\_out"}\FunctionTok{,} \DataTypeTok{"who"}\FunctionTok{:} \StringTok{"sidel"}\FunctionTok{,} \DataTypeTok{"x"}\FunctionTok{:} \FloatTok{7.5}\FunctionTok{\}}
\end{Highlighting}
\end{Shaded}

\begin{center}\rule{0.5\linewidth}{0.5pt}\end{center}

\paragraph{\texorpdfstring{3.2.5
\texttt{trot\_to}}{3.2.5 trot\_to}}\label{trot_to}

\emph{Parallel-safe:} yes (handled specially in \texttt{parallel}) ·
\emph{Animation:} timed · \emph{Required keys:} \texttt{who} or
\texttt{prop}, \texttt{x} · \emph{Applies to:}
\texttt{figure\_type:\ "dog"}

Dog-only locomotion. Path C dual-registration (v0.9.14) means a dog can
be addressed via \texttt{who:\ chekov} or \texttt{prop:\ chekov};
\texttt{trot\_to} resolves to the same \texttt{DogGraph} either way.

\textbf{Parameters:}

\begin{longtable}[]{@{}
  >{\raggedright\arraybackslash}p{(\columnwidth - 8\tabcolsep) * \real{0.2000}}
  >{\raggedright\arraybackslash}p{(\columnwidth - 8\tabcolsep) * \real{0.2000}}
  >{\raggedright\arraybackslash}p{(\columnwidth - 8\tabcolsep) * \real{0.2000}}
  >{\raggedright\arraybackslash}p{(\columnwidth - 8\tabcolsep) * \real{0.2000}}
  >{\raggedright\arraybackslash}p{(\columnwidth - 8\tabcolsep) * \real{0.2000}}@{}}
\toprule\noalign{}
\begin{minipage}[b]{\linewidth}\raggedright
Key
\end{minipage} & \begin{minipage}[b]{\linewidth}\raggedright
Type
\end{minipage} & \begin{minipage}[b]{\linewidth}\raggedright
Required
\end{minipage} & \begin{minipage}[b]{\linewidth}\raggedright
Default
\end{minipage} & \begin{minipage}[b]{\linewidth}\raggedright
Description
\end{minipage} \\
\midrule\noalign{}
\endhead
\bottomrule\noalign{}
\endlastfoot
\texttt{who} & \texttt{str} & one of & --- & Dog's character key. \\
\texttt{prop} & \texttt{str} & one of & --- & Dog's prop key (same name
under dual registration). \\
\texttt{x} & \texttt{float} & yes & --- & Target x in world
coordinates. \\
\texttt{stride} & \texttt{float} & no & \texttt{0.14} & Step length per
trot keyframe. Shorter for slower, more delicate trot; longer for faster
gait. \\
\end{longtable}

\textbf{Example:}

\begin{Shaded}
\begin{Highlighting}[]
\FunctionTok{\{}\DataTypeTok{"action"}\FunctionTok{:} \StringTok{"trot\_to"}\FunctionTok{,} \DataTypeTok{"who"}\FunctionTok{:} \StringTok{"chekov"}\FunctionTok{,} \DataTypeTok{"x"}\FunctionTok{:} \FloatTok{1.5}\FunctionTok{,} \DataTypeTok{"stride"}\FunctionTok{:} \DecValTok{0}\ErrorTok{.}\DecValTok{18}\FunctionTok{\}}
\end{Highlighting}
\end{Shaded}

\textbf{Note:} there is no equivalent \texttt{walk\_to\_dog} or
\texttt{run\_to\_dog} --- dogs use a single locomotion verb
(\texttt{trot\_to}) with stride control. An author needing a slow dog
should reduce \texttt{stride}; for a fast dog, increase it.

\begin{center}\rule{0.5\linewidth}{0.5pt}\end{center}

\paragraph{\texorpdfstring{3.2.6
\texttt{walk\_to\_prop}}{3.2.6 walk\_to\_prop}}\label{walk_to_prop}

\emph{Parallel-safe:} yes · \emph{Animation:} timed · \emph{Required
keys:} \texttt{who}, \texttt{prop}

Walk to the x position of a named prop. Convenience wrapper ---
equivalent to \texttt{walk\_to} with \texttt{x} computed from the prop's
\texttt{pam\_x}.

\begin{Shaded}
\begin{Highlighting}[]
\FunctionTok{\{}\DataTypeTok{"action"}\FunctionTok{:} \StringTok{"walk\_to\_prop"}\FunctionTok{,} \DataTypeTok{"who"}\FunctionTok{:} \StringTok{"sidel"}\FunctionTok{,} \DataTypeTok{"prop"}\FunctionTok{:} \StringTok{"desk"}\FunctionTok{\}}
\end{Highlighting}
\end{Shaded}

\textbf{Note:} the character ends up at the prop's x --- not ``next to''
it. For an offset stand-near-prop position, follow with a
\texttt{walk\_to} to a slightly different x.

\begin{center}\rule{0.5\linewidth}{0.5pt}\end{center}

\paragraph{\texorpdfstring{3.2.7
\texttt{run\_to\_prop}}{3.2.7 run\_to\_prop}}\label{run_to_prop}

\emph{Parallel-safe:} yes · \emph{Animation:} timed · \emph{Required
keys:} \texttt{who}, \texttt{prop}

Same as \texttt{walk\_to\_prop} but using the run cycle.

\begin{Shaded}
\begin{Highlighting}[]
\FunctionTok{\{}\DataTypeTok{"action"}\FunctionTok{:} \StringTok{"run\_to\_prop"}\FunctionTok{,} \DataTypeTok{"who"}\FunctionTok{:} \StringTok{"thalia"}\FunctionTok{,} \DataTypeTok{"prop"}\FunctionTok{:} \StringTok{"exit\_door"}\FunctionTok{\}}
\end{Highlighting}
\end{Shaded}

\begin{center}\rule{0.5\linewidth}{0.5pt}\end{center}

\paragraph{\texorpdfstring{3.2.8
\texttt{squeeze\_through}}{3.2.8 squeeze\_through}}\label{squeeze_through}

\emph{Parallel-safe:} no · \emph{Animation:} timed · \emph{Required
keys:} \texttt{who}, \texttt{x}

Narrow-stance walk through a tight space (doorway, crowd gap). Uses the
\texttt{squeeze\_walk\_a} / \texttt{squeeze\_walk\_b} poses with arms
held close to the body.

\textbf{Parameters:}

\begin{longtable}[]{@{}
  >{\raggedright\arraybackslash}p{(\columnwidth - 8\tabcolsep) * \real{0.2000}}
  >{\raggedright\arraybackslash}p{(\columnwidth - 8\tabcolsep) * \real{0.2000}}
  >{\raggedright\arraybackslash}p{(\columnwidth - 8\tabcolsep) * \real{0.2000}}
  >{\raggedright\arraybackslash}p{(\columnwidth - 8\tabcolsep) * \real{0.2000}}
  >{\raggedright\arraybackslash}p{(\columnwidth - 8\tabcolsep) * \real{0.2000}}@{}}
\toprule\noalign{}
\begin{minipage}[b]{\linewidth}\raggedright
Key
\end{minipage} & \begin{minipage}[b]{\linewidth}\raggedright
Type
\end{minipage} & \begin{minipage}[b]{\linewidth}\raggedright
Required
\end{minipage} & \begin{minipage}[b]{\linewidth}\raggedright
Default
\end{minipage} & \begin{minipage}[b]{\linewidth}\raggedright
Description
\end{minipage} \\
\midrule\noalign{}
\endhead
\bottomrule\noalign{}
\endlastfoot
\texttt{who} & \texttt{str} & yes & --- & Character key. \\
\texttt{x} & \texttt{float} & yes & --- & Target x. \\
\texttt{rt} & \texttt{float} & no & per cycle & Per-keyframe run time.
Tune for tightness pacing. \\
\end{longtable}

\begin{Shaded}
\begin{Highlighting}[]
\FunctionTok{\{}\DataTypeTok{"action"}\FunctionTok{:} \StringTok{"squeeze\_through"}\FunctionTok{,} \DataTypeTok{"who"}\FunctionTok{:} \StringTok{"freydoon"}\FunctionTok{,} \DataTypeTok{"x"}\FunctionTok{:} \FloatTok{2.5}\FunctionTok{\}}
\end{Highlighting}
\end{Shaded}

\begin{center}\rule{0.5\linewidth}{0.5pt}\end{center}

\paragraph{\texorpdfstring{3.2.9
\texttt{exit\_through}}{3.2.9 exit\_through}}\label{exit_through}

\emph{Parallel-safe:} no · \emph{Animation:} timed · \emph{Required
keys:} \texttt{who}, \texttt{prop}

Walk to a door (or other exit prop), then fade out. The character
disappears at the door's x position.

\begin{Shaded}
\begin{Highlighting}[]
\FunctionTok{\{}\DataTypeTok{"action"}\FunctionTok{:} \StringTok{"exit\_through"}\FunctionTok{,} \DataTypeTok{"who"}\FunctionTok{:} \StringTok{"bevers"}\FunctionTok{,} \DataTypeTok{"prop"}\FunctionTok{:} \StringTok{"exit\_door"}\FunctionTok{\}}
\end{Highlighting}
\end{Shaded}

\begin{center}\rule{0.5\linewidth}{0.5pt}\end{center}

\paragraph{\texorpdfstring{3.2.10
\texttt{exit\_through\_doors}}{3.2.10 exit\_through\_doors}}\label{exit_through_doors}

\emph{Parallel-safe:} no · \emph{Animation:} timed · \emph{Required
keys:} \texttt{who}, \texttt{prop}

Walk to a door (typically an elevator), pause for the door's open/close
animation, then disappear off-screen. Difference from
\texttt{exit\_through}: this version respects elevator-style props with
explicit door state.

\textbf{Parameters:}

\begin{longtable}[]{@{}
  >{\raggedright\arraybackslash}p{(\columnwidth - 8\tabcolsep) * \real{0.2000}}
  >{\raggedright\arraybackslash}p{(\columnwidth - 8\tabcolsep) * \real{0.2000}}
  >{\raggedright\arraybackslash}p{(\columnwidth - 8\tabcolsep) * \real{0.2000}}
  >{\raggedright\arraybackslash}p{(\columnwidth - 8\tabcolsep) * \real{0.2000}}
  >{\raggedright\arraybackslash}p{(\columnwidth - 8\tabcolsep) * \real{0.2000}}@{}}
\toprule\noalign{}
\begin{minipage}[b]{\linewidth}\raggedright
Key
\end{minipage} & \begin{minipage}[b]{\linewidth}\raggedright
Type
\end{minipage} & \begin{minipage}[b]{\linewidth}\raggedright
Required
\end{minipage} & \begin{minipage}[b]{\linewidth}\raggedright
Default
\end{minipage} & \begin{minipage}[b]{\linewidth}\raggedright
Description
\end{minipage} \\
\midrule\noalign{}
\endhead
\bottomrule\noalign{}
\endlastfoot
\texttt{who} & \texttt{str} & yes & --- & Character key. \\
\texttt{prop} & \texttt{str} & yes & --- & Elevator or door prop with
open/close animation. \\
\texttt{rt} & \texttt{float} & no & per stage & Per-stage run time. \\
\end{longtable}

\begin{Shaded}
\begin{Highlighting}[]
\FunctionTok{\{}\DataTypeTok{"action"}\FunctionTok{:} \StringTok{"exit\_through\_doors"}\FunctionTok{,} \DataTypeTok{"who"}\FunctionTok{:} \StringTok{"thalia"}\FunctionTok{,} \DataTypeTok{"prop"}\FunctionTok{:} \StringTok{"elevator\_1"}\FunctionTok{\}}
\end{Highlighting}
\end{Shaded}

\begin{center}\rule{0.5\linewidth}{0.5pt}\end{center}

\paragraph{\texorpdfstring{3.2.11
\texttt{jump\_up}}{3.2.11 jump\_up}}\label{jump_up}

\emph{Parallel-safe:} no · \emph{Animation:} timed · \emph{Required
keys:} \texttt{who}

Eager reactive jump: crouch -\textgreater{} peak -\textgreater{} land.
Used for affirmative reactions (``got it!'', ``yes!'').

\textbf{Parameters:}

\begin{longtable}[]{@{}
  >{\raggedright\arraybackslash}p{(\columnwidth - 8\tabcolsep) * \real{0.2000}}
  >{\raggedright\arraybackslash}p{(\columnwidth - 8\tabcolsep) * \real{0.2000}}
  >{\raggedright\arraybackslash}p{(\columnwidth - 8\tabcolsep) * \real{0.2000}}
  >{\raggedright\arraybackslash}p{(\columnwidth - 8\tabcolsep) * \real{0.2000}}
  >{\raggedright\arraybackslash}p{(\columnwidth - 8\tabcolsep) * \real{0.2000}}@{}}
\toprule\noalign{}
\begin{minipage}[b]{\linewidth}\raggedright
Key
\end{minipage} & \begin{minipage}[b]{\linewidth}\raggedright
Type
\end{minipage} & \begin{minipage}[b]{\linewidth}\raggedright
Required
\end{minipage} & \begin{minipage}[b]{\linewidth}\raggedright
Default
\end{minipage} & \begin{minipage}[b]{\linewidth}\raggedright
Description
\end{minipage} \\
\midrule\noalign{}
\endhead
\bottomrule\noalign{}
\endlastfoot
\texttt{who} & \texttt{str} & yes & --- & Character key. \\
\texttt{height} & \texttt{float} & no & \texttt{0.4} & Peak height in
world units. \\
\texttt{rt} & \texttt{float} & no & \texttt{0.4} & Total cycle run time
(crouch + peak + land). \\
\end{longtable}

\begin{Shaded}
\begin{Highlighting}[]
\FunctionTok{\{}\DataTypeTok{"action"}\FunctionTok{:} \StringTok{"jump\_up"}\FunctionTok{,} \DataTypeTok{"who"}\FunctionTok{:} \StringTok{"bart"}\FunctionTok{,} \DataTypeTok{"height"}\FunctionTok{:} \DecValTok{0}\ErrorTok{.}\DecValTok{6}\FunctionTok{\}}
\end{Highlighting}
\end{Shaded}

\begin{center}\rule{0.5\linewidth}{0.5pt}\end{center}

\paragraph{\texorpdfstring{3.2.12
\texttt{fall\_down}}{3.2.12 fall\_down}}\label{fall_down}

\emph{Parallel-safe:} yes · \emph{Animation:} timed · \emph{Required
keys:} \texttt{who}

Animate a character falling from standing to on-hands-and-knees.

\textbf{Parameters:}

\begin{longtable}[]{@{}
  >{\raggedright\arraybackslash}p{(\columnwidth - 8\tabcolsep) * \real{0.2000}}
  >{\raggedright\arraybackslash}p{(\columnwidth - 8\tabcolsep) * \real{0.2000}}
  >{\raggedright\arraybackslash}p{(\columnwidth - 8\tabcolsep) * \real{0.2000}}
  >{\raggedright\arraybackslash}p{(\columnwidth - 8\tabcolsep) * \real{0.2000}}
  >{\raggedright\arraybackslash}p{(\columnwidth - 8\tabcolsep) * \real{0.2000}}@{}}
\toprule\noalign{}
\begin{minipage}[b]{\linewidth}\raggedright
Key
\end{minipage} & \begin{minipage}[b]{\linewidth}\raggedright
Type
\end{minipage} & \begin{minipage}[b]{\linewidth}\raggedright
Required
\end{minipage} & \begin{minipage}[b]{\linewidth}\raggedright
Default
\end{minipage} & \begin{minipage}[b]{\linewidth}\raggedright
Description
\end{minipage} \\
\midrule\noalign{}
\endhead
\bottomrule\noalign{}
\endlastfoot
\texttt{who} & \texttt{str} & yes & --- & Character key. \\
\texttt{style} & \texttt{str} & no & \texttt{"stumble"} & Variant:
\texttt{"stumble"} (catch-and-go-down), \texttt{"fall\_catch"}
(catch-mid-fall, recover), or a direct pose name. \\
\texttt{rt} & \texttt{float} & no & per stage & Per-stage run time. \\
\end{longtable}

\begin{Shaded}
\begin{Highlighting}[]
\FunctionTok{\{}\DataTypeTok{"action"}\FunctionTok{:} \StringTok{"fall\_down"}\FunctionTok{,} \DataTypeTok{"who"}\FunctionTok{:} \StringTok{"brad"}\FunctionTok{,} \DataTypeTok{"style"}\FunctionTok{:} \StringTok{"stumble"}\FunctionTok{\}}
\end{Highlighting}
\end{Shaded}

\begin{center}\rule{0.5\linewidth}{0.5pt}\end{center}

\paragraph{\texorpdfstring{3.2.13
\texttt{dodge}}{3.2.13 dodge}}\label{dodge}

\emph{Parallel-safe:} yes · \emph{Animation:} timed · \emph{Required
keys:} \texttt{who}

Lateral sidestep away from another character's path. Short, sharp
displacement followed by return to upright.

\textbf{Parameters:}

\begin{longtable}[]{@{}
  >{\raggedright\arraybackslash}p{(\columnwidth - 8\tabcolsep) * \real{0.2000}}
  >{\raggedright\arraybackslash}p{(\columnwidth - 8\tabcolsep) * \real{0.2000}}
  >{\raggedright\arraybackslash}p{(\columnwidth - 8\tabcolsep) * \real{0.2000}}
  >{\raggedright\arraybackslash}p{(\columnwidth - 8\tabcolsep) * \real{0.2000}}
  >{\raggedright\arraybackslash}p{(\columnwidth - 8\tabcolsep) * \real{0.2000}}@{}}
\toprule\noalign{}
\begin{minipage}[b]{\linewidth}\raggedright
Key
\end{minipage} & \begin{minipage}[b]{\linewidth}\raggedright
Type
\end{minipage} & \begin{minipage}[b]{\linewidth}\raggedright
Required
\end{minipage} & \begin{minipage}[b]{\linewidth}\raggedright
Default
\end{minipage} & \begin{minipage}[b]{\linewidth}\raggedright
Description
\end{minipage} \\
\midrule\noalign{}
\endhead
\bottomrule\noalign{}
\endlastfoot
\texttt{who} & \texttt{str} & yes & --- & Character key. \\
\texttt{direction} & \texttt{str} & no & \texttt{"right"} &
\texttt{"left"} or \texttt{"right"}. \\
\texttt{distance} & \texttt{float} & no & \texttt{0.4} & Lateral
displacement in world units. \\
\texttt{rt} & \texttt{float} & no & \texttt{0.3} & Sidestep duration. \\
\end{longtable}

\begin{Shaded}
\begin{Highlighting}[]
\FunctionTok{\{}\DataTypeTok{"action"}\FunctionTok{:} \StringTok{"dodge"}\FunctionTok{,} \DataTypeTok{"who"}\FunctionTok{:} \StringTok{"thalia"}\FunctionTok{,} \DataTypeTok{"direction"}\FunctionTok{:} \StringTok{"left"}\FunctionTok{,} \DataTypeTok{"distance"}\FunctionTok{:} \DecValTok{0}\ErrorTok{.}\DecValTok{5}\FunctionTok{\}}
\end{Highlighting}
\end{Shaded}

\begin{center}\rule{0.5\linewidth}{0.5pt}\end{center}

\subsubsection{3.3 Posture transitions}\label{posture-transitions}

\paragraph{\texorpdfstring{3.3.1
\texttt{sit\_down}}{3.3.1 sit\_down}}\label{sit_down}

\emph{Parallel-safe:} no · \emph{Animation:} timed · \emph{Required
keys:} \texttt{who}

Lower the figure from standing to seated. Drives a three-keyframe cycle
through \texttt{sitting\_mid} to \texttt{sitting\_down}. The character
must be near a chair prop for the seated position to read correctly ---
\texttt{sit\_down} does not change x position.

\begin{Shaded}
\begin{Highlighting}[]
\FunctionTok{\{}\DataTypeTok{"action"}\FunctionTok{:} \StringTok{"sit\_down"}\FunctionTok{,} \DataTypeTok{"who"}\FunctionTok{:} \StringTok{"freydoon"}\FunctionTok{\}}
\end{Highlighting}
\end{Shaded}

\textbf{Use pattern:} walk to chair, then sit:

\begin{Shaded}
\begin{Highlighting}[]
\FunctionTok{\{}\DataTypeTok{"action"}\FunctionTok{:} \StringTok{"walk\_to\_prop"}\FunctionTok{,} \DataTypeTok{"who"}\FunctionTok{:} \StringTok{"freydoon"}\FunctionTok{,} \DataTypeTok{"prop"}\FunctionTok{:} \StringTok{"chair\_1"}\FunctionTok{\}}\ErrorTok{,}
\FunctionTok{\{}\DataTypeTok{"action"}\FunctionTok{:} \StringTok{"sit\_down"}\FunctionTok{,} \DataTypeTok{"who"}\FunctionTok{:} \StringTok{"freydoon"}\FunctionTok{\}}
\end{Highlighting}
\end{Shaded}

\begin{center}\rule{0.5\linewidth}{0.5pt}\end{center}

\paragraph{\texorpdfstring{3.3.2
\texttt{stand\_up}}{3.3.2 stand\_up}}\label{stand_up}

\emph{Parallel-safe:} no · \emph{Animation:} timed · \emph{Required
keys:} \texttt{who}

Inverse of \texttt{sit\_down}.

\begin{Shaded}
\begin{Highlighting}[]
\FunctionTok{\{}\DataTypeTok{"action"}\FunctionTok{:} \StringTok{"stand\_up"}\FunctionTok{,} \DataTypeTok{"who"}\FunctionTok{:} \StringTok{"freydoon"}\FunctionTok{\}}
\end{Highlighting}
\end{Shaded}

\begin{center}\rule{0.5\linewidth}{0.5pt}\end{center}

\subsubsection{3.4 Orientation}\label{orientation}

\paragraph{\texorpdfstring{3.4.1 \texttt{turn}}{3.4.1 turn}}\label{turn}

\emph{Parallel-safe:} yes · \emph{Animation:} timed · \emph{Required
keys:} \texttt{who}

Turn the figure to a side or named pose. Convenience wrapper around
\texttt{morph} for orientation-only changes.

\textbf{Parameters:}

\begin{longtable}[]{@{}
  >{\raggedright\arraybackslash}p{(\columnwidth - 8\tabcolsep) * \real{0.2000}}
  >{\raggedright\arraybackslash}p{(\columnwidth - 8\tabcolsep) * \real{0.2000}}
  >{\raggedright\arraybackslash}p{(\columnwidth - 8\tabcolsep) * \real{0.2000}}
  >{\raggedright\arraybackslash}p{(\columnwidth - 8\tabcolsep) * \real{0.2000}}
  >{\raggedright\arraybackslash}p{(\columnwidth - 8\tabcolsep) * \real{0.2000}}@{}}
\toprule\noalign{}
\begin{minipage}[b]{\linewidth}\raggedright
Key
\end{minipage} & \begin{minipage}[b]{\linewidth}\raggedright
Type
\end{minipage} & \begin{minipage}[b]{\linewidth}\raggedright
Required
\end{minipage} & \begin{minipage}[b]{\linewidth}\raggedright
Default
\end{minipage} & \begin{minipage}[b]{\linewidth}\raggedright
Description
\end{minipage} \\
\midrule\noalign{}
\endhead
\bottomrule\noalign{}
\endlastfoot
\texttt{who} & \texttt{str} & yes & --- & Character key. \\
\texttt{pose} & \texttt{str} & no & \texttt{"standing\_side"} & Target
pose. Common values: \texttt{"standing\_front"},
\texttt{"standing\_side"}. \\
\end{longtable}

\begin{Shaded}
\begin{Highlighting}[]
\FunctionTok{\{}\DataTypeTok{"action"}\FunctionTok{:} \StringTok{"turn"}\FunctionTok{,} \DataTypeTok{"who"}\FunctionTok{:} \StringTok{"thalia"}\FunctionTok{,} \DataTypeTok{"pose"}\FunctionTok{:} \StringTok{"standing\_side"}\FunctionTok{\}}
\end{Highlighting}
\end{Shaded}

\begin{center}\rule{0.5\linewidth}{0.5pt}\end{center}

\paragraph{\texorpdfstring{3.4.2 \texttt{face}}{3.4.2 face}}\label{face}

\emph{Parallel-safe:} yes · \emph{Animation:} timed · \emph{Required
keys:} \texttt{who}, \texttt{target}

Turn the figure to face a named prop or character. Computes the
appropriate side-view orientation based on the target's x position
relative to the character's.

\textbf{Parameters:}

\begin{longtable}[]{@{}
  >{\raggedright\arraybackslash}p{(\columnwidth - 8\tabcolsep) * \real{0.2000}}
  >{\raggedright\arraybackslash}p{(\columnwidth - 8\tabcolsep) * \real{0.2000}}
  >{\raggedright\arraybackslash}p{(\columnwidth - 8\tabcolsep) * \real{0.2000}}
  >{\raggedright\arraybackslash}p{(\columnwidth - 8\tabcolsep) * \real{0.2000}}
  >{\raggedright\arraybackslash}p{(\columnwidth - 8\tabcolsep) * \real{0.2000}}@{}}
\toprule\noalign{}
\begin{minipage}[b]{\linewidth}\raggedright
Key
\end{minipage} & \begin{minipage}[b]{\linewidth}\raggedright
Type
\end{minipage} & \begin{minipage}[b]{\linewidth}\raggedright
Required
\end{minipage} & \begin{minipage}[b]{\linewidth}\raggedright
Default
\end{minipage} & \begin{minipage}[b]{\linewidth}\raggedright
Description
\end{minipage} \\
\midrule\noalign{}
\endhead
\bottomrule\noalign{}
\endlastfoot
\texttt{who} & \texttt{str} & yes & --- & Character key. \\
\texttt{target} & \texttt{str} & yes & --- & Prop or character key. If
target's x \textgreater{} character's x, face right; else face left. \\
\end{longtable}

\begin{Shaded}
\begin{Highlighting}[]
\FunctionTok{\{}\DataTypeTok{"action"}\FunctionTok{:} \StringTok{"face"}\FunctionTok{,} \DataTypeTok{"who"}\FunctionTok{:} \StringTok{"thalia"}\FunctionTok{,} \DataTypeTok{"target"}\FunctionTok{:} \StringTok{"bevers"}\FunctionTok{\}}
\FunctionTok{\{}\DataTypeTok{"action"}\FunctionTok{:} \StringTok{"face"}\FunctionTok{,} \DataTypeTok{"who"}\FunctionTok{:} \StringTok{"athena"}\FunctionTok{,} \DataTypeTok{"target"}\FunctionTok{:} \StringTok{"desk"}\FunctionTok{\}}
\end{Highlighting}
\end{Shaded}

\begin{center}\rule{0.5\linewidth}{0.5pt}\end{center}

\subsubsection{3.5 Expressive gestures}\label{expressive-gestures}

\paragraph{\texorpdfstring{3.5.1 \texttt{wave}}{3.5.1 wave}}\label{wave}

\emph{Parallel-safe:} no · \emph{Animation:} timed · \emph{Required
keys:} \texttt{who}

Raise one arm and wag it.

\textbf{Parameters:}

\begin{longtable}[]{@{}
  >{\raggedright\arraybackslash}p{(\columnwidth - 8\tabcolsep) * \real{0.2000}}
  >{\raggedright\arraybackslash}p{(\columnwidth - 8\tabcolsep) * \real{0.2000}}
  >{\raggedright\arraybackslash}p{(\columnwidth - 8\tabcolsep) * \real{0.2000}}
  >{\raggedright\arraybackslash}p{(\columnwidth - 8\tabcolsep) * \real{0.2000}}
  >{\raggedright\arraybackslash}p{(\columnwidth - 8\tabcolsep) * \real{0.2000}}@{}}
\toprule\noalign{}
\begin{minipage}[b]{\linewidth}\raggedright
Key
\end{minipage} & \begin{minipage}[b]{\linewidth}\raggedright
Type
\end{minipage} & \begin{minipage}[b]{\linewidth}\raggedright
Required
\end{minipage} & \begin{minipage}[b]{\linewidth}\raggedright
Default
\end{minipage} & \begin{minipage}[b]{\linewidth}\raggedright
Description
\end{minipage} \\
\midrule\noalign{}
\endhead
\bottomrule\noalign{}
\endlastfoot
\texttt{who} & \texttt{str} & yes & --- & Character key. \\
\texttt{direction} & \texttt{str} & no & \texttt{"right"} &
\texttt{"right"} or \texttt{"left"} --- which arm waves. Synonym:
\texttt{"hand"}. \\
\texttt{hand} & \texttt{str} & no & \texttt{"right"} & Synonym for
\texttt{direction} (matches \texttt{figure.py}'s \texttt{wave(hand=...)}
parameter name). \\
\texttt{cycles} & \texttt{int} & no & \texttt{2} & Number of wag
oscillations. \\
\texttt{rt\_lift} & \texttt{float} & no & \texttt{0.4} & Arm raise/lower
run time. \\
\texttt{rt\_wag} & \texttt{float} & no & \texttt{0.24} &
Per-wag-keyframe run time. \\
\end{longtable}

\begin{Shaded}
\begin{Highlighting}[]
\FunctionTok{\{}\DataTypeTok{"action"}\FunctionTok{:} \StringTok{"wave"}\FunctionTok{,} \DataTypeTok{"who"}\FunctionTok{:} \StringTok{"bevers"}\FunctionTok{,} \DataTypeTok{"direction"}\FunctionTok{:} \StringTok{"left"}\FunctionTok{,} \DataTypeTok{"cycles"}\FunctionTok{:} \DecValTok{3}\FunctionTok{\}}
\end{Highlighting}
\end{Shaded}

\begin{center}\rule{0.5\linewidth}{0.5pt}\end{center}

\paragraph{\texorpdfstring{3.5.2 \texttt{wave\_left} /
\texttt{wave\_right}}{3.5.2 wave\_left / wave\_right}}\label{wave_left-wave_right}

\emph{Parallel-safe:} no · \emph{Animation:} timed · \emph{Required
keys:} \texttt{who}

Convenience aliases. \texttt{wave\_left} is equivalent to \texttt{wave}
with \texttt{direction:\ "left"}; \texttt{wave\_right} is equivalent to
\texttt{wave} with \texttt{direction:\ "right"}.

\begin{Shaded}
\begin{Highlighting}[]
\FunctionTok{\{}\DataTypeTok{"action"}\FunctionTok{:} \StringTok{"wave\_left"}\FunctionTok{,}  \DataTypeTok{"who"}\FunctionTok{:} \StringTok{"short\_emt"}\FunctionTok{\}}\ErrorTok{,}
\FunctionTok{\{}\DataTypeTok{"action"}\FunctionTok{:} \StringTok{"wave\_right"}\FunctionTok{,} \DataTypeTok{"who"}\FunctionTok{:} \StringTok{"tall\_emt"}\FunctionTok{\}}
\end{Highlighting}
\end{Shaded}

\begin{center}\rule{0.5\linewidth}{0.5pt}\end{center}

\paragraph{\texorpdfstring{3.5.3 \texttt{nod}}{3.5.3 nod}}\label{nod}

\emph{Parallel-safe:} no · \emph{Animation:} timed · \emph{Required
keys:} \texttt{who}

Nod the head up-and-down (affirmative).

\textbf{Parameters:}

\begin{longtable}[]{@{}
  >{\raggedright\arraybackslash}p{(\columnwidth - 8\tabcolsep) * \real{0.2000}}
  >{\raggedright\arraybackslash}p{(\columnwidth - 8\tabcolsep) * \real{0.2000}}
  >{\raggedright\arraybackslash}p{(\columnwidth - 8\tabcolsep) * \real{0.2000}}
  >{\raggedright\arraybackslash}p{(\columnwidth - 8\tabcolsep) * \real{0.2000}}
  >{\raggedright\arraybackslash}p{(\columnwidth - 8\tabcolsep) * \real{0.2000}}@{}}
\toprule\noalign{}
\begin{minipage}[b]{\linewidth}\raggedright
Key
\end{minipage} & \begin{minipage}[b]{\linewidth}\raggedright
Type
\end{minipage} & \begin{minipage}[b]{\linewidth}\raggedright
Required
\end{minipage} & \begin{minipage}[b]{\linewidth}\raggedright
Default
\end{minipage} & \begin{minipage}[b]{\linewidth}\raggedright
Description
\end{minipage} \\
\midrule\noalign{}
\endhead
\bottomrule\noalign{}
\endlastfoot
\texttt{who} & \texttt{str} & yes & --- & Character key. \\
\texttt{cycles} & \texttt{int} & no & \texttt{2} & Number of nod
oscillations. \\
\texttt{amp} & \texttt{float} & no & \texttt{0.08} & Vertical head
displacement amplitude. \\
\texttt{rt\_per\_step} & \texttt{float} & no & \texttt{0.18} &
Per-keyframe run time. \\
\end{longtable}

\begin{Shaded}
\begin{Highlighting}[]
\FunctionTok{\{}\DataTypeTok{"action"}\FunctionTok{:} \StringTok{"nod"}\FunctionTok{,} \DataTypeTok{"who"}\FunctionTok{:} \StringTok{"freydoon"}\FunctionTok{,} \DataTypeTok{"cycles"}\FunctionTok{:} \DecValTok{3}\FunctionTok{\}}
\end{Highlighting}
\end{Shaded}

\begin{center}\rule{0.5\linewidth}{0.5pt}\end{center}

\paragraph{\texorpdfstring{3.5.4
\texttt{shake\_head}}{3.5.4 shake\_head}}\label{shake_head}

\emph{Parallel-safe:} no · \emph{Animation:} timed · \emph{Required
keys:} \texttt{who}

Shake the head side-to-side (negation). Same parameters as \texttt{nod}
but the displacement is horizontal.

\begin{Shaded}
\begin{Highlighting}[]
\FunctionTok{\{}\DataTypeTok{"action"}\FunctionTok{:} \StringTok{"shake\_head"}\FunctionTok{,} \DataTypeTok{"who"}\FunctionTok{:} \StringTok{"thalia"}\FunctionTok{\}}
\end{Highlighting}
\end{Shaded}

\begin{center}\rule{0.5\linewidth}{0.5pt}\end{center}

\paragraph{\texorpdfstring{3.5.5
\texttt{shrug}}{3.5.5 shrug}}\label{shrug}

\emph{Parallel-safe:} no · \emph{Animation:} timed · \emph{Required
keys:} \texttt{who}

Raise both shoulders (arms follow), hold briefly, lower.

\textbf{Parameters:}

\begin{longtable}[]{@{}lllll@{}}
\toprule\noalign{}
Key & Type & Required & Default & Description \\
\midrule\noalign{}
\endhead
\bottomrule\noalign{}
\endlastfoot
\texttt{who} & \texttt{str} & yes & --- & Character key. \\
\texttt{amp} & \texttt{float} & no & \texttt{0.15} & Shoulder rise
amplitude. \\
\texttt{rt} & \texttt{float} & no & \texttt{0.4} & Rise/lower duration
per direction. \\
\texttt{hold} & \texttt{float} & no & \texttt{0.3} & Hold duration at
peak. \\
\end{longtable}

\begin{Shaded}
\begin{Highlighting}[]
\FunctionTok{\{}\DataTypeTok{"action"}\FunctionTok{:} \StringTok{"shrug"}\FunctionTok{,} \DataTypeTok{"who"}\FunctionTok{:} \StringTok{"bevers"}\FunctionTok{,} \DataTypeTok{"amp"}\FunctionTok{:} \DecValTok{0}\ErrorTok{.}\DecValTok{2}\FunctionTok{,} \DataTypeTok{"hold"}\FunctionTok{:} \DecValTok{0}\ErrorTok{.}\DecValTok{5}\FunctionTok{\}}
\end{Highlighting}
\end{Shaded}

\begin{center}\rule{0.5\linewidth}{0.5pt}\end{center}

\paragraph{\texorpdfstring{3.5.6
\texttt{point\_at}}{3.5.6 point\_at}}\label{point_at}

\emph{Parallel-safe:} yes · \emph{Animation:} timed · \emph{Required
keys:} \texttt{who}, \texttt{target}

Extend arm and point at a named prop or character.

\textbf{Parameters:}

\begin{longtable}[]{@{}
  >{\raggedright\arraybackslash}p{(\columnwidth - 8\tabcolsep) * \real{0.2000}}
  >{\raggedright\arraybackslash}p{(\columnwidth - 8\tabcolsep) * \real{0.2000}}
  >{\raggedright\arraybackslash}p{(\columnwidth - 8\tabcolsep) * \real{0.2000}}
  >{\raggedright\arraybackslash}p{(\columnwidth - 8\tabcolsep) * \real{0.2000}}
  >{\raggedright\arraybackslash}p{(\columnwidth - 8\tabcolsep) * \real{0.2000}}@{}}
\toprule\noalign{}
\begin{minipage}[b]{\linewidth}\raggedright
Key
\end{minipage} & \begin{minipage}[b]{\linewidth}\raggedright
Type
\end{minipage} & \begin{minipage}[b]{\linewidth}\raggedright
Required
\end{minipage} & \begin{minipage}[b]{\linewidth}\raggedright
Default
\end{minipage} & \begin{minipage}[b]{\linewidth}\raggedright
Description
\end{minipage} \\
\midrule\noalign{}
\endhead
\bottomrule\noalign{}
\endlastfoot
\texttt{who} & \texttt{str} & yes & --- & Character key. \\
\texttt{target} & \texttt{str} & yes & --- & Prop or character key. \\
\texttt{hold} & \texttt{float} & no & \texttt{0.6} & Hold time at full
extension before retracting. \\
\end{longtable}

\begin{Shaded}
\begin{Highlighting}[]
\FunctionTok{\{}\DataTypeTok{"action"}\FunctionTok{:} \StringTok{"point\_at"}\FunctionTok{,} \DataTypeTok{"who"}\FunctionTok{:} \StringTok{"nona"}\FunctionTok{,} \DataTypeTok{"target"}\FunctionTok{:} \StringTok{"factor"}\FunctionTok{\}}
\FunctionTok{\{}\DataTypeTok{"action"}\FunctionTok{:} \StringTok{"point\_at"}\FunctionTok{,} \DataTypeTok{"who"}\FunctionTok{:} \StringTok{"athena"}\FunctionTok{,} \DataTypeTok{"target"}\FunctionTok{:} \StringTok{"desk"}\FunctionTok{\}}
\end{Highlighting}
\end{Shaded}

\begin{center}\rule{0.5\linewidth}{0.5pt}\end{center}

\paragraph{\texorpdfstring{3.5.7
\texttt{express}}{3.5.7 express}}\label{express}

\emph{Parallel-safe:} yes · \emph{Animation:} timed · \emph{Required
keys:} \texttt{who}, \texttt{expression}

Flash a reaction glyph above the character's head. The glyph is a small
Unicode/text symbol drawn just above the head dot.

\textbf{Parameters:}

\begin{longtable}[]{@{}
  >{\raggedright\arraybackslash}p{(\columnwidth - 8\tabcolsep) * \real{0.2000}}
  >{\raggedright\arraybackslash}p{(\columnwidth - 8\tabcolsep) * \real{0.2000}}
  >{\raggedright\arraybackslash}p{(\columnwidth - 8\tabcolsep) * \real{0.2000}}
  >{\raggedright\arraybackslash}p{(\columnwidth - 8\tabcolsep) * \real{0.2000}}
  >{\raggedright\arraybackslash}p{(\columnwidth - 8\tabcolsep) * \real{0.2000}}@{}}
\toprule\noalign{}
\begin{minipage}[b]{\linewidth}\raggedright
Key
\end{minipage} & \begin{minipage}[b]{\linewidth}\raggedright
Type
\end{minipage} & \begin{minipage}[b]{\linewidth}\raggedright
Required
\end{minipage} & \begin{minipage}[b]{\linewidth}\raggedright
Default
\end{minipage} & \begin{minipage}[b]{\linewidth}\raggedright
Description
\end{minipage} \\
\midrule\noalign{}
\endhead
\bottomrule\noalign{}
\endlastfoot
\texttt{who} & \texttt{str} & yes & --- & Character key. \\
\texttt{expression} & \texttt{str} & yes & --- & Key from
\texttt{EXPRESSION\_GLYPHS} in \texttt{pam/poses.py}. Current registry:
\texttt{"smirk"}, \texttt{"roll\_eyes"}. \\
\texttt{hold} & \texttt{float} & no & per-glyph default & Display
duration. \\
\texttt{rt\_in} & \texttt{float} & no & \texttt{0.2} & Fade-in. \\
\texttt{rt\_out} & \texttt{float} & no & \texttt{0.2} & Fade-out. \\
\texttt{color} & \texttt{str} (hex) & no & per-glyph default & Override
the glyph color. \\
\end{longtable}

\begin{Shaded}
\begin{Highlighting}[]
\FunctionTok{\{}\DataTypeTok{"action"}\FunctionTok{:} \StringTok{"express"}\FunctionTok{,} \DataTypeTok{"who"}\FunctionTok{:} \StringTok{"bevers"}\FunctionTok{,} \DataTypeTok{"expression"}\FunctionTok{:} \StringTok{"smirk"}\FunctionTok{\}}
\FunctionTok{\{}\DataTypeTok{"action"}\FunctionTok{:} \StringTok{"express"}\FunctionTok{,} \DataTypeTok{"who"}\FunctionTok{:} \StringTok{"thalia"}\FunctionTok{,} \DataTypeTok{"expression"}\FunctionTok{:} \StringTok{"roll\_eyes"}\FunctionTok{,} \DataTypeTok{"hold"}\FunctionTok{:} \FloatTok{1.5}\FunctionTok{\}}
\end{Highlighting}
\end{Shaded}

\textbf{Note:} \texttt{EXPRESSION\_GLYPHS} is small (2 entries
currently). Each entry specifies its own default \texttt{glyph},
\texttt{dx}, \texttt{dy}, \texttt{font\_size}, \texttt{hold}, and
\texttt{color}. The action's \texttt{hold} and \texttt{color} parameters
override the glyph defaults.

\begin{center}\rule{0.5\linewidth}{0.5pt}\end{center}

\paragraph{\texorpdfstring{3.5.8
\texttt{react}}{3.5.8 react}}\label{react}

\emph{Parallel-safe:} yes · \emph{Animation:} timed · \emph{Required
keys:} \texttt{who}

Fire any tics in \texttt{fig.tic\_profile} whose triggers include
\texttt{"react"}. No-op for figures with no tic profile or no
react-triggered fragments.

\textbf{Parameters:}

\begin{longtable}[]{@{}
  >{\raggedright\arraybackslash}p{(\columnwidth - 8\tabcolsep) * \real{0.2000}}
  >{\raggedright\arraybackslash}p{(\columnwidth - 8\tabcolsep) * \real{0.2000}}
  >{\raggedright\arraybackslash}p{(\columnwidth - 8\tabcolsep) * \real{0.2000}}
  >{\raggedright\arraybackslash}p{(\columnwidth - 8\tabcolsep) * \real{0.2000}}
  >{\raggedright\arraybackslash}p{(\columnwidth - 8\tabcolsep) * \real{0.2000}}@{}}
\toprule\noalign{}
\begin{minipage}[b]{\linewidth}\raggedright
Key
\end{minipage} & \begin{minipage}[b]{\linewidth}\raggedright
Type
\end{minipage} & \begin{minipage}[b]{\linewidth}\raggedright
Required
\end{minipage} & \begin{minipage}[b]{\linewidth}\raggedright
Default
\end{minipage} & \begin{minipage}[b]{\linewidth}\raggedright
Description
\end{minipage} \\
\midrule\noalign{}
\endhead
\bottomrule\noalign{}
\endlastfoot
\texttt{who} & \texttt{str} & yes & --- & Character key. \\
\texttt{suppress\_tics} & \texttt{bool} & no & \texttt{false} & If
\texttt{true}, the action no-ops (useful for marking a beat without
actually firing the tic). \\
\end{longtable}

\begin{Shaded}
\begin{Highlighting}[]
\FunctionTok{\{}\DataTypeTok{"action"}\FunctionTok{:} \StringTok{"react"}\FunctionTok{,} \DataTypeTok{"who"}\FunctionTok{:} \StringTok{"bevers"}\FunctionTok{\}}
\end{Highlighting}
\end{Shaded}

\textbf{Tic fragment registry:} see \texttt{pam/tics.py} and §1.12 for
the full table. Current fragments: \texttt{hand\_twitch\_r},
\texttt{hand\_twitch\_l}, \texttt{foot\_tap\_r}, \texttt{foot\_tap\_l}
(biped, \texttt{human} and \texttt{alien}), and \texttt{tail\_wag}
(quadruped, \texttt{dog}). Body-type compatibility is checked at play
time --- a \texttt{tail\_wag} declared on a biped silently no-ops, as
does a \texttt{hand\_twitch\_*} or \texttt{foot\_tap\_*} declared on a
dog.

\textbf{Trigger interaction:} \texttt{say} and \texttt{prop\_say} also
fire tics under the \texttt{"sentence\_end"} trigger automatically. A
tic\_profile that triggers on both \texttt{react} and
\texttt{sentence\_end} fires on every \texttt{say} step \emph{and} every
explicit \texttt{react}.

\begin{center}\rule{0.5\linewidth}{0.5pt}\end{center}

\subsubsection{3.6 Prop interaction}\label{prop-interaction}

The following actions all take a character as \texttt{who} and a prop
(or surface) as the direct object. Common patterns:

\begin{itemize}
\tightlist
\item
  \texttt{target} --- the prop being approached or operated on (most
  actions).
\item
  \texttt{prop} --- the prop being carried, picked up, or set down
  (state-changing actions).
\item
  \texttt{arm} --- which arm to use: \texttt{"right"} (default) or
  \texttt{"left"}.
\end{itemize}

\paragraph{\texorpdfstring{3.6.1
\texttt{reach\_for}}{3.6.1 reach\_for}}\label{reach_for}

\emph{Parallel-safe:} yes · \emph{Animation:} timed · \emph{Required
keys:} \texttt{who}, \texttt{target}

Extend one arm toward a target prop, hold briefly, retract. Used as a
non-committal gesture --- character indicates the prop without acting on
it.

\textbf{Parameters:}

\begin{longtable}[]{@{}lllll@{}}
\toprule\noalign{}
Key & Type & Required & Default & Description \\
\midrule\noalign{}
\endhead
\bottomrule\noalign{}
\endlastfoot
\texttt{who} & \texttt{str} & yes & --- & Character key. \\
\texttt{target} & \texttt{str} & yes & --- & Prop key. \\
\texttt{arm} & \texttt{str} & no & \texttt{"right"} & \texttt{"right"}
or \texttt{"left"}. \\
\texttt{hold} & \texttt{float} & no & \texttt{0.4} & Hold time at full
extension. \\
\texttt{rt} & \texttt{float} & no & \texttt{0.4} & Extend/retract per
direction. \\
\end{longtable}

\begin{Shaded}
\begin{Highlighting}[]
\FunctionTok{\{}\DataTypeTok{"action"}\FunctionTok{:} \StringTok{"reach\_for"}\FunctionTok{,} \DataTypeTok{"who"}\FunctionTok{:} \StringTok{"freydoon"}\FunctionTok{,} \DataTypeTok{"target"}\FunctionTok{:} \StringTok{"doorknob"}\FunctionTok{\}}
\end{Highlighting}
\end{Shaded}

\begin{center}\rule{0.5\linewidth}{0.5pt}\end{center}

\paragraph{\texorpdfstring{3.6.2 \texttt{grab}}{3.6.2 grab}}\label{grab}

\emph{Parallel-safe:} no · \emph{Animation:} timed · \emph{Required
keys:} \texttt{who}, \texttt{prop} (or \texttt{target})

Decisive \texttt{reach\_for} + retract with the prop now held. Reads as
more urgent than \texttt{pick\_up} --- faster, no kneel. After
\texttt{grab}, the character is holding the prop in carry pose.

\textbf{Parameters:}

\begin{longtable}[]{@{}lllll@{}}
\toprule\noalign{}
Key & Type & Required & Default & Description \\
\midrule\noalign{}
\endhead
\bottomrule\noalign{}
\endlastfoot
\texttt{who} & \texttt{str} & yes & --- & Character key. \\
\texttt{prop} & \texttt{str} & one of & --- & Prop to grab. \\
\texttt{target} & \texttt{str} & one of & --- & Synonym for
\texttt{prop}. \\
\texttt{arm} & \texttt{str} & no & \texttt{"right"} & \texttt{"right"}
or \texttt{"left"}. \\
\texttt{rt} & \texttt{float} & no & \texttt{0.3} & Grab speed. \\
\end{longtable}

\begin{Shaded}
\begin{Highlighting}[]
\FunctionTok{\{}\DataTypeTok{"action"}\FunctionTok{:} \StringTok{"grab"}\FunctionTok{,} \DataTypeTok{"who"}\FunctionTok{:} \StringTok{"chava"}\FunctionTok{,} \DataTypeTok{"prop"}\FunctionTok{:} \StringTok{"folder"}\FunctionTok{\}}
\end{Highlighting}
\end{Shaded}

\begin{center}\rule{0.5\linewidth}{0.5pt}\end{center}

\paragraph{\texorpdfstring{3.6.3
\texttt{pick\_up}}{3.6.3 pick\_up}}\label{pick_up}

\emph{Parallel-safe:} yes · \emph{Animation:} timed · \emph{Required
keys:} \texttt{who}, \texttt{prop}

Reach down, pick up a prop, return to carry-hold pose. Slower and more
deliberate than \texttt{grab} --- used for ground-level or table-level
objects.

\begin{Shaded}
\begin{Highlighting}[]
\FunctionTok{\{}\DataTypeTok{"action"}\FunctionTok{:} \StringTok{"pick\_up"}\FunctionTok{,} \DataTypeTok{"who"}\FunctionTok{:} \StringTok{"sidel"}\FunctionTok{,} \DataTypeTok{"prop"}\FunctionTok{:} \StringTok{"briefcase"}\FunctionTok{\}}
\end{Highlighting}
\end{Shaded}

\begin{center}\rule{0.5\linewidth}{0.5pt}\end{center}

\paragraph{\texorpdfstring{3.6.4
\texttt{put\_down}}{3.6.4 put\_down}}\label{put_down}

\emph{Parallel-safe:} yes · \emph{Animation:} timed · \emph{Required
keys:} \texttt{who}, \texttt{prop}

Lower a held prop onto a target surface or the floor.

\textbf{Parameters:}

\begin{longtable}[]{@{}
  >{\raggedright\arraybackslash}p{(\columnwidth - 8\tabcolsep) * \real{0.2000}}
  >{\raggedright\arraybackslash}p{(\columnwidth - 8\tabcolsep) * \real{0.2000}}
  >{\raggedright\arraybackslash}p{(\columnwidth - 8\tabcolsep) * \real{0.2000}}
  >{\raggedright\arraybackslash}p{(\columnwidth - 8\tabcolsep) * \real{0.2000}}
  >{\raggedright\arraybackslash}p{(\columnwidth - 8\tabcolsep) * \real{0.2000}}@{}}
\toprule\noalign{}
\begin{minipage}[b]{\linewidth}\raggedright
Key
\end{minipage} & \begin{minipage}[b]{\linewidth}\raggedright
Type
\end{minipage} & \begin{minipage}[b]{\linewidth}\raggedright
Required
\end{minipage} & \begin{minipage}[b]{\linewidth}\raggedright
Default
\end{minipage} & \begin{minipage}[b]{\linewidth}\raggedright
Description
\end{minipage} \\
\midrule\noalign{}
\endhead
\bottomrule\noalign{}
\endlastfoot
\texttt{who} & \texttt{str} & yes & --- & Character key. \\
\texttt{prop} & \texttt{str} & yes & --- & Currently-held prop. \\
\texttt{on} & \texttt{str} & no & --- & Target surface prop (desk,
table). Omit to put down on the floor. \\
\end{longtable}

\begin{Shaded}
\begin{Highlighting}[]
\FunctionTok{\{}\DataTypeTok{"action"}\FunctionTok{:} \StringTok{"put\_down"}\FunctionTok{,} \DataTypeTok{"who"}\FunctionTok{:} \StringTok{"sidel"}\FunctionTok{,} \DataTypeTok{"prop"}\FunctionTok{:} \StringTok{"briefcase"}\FunctionTok{,} \DataTypeTok{"on"}\FunctionTok{:} \StringTok{"desk"}\FunctionTok{\}}
\FunctionTok{\{}\DataTypeTok{"action"}\FunctionTok{:} \StringTok{"put\_down"}\FunctionTok{,} \DataTypeTok{"who"}\FunctionTok{:} \StringTok{"freydoon"}\FunctionTok{,} \DataTypeTok{"prop"}\FunctionTok{:} \StringTok{"phone"}\FunctionTok{\}}
\end{Highlighting}
\end{Shaded}

\begin{center}\rule{0.5\linewidth}{0.5pt}\end{center}

\paragraph{\texorpdfstring{3.6.5
\texttt{place\_on}}{3.6.5 place\_on}}\label{place_on}

\emph{Parallel-safe:} yes · \emph{Animation:} timed · \emph{Required
keys:} \texttt{who}, \texttt{prop}, \texttt{target}

Set a carried prop onto a target prop's surface. Difference from
\texttt{put\_down}: \texttt{place\_on} always requires a target surface
and animates a more deliberate placement (slightly higher arc).

\textbf{Parameters:}

\begin{longtable}[]{@{}lllll@{}}
\toprule\noalign{}
Key & Type & Required & Default & Description \\
\midrule\noalign{}
\endhead
\bottomrule\noalign{}
\endlastfoot
\texttt{who} & \texttt{str} & yes & --- & Character key. \\
\texttt{prop} & \texttt{str} & yes & --- & Currently-held prop. \\
\texttt{target} & \texttt{str} & yes & --- & Target surface prop. \\
\texttt{rt} & \texttt{float} & no & \texttt{0.5} & Placement
duration. \\
\end{longtable}

\begin{Shaded}
\begin{Highlighting}[]
\FunctionTok{\{}\DataTypeTok{"action"}\FunctionTok{:} \StringTok{"place\_on"}\FunctionTok{,} \DataTypeTok{"who"}\FunctionTok{:} \StringTok{"athena"}\FunctionTok{,} \DataTypeTok{"prop"}\FunctionTok{:} \StringTok{"laptop"}\FunctionTok{,} \DataTypeTok{"target"}\FunctionTok{:} \StringTok{"desk"}\FunctionTok{\}}
\end{Highlighting}
\end{Shaded}

\begin{center}\rule{0.5\linewidth}{0.5pt}\end{center}

\paragraph{\texorpdfstring{3.6.6
\texttt{peel\_from\_hand}}{3.6.6 peel\_from\_hand}}\label{peel_from_hand}

\emph{Parallel-safe:} no · \emph{Animation:} timed · \emph{Required
keys:} \texttt{who}, \texttt{prop}

Peel a tiny prop (e.g.~an audio/video bug) off the character's palm. The
prop detaches from the carry pose and remains at the peeled-off
position.

\textbf{Parameters:}

\begin{longtable}[]{@{}lllll@{}}
\toprule\noalign{}
Key & Type & Required & Default & Description \\
\midrule\noalign{}
\endhead
\bottomrule\noalign{}
\endlastfoot
\texttt{who} & \texttt{str} & yes & --- & Character key. \\
\texttt{prop} & \texttt{str} & yes & --- & Held prop to peel off. \\
\texttt{arm} & \texttt{str} & no & \texttt{"right"} & \texttt{"right"}
or \texttt{"left"}. \\
\texttt{rt} & \texttt{float} & no & \texttt{0.4} & Peel duration. \\
\end{longtable}

\begin{Shaded}
\begin{Highlighting}[]
\FunctionTok{\{}\DataTypeTok{"action"}\FunctionTok{:} \StringTok{"peel\_from\_hand"}\FunctionTok{,} \DataTypeTok{"who"}\FunctionTok{:} \StringTok{"agent"}\FunctionTok{,} \DataTypeTok{"prop"}\FunctionTok{:} \StringTok{"av\_bug"}\FunctionTok{\}}
\end{Highlighting}
\end{Shaded}

\begin{center}\rule{0.5\linewidth}{0.5pt}\end{center}

\paragraph{\texorpdfstring{3.6.7
\texttt{punch\_button}}{3.6.7 punch\_button}}\label{punch_button}

\emph{Parallel-safe:} no · \emph{Animation:} timed · \emph{Required
keys:} \texttt{who}, \texttt{target}

Sharp jab at a prop (button panel, elevator call, doorbell) then
retract. Faster and more aggressive than \texttt{reach\_for}.

\textbf{Parameters:}

\begin{longtable}[]{@{}
  >{\raggedright\arraybackslash}p{(\columnwidth - 8\tabcolsep) * \real{0.2000}}
  >{\raggedright\arraybackslash}p{(\columnwidth - 8\tabcolsep) * \real{0.2000}}
  >{\raggedright\arraybackslash}p{(\columnwidth - 8\tabcolsep) * \real{0.2000}}
  >{\raggedright\arraybackslash}p{(\columnwidth - 8\tabcolsep) * \real{0.2000}}
  >{\raggedright\arraybackslash}p{(\columnwidth - 8\tabcolsep) * \real{0.2000}}@{}}
\toprule\noalign{}
\begin{minipage}[b]{\linewidth}\raggedright
Key
\end{minipage} & \begin{minipage}[b]{\linewidth}\raggedright
Type
\end{minipage} & \begin{minipage}[b]{\linewidth}\raggedright
Required
\end{minipage} & \begin{minipage}[b]{\linewidth}\raggedright
Default
\end{minipage} & \begin{minipage}[b]{\linewidth}\raggedright
Description
\end{minipage} \\
\midrule\noalign{}
\endhead
\bottomrule\noalign{}
\endlastfoot
\texttt{who} & \texttt{str} & yes & --- & Character key. \\
\texttt{target} & \texttt{str} & yes & --- & Target prop. \\
\texttt{arm} & \texttt{str} & no & \texttt{"right"} & \texttt{"right"}
or \texttt{"left"}. \\
\texttt{offset\_x} & \texttt{float} & no & \texttt{0.0} & Fine-tune jab
landing point relative to target's center. \\
\texttt{offset\_y} & \texttt{float} & no & \texttt{0.0} & As above,
vertical. \\
\texttt{hold} & \texttt{float} & no & \texttt{0.1} & Brief hold at
impact. \\
\texttt{rt} & \texttt{float} & no & \texttt{0.2} & Jab and retract
duration. \\
\end{longtable}

\begin{Shaded}
\begin{Highlighting}[]
\FunctionTok{\{}\DataTypeTok{"action"}\FunctionTok{:} \StringTok{"punch\_button"}\FunctionTok{,} \DataTypeTok{"who"}\FunctionTok{:} \StringTok{"athena"}\FunctionTok{,} \DataTypeTok{"target"}\FunctionTok{:} \StringTok{"call\_button"}\FunctionTok{\}}
\FunctionTok{\{}\DataTypeTok{"action"}\FunctionTok{:} \StringTok{"punch\_button"}\FunctionTok{,} \DataTypeTok{"who"}\FunctionTok{:} \StringTok{"freydoon"}\FunctionTok{,} \DataTypeTok{"target"}\FunctionTok{:} \StringTok{"elevator\_panel"}\FunctionTok{,}
 \DataTypeTok{"offset\_x"}\FunctionTok{:} \DecValTok{0}\ErrorTok{.}\DecValTok{1}\FunctionTok{,} \DataTypeTok{"offset\_y"}\FunctionTok{:} \DecValTok{{-}0}\ErrorTok{.}\DecValTok{2}\FunctionTok{\}}
\end{Highlighting}
\end{Shaded}

\begin{center}\rule{0.5\linewidth}{0.5pt}\end{center}

\paragraph{\texorpdfstring{3.6.8
\texttt{search\_drawers}}{3.6.8 search\_drawers}}\label{search_drawers}

\emph{Parallel-safe:} no · \emph{Animation:} timed · \emph{Required
keys:} \texttt{who}, \texttt{target}

Rummaging macro: repeated reach-for with downward pose variants. Reads
as opening and searching drawers in a desk.

\textbf{Parameters:}

\begin{longtable}[]{@{}lllll@{}}
\toprule\noalign{}
Key & Type & Required & Default & Description \\
\midrule\noalign{}
\endhead
\bottomrule\noalign{}
\endlastfoot
\texttt{who} & \texttt{str} & yes & --- & Character key. \\
\texttt{target} & \texttt{str} & yes & --- & Desk or chest prop. \\
\texttt{cycles} & \texttt{int} & no & \texttt{3} & Number of search
cycles. \\
\texttt{rt} & \texttt{float} & no & \texttt{0.4} & Per-cycle run
time. \\
\end{longtable}

\begin{Shaded}
\begin{Highlighting}[]
\FunctionTok{\{}\DataTypeTok{"action"}\FunctionTok{:} \StringTok{"search\_drawers"}\FunctionTok{,} \DataTypeTok{"who"}\FunctionTok{:} \StringTok{"thalia"}\FunctionTok{,} \DataTypeTok{"target"}\FunctionTok{:} \StringTok{"desk"}\FunctionTok{,} \DataTypeTok{"cycles"}\FunctionTok{:} \DecValTok{4}\FunctionTok{\}}
\end{Highlighting}
\end{Shaded}

\begin{center}\rule{0.5\linewidth}{0.5pt}\end{center}

\paragraph{\texorpdfstring{3.6.9
\texttt{snap\_photo}}{3.6.9 snap\_photo}}\label{snap_photo}

\emph{Parallel-safe:} yes · \emph{Animation:} timed · \emph{Required
keys:} \texttt{who}, \texttt{target}

Point a held phone at a target and flash. Requires the character to be
holding a phone (via \texttt{pick\_up} or \texttt{grab}).

\textbf{Parameters:}

\begin{longtable}[]{@{}
  >{\raggedright\arraybackslash}p{(\columnwidth - 8\tabcolsep) * \real{0.2000}}
  >{\raggedright\arraybackslash}p{(\columnwidth - 8\tabcolsep) * \real{0.2000}}
  >{\raggedright\arraybackslash}p{(\columnwidth - 8\tabcolsep) * \real{0.2000}}
  >{\raggedright\arraybackslash}p{(\columnwidth - 8\tabcolsep) * \real{0.2000}}
  >{\raggedright\arraybackslash}p{(\columnwidth - 8\tabcolsep) * \real{0.2000}}@{}}
\toprule\noalign{}
\begin{minipage}[b]{\linewidth}\raggedright
Key
\end{minipage} & \begin{minipage}[b]{\linewidth}\raggedright
Type
\end{minipage} & \begin{minipage}[b]{\linewidth}\raggedright
Required
\end{minipage} & \begin{minipage}[b]{\linewidth}\raggedright
Default
\end{minipage} & \begin{minipage}[b]{\linewidth}\raggedright
Description
\end{minipage} \\
\midrule\noalign{}
\endhead
\bottomrule\noalign{}
\endlastfoot
\texttt{who} & \texttt{str} & yes & --- & Character key. \\
\texttt{target} & \texttt{str} & yes & --- & Subject of the photo (prop
or character). \\
\texttt{flash\_rt} & \texttt{float} & no & \texttt{0.15} & Flash overlay
duration. \\
\texttt{hold} & \texttt{float} & no & \texttt{0.3} & Total hold time at
the photo-taking pose. \\
\end{longtable}

\begin{Shaded}
\begin{Highlighting}[]
\FunctionTok{\{}\DataTypeTok{"action"}\FunctionTok{:} \StringTok{"snap\_photo"}\FunctionTok{,} \DataTypeTok{"who"}\FunctionTok{:} \StringTok{"thalia"}\FunctionTok{,} \DataTypeTok{"target"}\FunctionTok{:} \StringTok{"bevers"}\FunctionTok{\}}
\end{Highlighting}
\end{Shaded}

\begin{center}\rule{0.5\linewidth}{0.5pt}\end{center}

\paragraph{\texorpdfstring{3.6.10
\texttt{stick\_to}}{3.6.10 stick\_to}}\label{stick_to}

\emph{Parallel-safe:} yes · \emph{Animation:} timed · \emph{Required
keys:} \texttt{who}, \texttt{prop}, \texttt{target}

Attach a small prop (e.g.~audio/video bug) to a target prop surface.
After the action, the small prop is parented to the target and tracks
with it.

\textbf{Parameters:}

\begin{longtable}[]{@{}
  >{\raggedright\arraybackslash}p{(\columnwidth - 8\tabcolsep) * \real{0.2000}}
  >{\raggedright\arraybackslash}p{(\columnwidth - 8\tabcolsep) * \real{0.2000}}
  >{\raggedright\arraybackslash}p{(\columnwidth - 8\tabcolsep) * \real{0.2000}}
  >{\raggedright\arraybackslash}p{(\columnwidth - 8\tabcolsep) * \real{0.2000}}
  >{\raggedright\arraybackslash}p{(\columnwidth - 8\tabcolsep) * \real{0.2000}}@{}}
\toprule\noalign{}
\begin{minipage}[b]{\linewidth}\raggedright
Key
\end{minipage} & \begin{minipage}[b]{\linewidth}\raggedright
Type
\end{minipage} & \begin{minipage}[b]{\linewidth}\raggedright
Required
\end{minipage} & \begin{minipage}[b]{\linewidth}\raggedright
Default
\end{minipage} & \begin{minipage}[b]{\linewidth}\raggedright
Description
\end{minipage} \\
\midrule\noalign{}
\endhead
\bottomrule\noalign{}
\endlastfoot
\texttt{who} & \texttt{str} & yes & --- & Character key. \\
\texttt{prop} & \texttt{str} & yes & --- & Held prop to stick. \\
\texttt{target} & \texttt{str} & yes & --- & Target prop (must declare a
\texttt{shade}, \texttt{surface}, or similar attachment point). \\
\texttt{rt} & \texttt{float} & no & \texttt{0.4} & Stick duration. \\
\end{longtable}

\begin{Shaded}
\begin{Highlighting}[]
\FunctionTok{\{}\DataTypeTok{"action"}\FunctionTok{:} \StringTok{"stick\_to"}\FunctionTok{,} \DataTypeTok{"who"}\FunctionTok{:} \StringTok{"agent"}\FunctionTok{,} \DataTypeTok{"prop"}\FunctionTok{:} \StringTok{"av\_bug"}\FunctionTok{,}
 \DataTypeTok{"target"}\FunctionTok{:} \StringTok{"desk\_lamp"}\FunctionTok{\}}
\end{Highlighting}
\end{Shaded}

\begin{center}\rule{0.5\linewidth}{0.5pt}\end{center}

\paragraph{\texorpdfstring{3.6.11
\texttt{move\_aside}}{3.6.11 move\_aside}}\label{move_aside}

\emph{Parallel-safe:} yes · \emph{Animation:} timed · \emph{Required
keys:} \texttt{who}, \texttt{prop}

Push a prop laterally without picking it up.

\textbf{Parameters:}

\begin{longtable}[]{@{}lllll@{}}
\toprule\noalign{}
Key & Type & Required & Default & Description \\
\midrule\noalign{}
\endhead
\bottomrule\noalign{}
\endlastfoot
\texttt{who} & \texttt{str} & yes & --- & Character key. \\
\texttt{prop} & \texttt{str} & yes & --- & Prop to push. \\
\texttt{direction} & \texttt{str} & no & \texttt{"right"} &
\texttt{"left"} or \texttt{"right"}. \\
\texttt{distance} & \texttt{float} & no & \texttt{0.5} & Lateral
displacement. \\
\texttt{rt} & \texttt{float} & no & \texttt{0.4} & Push duration. \\
\end{longtable}

\begin{Shaded}
\begin{Highlighting}[]
\FunctionTok{\{}\DataTypeTok{"action"}\FunctionTok{:} \StringTok{"move\_aside"}\FunctionTok{,} \DataTypeTok{"who"}\FunctionTok{:} \StringTok{"freydoon"}\FunctionTok{,} \DataTypeTok{"prop"}\FunctionTok{:} \StringTok{"chair\_2"}\FunctionTok{,}
 \DataTypeTok{"direction"}\FunctionTok{:} \StringTok{"left"}\FunctionTok{,} \DataTypeTok{"distance"}\FunctionTok{:} \DecValTok{0}\ErrorTok{.}\DecValTok{8}\FunctionTok{\}}
\end{Highlighting}
\end{Shaded}

\begin{center}\rule{0.5\linewidth}{0.5pt}\end{center}

\paragraph{\texorpdfstring{3.6.12
\texttt{pick\_up\_phone}}{3.6.12 pick\_up\_phone}}\label{pick_up_phone}

\emph{Parallel-safe:} no · \emph{Animation:} timed · \emph{Required
keys:} \texttt{who}, \texttt{prop}

Specialized variant of \texttt{pick\_up} for landline desk phones.
Raises the handset to ear height in a specific carry pose (not the
generic \texttt{carry\_hold}).

\textbf{Parameters:}

\begin{longtable}[]{@{}
  >{\raggedright\arraybackslash}p{(\columnwidth - 8\tabcolsep) * \real{0.2000}}
  >{\raggedright\arraybackslash}p{(\columnwidth - 8\tabcolsep) * \real{0.2000}}
  >{\raggedright\arraybackslash}p{(\columnwidth - 8\tabcolsep) * \real{0.2000}}
  >{\raggedright\arraybackslash}p{(\columnwidth - 8\tabcolsep) * \real{0.2000}}
  >{\raggedright\arraybackslash}p{(\columnwidth - 8\tabcolsep) * \real{0.2000}}@{}}
\toprule\noalign{}
\begin{minipage}[b]{\linewidth}\raggedright
Key
\end{minipage} & \begin{minipage}[b]{\linewidth}\raggedright
Type
\end{minipage} & \begin{minipage}[b]{\linewidth}\raggedright
Required
\end{minipage} & \begin{minipage}[b]{\linewidth}\raggedright
Default
\end{minipage} & \begin{minipage}[b]{\linewidth}\raggedright
Description
\end{minipage} \\
\midrule\noalign{}
\endhead
\bottomrule\noalign{}
\endlastfoot
\texttt{who} & \texttt{str} & yes & --- & Character key. \\
\texttt{prop} & \texttt{str} & yes & --- & Phone prop (must be
\texttt{landline\_desk} style). \\
\texttt{arm} & \texttt{str} & no & \texttt{"right"} & \texttt{"right"}
or \texttt{"left"}. \\
\texttt{rt} & \texttt{float} & no & \texttt{0.4} & Raise duration. \\
\end{longtable}

\begin{Shaded}
\begin{Highlighting}[]
\FunctionTok{\{}\DataTypeTok{"action"}\FunctionTok{:} \StringTok{"pick\_up\_phone"}\FunctionTok{,} \DataTypeTok{"who"}\FunctionTok{:} \StringTok{"sidel"}\FunctionTok{,} \DataTypeTok{"prop"}\FunctionTok{:} \StringTok{"office\_phone"}\FunctionTok{\}}
\end{Highlighting}
\end{Shaded}

\begin{center}\rule{0.5\linewidth}{0.5pt}\end{center}

\paragraph{\texorpdfstring{3.6.13
\texttt{hang\_up}}{3.6.13 hang\_up}}\label{hang_up}

\emph{Parallel-safe:} no · \emph{Animation:} timed · \emph{Required
keys:} \texttt{who}, \texttt{prop}

Return a held phone handset to its cradle. Inverse of
\texttt{pick\_up\_phone}.

\textbf{Parameters:}

\begin{longtable}[]{@{}lllll@{}}
\toprule\noalign{}
Key & Type & Required & Default & Description \\
\midrule\noalign{}
\endhead
\bottomrule\noalign{}
\endlastfoot
\texttt{who} & \texttt{str} & yes & --- & Character key. \\
\texttt{prop} & \texttt{str} & yes & --- & Held phone handset. Synonym:
\texttt{target}. \\
\texttt{target} & \texttt{str} & optional & --- & Synonym for
\texttt{prop}. \\
\texttt{rt} & \texttt{float} & no & \texttt{0.4} & Lower duration. \\
\end{longtable}

\begin{Shaded}
\begin{Highlighting}[]
\FunctionTok{\{}\DataTypeTok{"action"}\FunctionTok{:} \StringTok{"hang\_up"}\FunctionTok{,} \DataTypeTok{"who"}\FunctionTok{:} \StringTok{"sidel"}\FunctionTok{,} \DataTypeTok{"prop"}\FunctionTok{:} \StringTok{"office\_phone"}\FunctionTok{\}}
\end{Highlighting}
\end{Shaded}

\begin{center}\rule{0.5\linewidth}{0.5pt}\end{center}

\paragraph{\texorpdfstring{3.6.14
\texttt{carry}}{3.6.14 carry}}\label{carry}

\emph{Parallel-safe:} no · \emph{Animation:} timed · \emph{Required
keys:} \texttt{who}, \texttt{prop}, \texttt{x}

Walk while carrying a named prop, with optional head companions (other
props that ride along on the head).

\textbf{Parameters:}

\begin{longtable}[]{@{}
  >{\raggedright\arraybackslash}p{(\columnwidth - 8\tabcolsep) * \real{0.2000}}
  >{\raggedright\arraybackslash}p{(\columnwidth - 8\tabcolsep) * \real{0.2000}}
  >{\raggedright\arraybackslash}p{(\columnwidth - 8\tabcolsep) * \real{0.2000}}
  >{\raggedright\arraybackslash}p{(\columnwidth - 8\tabcolsep) * \real{0.2000}}
  >{\raggedright\arraybackslash}p{(\columnwidth - 8\tabcolsep) * \real{0.2000}}@{}}
\toprule\noalign{}
\begin{minipage}[b]{\linewidth}\raggedright
Key
\end{minipage} & \begin{minipage}[b]{\linewidth}\raggedright
Type
\end{minipage} & \begin{minipage}[b]{\linewidth}\raggedright
Required
\end{minipage} & \begin{minipage}[b]{\linewidth}\raggedright
Default
\end{minipage} & \begin{minipage}[b]{\linewidth}\raggedright
Description
\end{minipage} \\
\midrule\noalign{}
\endhead
\bottomrule\noalign{}
\endlastfoot
\texttt{who} & \texttt{str} & yes & --- & Character key. \\
\texttt{prop} & \texttt{str} & yes & --- & Prop being carried (must be
already held --- call \texttt{pick\_up} first). \\
\texttt{x} & \texttt{float} & yes & --- & Walk destination. \\
\texttt{companions} & \texttt{list{[}str{]}} & no & \texttt{{[}{]}} &
Other prop keys that follow the head joint while walking. \\
\texttt{torso\_companions} & \texttt{list{[}str{]}} & no &
\texttt{{[}{]}} & Other prop keys that follow the torso joint. \\
\texttt{size} & \texttt{float} & no & \texttt{1.0} & Visual scale
multiplier applied during the carry walk. \\
\texttt{color} & \texttt{str} (hex) & no & --- & Optional uniform/torso
color recolor during the carry. \\
\end{longtable}

\begin{Shaded}
\begin{Highlighting}[]
\FunctionTok{\{}\DataTypeTok{"action"}\FunctionTok{:} \StringTok{"carry"}\FunctionTok{,} \DataTypeTok{"who"}\FunctionTok{:} \StringTok{"chava"}\FunctionTok{,} \DataTypeTok{"prop"}\FunctionTok{:} \StringTok{"briefcase"}\FunctionTok{,} \DataTypeTok{"x"}\FunctionTok{:} \FloatTok{2.5}\FunctionTok{,}
 \DataTypeTok{"torso\_companions"}\FunctionTok{:} \OtherTok{[}\StringTok{"folder"}\OtherTok{]}\FunctionTok{\}}
\end{Highlighting}
\end{Shaded}

\begin{center}\rule{0.5\linewidth}{0.5pt}\end{center}

\subsubsection{3.7 Two-figure interaction}\label{two-figure-interaction}

The following actions take two character keys: \texttt{who} (the acting
character) and \texttt{target} (the other character). Both must be in
the cast and faded in. Most actions read the target via
\texttt{cast{[}target{]}{[}"fig"{]}} and skip with a console warning if
the target is not live.

\paragraph{\texorpdfstring{3.7.1 \texttt{pat}}{3.7.1 pat}}\label{pat}

\emph{Parallel-safe:} no · \emph{Animation:} timed · \emph{Required
keys:} \texttt{who}

Short repeated tapping gesture toward a prop or body area.
Single-character action --- taps the air at a target position.

\textbf{Parameters:}

\begin{longtable}[]{@{}lllll@{}}
\toprule\noalign{}
Key & Type & Required & Default & Description \\
\midrule\noalign{}
\endhead
\bottomrule\noalign{}
\endlastfoot
\texttt{who} & \texttt{str} & yes & --- & Character key. \\
\texttt{cycles} & \texttt{int} & no & \texttt{3} & Number of pat
oscillations. \\
\texttt{rt} & \texttt{float} & no & per cycle & Per-cycle run time. \\
\end{longtable}

\begin{Shaded}
\begin{Highlighting}[]
\FunctionTok{\{}\DataTypeTok{"action"}\FunctionTok{:} \StringTok{"pat"}\FunctionTok{,} \DataTypeTok{"who"}\FunctionTok{:} \StringTok{"thalia"}\FunctionTok{,} \DataTypeTok{"cycles"}\FunctionTok{:} \DecValTok{3}\FunctionTok{\}}
\end{Highlighting}
\end{Shaded}

\begin{center}\rule{0.5\linewidth}{0.5pt}\end{center}

\paragraph{\texorpdfstring{3.7.2
\texttt{pat\_head}}{3.7.2 pat\_head}}\label{pat_head}

\emph{Parallel-safe:} yes · \emph{Animation:} timed · \emph{Required
keys:} \texttt{who}, \texttt{target}

Pat another character or dog on the head. The acting figure extends an
arm to the target's head position and does a gentle up-down oscillation
(like \texttt{pat}, but targeting another character's head joint).

\textbf{Parameters:}

\begin{longtable}[]{@{}
  >{\raggedright\arraybackslash}p{(\columnwidth - 8\tabcolsep) * \real{0.2000}}
  >{\raggedright\arraybackslash}p{(\columnwidth - 8\tabcolsep) * \real{0.2000}}
  >{\raggedright\arraybackslash}p{(\columnwidth - 8\tabcolsep) * \real{0.2000}}
  >{\raggedright\arraybackslash}p{(\columnwidth - 8\tabcolsep) * \real{0.2000}}
  >{\raggedright\arraybackslash}p{(\columnwidth - 8\tabcolsep) * \real{0.2000}}@{}}
\toprule\noalign{}
\begin{minipage}[b]{\linewidth}\raggedright
Key
\end{minipage} & \begin{minipage}[b]{\linewidth}\raggedright
Type
\end{minipage} & \begin{minipage}[b]{\linewidth}\raggedright
Required
\end{minipage} & \begin{minipage}[b]{\linewidth}\raggedright
Default
\end{minipage} & \begin{minipage}[b]{\linewidth}\raggedright
Description
\end{minipage} \\
\midrule\noalign{}
\endhead
\bottomrule\noalign{}
\endlastfoot
\texttt{who} & \texttt{str} & yes & --- & Acting character. \\
\texttt{target} & \texttt{str} & yes & --- & Character or dog prop whose
head receives the pat. \\
\texttt{cycles} & \texttt{int} & no & \texttt{3} & Number of pat
oscillations. \\
\texttt{rt} & \texttt{float} & no & \texttt{0.2} & Per-oscillation run
time. \\
\texttt{arm} & \texttt{str} & no & \texttt{"auto"} & \texttt{"r"},
\texttt{"l"}, or \texttt{"auto"}. \texttt{"auto"} picks the arm on the
side closest to the target. \\
\end{longtable}

\begin{Shaded}
\begin{Highlighting}[]
\FunctionTok{\{}\DataTypeTok{"action"}\FunctionTok{:} \StringTok{"pat\_head"}\FunctionTok{,} \DataTypeTok{"who"}\FunctionTok{:} \StringTok{"thalia"}\FunctionTok{,} \DataTypeTok{"target"}\FunctionTok{:} \StringTok{"chekov"}\FunctionTok{\}}
\FunctionTok{\{}\DataTypeTok{"action"}\FunctionTok{:} \StringTok{"pat\_head"}\FunctionTok{,} \DataTypeTok{"who"}\FunctionTok{:} \StringTok{"bevers"}\FunctionTok{,} \DataTypeTok{"target"}\FunctionTok{:} \StringTok{"chekov"}\FunctionTok{,} \DataTypeTok{"cycles"}\FunctionTok{:} \DecValTok{5}\FunctionTok{,} \DataTypeTok{"arm"}\FunctionTok{:} \StringTok{"l"}\FunctionTok{\}}
\end{Highlighting}
\end{Shaded}

\begin{center}\rule{0.5\linewidth}{0.5pt}\end{center}

\paragraph{\texorpdfstring{3.7.3
\texttt{hand\_to}}{3.7.3 hand\_to}}\label{hand_to}

\emph{Parallel-safe:} yes · \emph{Animation:} timed · \emph{Required
keys:} \texttt{who}, \texttt{target}, \texttt{prop}

Hand a prop to another character: the giver extends an arm with the prop
toward the target, the target reaches to receive, and the prop transfers
ownership. After the action, the prop is parented to the target.

\textbf{Parameters:}

\begin{longtable}[]{@{}
  >{\raggedright\arraybackslash}p{(\columnwidth - 8\tabcolsep) * \real{0.2000}}
  >{\raggedright\arraybackslash}p{(\columnwidth - 8\tabcolsep) * \real{0.2000}}
  >{\raggedright\arraybackslash}p{(\columnwidth - 8\tabcolsep) * \real{0.2000}}
  >{\raggedright\arraybackslash}p{(\columnwidth - 8\tabcolsep) * \real{0.2000}}
  >{\raggedright\arraybackslash}p{(\columnwidth - 8\tabcolsep) * \real{0.2000}}@{}}
\toprule\noalign{}
\begin{minipage}[b]{\linewidth}\raggedright
Key
\end{minipage} & \begin{minipage}[b]{\linewidth}\raggedright
Type
\end{minipage} & \begin{minipage}[b]{\linewidth}\raggedright
Required
\end{minipage} & \begin{minipage}[b]{\linewidth}\raggedright
Default
\end{minipage} & \begin{minipage}[b]{\linewidth}\raggedright
Description
\end{minipage} \\
\midrule\noalign{}
\endhead
\bottomrule\noalign{}
\endlastfoot
\texttt{who} & \texttt{str} & yes & --- & Giver. \\
\texttt{target} & \texttt{str} & yes & --- & Receiver. \\
\texttt{prop} & \texttt{str} & yes & --- & Prop being transferred. \\
\texttt{rt} & \texttt{float} & no & \texttt{0.35} & Morph speed for both
characters' arms. \\
\texttt{hold} & \texttt{float} & no & \texttt{0.3} & Dwell at the
hand-off point before transfer. \\
\end{longtable}

\begin{Shaded}
\begin{Highlighting}[]
\FunctionTok{\{}\DataTypeTok{"action"}\FunctionTok{:} \StringTok{"hand\_to"}\FunctionTok{,} \DataTypeTok{"who"}\FunctionTok{:} \StringTok{"freydoon"}\FunctionTok{,} \DataTypeTok{"prop"}\FunctionTok{:} \StringTok{"phone01"}\FunctionTok{,}
 \DataTypeTok{"target"}\FunctionTok{:} \StringTok{"chava"}\FunctionTok{,} \DataTypeTok{"rt"}\FunctionTok{:} \DecValTok{0}\ErrorTok{.}\DecValTok{35}\FunctionTok{\}}
\end{Highlighting}
\end{Shaded}

\begin{center}\rule{0.5\linewidth}{0.5pt}\end{center}

\paragraph{\texorpdfstring{3.7.4 \texttt{kiss}}{3.7.4 kiss}}\label{kiss}

\emph{Parallel-safe:} yes · \emph{Animation:} timed · \emph{Required
keys:} \texttt{who}, \texttt{target}

Two characters lean forward, touch heads, and show a \textless3 glyph
with a pink dashed \texttt{(kiss)} bubble. Both characters morph to a
lean-forward pose (heads tilted toward each other), hold, then morph
back.

\textbf{Parameters:}

\begin{longtable}[]{@{}
  >{\raggedright\arraybackslash}p{(\columnwidth - 8\tabcolsep) * \real{0.2000}}
  >{\raggedright\arraybackslash}p{(\columnwidth - 8\tabcolsep) * \real{0.2000}}
  >{\raggedright\arraybackslash}p{(\columnwidth - 8\tabcolsep) * \real{0.2000}}
  >{\raggedright\arraybackslash}p{(\columnwidth - 8\tabcolsep) * \real{0.2000}}
  >{\raggedright\arraybackslash}p{(\columnwidth - 8\tabcolsep) * \real{0.2000}}@{}}
\toprule\noalign{}
\begin{minipage}[b]{\linewidth}\raggedright
Key
\end{minipage} & \begin{minipage}[b]{\linewidth}\raggedright
Type
\end{minipage} & \begin{minipage}[b]{\linewidth}\raggedright
Required
\end{minipage} & \begin{minipage}[b]{\linewidth}\raggedright
Default
\end{minipage} & \begin{minipage}[b]{\linewidth}\raggedright
Description
\end{minipage} \\
\midrule\noalign{}
\endhead
\bottomrule\noalign{}
\endlastfoot
\texttt{who} & \texttt{str} & yes & --- & Acting character. \\
\texttt{target} & \texttt{str} & yes & --- & Recipient. \\
\texttt{hold} & \texttt{float} & no & \texttt{1.5} & Kiss duration in
seconds. \\
\texttt{rt} & \texttt{float} & no & \texttt{0.4} & Lean-in / lean-back
morph speed. \\
\texttt{glyph} & \texttt{str} & no & \texttt{"\textless{}3"} & The glyph
character displayed during the kiss. \\
\texttt{color} & \texttt{str} (hex) & no & \texttt{"\#ff69b4"} & Glyph
and bubble color. \\
\texttt{bubble} & \texttt{bool} & no & \texttt{true} & Show the
\texttt{(kiss)} bubble. Set \texttt{false} to suppress for a silent
kiss. \\
\end{longtable}

\begin{Shaded}
\begin{Highlighting}[]
\FunctionTok{\{}\DataTypeTok{"action"}\FunctionTok{:} \StringTok{"kiss"}\FunctionTok{,} \DataTypeTok{"who"}\FunctionTok{:} \StringTok{"nona"}\FunctionTok{,} \DataTypeTok{"target"}\FunctionTok{:} \StringTok{"factor"}\FunctionTok{,} \DataTypeTok{"hold"}\FunctionTok{:} \FloatTok{2.0}\FunctionTok{\}}
\FunctionTok{\{}\DataTypeTok{"action"}\FunctionTok{:} \StringTok{"kiss"}\FunctionTok{,} \DataTypeTok{"who"}\FunctionTok{:} \StringTok{"freydoon"}\FunctionTok{,} \DataTypeTok{"target"}\FunctionTok{:} \StringTok{"athena"}\FunctionTok{,}
 \DataTypeTok{"glyph"}\FunctionTok{:} \StringTok{"*"}\FunctionTok{,} \DataTypeTok{"color"}\FunctionTok{:} \StringTok{"\#ffdd55"}\FunctionTok{,} \DataTypeTok{"bubble"}\FunctionTok{:} \KeywordTok{false}\FunctionTok{\}}
\end{Highlighting}
\end{Shaded}

\begin{center}\rule{0.5\linewidth}{0.5pt}\end{center}

\paragraph{\texorpdfstring{3.7.5
\texttt{hold\_hands}}{3.7.5 hold\_hands}}\label{hold_hands}

\emph{Parallel-safe:} yes · \emph{Animation:} timed · \emph{Required
keys:} \texttt{who}, \texttt{target}

Two adjacent characters extend their inner arms until wrists meet at a
midpoint. The action figures out which character is left vs.~right of
the other, then morphs each figure's inner arm (the left character's
right arm; the right character's left arm) to a midpoint handhold pose.

After the \texttt{hold} duration, both characters return to their
pre-handhold rest poses. Not a persistent state --- for a sustained
handhold, run a longer \texttt{hold} or repeat the action.

\textbf{Parameters:}

\begin{longtable}[]{@{}
  >{\raggedright\arraybackslash}p{(\columnwidth - 8\tabcolsep) * \real{0.2000}}
  >{\raggedright\arraybackslash}p{(\columnwidth - 8\tabcolsep) * \real{0.2000}}
  >{\raggedright\arraybackslash}p{(\columnwidth - 8\tabcolsep) * \real{0.2000}}
  >{\raggedright\arraybackslash}p{(\columnwidth - 8\tabcolsep) * \real{0.2000}}
  >{\raggedright\arraybackslash}p{(\columnwidth - 8\tabcolsep) * \real{0.2000}}@{}}
\toprule\noalign{}
\begin{minipage}[b]{\linewidth}\raggedright
Key
\end{minipage} & \begin{minipage}[b]{\linewidth}\raggedright
Type
\end{minipage} & \begin{minipage}[b]{\linewidth}\raggedright
Required
\end{minipage} & \begin{minipage}[b]{\linewidth}\raggedright
Default
\end{minipage} & \begin{minipage}[b]{\linewidth}\raggedright
Description
\end{minipage} \\
\midrule\noalign{}
\endhead
\bottomrule\noalign{}
\endlastfoot
\texttt{who} & \texttt{str} & yes & --- & Acting character. \\
\texttt{target} & \texttt{str} & yes & --- & Other character. Order
doesn't matter; the action picks left/right based on x positions. \\
\texttt{hold} & \texttt{float} & no & \texttt{2.0} & Hold duration
before release. \\
\texttt{rt} & \texttt{float} & no & \texttt{0.35} & Arm-extend morph
speed. Release uses \texttt{rt\ *\ 0.8}. \\
\end{longtable}

\begin{Shaded}
\begin{Highlighting}[]
\FunctionTok{\{}\DataTypeTok{"action"}\FunctionTok{:} \StringTok{"hold\_hands"}\FunctionTok{,} \DataTypeTok{"who"}\FunctionTok{:} \StringTok{"thalia"}\FunctionTok{,} \DataTypeTok{"target"}\FunctionTok{:} \StringTok{"bevers"}\FunctionTok{,} \DataTypeTok{"hold"}\FunctionTok{:} \FloatTok{2.0}\FunctionTok{\}}
\end{Highlighting}
\end{Shaded}

\begin{center}\rule{0.5\linewidth}{0.5pt}\end{center}

\paragraph{\texorpdfstring{3.7.6
\texttt{grab\_arm}}{3.7.6 grab\_arm}}\label{grab_arm}

\emph{Parallel-safe:} no · \emph{Animation:} timed · \emph{Required
keys:} \texttt{who}, \texttt{target}

Grab another character's arm from behind: the actor reaches forward and
seizes the target's wrist, locking that arm in place.

\textbf{The action does NOT auto-return either character to a rest pose}
--- the constraint persists until \texttt{release\_arm} is called. The
target's seized arm is recorded on
\texttt{target\_fig.\_restrained\_arm} (\texttt{"r"} or \texttt{"l"}),
which \texttt{twist\_arm\_behind} reads to know which arm to fold and
which other actions can check to skip the constrained character.

\textbf{Parameters:}

\begin{longtable}[]{@{}
  >{\raggedright\arraybackslash}p{(\columnwidth - 8\tabcolsep) * \real{0.2000}}
  >{\raggedright\arraybackslash}p{(\columnwidth - 8\tabcolsep) * \real{0.2000}}
  >{\raggedright\arraybackslash}p{(\columnwidth - 8\tabcolsep) * \real{0.2000}}
  >{\raggedright\arraybackslash}p{(\columnwidth - 8\tabcolsep) * \real{0.2000}}
  >{\raggedright\arraybackslash}p{(\columnwidth - 8\tabcolsep) * \real{0.2000}}@{}}
\toprule\noalign{}
\begin{minipage}[b]{\linewidth}\raggedright
Key
\end{minipage} & \begin{minipage}[b]{\linewidth}\raggedright
Type
\end{minipage} & \begin{minipage}[b]{\linewidth}\raggedright
Required
\end{minipage} & \begin{minipage}[b]{\linewidth}\raggedright
Default
\end{minipage} & \begin{minipage}[b]{\linewidth}\raggedright
Description
\end{minipage} \\
\midrule\noalign{}
\endhead
\bottomrule\noalign{}
\endlastfoot
\texttt{who} & \texttt{str} & yes & --- & Acting character. \\
\texttt{target} & \texttt{str} & yes & --- & Character whose arm is
grabbed. \\
\texttt{arm} & \texttt{str} & no & \texttt{"auto"} & Which of the
\emph{target's} arms to seize: \texttt{"r"}, \texttt{"l"}, or
\texttt{"auto"} (picks the arm on the side closest to the actor). \\
\texttt{rt} & \texttt{float} & no & \texttt{0.3} & Morph speed for the
grab. \\
\end{longtable}

\begin{Shaded}
\begin{Highlighting}[]
\FunctionTok{\{}\DataTypeTok{"action"}\FunctionTok{:} \StringTok{"grab\_arm"}\FunctionTok{,} \DataTypeTok{"who"}\FunctionTok{:} \StringTok{"chava"}\FunctionTok{,} \DataTypeTok{"target"}\FunctionTok{:} \StringTok{"brad"}\FunctionTok{,} \DataTypeTok{"arm"}\FunctionTok{:} \StringTok{"r"}\FunctionTok{\}}
\FunctionTok{\{}\DataTypeTok{"action"}\FunctionTok{:} \StringTok{"grab\_arm"}\FunctionTok{,} \DataTypeTok{"who"}\FunctionTok{:} \StringTok{"chava"}\FunctionTok{,} \DataTypeTok{"target"}\FunctionTok{:} \StringTok{"brad"}\FunctionTok{\}}
\end{Highlighting}
\end{Shaded}

\begin{center}\rule{0.5\linewidth}{0.5pt}\end{center}

\paragraph{\texorpdfstring{3.7.7
\texttt{twist\_arm\_behind}}{3.7.7 twist\_arm\_behind}}\label{twist_arm_behind}

\emph{Parallel-safe:} no · \emph{Animation:} timed · \emph{Required
keys:} \texttt{who}, \texttt{target}

Escalate a \texttt{grab\_arm} into an arm-lock: fold the target's seized
arm behind their back and tilt their torso forward under the pressure.
Pairs with the \texttt{restrained\_side} / \texttt{grip\_behind} poses
(see §1.7).

\textbf{Must be called \emph{after} \texttt{grab\_arm} on the same
target.} Reads \texttt{target\_fig.\_restrained\_arm} to know which arm
to fold; if no arm is currently grabbed, the action prints a warning and
skips.

\textbf{Parameters:}

\begin{longtable}[]{@{}
  >{\raggedright\arraybackslash}p{(\columnwidth - 8\tabcolsep) * \real{0.2000}}
  >{\raggedright\arraybackslash}p{(\columnwidth - 8\tabcolsep) * \real{0.2000}}
  >{\raggedright\arraybackslash}p{(\columnwidth - 8\tabcolsep) * \real{0.2000}}
  >{\raggedright\arraybackslash}p{(\columnwidth - 8\tabcolsep) * \real{0.2000}}
  >{\raggedright\arraybackslash}p{(\columnwidth - 8\tabcolsep) * \real{0.2000}}@{}}
\toprule\noalign{}
\begin{minipage}[b]{\linewidth}\raggedright
Key
\end{minipage} & \begin{minipage}[b]{\linewidth}\raggedright
Type
\end{minipage} & \begin{minipage}[b]{\linewidth}\raggedright
Required
\end{minipage} & \begin{minipage}[b]{\linewidth}\raggedright
Default
\end{minipage} & \begin{minipage}[b]{\linewidth}\raggedright
Description
\end{minipage} \\
\midrule\noalign{}
\endhead
\bottomrule\noalign{}
\endlastfoot
\texttt{who} & \texttt{str} & yes & --- & Acting character. \\
\texttt{target} & \texttt{str} & yes & --- & Character being restrained
(must already be grabbed). \\
\texttt{tilt} & \texttt{float} & no & \texttt{0.12} & How far the
target's torso tilts forward, in unscaled pose units. Higher values read
as more painful / more aggressive lock. \\
\texttt{rt} & \texttt{float} & no & \texttt{0.35} & Morph speed. \\
\end{longtable}

\begin{Shaded}
\begin{Highlighting}[]
\FunctionTok{\{}\DataTypeTok{"action"}\FunctionTok{:} \StringTok{"twist\_arm\_behind"}\FunctionTok{,} \DataTypeTok{"who"}\FunctionTok{:} \StringTok{"chava"}\FunctionTok{,} \DataTypeTok{"target"}\FunctionTok{:} \StringTok{"brad"}\FunctionTok{,} \DataTypeTok{"tilt"}\FunctionTok{:} \DecValTok{0}\ErrorTok{.}\DecValTok{15}\FunctionTok{\}}
\end{Highlighting}
\end{Shaded}

\begin{center}\rule{0.5\linewidth}{0.5pt}\end{center}

\paragraph{\texorpdfstring{3.7.8
\texttt{release\_arm}}{3.7.8 release\_arm}}\label{release_arm}

\emph{Parallel-safe:} no · \emph{Animation:} timed · \emph{Required
keys:} \texttt{who}, \texttt{target}

Release a grabbed arm: both characters return to their pre-grab rest
poses, and the \texttt{target\_fig.\_restrained\_arm} constraint state
is cleared.

\textbf{Target must currently be restrained.} If \texttt{release\_arm}
is called on a target with no active grab, the action prints a warning
and skips. The full restraint sequence is therefore: \texttt{grab\_arm}
-\textgreater{} (optional \texttt{twist\_arm\_behind}) -\textgreater{}
\texttt{release\_arm}.

\textbf{Parameters:}

\begin{longtable}[]{@{}lllll@{}}
\toprule\noalign{}
Key & Type & Required & Default & Description \\
\midrule\noalign{}
\endhead
\bottomrule\noalign{}
\endlastfoot
\texttt{who} & \texttt{str} & yes & --- & Actor who originally
grabbed. \\
\texttt{target} & \texttt{str} & yes & --- & Character being
released. \\
\texttt{rt} & \texttt{float} & no & \texttt{0.3} & Morph speed for the
release. \\
\end{longtable}

\begin{Shaded}
\begin{Highlighting}[]
\FunctionTok{\{}\DataTypeTok{"action"}\FunctionTok{:} \StringTok{"release\_arm"}\FunctionTok{,} \DataTypeTok{"who"}\FunctionTok{:} \StringTok{"chava"}\FunctionTok{,} \DataTypeTok{"target"}\FunctionTok{:} \StringTok{"brad"}\FunctionTok{\}}
\end{Highlighting}
\end{Shaded}

\begin{center}\rule{0.5\linewidth}{0.5pt}\end{center}

\paragraph{\texorpdfstring{3.7.9
\texttt{reach\_character}}{3.7.9 reach\_character}}\label{reach_character}

\emph{Parallel-safe:} no · \emph{Animation:} timed · \emph{Required
keys:} \texttt{who}, \texttt{target}

Jab one arm toward a body zone on another cast member, then retract.
Similar mechanics to \texttt{punch\_button}, but with a character as the
target.

\textbf{Parameters:}

\begin{longtable}[]{@{}
  >{\raggedright\arraybackslash}p{(\columnwidth - 8\tabcolsep) * \real{0.2000}}
  >{\raggedright\arraybackslash}p{(\columnwidth - 8\tabcolsep) * \real{0.2000}}
  >{\raggedright\arraybackslash}p{(\columnwidth - 8\tabcolsep) * \real{0.2000}}
  >{\raggedright\arraybackslash}p{(\columnwidth - 8\tabcolsep) * \real{0.2000}}
  >{\raggedright\arraybackslash}p{(\columnwidth - 8\tabcolsep) * \real{0.2000}}@{}}
\toprule\noalign{}
\begin{minipage}[b]{\linewidth}\raggedright
Key
\end{minipage} & \begin{minipage}[b]{\linewidth}\raggedright
Type
\end{minipage} & \begin{minipage}[b]{\linewidth}\raggedright
Required
\end{minipage} & \begin{minipage}[b]{\linewidth}\raggedright
Default
\end{minipage} & \begin{minipage}[b]{\linewidth}\raggedright
Description
\end{minipage} \\
\midrule\noalign{}
\endhead
\bottomrule\noalign{}
\endlastfoot
\texttt{who} & \texttt{str} & yes & --- & Acting character. \\
\texttt{target} & \texttt{str} & yes & --- & Target character. \\
\texttt{zone} & \texttt{str} & no & \texttt{"torso"} & Body zone of
target: \texttt{"torso"}, \texttt{"head"}, \texttt{"larm"},
\texttt{"rarm"}, etc. Mapped to joint name on the target. \\
\texttt{arm} & \texttt{str} & no & \texttt{"right"} & Acting character's
arm. \\
\texttt{offset\_x} & \texttt{float} & no & \texttt{0.0} & Fine-tune
landing point. \\
\texttt{offset\_y} & \texttt{float} & no & \texttt{0.0} & As above. \\
\texttt{hold} & \texttt{float} & no & \texttt{0.2} & Hold at full
extension. \\
\texttt{rt} & \texttt{float} & no & \texttt{0.3} & Reach and retract
duration. \\
\end{longtable}

\begin{Shaded}
\begin{Highlighting}[]
\FunctionTok{\{}\DataTypeTok{"action"}\FunctionTok{:} \StringTok{"reach\_character"}\FunctionTok{,} \DataTypeTok{"who"}\FunctionTok{:} \StringTok{"freydoon"}\FunctionTok{,} \DataTypeTok{"target"}\FunctionTok{:} \StringTok{"brad"}\FunctionTok{,}
 \DataTypeTok{"zone"}\FunctionTok{:} \StringTok{"torso"}\FunctionTok{,} \DataTypeTok{"arm"}\FunctionTok{:} \StringTok{"right"}\FunctionTok{\}}
\end{Highlighting}
\end{Shaded}

\textbf{Note:} the cast registry lookup for \texttt{target} goes through
\texttt{cast{[}target\_name{]}{[}"fig"{]}} (the live \texttt{HumanGraph}
object), not the cast spec dict. If \texttt{target} isn't in cast or
hasn't been faded in, the action skips with a warning.

\begin{center}\rule{0.5\linewidth}{0.5pt}\end{center}

\subsubsection{3.8 Multi-target}\label{multi-target}

\paragraph{\texorpdfstring{3.8.1
\texttt{group\_translate}}{3.8.1 group\_translate}}\label{group_translate}

\emph{Parallel-safe:} no · \emph{Animation:} timed · \emph{Required
keys:} \texttt{who}

Move multiple characters (and optionally props) simultaneously along a
vector. Used for camera-like ``treadmill'' effects in walk-and-talk: the
world translates underneath the characters, suggesting continuous motion
without the characters changing scene position.

\textbf{Parameters:}

\begin{longtable}[]{@{}
  >{\raggedright\arraybackslash}p{(\columnwidth - 8\tabcolsep) * \real{0.2000}}
  >{\raggedright\arraybackslash}p{(\columnwidth - 8\tabcolsep) * \real{0.2000}}
  >{\raggedright\arraybackslash}p{(\columnwidth - 8\tabcolsep) * \real{0.2000}}
  >{\raggedright\arraybackslash}p{(\columnwidth - 8\tabcolsep) * \real{0.2000}}
  >{\raggedright\arraybackslash}p{(\columnwidth - 8\tabcolsep) * \real{0.2000}}@{}}
\toprule\noalign{}
\begin{minipage}[b]{\linewidth}\raggedright
Key
\end{minipage} & \begin{minipage}[b]{\linewidth}\raggedright
Type
\end{minipage} & \begin{minipage}[b]{\linewidth}\raggedright
Required
\end{minipage} & \begin{minipage}[b]{\linewidth}\raggedright
Default
\end{minipage} & \begin{minipage}[b]{\linewidth}\raggedright
Description
\end{minipage} \\
\midrule\noalign{}
\endhead
\bottomrule\noalign{}
\endlastfoot
\texttt{who} & \texttt{list{[}str{]}} or \texttt{str} & yes & --- &
Character keys to translate. Either a list of names or
\texttt{"all"}. \\
\texttt{dx} & \texttt{float} & no & \texttt{0.0} & Horizontal
translation. \\
\texttt{dy} & \texttt{float} & no & \texttt{0.0} & Vertical
translation. \\
\texttt{props} & \texttt{list{[}str{]}} & no & \texttt{{[}{]}} & Prop
keys to also translate. \\
\texttt{rt} & \texttt{float} & no & \texttt{1.0} & Translation
duration. \\
\end{longtable}

\begin{Shaded}
\begin{Highlighting}[]
\FunctionTok{\{}\DataTypeTok{"action"}\FunctionTok{:} \StringTok{"group\_translate"}\FunctionTok{,}
 \DataTypeTok{"who"}\FunctionTok{:} \OtherTok{[}\StringTok{"bevers"}\OtherTok{,} \StringTok{"thalia"}\OtherTok{]}\FunctionTok{,}
 \DataTypeTok{"props"}\FunctionTok{:} \OtherTok{[}\StringTok{"bench\_1"}\OtherTok{,} \StringTok{"bench\_2"}\OtherTok{]}\FunctionTok{,}
 \DataTypeTok{"dx"}\FunctionTok{:} \FloatTok{{-}2.0}\FunctionTok{,} \DataTypeTok{"rt"}\FunctionTok{:} \FloatTok{1.5}\FunctionTok{\}}
\end{Highlighting}
\end{Shaded}

\textbf{Parallel-safety:} currently marked unsafe. A clean fix is on the
back-burner list (walk-and-talk wants this inside \texttt{parallel} for
the talk-during-translate pattern). Workaround: split the talk into
stationary beats interleaved with translation beats.

\begin{center}\rule{0.5\linewidth}{0.5pt}\end{center}

\emph{Section 4 (meta --- name resolution, the \texttt{hidden} flag,
z-order, color authority chain, registry behavior) is the final
substep.}

\subsection{4. Meta}\label{meta}

Cross-cutting concerns that don't fit cleanly inside any single field or
action: how names resolve to figures and props, where the
\texttt{hidden} flag does and doesn't apply, how z-order works, the
color authority chain across the registry / scene / annotation layers,
registry maintenance, and a consolidated gotchas appendix.

\subsubsection{4.1 Name resolution}\label{name-resolution}

A step's target is identified by either \texttt{who} or \texttt{prop}.
Both are string keys into separate registries --- \emph{characters} in
the cast registry, \emph{props} in the prop registry --- but for some
character types the same name resolves through both.

\textbf{The cast registry} is the dict \texttt{cast{[}name{]}} keyed on
character names from the \texttt{cast} action block. Each entry holds
the canonical spec plus a \texttt{"fig"} field that points to the live
figure instance after fade-in:

\begin{Shaded}
\begin{Highlighting}[]
\NormalTok{cast[}\StringTok{"bevers"}\NormalTok{] }\OperatorTok{=}\NormalTok{ \{}
    \StringTok{"figure\_type"}\NormalTok{: }\StringTok{"alien"}\NormalTok{,}
    \StringTok{"build"}\NormalTok{: }\StringTok{"alien"}\NormalTok{,}
    \StringTok{"fig"}\NormalTok{: }\OperatorTok{\textless{}}\NormalTok{AlienGraph instance}\OperatorTok{\textgreater{}}\NormalTok{,  }\CommentTok{\# populated by fade\_in}
    \CommentTok{\# ... all the creation{-}time fields from §1}
\NormalTok{\}}
\end{Highlighting}
\end{Shaded}

\texttt{cast{[}name{]}{[}"fig"{]}} is \texttt{None} between fade-out and
the next fade-in for the same character.

\textbf{The prop registry} is the \texttt{PropRegistry} instance
(\texttt{props} in handler signatures). Each entry is the prop's Manim
\texttt{VGroup} decorated with \texttt{pam\_*} attributes
(\texttt{pam\_name}, \texttt{pam\_type}, \texttt{pam\_x},
\texttt{pam\_y}, etc.).

\paragraph{Path C dual registration}\label{path-c-dual-registration}

DogGraph instances live in both registries simultaneously (v0.9.14+).
When a dog is declared in a \texttt{cast} block, \texttt{pam\_player}'s
\texttt{fade\_in} branch:

\begin{enumerate}
\def\labelenumi{\arabic{enumi}.}
\tightlist
\item
  Builds the \texttt{DogGraph} and assigns it to
  \texttt{cast{[}name{]}{[}"fig"{]}}.
\item
  Wraps it as a prop: binds \texttt{dog.group} to a local, stamps
  \texttt{pam\_dog\ =\ fig} plus the other \texttt{pam\_*} attributes on
  the group, and calls \texttt{props.add(name,\ dog\_group)}.
\item
  Skips step 2 if a prop with that name already exists
  (e.g.~\texttt{spawn\_prop} ran first).
\end{enumerate}

The result: \texttt{\{"who":\ "chekov"\}} and
\texttt{\{"prop":\ "chekov"\}} both refer to the same \texttt{DogGraph}.
Locomotion (\texttt{trot\_to}) and speech (\texttt{prop\_say}) resolve
correctly from either side.

\textbf{Why the \texttt{dog.group} local-bind matters:}
\texttt{DogGraph.group} is a \texttt{@property} that returns a fresh
\texttt{VGroup} on every access. Setting attributes on
\texttt{dog.group} directly would stamp them onto a throwaway object and
leave the registered group bare. The pattern is:

\begin{Shaded}
\begin{Highlighting}[]
\NormalTok{dog\_group }\OperatorTok{=}\NormalTok{ fig.group     }\CommentTok{\# bind once}
\NormalTok{dog\_group.pam\_name }\OperatorTok{=}\NormalTok{ name}
\NormalTok{dog\_group.pam\_dog  }\OperatorTok{=}\NormalTok{ fig}
\CommentTok{\# ...}
\NormalTok{props.add(name, dog\_group)}
\end{Highlighting}
\end{Shaded}

Same pattern is used in both \texttt{spawn\_prop}'s dog branch and
\texttt{fade\_in}'s dual-registration branch --- symmetric registry
maintenance.

\paragraph{\texorpdfstring{The \texttt{\_resolve\_speaker} canonical
helper}{The \_resolve\_speaker canonical helper}}\label{the-_resolve_speaker-canonical-helper}

\texttt{pam.actions.\_resolve\_speaker(step,\ cast=cast,\ props=props)}
is the canonical name-to-figure lookup. It reads
\texttt{step.get("who")} first, then falls back to
\texttt{step.get("prop")}, and returns a \texttt{ResolvedSpeaker}
namedtuple:

\begin{verbatim}
ResolvedSpeaker(
    name: str | None,
    fig:  HumanGraph | AlienGraph | DogGraph | GovernorGraph | None,
    prop: VGroup | None,
)
\end{verbatim}

Both \texttt{fig} and \texttt{prop} may be populated simultaneously for
dual-registered entries. Callers select whichever they need.

\textbf{The legacy fallthrough:} if the cast lookup misses but a prop
with \texttt{pam\_dog} exists, \texttt{\_resolve\_speaker} surfaces the
\texttt{DogGraph} as \texttt{fig}. This handles \texttt{spawn\_prop}'d
dogs cleanly even before they're cast-registered. The fallthrough
preserves backward compatibility with pre-Path-C screenplays.

\textbf{Resolution failure:} if neither \texttt{who} nor \texttt{prop}
resolves, the action handler typically prints a console warning and
returns \texttt{None} (silent no-op). This is the documented behavior
for safety --- a malformed step should not crash a render.

\paragraph{Resolution priority
(action-specific)}\label{resolution-priority-action-specific}

Some actions accept both \texttt{who} and \texttt{prop} and have
explicit priority rules:

\begin{itemize}
\tightlist
\item
  \textbf{\texttt{rotate}} --- if both are present, \texttt{prop} wins.
  Documented in §2.5.
\item
  \textbf{\texttt{grab}} --- accepts \texttt{prop} or \texttt{target} as
  synonyms for the held-item key. \texttt{prop} wins if both are
  present.
\item
  \textbf{\texttt{hang\_up}} --- accepts \texttt{prop} or
  \texttt{target} as synonyms for the phone-handset key.
\end{itemize}

These rules are per-action and not centralized; check the action's
docstring in §3 for the specific resolution path.

\begin{center}\rule{0.5\linewidth}{0.5pt}\end{center}

\subsubsection{\texorpdfstring{4.2 The \texttt{hidden}
flag}{4.2 The hidden flag}}\label{the-hidden-flag}

\textbf{The \texttt{hidden} flag is a prop concept, not a character
concept.}

For props, \texttt{"hidden":\ true} in a \texttt{props} (or
\texttt{scene\_props}) action declares the prop but doesn't add it to
the scene visibly. The prop is registered (so \texttt{parent} /
\texttt{attach} references can resolve it) but rendered at opacity 0. A
subsequent \texttt{spawn\_prop} action reveals it via a fade-in.

\begin{Shaded}
\begin{Highlighting}[]
\FunctionTok{\{}
  \DataTypeTok{"action"}\FunctionTok{:} \StringTok{"props"}\FunctionTok{,}
  \DataTypeTok{"items"}\FunctionTok{:} \FunctionTok{\{}
    \DataTypeTok{"athena\_laptop"}\FunctionTok{:} \FunctionTok{\{}\DataTypeTok{"type"}\FunctionTok{:} \StringTok{"laptop"}\FunctionTok{,} \DataTypeTok{"x"}\FunctionTok{:} \FloatTok{4.5}\FunctionTok{,} \DataTypeTok{"y"}\FunctionTok{:} \DecValTok{0}\ErrorTok{.}\DecValTok{4}\FunctionTok{,}
                      \DataTypeTok{"parent"}\FunctionTok{:} \StringTok{"athena"}\FunctionTok{,} \DataTypeTok{"hidden"}\FunctionTok{:} \KeywordTok{true}\FunctionTok{\}}
  \FunctionTok{\}}
\FunctionTok{\}}
\end{Highlighting}
\end{Shaded}

Later in the screenplay:

\begin{Shaded}
\begin{Highlighting}[]
\FunctionTok{\{}\DataTypeTok{"action"}\FunctionTok{:} \StringTok{"spawn\_prop"}\FunctionTok{,} \DataTypeTok{"prop"}\FunctionTok{:} \StringTok{"athena\_laptop"}\FunctionTok{\}}
\end{Highlighting}
\end{Shaded}

For \emph{characters}, visibility is controlled by \texttt{fade\_in} /
\texttt{fade\_out}. There is no character-side \texttt{hidden} flag ---
a character is either in the cast block (and faded in to become visible)
or not. Setting \texttt{"hidden":\ true} on a cast entry is silently
ignored by \texttt{pam\_player}.

\textbf{Why this is worth flagging:} the v1.0.0 back-burner note
tentatively listed a character-side \texttt{hidden} flag. The
implementation didn't land --- character visibility remained tied to
\texttt{fade\_in} / \texttt{fade\_out}. Authors who want a character
on-stage but invisible (e.g.~for an instant reveal) should
\texttt{fade\_in} with a very short \texttt{rt} rather than reach for a
non-existent flag.

\begin{center}\rule{0.5\linewidth}{0.5pt}\end{center}

\subsubsection{4.3 Z-order}\label{z-order}

\textbf{PAM's z-order is implicit.} The order in which mobjects are
added to the scene determines their painter's-algorithm layering: later
additions draw on top of earlier ones. There is no explicit
\texttt{z-order} field on cast entries or actions.

\textbf{Order observed in a typical scene:}

\begin{enumerate}
\def\labelenumi{\arabic{enumi}.}
\tightlist
\item
  \textbf{Scene props} --- declared in the \texttt{scene\_props} or
  \texttt{props} action block at the top. Drawn first; sit in the
  background.
\item
  \textbf{Characters} --- added during \texttt{fade\_in}. Drawn on top
  of props. Multiple characters layer in the order they're faded in.
\item
  \textbf{Mid-scene-spawned props} (\texttt{spawn\_prop}) --- drawn on
  top of characters that already faded in. This is how a held laptop or
  carried briefcase renders in front of the character's body.
\item
  \textbf{Attached faces} (\texttt{attach\_face}, v0.9.15) --- added
  per-character at attach time. Drawn on top of the head dot.
\item
  \textbf{Speech bubbles} (\texttt{say}, \texttt{prop\_say}) --- added
  at speech time. Drawn on top of everything below them at that moment.
\end{enumerate}

\textbf{Where z-order matters:}

\begin{itemize}
\tightlist
\item
  \textbf{Face on top of head dot.} \texttt{attach\_face} adds the face
  mobject \emph{after} the figure was faded in, so the face draws on
  top. The head dot is set to opacity 0 underneath; the face replaces it
  visually.
\item
  \textbf{Bubble above face above head dot.} A \texttt{say} step adds
  the bubble after \texttt{attach\_face}, so the bubble draws above the
  face. Visually: dot (hidden) -\textgreater{} face -\textgreater{}
  bubble, stacked correctly.
\item
  \textbf{Foreground-character occlusion.} If character A (with an
  attached face) walks behind character B (no face), the face draws on
  top of B because it was added later. In practice this is rare in TNTD
  blocking; flagged for future awareness in §2.7.
\item
  \textbf{The \texttt{focus\_reset} gotcha.} See §2.13: persistent
  bubbles can be drawn over by figures restored by \texttt{focus\_reset}
  because Manim's opacity restoration interacts with z-order. Workaround
  documented there.
\end{itemize}

\textbf{No explicit override:} there is currently no action like
\texttt{bring\_to\_front} or \texttt{send\_to\_back} for characters. An
author who needs a specific layering should \texttt{fade\_in} characters
in the desired back-to-front order. This is on the v1.0.0 wishlist.

\begin{center}\rule{0.5\linewidth}{0.5pt}\end{center}

\subsubsection{4.4 Color authority chain}\label{color-authority-chain}

A character's color can be specified in three places:

\begin{enumerate}
\def\labelenumi{\arabic{enumi}.}
\tightlist
\item
  The canonical registry file (e.g.~\texttt{tntd\_characters.json}).
\item
  The scene's cast block (the
  \texttt{\{"action":\ "cast",\ "characters":\ \{...\}\}} step).
\item
  A \texttt{{[}{[}\ CHARACTER:\ ...\ {]}{]}} Fountain+ annotation in the
  \texttt{.fountain} source.
\end{enumerate}

These are merged by \texttt{fountain2pam.py} at conversion time. The
merge order and authority are:

\begin{verbatim}
[[ CHARACTER: ]] annotation     (lowest authority; applied first)
        down then
canonical registry file         (highest authority for canonical fields)
        down then
scene cast block                (scene-specific overrides only)
\end{verbatim}

\paragraph{What's canonical, what's
scene-specific}\label{whats-canonical-whats-scene-specific}

From §1.13 (repeated here for convenience):

\begin{longtable}[]{@{}ll@{}}
\toprule\noalign{}
Canonical (registry-stored) & Scene-specific (never registry-stored) \\
\midrule\noalign{}
\endhead
\bottomrule\noalign{}
\endlastfoot
\texttt{figure\_type} & \texttt{offset} \\
\texttt{build} & \texttt{scale} \\
\texttt{gender} & \texttt{pose} \\
\texttt{color} & \texttt{facing} (dog only) \\
\texttt{torso\_color} & \\
\texttt{style} & \\
\texttt{default\_scale} & \\
\texttt{tic\_profile} & \\
\texttt{uniforms} & \\
\end{longtable}

The canonical registry is authoritative for color. If the registry has
\texttt{"bevers":\ \{"color":\ "\#44bb66",\ ...\}}, then every scene
Bevers appears in renders him at that color, regardless of what the
per-scene cast block or the \texttt{{[}{[}\ CHARACTER:\ {]}{]}}
annotation says.

\paragraph{The v0.9.8 fix in detail}\label{the-v0.9.8-fix-in-detail}

Before v0.9.8, \texttt{fountain2pam.py} would apply a round-robin color
palette to any character missing an explicit color, then refuse to
overwrite it from the registry. The result: Bevers might appear in green
in one scene and amber in another, depending on alphabetical order of
the cast block.

The fix re-derives the full six-key palette from the canonical
\texttt{color} at conversion time and stamps it into the character's
\texttt{style} dict, overwriting whatever the round-robin assigned. This
is why:

\begin{itemize}
\tightlist
\item
  Editing the registry is the only durable way to change a character's
  color.
\item
  Per-scene \texttt{color} / \texttt{style.edge\_color} overrides are
  not durable --- the next conversion stomps them.
\item
  A \texttt{{[}{[}\ CHARACTER:\ color=...\ {]}{]}} annotation in the
  Fountain file is applied during conversion \emph{before} the registry
  merge, so it can seed a new character but cannot durably override an
  existing canonical color.
\end{itemize}

\paragraph{Practical workflow}\label{practical-workflow}

\begin{itemize}
\tightlist
\item
  \textbf{To add a new character:} write a
  \texttt{{[}{[}\ CHARACTER:\ {]}{]}} annotation in the Fountain file
  with \texttt{color=}, \texttt{torso\_color=}, \texttt{gender=}, etc.
  After the first conversion, the canonical entry exists in the
  registry; subsequent scenes inherit it automatically.
\item
  \textbf{To change an existing character's color:} edit the registry
  file directly. Don't rely on per-scene overrides.
\item
  \textbf{To override color in one scene only:} edit the scene's
  converted JSON cast block after conversion, and don't re-run the
  converter on that scene. Or accept the canonical color.
\end{itemize}

\begin{center}\rule{0.5\linewidth}{0.5pt}\end{center}

\subsubsection{4.5 Registry behavior}\label{registry-behavior}

\paragraph{The cast registry}\label{the-cast-registry}

The \texttt{cast} dict is constructed at scene start from the
\texttt{cast} action block. Each entry follows the structure:

\begin{Shaded}
\begin{Highlighting}[]
\NormalTok{cast[name] }\OperatorTok{=}\NormalTok{ \{}
    \StringTok{"fig"}\NormalTok{:         }\OperatorTok{\textless{}}\NormalTok{live figure instance}\OperatorTok{\textgreater{}}\NormalTok{,  }\CommentTok{\# None before fade\_in}
    \StringTok{"figure\_type"}\NormalTok{: }\StringTok{"..."}\NormalTok{,}
    \StringTok{"build"}\NormalTok{:       }\StringTok{"..."}\NormalTok{,}
    \StringTok{"gender"}\NormalTok{:      }\StringTok{"..."}\NormalTok{,}
    \StringTok{"pose"}\NormalTok{:        }\StringTok{"..."}\NormalTok{,}
    \StringTok{"offset"}\NormalTok{:      [x, y, z],}
    \StringTok{"scale"}\NormalTok{:       \{}\StringTok{"sy"}\NormalTok{: ..., }\StringTok{"sx"}\NormalTok{: ..., }\StringTok{"anchor"}\NormalTok{: }\StringTok{"..."}\NormalTok{\},}
    \StringTok{"style"}\NormalTok{:       \{...\},}
    \StringTok{"color"}\NormalTok{:       }\StringTok{"..."}\NormalTok{,}
    \StringTok{"torso\_color"}\NormalTok{: }\StringTok{"..."}\NormalTok{,}
    \StringTok{"tic\_profile"}\NormalTok{: [...],}
    \StringTok{"uniforms"}\NormalTok{:    \{...\},}
    \StringTok{"facing"}\NormalTok{:      }\StringTok{"..."}\NormalTok{,     }\CommentTok{\# dog only}
    \StringTok{"opacity"}\NormalTok{:     }\FloatTok{1.0}\NormalTok{,       }\CommentTok{\# set by focus / focus\_reset}
\NormalTok{\}}
\end{Highlighting}
\end{Shaded}

\texttt{fade\_in} writes the live figure into
\texttt{cast{[}name{]}{[}"fig"{]}}; \texttt{fade\_out} clears it to
\texttt{None}. Actions that need the live figure (almost all of them)
read \texttt{cast{[}name{]}{[}"fig"{]}} via \texttt{\_resolve\_speaker}
(§4.1).

\paragraph{The prop registry}\label{the-prop-registry}

\texttt{props} is a \texttt{PropRegistry} instance. Most relevant
methods for character-side authoring:

\begin{itemize}
\tightlist
\item
  \texttt{props.get(name)} --- fetches the prop's \texttt{VGroup},
  prints a warning on miss.
\item
  \texttt{props.get\_raw(name)} --- same, but silent on miss (used for
  membership probes).
\item
  \texttt{props.add(name,\ vgroup)} --- registers a prop. Used by
  \texttt{spawn\_prop} and by \texttt{fade\_in}'s dual-registration
  branch.
\item
  \texttt{props.world\_pos(name)} --- returns the \texttt{(x,\ y)} world
  position of the prop, accounting for any \texttt{pam\_follows} parent
  chain.
\end{itemize}

\paragraph{Dual-registered dogs}\label{dual-registered-dogs}

A \texttt{DogGraph} faded in via the cast block appears in both
\texttt{cast{[}"chekov"{]}{[}"fig"{]}} and \texttt{props.get("chekov")}
(the latter being the \texttt{dog.group} \texttt{VGroup} with
\texttt{pam\_dog\ =\ fig}). Actions that resolve via
\texttt{\_resolve\_speaker} see the same \texttt{DogGraph} from either
side.

\paragraph{\texorpdfstring{The \texttt{remove\_prop} symmetric-clear
lesson}{The remove\_prop symmetric-clear lesson}}\label{the-remove_prop-symmetric-clear-lesson}

\texttt{remove\_prop} removes a prop from the prop registry \emph{and}
clears any cast entry of the same name (for Path C dogs). Before the
symmetric-clear fix, removing a dog cleared only the prop side, leaving
\texttt{cast{[}"chekov"{]}{[}"fig"{]}} pointing at a detached figure
that no longer rendered. Subsequent \texttt{say} / \texttt{trot\_to}
would resolve to the orphan and crash or silently no-op.

\textbf{Lesson for any future dual-registered figure type:} registry
maintenance must be symmetric. Add to both registries on creation;
remove from both on cleanup. A helper like
\texttt{\_register\_dual(cast,\ props,\ name,\ fig,\ group)} and
\texttt{\_unregister\_dual(cast,\ props,\ name)} is on the back-burner
list.

\begin{center}\rule{0.5\linewidth}{0.5pt}\end{center}

\subsubsection{4.6 Common gotchas}\label{common-gotchas}

Consolidated footguns from the prior sections. Each entry
cross-references the section where the full detail lives.

\paragraph{Identity}\label{identity}

\begin{itemize}
\tightlist
\item
  \textbf{\texttt{figure\_type} defaults to \texttt{"human"}.} Always
  set explicitly for alien and dog characters; the default-to-human
  fallthrough silently rebuilds an alien character as a human skeleton.
  See §1.1.
\item
  \textbf{\texttt{build} must match \texttt{figure\_type} for aliens.}
  \texttt{figure\_type:\ "alien"} + \texttt{build:\ "default"} produces
  a human-proportioned alien that breaks on morphs to alien-only poses.
  Use \texttt{build:\ "alien"} or \texttt{build:\ "alien\_female"}. See
  §1.2.
\item
  \textbf{\texttt{gender} defaults to \texttt{"male"} with a warning.}
  Set explicitly to avoid the warning, especially for female and child
  characters. See §1.3.
\end{itemize}

\paragraph{Color}\label{color-1}

\begin{itemize}
\tightlist
\item
  \textbf{The canonical registry is authoritative for color.} Per-scene
  \texttt{color} and \texttt{style.edge\_color} overrides are stomped on
  the next \texttt{fountain2pam.py} conversion. Edit the registry file
  directly. See §4.4.
\end{itemize}

\paragraph{Pose and scale}\label{pose-and-scale}

\begin{itemize}
\tightlist
\item
  \textbf{\texttt{lying\_flat\_right} / \texttt{lying\_flat\_left} are
  required for correct bubble anchoring on horizontal characters.} An
  ad-hoc \texttt{rotate} to horizontal mis-anchors \texttt{say} bubbles.
  See §1.7 and §3.1.1.
\item
  \textbf{The \texttt{scale} action replaces, doesn't multiply.}
  \texttt{\{"action":\ "scale",\ "who":\ "chava",\ "sy":\ 0.5\}} on a
  character currently at 0.7 results in 0.5, not 0.35. See §2.4.
\item
  \textbf{\texttt{rotate} on a character is visual-only.} The figure's
  internal \texttt{pose} and \texttt{offset} don't update. A subsequent
  \texttt{morph} snaps the character back to upright unless you
  counter-rotate first. See §2.5.
\end{itemize}

\paragraph{Locomotion}\label{locomotion}

\begin{itemize}
\tightlist
\item
  \textbf{\texttt{walk\_to} doesn't update the cast spec.} After
  fade-out + fade-in, the character returns to the cast-spec offset, not
  the walked-to x. See §1.8 and §3.2.1.
\item
  \textbf{\texttt{walk\_to\_prop} lands \emph{at} the prop x, not next
  to it.} Follow with a small \texttt{walk\_to} for offset positioning.
  See §3.2.6.
\end{itemize}

\paragraph{Speech bubbles}\label{speech-bubbles}

\begin{itemize}
\tightlist
\item
  \textbf{\texttt{say} is not parallel-safe at the screenplay level.}
  The player refuses \texttt{say} inside a \texttt{parallel} block. See
  §3.1.1.
\item
  \textbf{Persistent speech bubbles anchor world-space, not head-space.}
  A character who walks while a \texttt{persist} bubble is up leaves it
  behind. Use interleaved short non-persist \texttt{say} steps for
  walk-and-talk. See §3.1.1.
\item
  \textbf{\texttt{focus\_reset} can paint over persistent bubbles}
  (v0.9.10 gotcha). Issue \texttt{clear\_bubble} or
  \texttt{clear\_all\_bubbles} \emph{before} \texttt{focus\_reset}. See
  §2.13.
\end{itemize}

\paragraph{Face attachment (v0.9.15)}\label{face-attachment-v0.9.15}

\begin{itemize}
\tightlist
\item
  \textbf{Faces do not persist across \texttt{fade\_out} +
  \texttt{fade\_in}.} A character that fades out loses their face
  cleanly; re-attach explicitly after fade-in. See §2.7.
\item
  \textbf{The face stays upright during \texttt{rotate}.} A character
  laid flat keeps an upright face by design. See §2.5 and §2.7.
\item
  \textbf{Manim CE 0.20.1 (Cairo backend) doesn't composite RGBA alpha
  correctly} --- placeholder PNGs need opaque backgrounds matching the
  scene \texttt{BG\_COLOR}. See §2.7.
\end{itemize}

\paragraph{Tics}\label{tics}

\begin{itemize}
\tightlist
\item
  \textbf{Tics declared on incompatible body types are silent no-ops} at
  play time. \texttt{tail\_wag} on a biped fires harmlessly with no
  visual. See §1.12 and §3.5.8.
\item
  \textbf{\texttt{sentence\_end} triggers fire on every \texttt{say} /
  \texttt{prop\_say}.} Set \texttt{"suppress\_tics":\ true} on a step to
  opt out. See §1.12 and §3.1.1.
\end{itemize}

\paragraph{Resolution}\label{resolution}

\begin{itemize}
\tightlist
\item
  \textbf{\texttt{who} and \texttt{prop} are separate registries except
  for Path C dogs.} A human character is never reachable via
  \texttt{prop:\ name}; a \texttt{tv\_monitor} prop is never reachable
  via \texttt{who:\ name}. Dogs are reachable from both sides. See §4.1.
\item
  \textbf{Action handlers print a warning and no-op on resolution
  failure.} A typo in \texttt{who} or \texttt{prop} produces a console
  warning but does not crash the render. Watch the log during testing.
\end{itemize}

\paragraph{Persistent bubble registry}\label{persistent-bubble-registry}

\begin{itemize}
\tightlist
\item
  \textbf{Subscene transitions do not auto-clear persistent bubbles.}
  Explicit \texttt{clear\_all\_bubbles} at beat boundaries is the
  documented cleanup pattern. See §2.12.
\item
  \textbf{\texttt{clear\_bubble} is a silent no-op on stale clears.}
  Calling \texttt{clear\_bubble} for a bubble that's already faded is
  safe. See §2.11.
\end{itemize}

\paragraph{Registry maintenance}\label{registry-maintenance}

\begin{itemize}
\tightlist
\item
  \textbf{\texttt{DogGraph.group} returns a fresh VGroup on every
  access.} Bind to a local before stamping \texttt{pam\_*} attributes;
  otherwise the attributes attach to a throwaway object. See §4.1.
\item
  \textbf{Registry cleanup must be symmetric for dual-registered
  entries.} \texttt{remove\_prop} clears both the prop registry and the
  cast entry for Path C dogs. See §4.5.
\end{itemize}

\begin{center}\rule{0.5\linewidth}{0.5pt}\end{center}

\subsubsection{Status and next steps}\label{status-and-next-steps}

This document is intended to evolve into the v1.0.0 frozen reference. At
v0.9.15, the following items are still provisional and may shift before
the v1.0.0 freeze:

\begin{itemize}
\tightlist
\item
  The \texttt{to\_x} / \texttt{x} locomotion-parameter normalization
  (the v1.0.0 back-burner item). Currently \texttt{walk\_to} uses
  \texttt{x}; some older \texttt{walk} handlers accept \texttt{to\_x}.
  When normalized, the documentation in §3.2 will use whichever key
  wins.
\item
  The \texttt{group\_translate} parallel-safety fix (v1.0.0
  back-burner). Currently in \texttt{\_CANNOT\_PARALLEL}; if fixed, the
  §3.8 flag will flip to parallel-safe.
\item
  The \texttt{parallel} + \texttt{say} constraint (currently tribal
  knowledge / refused by the player). Either fixed (head-anchored
  bubbles allowing walk-and-talk) or formally documented with the
  refusal error.
\item
  The bubble / \texttt{focus\_reset} ordering gotcha (§2.13). Either
  fixed in the z-order math or kept as a documented pattern requiring
  explicit \texttt{clear\_bubble}.
\item
  Quadruped face attachment (§2.7). Currently \texttt{attach\_face}
  no-ops on figures without a \texttt{"head"} dot. If \texttt{DogGraph}
  grows a head-positioned face anchor, the action becomes meaningful for
  dogs.
\end{itemize}

The companion document \texttt{PROP\_PROPERTIES.md} (also on the v1.0.0
back-burner list) will mirror this document's structure for prop-side
authoring. Topics deferred to that document: prop spawning / despawning
lifecycle, the \texttt{hidden} flag for props, \texttt{DogGraph} as
prop, \texttt{GovernorGraph}'s autonomous rotation, the \texttt{phone}
builder's three styles, prop-following lifecycle
(\texttt{\_drag\_attached\_props}), the \texttt{tv\_monitor} /
\texttt{news\_anchor} prop-character pattern, and prop-side z-order
considerations.

\begin{center}\rule{0.5\linewidth}{0.5pt}\end{center}

\emph{End of Section 4. Section 5 follows as the Face \& Appearance
reference.}

\subsection{5. Face \& Appearance}\label{face-appearance}

This section documents PAM's flat-cartoon face system: how character
face data is declared in the screenplay JSON, how \texttt{attach\_face}
builds and attaches a face VGroup over a character's head dot, how
expression variants are registered and swapped at runtime, and the
geometry / layering reference needed to debug face render anomalies.

The face system is implemented in \texttt{pam/face\_builder.py}. Face
data is \textbf{separate from skeleton/physics data} --- the
\texttt{"cast"} block (§1) defines the figure (skeleton, build, pose,
scale); the \texttt{"faces"} block (§5.1) defines the cartoon face that
overlays the head dot. The two blocks are independent and may appear in
either order in the screenplay JSON, provided both appear before any
\texttt{attach\_face} step that references them.

Faces are built entirely from Manim VMobjects --- \texttt{Circle},
\texttt{Ellipse}, \texttt{Rectangle}, \texttt{Polygon}, \texttt{Line},
\texttt{ArcBetweenPoints}. No bitmap or SVG asset loading is involved
for character-key faces (the legacy \texttt{pam/assets/} path is still
supported for file-extension \texttt{image} values; see §5.4.1).

\begin{itemize}
\tightlist
\item
  \textbf{§5.1} --- The \texttt{"faces"} block: how project character
  face data is declared in tntd\_characters.json and loaded by
  pam\_player at runtime.
\item
  \textbf{§5.2} --- The complete FACE\_DATA schema: skin, hair, eyes,
  brows, mouth, clothing, extras, hats, bun geometry.
\item
  \textbf{§5.3} --- Expression variants: the \texttt{"expressions"} key
  on cast entries, the \texttt{\{char\}\_\{mouth\}} naming convention,
  compound variants, and \texttt{register\_variants()}.
\item
  \textbf{§5.4} --- The \texttt{attach\_face} action: JSON syntax, scale
  convention, expression swapping, the \texttt{pam\_head\_ref} anchor
  fix.
\item
  \textbf{§5.5} --- Hat geometry reference: triangle / triangle\_inv
  shapes, module-level constants, bun coverage geometry.
\item
  \textbf{§5.6} --- Layering order reference: front-view and side-view
  z-order tables.
\item
  \textbf{§5.7} --- Debugging checklist for missing or mis-anchored
  faces.
\end{itemize}

\begin{center}\rule{0.5\linewidth}{0.5pt}\end{center}

\subsubsection{5.1 The ``faces'' block}\label{the-faces-block}

Face data for project characters lives in tntd\_characters.json (or any
screenplay JSON that pam\_player loads), in a top-level action block
with \texttt{"action":\ "faces"}. The block is a sibling of
\texttt{"action":\ "cast"} and structurally mirrors it: a
\texttt{"characters"} key mapping character keys to per-character face
specs.

\texttt{face\_builder.py} itself ships with only three illustrative
entries in \texttt{FACE\_DATA}: \texttt{example\_human},
\texttt{example\_alien}, \texttt{example\_dog}. Project-specific
characters belong in the screenplay JSON, \textbf{not} in
face\_builder.py source code.

\paragraph{JSON shape}\label{json-shape}

\begin{Shaded}
\begin{Highlighting}[]
\FunctionTok{\{}
  \DataTypeTok{"action"}\FunctionTok{:} \StringTok{"faces"}\FunctionTok{,}
  \DataTypeTok{"characters"}\FunctionTok{:} \FunctionTok{\{}
    \DataTypeTok{"bevers"}\FunctionTok{:} \FunctionTok{\{}
      \DataTypeTok{"name"}\FunctionTok{:}        \StringTok{"Bevers"}\FunctionTok{,}
      \DataTypeTok{"note"}\FunctionTok{:}        \StringTok{"alien · avatar cadet"}\FunctionTok{,}
      \DataTypeTok{"skin"}\FunctionTok{:}        \StringTok{"\#8cc47c"}\FunctionTok{,}
      \DataTypeTok{"hair"}\FunctionTok{:}        \StringTok{"\#3a56a8"}\FunctionTok{,}
      \DataTypeTok{"hair\_style"}\FunctionTok{:}  \StringTok{"flat\_top"}\FunctionTok{,}
      \DataTypeTok{"eye\_type"}\FunctionTok{:}    \StringTok{"alien"}\FunctionTok{,}
      \DataTypeTok{"eye\_color"}\FunctionTok{:}   \StringTok{"\#c87820"}\FunctionTok{,}
      \DataTypeTok{"brow\_thick"}\FunctionTok{:}  \KeywordTok{false}\FunctionTok{,}
      \DataTypeTok{"mouth"}\FunctionTok{:}       \StringTok{"neutral"}\FunctionTok{,}
      \DataTypeTok{"cloth\_color"}\FunctionTok{:} \StringTok{"\#202840"}\FunctionTok{,}
      \DataTypeTok{"cloth\_style"}\FunctionTok{:} \StringTok{"uniform"}\FunctionTok{,}
      \DataTypeTok{"ep\_color"}\FunctionTok{:}    \StringTok{"\#c8803a"}\FunctionTok{,}
      \DataTypeTok{"extras"}\FunctionTok{:}      \OtherTok{[]}
    \FunctionTok{\},}
    \DataTypeTok{"freydoon"}\FunctionTok{:} \FunctionTok{\{}
      \DataTypeTok{"name"}\FunctionTok{:}        \StringTok{"Freydoon"}\FunctionTok{,}
      \DataTypeTok{"note"}\FunctionTok{:}        \StringTok{"human · hospital patient"}\FunctionTok{,}
      \DataTypeTok{"skin"}\FunctionTok{:}        \StringTok{"\#c07848"}\FunctionTok{,}
      \DataTypeTok{"hair"}\FunctionTok{:}        \StringTok{"\#181818"}\FunctionTok{,}
      \DataTypeTok{"hair\_style"}\FunctionTok{:}  \StringTok{"short\_male"}\FunctionTok{,}
      \DataTypeTok{"eye\_type"}\FunctionTok{:}    \StringTok{"human"}\FunctionTok{,}
      \DataTypeTok{"eye\_color"}\FunctionTok{:}   \StringTok{"\#6a3808"}\FunctionTok{,}
      \DataTypeTok{"brow\_thick"}\FunctionTok{:}  \KeywordTok{true}\FunctionTok{,}
      \DataTypeTok{"mouth"}\FunctionTok{:}       \StringTok{"neutral"}\FunctionTok{,}
      \DataTypeTok{"cloth\_color"}\FunctionTok{:} \StringTok{"\#8090b0"}\FunctionTok{,}
      \DataTypeTok{"cloth\_style"}\FunctionTok{:} \StringTok{"standard"}\FunctionTok{,}
      \DataTypeTok{"extras"}\FunctionTok{:}      \OtherTok{[]}
    \FunctionTok{\}}
  \FunctionTok{\}}
\FunctionTok{\}}
\end{Highlighting}
\end{Shaded}

Each entry in the \texttt{"characters"} dict maps a character key to a
face spec dict using the schema in §5.2.

\paragraph{\texorpdfstring{Loading:
\texttt{load\_faces()}}{Loading: load\_faces()}}\label{loading-load_faces}

When pam\_player encounters an \texttt{"action":\ "faces"} step, it
calls \texttt{face\_builder.load\_faces(step{[}"characters"{]})}, which
iterates the dict and writes each entry into the module-level
\texttt{FACE\_DATA} dict:

\begin{Shaded}
\begin{Highlighting}[]
\KeywordTok{def}\NormalTok{ load\_faces(characters: }\BuiltInTok{dict}\NormalTok{) }\OperatorTok{{-}\textgreater{}} \VariableTok{None}\NormalTok{:}
    \ControlFlowTok{for}\NormalTok{ char\_key, spec }\KeywordTok{in}\NormalTok{ characters.items():}
\NormalTok{        FACE\_DATA[char\_key] }\OperatorTok{=}\NormalTok{ spec}
\end{Highlighting}
\end{Shaded}

Key behaviors:

\begin{itemize}
\tightlist
\item
  \textbf{Overwrites} any existing \texttt{FACE\_DATA} entry with the
  same key. Re-running \texttt{load\_faces} with updated data refreshes
  the entry --- useful for iterative tuning during development.
\item
  \textbf{Silent no-op} if the \texttt{characters} dict is empty.
\item
  \textbf{Expression variants are NOT loaded here.} Variants live under
  the \texttt{"expressions"} key in each cast entry and are registered
  lazily by \texttt{register\_variants()} on the first
  \texttt{attach\_face} for that character (see §5.3.4).
\end{itemize}

\paragraph{Placement in
tntd\_characters.json}\label{placement-in-tntd_characters.json}

The \texttt{"faces"} block must appear \textbf{before any scene steps
that use \texttt{attach\_face}}. It can appear before or after the
\texttt{"cast"} block --- the two are independent. Recommended order
keeps the character bootstrap at the top of the file:

\begin{Shaded}
\begin{Highlighting}[]
\OtherTok{[}
  \FunctionTok{\{}\DataTypeTok{"action"}\FunctionTok{:} \StringTok{"title"}\FunctionTok{,} \DataTypeTok{"text"}\FunctionTok{:} \StringTok{"..."}\FunctionTok{\}}\OtherTok{,}
  \FunctionTok{\{}\DataTypeTok{"action"}\FunctionTok{:} \StringTok{"cast"}\FunctionTok{,}  \DataTypeTok{"characters"}\FunctionTok{:} \FunctionTok{\{} \ErrorTok{...} \FunctionTok{\}\}}\OtherTok{,}
  \FunctionTok{\{}\DataTypeTok{"action"}\FunctionTok{:} \StringTok{"faces"}\FunctionTok{,} \DataTypeTok{"characters"}\FunctionTok{:} \FunctionTok{\{} \ErrorTok{...} \FunctionTok{\}\}}\OtherTok{,}
  \ErrorTok{...} \ErrorTok{scene} \ErrorTok{steps} \ErrorTok{...}
\OtherTok{]}
\end{Highlighting}
\end{Shaded}

\begin{center}\rule{0.5\linewidth}{0.5pt}\end{center}

\subsubsection{5.2 The FACE\_DATA schema}\label{the-face_data-schema}

Each face spec is a dict. Only \texttt{skin} and \texttt{hair} are
\textbf{strictly required} --- both are accessed directly via
\texttt{char{[}"skin"{]}} and \texttt{char{[}"hair"{]}} in
\texttt{build\_face()} and raise \texttt{KeyError} if absent. Every
other key uses \texttt{char.get(key,\ default)} and falls back to the
documented default when omitted.

\paragraph{5.2.1 Identity}\label{identity2}

\begin{longtable}[]{@{}
  >{\raggedright\arraybackslash}p{(\columnwidth - 8\tabcolsep) * \real{0.2000}}
  >{\raggedright\arraybackslash}p{(\columnwidth - 8\tabcolsep) * \real{0.2000}}
  >{\raggedright\arraybackslash}p{(\columnwidth - 8\tabcolsep) * \real{0.2000}}
  >{\raggedright\arraybackslash}p{(\columnwidth - 8\tabcolsep) * \real{0.2000}}
  >{\raggedright\arraybackslash}p{(\columnwidth - 8\tabcolsep) * \real{0.2000}}@{}}
\toprule\noalign{}
\begin{minipage}[b]{\linewidth}\raggedright
Key
\end{minipage} & \begin{minipage}[b]{\linewidth}\raggedright
Type
\end{minipage} & \begin{minipage}[b]{\linewidth}\raggedright
Required
\end{minipage} & \begin{minipage}[b]{\linewidth}\raggedright
Default
\end{minipage} & \begin{minipage}[b]{\linewidth}\raggedright
Description
\end{minipage} \\
\midrule\noalign{}
\endhead
\bottomrule\noalign{}
\endlastfoot
\texttt{name} & \texttt{str} & no & \texttt{"?"} & Display name shown in
\texttt{list\_characters()} output. \\
\texttt{note} & \texttt{str} & no & \texttt{""} & Free text
role/description, also shown in \texttt{list\_characters()}. \\
\end{longtable}

These are documentation fields only --- they do not affect rendering.

\paragraph{5.2.2 Skin}\label{skin}

\begin{longtable}[]{@{}
  >{\raggedright\arraybackslash}p{(\columnwidth - 8\tabcolsep) * \real{0.2000}}
  >{\raggedright\arraybackslash}p{(\columnwidth - 8\tabcolsep) * \real{0.2000}}
  >{\raggedright\arraybackslash}p{(\columnwidth - 8\tabcolsep) * \real{0.2000}}
  >{\raggedright\arraybackslash}p{(\columnwidth - 8\tabcolsep) * \real{0.2000}}
  >{\raggedright\arraybackslash}p{(\columnwidth - 8\tabcolsep) * \real{0.2000}}@{}}
\toprule\noalign{}
\begin{minipage}[b]{\linewidth}\raggedright
Key
\end{minipage} & \begin{minipage}[b]{\linewidth}\raggedright
Type
\end{minipage} & \begin{minipage}[b]{\linewidth}\raggedright
Required
\end{minipage} & \begin{minipage}[b]{\linewidth}\raggedright
Default
\end{minipage} & \begin{minipage}[b]{\linewidth}\raggedright
Description
\end{minipage} \\
\midrule\noalign{}
\endhead
\bottomrule\noalign{}
\endlastfoot
\texttt{skin} & hex & \textbf{yes} & --- & Face, neck, and ear fill
color. Also used to derive nose and mouth stroke colors (auto-darkened
by \texttt{\_darken()}). \\
\end{longtable}

\paragraph{5.2.3 Hair}\label{hair}

\begin{longtable}[]{@{}
  >{\raggedright\arraybackslash}p{(\columnwidth - 8\tabcolsep) * \real{0.2000}}
  >{\raggedright\arraybackslash}p{(\columnwidth - 8\tabcolsep) * \real{0.2000}}
  >{\raggedright\arraybackslash}p{(\columnwidth - 8\tabcolsep) * \real{0.2000}}
  >{\raggedright\arraybackslash}p{(\columnwidth - 8\tabcolsep) * \real{0.2000}}
  >{\raggedright\arraybackslash}p{(\columnwidth - 8\tabcolsep) * \real{0.2000}}@{}}
\toprule\noalign{}
\begin{minipage}[b]{\linewidth}\raggedright
Key
\end{minipage} & \begin{minipage}[b]{\linewidth}\raggedright
Type
\end{minipage} & \begin{minipage}[b]{\linewidth}\raggedright
Required
\end{minipage} & \begin{minipage}[b]{\linewidth}\raggedright
Default
\end{minipage} & \begin{minipage}[b]{\linewidth}\raggedright
Description
\end{minipage} \\
\midrule\noalign{}
\endhead
\bottomrule\noalign{}
\endlastfoot
\texttt{hair} & hex & \textbf{yes} & --- & Hair fill color. Also used to
derive eyebrow color (auto-darkened by \texttt{0.82}) unless
\texttt{brow\_color} is set explicitly. \\
\texttt{hair\_style} & \texttt{str} & no & \texttt{"short\_male"} & One
of the named styles in the table below. \\
\texttt{hair\_width} & \texttt{float} & no & (per style) & \emph{(v0.9.18)}
Front-view cap width override in Manim units. Replaces the per-style
default. Use to make the hairdo thinner or wider. Not applied in side
view (cap width is locked to head depth). See
\hyperref[hat-geometry-reference]{§5.5.5} for per-style defaults and
examples. \\
\texttt{hair\_height} & \texttt{float} & no & (per style) & \emph{(v0.9.18)}
Cap height override in Manim units, applied in both front and side
views. Use to make the hairdo shorter or taller. Note: cap extends
above and below its center, so taller caps lower the hairline unless
\texttt{hair\_offset\_y} is also raised. \\
\texttt{hair\_offset\_y} & \texttt{float} & no & (per style) & \emph{(v0.9.18)}
Vertical offset from \texttt{HEAD\_RY} to the cap CENTER, both views.
Replaces the per-style default offset. Use to raise or lower the
hairline independently of height. \\
\end{longtable}

\textbf{Hair style table:}

\begin{longtable}[]{@{}
  >{\raggedright\arraybackslash}p{(\columnwidth - 2\tabcolsep) * \real{0.5000}}
  >{\raggedright\arraybackslash}p{(\columnwidth - 2\tabcolsep) * \real{0.5000}}@{}}
\toprule\noalign{}
\begin{minipage}[b]{\linewidth}\raggedright
Value
\end{minipage} & \begin{minipage}[b]{\linewidth}\raggedright
Description
\end{minipage} \\
\midrule\noalign{}
\endhead
\bottomrule\noalign{}
\endlastfoot
\texttt{"flat\_top"} & Flat-topped rectangle. Bevers's signature
look. \\
\texttt{"stubble"} & Very thin ellipse cap --- barely-there buzz. \\
\texttt{"short\_male"} & Low half-ellipse cap with ear visibility. \\
\texttt{"slicked"} & Flat cap with a highlight arc showing the slick. \\
\texttt{"medium\_female"} & Cap + side panels hanging to shoulder
level. \\
\texttt{"grey\_wavy"} & \texttt{medium\_female} variant with a wavy
highlight arc. \\
\texttt{"bob"} & Cap + shorter side panels (chin-length bob). \\
\texttt{"updo"} & Tall dome + top-knot. Thalia's braided updo. \\
\texttt{"bun"} & Cap + round bun circle above. Nona's grey bun. \\
\end{longtable}

Styles in the \texttt{\_HAIR\_SHOWS\_EARS} set --- \texttt{short\_male},
\texttt{stubble}, \texttt{slicked}, \texttt{flat\_top}, \texttt{bun},
\texttt{updo} --- leave the ears visible. The longer styles
\texttt{medium\_female}, \texttt{grey\_wavy}, \texttt{bob} draw side
panels that cover the ears. Earring rendering is conditional on this set
(see §5.2.8).

\paragraph{5.2.4 Eyes}\label{eyes}

\begin{longtable}[]{@{}
  >{\raggedright\arraybackslash}p{(\columnwidth - 8\tabcolsep) * \real{0.2000}}
  >{\raggedright\arraybackslash}p{(\columnwidth - 8\tabcolsep) * \real{0.2000}}
  >{\raggedright\arraybackslash}p{(\columnwidth - 8\tabcolsep) * \real{0.2000}}
  >{\raggedright\arraybackslash}p{(\columnwidth - 8\tabcolsep) * \real{0.2000}}
  >{\raggedright\arraybackslash}p{(\columnwidth - 8\tabcolsep) * \real{0.2000}}@{}}
\toprule\noalign{}
\begin{minipage}[b]{\linewidth}\raggedright
Key
\end{minipage} & \begin{minipage}[b]{\linewidth}\raggedright
Type
\end{minipage} & \begin{minipage}[b]{\linewidth}\raggedright
Required
\end{minipage} & \begin{minipage}[b]{\linewidth}\raggedright
Default
\end{minipage} & \begin{minipage}[b]{\linewidth}\raggedright
Description
\end{minipage} \\
\midrule\noalign{}
\endhead
\bottomrule\noalign{}
\endlastfoot
\texttt{eye\_type} & \texttt{str} & no & \texttt{"alien"} &
\texttt{"alien"} or \texttt{"human"}. \\
\texttt{eye\_color} & hex & no & \texttt{"\#c87820"} & Alien: full iris
disc color. Human: iris ring color (pupil is always dark). \\
\end{longtable}

\begin{itemize}
\tightlist
\item
  \textbf{\texttt{"alien"}} --- single amber/colored disc with a dark
  pupil and small white catchlight. No white sclera.
\item
  \textbf{\texttt{"human"}} --- white sclera circle + small colored iris
  ring + dark pupil + catchlight.
\end{itemize}

\paragraph{5.2.5 Eyebrows}\label{eyebrows}

\begin{longtable}[]{@{}
  >{\raggedright\arraybackslash}p{(\columnwidth - 8\tabcolsep) * \real{0.2000}}
  >{\raggedright\arraybackslash}p{(\columnwidth - 8\tabcolsep) * \real{0.2000}}
  >{\raggedright\arraybackslash}p{(\columnwidth - 8\tabcolsep) * \real{0.2000}}
  >{\raggedright\arraybackslash}p{(\columnwidth - 8\tabcolsep) * \real{0.2000}}
  >{\raggedright\arraybackslash}p{(\columnwidth - 8\tabcolsep) * \real{0.2000}}@{}}
\toprule\noalign{}
\begin{minipage}[b]{\linewidth}\raggedright
Key
\end{minipage} & \begin{minipage}[b]{\linewidth}\raggedright
Type
\end{minipage} & \begin{minipage}[b]{\linewidth}\raggedright
Required
\end{minipage} & \begin{minipage}[b]{\linewidth}\raggedright
Default
\end{minipage} & \begin{minipage}[b]{\linewidth}\raggedright
Description
\end{minipage} \\
\midrule\noalign{}
\endhead
\bottomrule\noalign{}
\endlastfoot
\texttt{brow\_color} & hex & no & hair darkened by \texttt{0.82} &
Stroke color. Set explicitly for alien characters whose hair-derived
brow would be too subtle (e.g.~Thalia's \texttt{"\#3a1858"} against hair
\texttt{"\#5a3080"}). \\
\texttt{brow\_thick} & \texttt{bool} & no & \texttt{false} &
\texttt{false} = standard arch stroke weight. \texttt{true} = heavier
brow, reads as more masculine or emphatic. Used for Freydoon. \\
\end{longtable}

\paragraph{5.2.6 Mouth}\label{mouth}

\begin{longtable}[]{@{}lllll@{}}
\toprule\noalign{}
Key & Type & Required & Default & Description \\
\midrule\noalign{}
\endhead
\bottomrule\noalign{}
\endlastfoot
\texttt{mouth} & \texttt{str} & no & \texttt{"neutral"} & See values
below. \\
\end{longtable}

\begin{longtable}[]{@{}
  >{\raggedright\arraybackslash}p{(\columnwidth - 2\tabcolsep) * \real{0.5000}}
  >{\raggedright\arraybackslash}p{(\columnwidth - 2\tabcolsep) * \real{0.5000}}@{}}
\toprule\noalign{}
\begin{minipage}[b]{\linewidth}\raggedright
Value
\end{minipage} & \begin{minipage}[b]{\linewidth}\raggedright
Description
\end{minipage} \\
\midrule\noalign{}
\endhead
\bottomrule\noalign{}
\endlastfoot
\texttt{"neutral"} & Very slight upward arc. Resting face. \\
\texttt{"smile"} & Pronounced upward arc. Friendly. \\
\texttt{"flat"} & Straight horizontal line. Deadpan. \\
\texttt{"wry"} & Asymmetric: left corner lower. Ironic or skeptical. \\
\texttt{"wry\_down"} & Same left-corner drop as \texttt{"wry"}, but the
arc bows downward. Reads as a grimace, resigned displeasure, or
``really?''. \\
\end{longtable}

\textbf{Side-view note:} \texttt{"wry"} maps to \texttt{"neutral"} in
profile view --- the asymmetry does not read from the side, so the
side-view builder silently substitutes a neutral arc. \texttt{"wry\_down"}
keeps a downward-bowed grimace in profile, so use it when the beat needs
visible side-view displeasure rather than a neutralized skeptical mouth.

\paragraph{5.2.7 Clothing}\label{clothing}

\begin{longtable}[]{@{}
  >{\raggedright\arraybackslash}p{(\columnwidth - 8\tabcolsep) * \real{0.2000}}
  >{\raggedright\arraybackslash}p{(\columnwidth - 8\tabcolsep) * \real{0.2000}}
  >{\raggedright\arraybackslash}p{(\columnwidth - 8\tabcolsep) * \real{0.2000}}
  >{\raggedright\arraybackslash}p{(\columnwidth - 8\tabcolsep) * \real{0.2000}}
  >{\raggedright\arraybackslash}p{(\columnwidth - 8\tabcolsep) * \real{0.2000}}@{}}
\toprule\noalign{}
\begin{minipage}[b]{\linewidth}\raggedright
Key
\end{minipage} & \begin{minipage}[b]{\linewidth}\raggedright
Type
\end{minipage} & \begin{minipage}[b]{\linewidth}\raggedright
Required
\end{minipage} & \begin{minipage}[b]{\linewidth}\raggedright
Default
\end{minipage} & \begin{minipage}[b]{\linewidth}\raggedright
Description
\end{minipage} \\
\midrule\noalign{}
\endhead
\bottomrule\noalign{}
\endlastfoot
\texttt{cloth\_color} & hex & no & \texttt{"\#202840"} & Main
torso/clothing fill. \\
\texttt{cloth\_style} & \texttt{str} & no & \texttt{"standard"} & See
values below. \\
\texttt{ep\_color} & hex & no & \texttt{"\#c8a020"} & Epaulette fill.
Used when \texttt{extras} includes \texttt{"epaulettes"} OR when
\texttt{cloth\_style} is \texttt{"military"}. \\
\texttt{panel\_length} & float & no & \texttt{0.38} & Height of a
\texttt{cloth\_style:\ "panel"} bib. Larger values hang the bib lower
while preserving the top anchor. \\
\texttt{panel\_width} & float & no & \texttt{0.32} & Width of a panel
bib. \\
\texttt{panel\_color} & hex & no & \texttt{cloth\_color} & Optional
panel fill override. \\
\texttt{panel\_border\_color} & hex & no & \texttt{OUTLINE\_COLOR} &
Optional panel outline color; useful for ceremonial rims or rank
coding. \\
\texttt{panel\_border\_width} & float & no & \texttt{STROKE\_WIDTH} &
Optional panel outline width; use 4--6 for a visible decorative
border. \\
\end{longtable}

\begin{longtable}[]{@{}
  >{\raggedright\arraybackslash}p{(\columnwidth - 2\tabcolsep) * \real{0.5000}}
  >{\raggedright\arraybackslash}p{(\columnwidth - 2\tabcolsep) * \real{0.5000}}@{}}
\toprule\noalign{}
\begin{minipage}[b]{\linewidth}\raggedright
\texttt{cloth\_style} value
\end{minipage} & \begin{minipage}[b]{\linewidth}\raggedright
Description
\end{minipage} \\
\midrule\noalign{}
\endhead
\bottomrule\noalign{}
\endlastfoot
\texttt{"standard"} & Plain trapezoid torso, no detail. \\
\texttt{"blazer"} & Adds a V-neck collar crease (two short angled
strokes). \\
\texttt{"military"} & Adds rectangular epaulette slots (filled with
\texttt{ep\_color}). Implicitly enables epaulettes. \\
\texttt{"uniform"} & Adds a center-front line (PAM cadet uniform
suggestion). \\
\texttt{"panel"} (v0.9.X) & Rectangular chest bib replacing the
trapezoid. See exclusivity note and §5.5.4. \\
\end{longtable}

Epaulettes are triggered by either route --- explicit
\texttt{"epaulettes"} in extras, or \texttt{cloth\_style:\ "military"}.
The internal check is
\texttt{ep\_color\ if\ ("epaulettes"\ in\ extras\ or\ cloth\_style\ ==\ "military")\ else\ None}.

\textbf{Panel exclusivity:} \texttt{cloth\_style:\ "panel"} is exclusive
with all the other decorations. When \texttt{cloth\_style\ ==\ "panel"},
the blazer V-neck crease, military/uniform center seam, and shoulder
epaulettes are all skipped --- even if \texttt{"epaulettes"} appears in
\texttt{extras} or another decoration would otherwise apply. For rank
insignia or a name tag on a panel, attach a badge prop using
\texttt{get\_panel\_badge\_anchor()} (see §5.5.4) rather than relying on
\texttt{extras:\ {[}"epaulettes"{]}}.

\paragraph{5.2.8 Extras}\label{extras}

\begin{longtable}[]{@{}
  >{\raggedright\arraybackslash}p{(\columnwidth - 8\tabcolsep) * \real{0.2000}}
  >{\raggedright\arraybackslash}p{(\columnwidth - 8\tabcolsep) * \real{0.2000}}
  >{\raggedright\arraybackslash}p{(\columnwidth - 8\tabcolsep) * \real{0.2000}}
  >{\raggedright\arraybackslash}p{(\columnwidth - 8\tabcolsep) * \real{0.2000}}
  >{\raggedright\arraybackslash}p{(\columnwidth - 8\tabcolsep) * \real{0.2000}}@{}}
\toprule\noalign{}
\begin{minipage}[b]{\linewidth}\raggedright
Key
\end{minipage} & \begin{minipage}[b]{\linewidth}\raggedright
Type
\end{minipage} & \begin{minipage}[b]{\linewidth}\raggedright
Required
\end{minipage} & \begin{minipage}[b]{\linewidth}\raggedright
Default
\end{minipage} & \begin{minipage}[b]{\linewidth}\raggedright
Description
\end{minipage} \\
\midrule\noalign{}
\endhead
\bottomrule\noalign{}
\endlastfoot
\texttt{extras} & \texttt{list{[}str{]}} & no & \texttt{{[}{]}} & Zero
or more of the values below. \\
\texttt{blush\_color} & hex & no & \texttt{"\#d04040"} & Blush fill. Use
a muted red-pink for aliens; warmer peach for humans. \\
\end{longtable}

\begin{longtable}[]{@{}
  >{\raggedright\arraybackslash}p{(\columnwidth - 2\tabcolsep) * \real{0.5000}}
  >{\raggedright\arraybackslash}p{(\columnwidth - 2\tabcolsep) * \real{0.5000}}@{}}
\toprule\noalign{}
\begin{minipage}[b]{\linewidth}\raggedright
\texttt{extras} value
\end{minipage} & \begin{minipage}[b]{\linewidth}\raggedright
Description
\end{minipage} \\
\midrule\noalign{}
\endhead
\bottomrule\noalign{}
\endlastfoot
\texttt{"blush"} & Soft ellipses on cheeks. \\
\texttt{"epaulettes"} & Shoulder rank bars (same visual effect as
\texttt{cloth\_style:\ "military"}). \\
\texttt{"gold\_earrings"} & Gold oval earrings at earlobes. \\
\texttt{"pearl\_earrings"} & White/cream disc earrings. \\
\end{longtable}

Earrings render in front view at both earlobes only for hair styles in
\texttt{\_HAIR\_SHOWS\_EARS} (see §5.2.3). In side view, the earring
appears at the far-side ear (back of head) for the same set of
short-hair styles; longer styles (\texttt{medium\_female},
\texttt{grey\_wavy}, \texttt{bob}) suppress the earring entirely because
the side panels would cover it.

Pearl earrings take precedence over gold if both are listed (mutually
exclusive via \texttt{if/elif} in the builder).

\paragraph{5.2.9 Hats}\label{hats}

Hats are optional. Omit \texttt{hat\_style} entirely for a bare-headed
character.

\begin{longtable}[]{@{}
  >{\raggedright\arraybackslash}p{(\columnwidth - 8\tabcolsep) * \real{0.2000}}
  >{\raggedright\arraybackslash}p{(\columnwidth - 8\tabcolsep) * \real{0.2000}}
  >{\raggedright\arraybackslash}p{(\columnwidth - 8\tabcolsep) * \real{0.2000}}
  >{\raggedright\arraybackslash}p{(\columnwidth - 8\tabcolsep) * \real{0.2000}}
  >{\raggedright\arraybackslash}p{(\columnwidth - 8\tabcolsep) * \real{0.2000}}@{}}
\toprule\noalign{}
\begin{minipage}[b]{\linewidth}\raggedright
Key
\end{minipage} & \begin{minipage}[b]{\linewidth}\raggedright
Type
\end{minipage} & \begin{minipage}[b]{\linewidth}\raggedright
Required
\end{minipage} & \begin{minipage}[b]{\linewidth}\raggedright
Default
\end{minipage} & \begin{minipage}[b]{\linewidth}\raggedright
Description
\end{minipage} \\
\midrule\noalign{}
\endhead
\bottomrule\noalign{}
\endlastfoot
\texttt{hat\_style} & \texttt{str} & no & (no hat) & \texttt{"triangle"}
or \texttt{"triangle\_inv"}. See §5.5 for full geometry. \\
\texttt{hat\_color} & hex & no & \texttt{"\#303030"} & Hat fill color.
Outline uses \texttt{OUTLINE\_COLOR} at \texttt{STROKE\_WIDTH}. \\
\end{longtable}

The hat is drawn \textbf{last} --- on top of all hair and face elements.
Hair underneath is \textbf{not} suppressed; it remains visible at the
sides where it extends beyond the hat outline. Because the hat is
embedded in the face VGroup, it persists automatically through
expression swaps that target the same base character (an
\texttt{attach\_face} with \texttt{image:\ "factor\_flat"} keeps
Factor's triangle hat).

To show a character \textbf{without} their hat mid-scene, define an
expression variant in the \texttt{"expressions"} block that points to a
separate FACE\_DATA entry omitting \texttt{hat\_style}, or define a
parallel hat-less base entry under a different key. There is no
\texttt{"attach\_hat"} action; hats are not separate mobjects.

\paragraph{5.2.10 Bun geometry}\label{bun-geometry}

Only relevant when \texttt{hair\_style\ ==\ "bun"}.

\begin{longtable}[]{@{}
  >{\raggedright\arraybackslash}p{(\columnwidth - 8\tabcolsep) * \real{0.2000}}
  >{\raggedright\arraybackslash}p{(\columnwidth - 8\tabcolsep) * \real{0.2000}}
  >{\raggedright\arraybackslash}p{(\columnwidth - 8\tabcolsep) * \real{0.2000}}
  >{\raggedright\arraybackslash}p{(\columnwidth - 8\tabcolsep) * \real{0.2000}}
  >{\raggedright\arraybackslash}p{(\columnwidth - 8\tabcolsep) * \real{0.2000}}@{}}
\toprule\noalign{}
\begin{minipage}[b]{\linewidth}\raggedright
Key
\end{minipage} & \begin{minipage}[b]{\linewidth}\raggedright
Type
\end{minipage} & \begin{minipage}[b]{\linewidth}\raggedright
Required
\end{minipage} & \begin{minipage}[b]{\linewidth}\raggedright
Default
\end{minipage} & \begin{minipage}[b]{\linewidth}\raggedright
Description
\end{minipage} \\
\midrule\noalign{}
\endhead
\bottomrule\noalign{}
\endlastfoot
\texttt{bun\_offset\_y} & \texttt{float} & no & \texttt{0.155} &
Vertical offset from \texttt{HEAD\_RY} to bun center in front view.
Side-view offset is derived automatically as
\texttt{bun\_offset\_y\ *\ 0.71}. \\
\end{longtable}

Larger values push the bun higher above the head. The default
\texttt{0.155} places the bun close to the cap; if a
\texttt{triangle\_inv} hat is worn, the default leaves enough clearance
to cover the bun without clipping (see §5.5.3 for the clearance
arithmetic).

\begin{center}\rule{0.5\linewidth}{0.5pt}\end{center}

\paragraph{Complete schema example: Factor (alien · professor, with
triangle
hat)}\label{complete-schema-example-factor-alien-professor-with-triangle-hat}

\begin{Shaded}
\begin{Highlighting}[]
\FunctionTok{\{}
  \DataTypeTok{"factor"}\FunctionTok{:} \FunctionTok{\{}
    \DataTypeTok{"name"}\FunctionTok{:}        \StringTok{"Factor"}\FunctionTok{,}
    \DataTypeTok{"note"}\FunctionTok{:}        \StringTok{"alien · professor"}\FunctionTok{,}
    \DataTypeTok{"skin"}\FunctionTok{:}        \StringTok{"\#90aa58"}\FunctionTok{,}
    \DataTypeTok{"hair"}\FunctionTok{:}        \StringTok{"\#787878"}\FunctionTok{,}
    \DataTypeTok{"hair\_style"}\FunctionTok{:}  \StringTok{"slicked"}\FunctionTok{,}
    \DataTypeTok{"eye\_type"}\FunctionTok{:}    \StringTok{"alien"}\FunctionTok{,}
    \DataTypeTok{"eye\_color"}\FunctionTok{:}   \StringTok{"\#c87820"}\FunctionTok{,}
    \DataTypeTok{"brow\_thick"}\FunctionTok{:}  \KeywordTok{false}\FunctionTok{,}
    \DataTypeTok{"mouth"}\FunctionTok{:}       \StringTok{"flat"}\FunctionTok{,}
    \DataTypeTok{"cloth\_color"}\FunctionTok{:} \StringTok{"\#282828"}\FunctionTok{,}
    \DataTypeTok{"cloth\_style"}\FunctionTok{:} \StringTok{"blazer"}\FunctionTok{,}
    \DataTypeTok{"extras"}\FunctionTok{:}      \OtherTok{[]}\FunctionTok{,}
    \DataTypeTok{"hat\_style"}\FunctionTok{:}   \StringTok{"triangle"}\FunctionTok{,}
    \DataTypeTok{"hat\_color"}\FunctionTok{:}   \StringTok{"\#2a6830"}
  \FunctionTok{\}}
\FunctionTok{\}}
\end{Highlighting}
\end{Shaded}

\paragraph{Complete schema example: Nona (alien · elder, bun +
triangle\_inv hat +
blush)}\label{complete-schema-example-nona-alien-elder-bun-triangle_inv-hat-blush}

\begin{Shaded}
\begin{Highlighting}[]
\FunctionTok{\{}
  \DataTypeTok{"nona"}\FunctionTok{:} \FunctionTok{\{}
    \DataTypeTok{"name"}\FunctionTok{:}         \StringTok{"Nona"}\FunctionTok{,}
    \DataTypeTok{"note"}\FunctionTok{:}         \StringTok{"alien · elder"}\FunctionTok{,}
    \DataTypeTok{"skin"}\FunctionTok{:}         \StringTok{"\#809040"}\FunctionTok{,}
    \DataTypeTok{"hair"}\FunctionTok{:}         \StringTok{"\#909898"}\FunctionTok{,}
    \DataTypeTok{"hair\_style"}\FunctionTok{:}   \StringTok{"bun"}\FunctionTok{,}
    \DataTypeTok{"bun\_offset\_y"}\FunctionTok{:} \DecValTok{0}\ErrorTok{.}\DecValTok{155}\FunctionTok{,}
    \DataTypeTok{"eye\_type"}\FunctionTok{:}     \StringTok{"alien"}\FunctionTok{,}
    \DataTypeTok{"eye\_color"}\FunctionTok{:}    \StringTok{"\#3a5028"}\FunctionTok{,}
    \DataTypeTok{"brow\_thick"}\FunctionTok{:}   \KeywordTok{false}\FunctionTok{,}
    \DataTypeTok{"mouth"}\FunctionTok{:}        \StringTok{"smile"}\FunctionTok{,}
    \DataTypeTok{"cloth\_color"}\FunctionTok{:}  \StringTok{"\#8a6030"}\FunctionTok{,}
    \DataTypeTok{"cloth\_style"}\FunctionTok{:}  \StringTok{"standard"}\FunctionTok{,}
    \DataTypeTok{"extras"}\FunctionTok{:}       \OtherTok{[}\StringTok{"blush"}\OtherTok{]}\FunctionTok{,}
    \DataTypeTok{"blush\_color"}\FunctionTok{:}  \StringTok{"\#c03828"}\FunctionTok{,}
    \DataTypeTok{"hat\_style"}\FunctionTok{:}    \StringTok{"triangle\_inv"}\FunctionTok{,}
    \DataTypeTok{"hat\_color"}\FunctionTok{:}    \StringTok{"\#3a4858"}
  \FunctionTok{\}}
\FunctionTok{\}}
\end{Highlighting}
\end{Shaded}

\begin{center}\rule{0.5\linewidth}{0.5pt}\end{center}

\subsubsection{5.3 Expression variants}\label{expression-variants}

Expression variants are per-character face overrides --- typically a
mouth swap, sometimes a brow change or a blush toggle --- used to drive
moment-to-moment emotional beats without redefining the full face spec.
They live in the \texttt{"expressions"} key on each character's
\textbf{cast} entry (not the \texttt{"faces"} block) and are registered
into \texttt{FACE\_DATA} lazily by \texttt{register\_variants()} on the
first \texttt{attach\_face} step for that character.

You do \textbf{not} edit face\_builder.py to add expressions, and you do
not declare them in the \texttt{"faces"} block. All expression data
lives in tntd\_characters.json under the cast entry for the character.

\paragraph{\texorpdfstring{5.3.1 The \texttt{"expressions"} key on cast
entries}{5.3.1 The "expressions" key on cast entries}}\label{the-expressions-key-on-cast-entries}

In a character's cast spec, add an \texttt{"expressions"} dict alongside
\texttt{style}, \texttt{scale}, \texttt{pose}, etc.:

\begin{Shaded}
\begin{Highlighting}[]
\FunctionTok{\{}
  \DataTypeTok{"action"}\FunctionTok{:} \StringTok{"cast"}\FunctionTok{,}
  \DataTypeTok{"characters"}\FunctionTok{:} \FunctionTok{\{}
    \DataTypeTok{"sidel"}\FunctionTok{:} \FunctionTok{\{}
      \DataTypeTok{"figure\_type"}\FunctionTok{:} \StringTok{"alien"}\FunctionTok{,}
      \DataTypeTok{"build"}\FunctionTok{:}       \StringTok{"alien"}\FunctionTok{,}
      \DataTypeTok{"gender"}\FunctionTok{:}      \StringTok{"male"}\FunctionTok{,}
      \DataTypeTok{"pose"}\FunctionTok{:}        \StringTok{"standing\_front"}\FunctionTok{,}
      \DataTypeTok{"scale"}\FunctionTok{:}       \FunctionTok{\{} \DataTypeTok{"sy"}\FunctionTok{:} \DecValTok{0}\ErrorTok{.}\DecValTok{7}\FunctionTok{,} \DataTypeTok{"sx"}\FunctionTok{:} \DecValTok{0}\ErrorTok{.}\DecValTok{7}\FunctionTok{,} \DataTypeTok{"anchor"}\FunctionTok{:} \StringTok{"lankle"} \FunctionTok{\},}
      \DataTypeTok{"style"}\FunctionTok{:}       \FunctionTok{\{} \DataTypeTok{"head\_label"}\FunctionTok{:} \StringTok{"Sidel"}\FunctionTok{,} \DataTypeTok{"edge\_color"}\FunctionTok{:} \StringTok{"\#c04010"} \FunctionTok{\},}
      \DataTypeTok{"expressions"}\FunctionTok{:} \FunctionTok{\{}
        \DataTypeTok{"sidel\_smile"}\FunctionTok{:}   \FunctionTok{\{}\DataTypeTok{"mouth"}\FunctionTok{:} \StringTok{"smile"}\FunctionTok{\},}
        \DataTypeTok{"sidel\_flat"}\FunctionTok{:}    \FunctionTok{\{}\DataTypeTok{"mouth"}\FunctionTok{:} \StringTok{"flat"}\FunctionTok{\},}
        \DataTypeTok{"sidel\_neutral"}\FunctionTok{:} \FunctionTok{\{}\DataTypeTok{"mouth"}\FunctionTok{:} \StringTok{"neutral"}\FunctionTok{\},}
        \DataTypeTok{"sidel\_wry"}\FunctionTok{:}     \FunctionTok{\{}\DataTypeTok{"mouth"}\FunctionTok{:} \StringTok{"wry"}\FunctionTok{\},}
        \DataTypeTok{"sidel\_worried"}\FunctionTok{:} \FunctionTok{\{}\DataTypeTok{"mouth"}\FunctionTok{:} \StringTok{"flat"}\FunctionTok{,} \DataTypeTok{"brow\_thick"}\FunctionTok{:} \KeywordTok{true}\FunctionTok{\}}
      \FunctionTok{\}}
    \FunctionTok{\}}
  \FunctionTok{\}}
\FunctionTok{\}}
\end{Highlighting}
\end{Shaded}

Each variant key (\texttt{"sidel\_smile"}, \texttt{"sidel\_worried"},
\ldots) is the literal string used as the \texttt{"image"} value in
subsequent \texttt{attach\_face} steps. The override dict lists only the
fields that change relative to the base face spec --- everything else
(skin, hair, eyes, clothing, hat) is inherited automatically.

\textbf{pam\_player must copy \texttt{"expressions"} from the cast spec
into the live cast dict.} If \texttt{pam\_player.py} is missing the line

\begin{Shaded}
\begin{Highlighting}[]
\CommentTok{"expressions"}\NormalTok{: spec.get(}\StringTok{"expressions"}\NormalTok{, \{\}),}
\end{Highlighting}
\end{Shaded}

in its cast handler, the expressions block never reaches
\texttt{act\_attach\_face} and variants will not register. See §5.7,
item 1.

\paragraph{5.3.2 Naming convention}\label{naming-convention}

\textbf{Simple variants} (one mouth value, no other overrides) use the
pattern \texttt{\{character\}\_\{mouth\_value\}}:

\begin{verbatim}
sidel_smile    sidel_flat    sidel_neutral    sidel_wry
nona_smile     nona_flat     nona_neutral     nona_wry
bevers_smile   bevers_flat   bevers_neutral   bevers_wry
\end{verbatim}

The variant key suffix matches the \texttt{"mouth"} parameter value
directly --- no translation needed.

\textbf{Compound variants} combine two parameters and have no single
mouth-value equivalent, so they use a descriptive suffix:

\begin{longtable}[]{@{}ll@{}}
\toprule\noalign{}
Variant suffix & Composition \\
\midrule\noalign{}
\endhead
\bottomrule\noalign{}
\endlastfoot
\texttt{\_worried} & \texttt{"mouth":\ "flat",\ "brow\_thick":\ true} \\
\texttt{\_flustered} & \texttt{"mouth":\ "flat",\ "blush":\ true} \\
\end{longtable}

Add new compound names ad-hoc as scenes require them; the naming is
convention only, not enforced anywhere in code. Once you settle on a new
compound name (e.g.~\texttt{\_smug} = \texttt{wry} +
\texttt{brow\_thick}), use it consistently across all characters so a
scene's expression beats read uniformly.

\paragraph{5.3.3 The override knobs}\label{the-override-knobs}

Only these three fields differ between a base face and its variants.
Everything else is inherited from the base FACE\_DATA entry ---
including hats, hair style, clothing, and \texttt{bun\_offset\_y}.

\begin{longtable}[]{@{}
  >{\raggedright\arraybackslash}p{(\columnwidth - 4\tabcolsep) * \real{0.3333}}
  >{\raggedright\arraybackslash}p{(\columnwidth - 4\tabcolsep) * \real{0.3333}}
  >{\raggedright\arraybackslash}p{(\columnwidth - 4\tabcolsep) * \real{0.3333}}@{}}
\toprule\noalign{}
\begin{minipage}[b]{\linewidth}\raggedright
Key
\end{minipage} & \begin{minipage}[b]{\linewidth}\raggedright
Type
\end{minipage} & \begin{minipage}[b]{\linewidth}\raggedright
Description
\end{minipage} \\
\midrule\noalign{}
\endhead
\bottomrule\noalign{}
\endlastfoot
\texttt{mouth} & \texttt{str} & \texttt{"smile"}, \texttt{"flat"},
\texttt{"neutral"}, or \texttt{"wry"}. (\texttt{"wry"} maps to
\texttt{"neutral"} in side view; see §5.2.6.) \\
\texttt{brow\_thick} & \texttt{bool} & \texttt{true} → heavier brow
stroke. Reads as concern or anger depending on mouth pairing. Omit
(default \texttt{false}) for simple mouth-only variants. \\
\texttt{blush} & \texttt{bool} & \texttt{true} → add \texttt{"blush"} to
extras. \texttt{false} → remove \texttt{"blush"} from extras even if the
base face has it. Do \textbf{not} pass \texttt{"extras"} directly in an
expression --- \texttt{register\_variants()} handles the list merge. \\
\end{longtable}

Any other keys present in an override dict are silently ignored. If you
find yourself wanting to override \texttt{cloth\_color} or
\texttt{hair\_style} mid-scene, you've left the expression-variant lane
--- define a separate base face under a new key in the \texttt{"faces"}
block instead.

\paragraph{\texorpdfstring{5.3.4 How \texttt{register\_variants()}
works}{5.3.4 How register\_variants() works}}\label{how-register_variants-works}

\texttt{act\_attach\_face} calls
\texttt{register\_variants(char\_key,\ expressions)} the first time a
face is attached for a character, \textbf{before} the FACE\_DATA lookup
runs. The function logic:

\begin{enumerate}
\def\labelenumi{\arabic{enumi}.}
\tightlist
\item
  If \texttt{char\_key} is not in \texttt{FACE\_DATA}, return early
  (silent skip --- the base face was never loaded).
\item
  Read the base spec: \texttt{base\ =\ FACE\_DATA{[}char\_key{]}}.
\item
  For each \texttt{variant\_key\ →\ overrides} pair:

  \begin{itemize}
  \tightlist
  \item
    \textbf{Idempotent skip:} if \texttt{variant\_key} is already in
    \texttt{FACE\_DATA}, leave it alone.
  \item
    Build \texttt{merged\ =\ \{**base\}} (a shallow copy of the base).
  \item
    Apply \texttt{mouth} and \texttt{brow\_thick} overrides verbatim.
  \item
    For the \texttt{blush} convenience bool: copy
    \texttt{base{[}"extras"{]}} into a fresh list and add or remove
    \texttt{"blush"} to match the override.
  \item
    Write \texttt{merged} into \texttt{FACE\_DATA{[}variant\_key{]}}.
  \end{itemize}
\end{enumerate}

Because the function is idempotent, \texttt{act\_attach\_face} calls it
on every face swap without harm --- second and later calls do nothing
for already-registered variants.

\texttt{register\_variants()} is also exposed as public API for
standalone testing:

\begin{Shaded}
\begin{Highlighting}[]
\ImportTok{from}\NormalTok{ pam.face\_builder }\ImportTok{import}\NormalTok{ register\_variants}
\NormalTok{register\_variants(}\StringTok{"sidel"}\NormalTok{, \{}
    \StringTok{"sidel\_smile"}\NormalTok{:   \{}\StringTok{"mouth"}\NormalTok{: }\StringTok{"smile"}\NormalTok{\},}
    \StringTok{"sidel\_worried"}\NormalTok{: \{}\StringTok{"mouth"}\NormalTok{: }\StringTok{"flat"}\NormalTok{, }\StringTok{"brow\_thick"}\NormalTok{: }\VariableTok{True}\NormalTok{\},}
\NormalTok{\})}
\end{Highlighting}
\end{Shaded}

\begin{center}\rule{0.5\linewidth}{0.5pt}\end{center}

\subsubsection{\texorpdfstring{5.4 \texttt{attach\_face} --- the
action}{5.4 attach\_face --- the action}}\label{attach_face-the-action}

\emph{Parallel-safe:} yes · \emph{Animation:} instant (no tween) ·
\emph{Required keys:} \texttt{who}, \texttt{image}

Attach a flat-cartoon face VGroup over a character's head dot. The face
is built fresh from FACE\_DATA on every call (no caching), positioned
via the \texttt{pam\_head\_ref} anchor (§5.4.4), and attached with an
updater so it follows the head dot through subsequent motion.

\paragraph{5.4.1 JSON syntax}\label{json-syntax}

\begin{Shaded}
\begin{Highlighting}[]
\FunctionTok{\{}\DataTypeTok{"action"}\FunctionTok{:} \StringTok{"attach\_face"}\FunctionTok{,} \DataTypeTok{"who"}\FunctionTok{:} \StringTok{"freydoon"}\FunctionTok{,}
 \DataTypeTok{"image"}\FunctionTok{:} \StringTok{"freydoon"}\FunctionTok{,} \DataTypeTok{"scale"}\FunctionTok{:} \DecValTok{0}\ErrorTok{.}\DecValTok{35}\FunctionTok{\}}
\end{Highlighting}
\end{Shaded}

\begin{longtable}[]{@{}
  >{\raggedright\arraybackslash}p{(\columnwidth - 8\tabcolsep) * \real{0.2000}}
  >{\raggedright\arraybackslash}p{(\columnwidth - 8\tabcolsep) * \real{0.2000}}
  >{\raggedright\arraybackslash}p{(\columnwidth - 8\tabcolsep) * \real{0.2000}}
  >{\raggedright\arraybackslash}p{(\columnwidth - 8\tabcolsep) * \real{0.2000}}
  >{\raggedright\arraybackslash}p{(\columnwidth - 8\tabcolsep) * \real{0.2000}}@{}}
\toprule\noalign{}
\begin{minipage}[b]{\linewidth}\raggedright
Key
\end{minipage} & \begin{minipage}[b]{\linewidth}\raggedright
Type
\end{minipage} & \begin{minipage}[b]{\linewidth}\raggedright
Required
\end{minipage} & \begin{minipage}[b]{\linewidth}\raggedright
Default
\end{minipage} & \begin{minipage}[b]{\linewidth}\raggedright
Description
\end{minipage} \\
\midrule\noalign{}
\endhead
\bottomrule\noalign{}
\endlastfoot
\texttt{who} & \texttt{str} & yes & --- & Character key from the cast
block. \\
\texttt{image} & \texttt{str} & yes & --- & Either a \texttt{FACE\_DATA}
key (no file extension) or a filename in \texttt{pam/assets/}
(\texttt{.png} or \texttt{.svg} extension). \\
\texttt{scale} & \texttt{float} & no & per-build default (typically
\texttt{0.35}) & Scale applied to the built VGroup before alignment. \\
\texttt{view} & \texttt{str} & no & \texttt{"front"} & \texttt{"front"},
\texttt{"lside"}, or \texttt{"rside"} (see §5.4.5). \\
\end{longtable}

\textbf{Source dispatch on \texttt{"image"}:}

\begin{itemize}
\tightlist
\item
  An \texttt{image} value with \textbf{no file extension} is treated as
  a \texttt{FACE\_DATA} key. \texttt{act\_attach\_face} calls
  \texttt{face\_builder.build\_face(image,\ view=view)} to build a fresh
  VGroup, then attaches it. Unknown keys raise \texttt{KeyError} from
  \texttt{build\_face()} with a list of valid options.
\item
  An \texttt{image} value with a \textbf{\texttt{.png} or \texttt{.svg}
  extension} is loaded from \texttt{pam/assets/} as a bitmap or SVG
  mobject. This is the legacy path; no FACE\_DATA lookup occurs, and
  \texttt{pam\_head\_ref} does not apply (see §5.4.4).
\end{itemize}

A character key is valid if it has been registered in
\texttt{FACE\_DATA} via the \texttt{"faces"} block (§5.1), \textbf{or}
is an expression variant lazily registered by
\texttt{register\_variants()} (§5.3.4), \textbf{or} is one of the three
illustrative example entries (\texttt{example\_human},
\texttt{example\_alien}, \texttt{example\_dog}). Project characters not
yet declared in the \texttt{"faces"} block fail with \texttt{KeyError}
on first attach.

\paragraph{5.4.2 Scale convention}\label{scale-convention}

A scale of \textbf{0.30--0.40} places the face at approximately the
right visual weight relative to the skeleton figure at default character
scale. Tune per scene if a character is unusually large or small on
screen.

\textbf{Keep \texttt{scale} identical across all \texttt{attach\_face}
calls for the same character.} A scale change between calls produces a
visible size jump because each \texttt{attach\_face} builds a fresh
VGroup and rescales it from the default. If a per-scene scale change is
genuinely needed, accept the jump and consider hiding it behind a cut or
a fade.

\paragraph{5.4.3 Expression swapping}\label{expression-swapping}

\texttt{act\_attach\_face} tears down the prior face automatically
before attaching the new one --- no \texttt{detach\_face} step is needed
between swaps. Typical mid-scene swap pattern:

\begin{Shaded}
\begin{Highlighting}[]
\FunctionTok{\{}\DataTypeTok{"action"}\FunctionTok{:} \StringTok{"attach\_face"}\FunctionTok{,} \DataTypeTok{"who"}\FunctionTok{:} \StringTok{"sidel"}\FunctionTok{,}
 \DataTypeTok{"image"}\FunctionTok{:} \StringTok{"sidel"}\FunctionTok{,} \DataTypeTok{"scale"}\FunctionTok{:} \DecValTok{0}\ErrorTok{.}\DecValTok{35}\FunctionTok{\}}\ErrorTok{,}

\FunctionTok{\{}\DataTypeTok{"action"}\FunctionTok{:} \StringTok{"say"}\FunctionTok{,} \DataTypeTok{"who"}\FunctionTok{:} \StringTok{"sidel"}\FunctionTok{,}
 \DataTypeTok{"text"}\FunctionTok{:} \StringTok{"The Council\textquotesingle{}s decision is final."}\FunctionTok{\}}\ErrorTok{,}

\FunctionTok{\{}\DataTypeTok{"action"}\FunctionTok{:} \StringTok{"attach\_face"}\FunctionTok{,} \DataTypeTok{"who"}\FunctionTok{:} \StringTok{"sidel"}\FunctionTok{,}
 \DataTypeTok{"image"}\FunctionTok{:} \StringTok{"sidel\_worried"}\FunctionTok{,} \DataTypeTok{"scale"}\FunctionTok{:} \DecValTok{0}\ErrorTok{.}\DecValTok{35}\FunctionTok{\}}\ErrorTok{,}

\FunctionTok{\{}\DataTypeTok{"action"}\FunctionTok{:} \StringTok{"say"}\FunctionTok{,} \DataTypeTok{"who"}\FunctionTok{:} \StringTok{"sidel"}\FunctionTok{,}
 \DataTypeTok{"text"}\FunctionTok{:} \StringTok{"This will not go well for us."}\FunctionTok{\}}\ErrorTok{,}

\FunctionTok{\{}\DataTypeTok{"action"}\FunctionTok{:} \StringTok{"attach\_face"}\FunctionTok{,} \DataTypeTok{"who"}\FunctionTok{:} \StringTok{"sidel"}\FunctionTok{,}
 \DataTypeTok{"image"}\FunctionTok{:} \StringTok{"sidel\_smile"}\FunctionTok{,} \DataTypeTok{"scale"}\FunctionTok{:} \DecValTok{0}\ErrorTok{.}\DecValTok{35}\FunctionTok{\}}
\end{Highlighting}
\end{Shaded}

Because hats are embedded in the face VGroup (§5.2.9), they persist
across swaps that resolve to variants of the same base character.
Factor's \texttt{triangle} hat stays in place through every
\texttt{factor\_*} swap.

\paragraph{\texorpdfstring{5.4.4 The \texttt{pam\_head\_ref} anchor
fix}{5.4.4 The pam\_head\_ref anchor fix}}\label{the-pam_head_ref-anchor-fix}

This is the alignment trick that makes tall hair and hats render
correctly. Understanding it is essential for debugging mis-anchored
faces.

\texttt{build\_face()} tags the returned VGroup with an attribute
pointing at the head oval Ellipse (layer 5 in front view, layer 6 in
side view):

\begin{Shaded}
\begin{Highlighting}[]
\NormalTok{face.pam\_head\_ref }\OperatorTok{=}\NormalTok{ head   }\CommentTok{\# the head oval Ellipse}
\end{Highlighting}
\end{Shaded}

\texttt{act\_attach\_face} uses this attribute to align the face by the
\textbf{head oval center}, not the VGroup bounding-box center. The shift
is:

\begin{Shaded}
\begin{Highlighting}[]
\NormalTok{face.shift(head\_dot\_center }\OperatorTok{{-}}\NormalTok{ face.pam\_head\_ref.get\_center())}
\end{Highlighting}
\end{Shaded}

\textbf{Why this exists:} without \texttt{pam\_head\_ref}, alignment
uses \texttt{face.move\_to(head\_dot\_center)}, which moves the VGroup's
bounding box center to the head dot. Tall hair (buns, updos) and hats
extend the bounding box upward, so \texttt{move\_to()} pushes the face
\textbf{downward} until the bounding box center coincides with the head
dot --- leaving the head oval below the head dot, with the hair or hat
occupying the space where the face should be.

\texttt{pam\_head\_ref.get\_center()} returns the \textbf{live world
position} of the head oval Ellipse after any prior scale or shift, so
the formula stays correct after scaling --- the scale operation moves
the Ellipse along with the rest of the VGroup, and
\texttt{get\_center()} reads the post-scale position.

\textbf{Fallback for bitmap / SVG faces:} PNG and SVG assets loaded from
\texttt{pam/assets/} have no \texttt{pam\_head\_ref} attribute.
\texttt{act\_attach\_face} falls back to the legacy
\texttt{move\_to(head\_dot\_center)} behavior in that case. If a bitmap
face mis-anchors, the fix is to crop or compose the source image so its
center matches the intended head center --- there is no
\texttt{pam\_head\_ref} equivalent for bitmap sources.

\textbf{Related anchor:} \texttt{cloth\_style:\ "panel"} (v0.9.X)
exposes a parallel attribute \texttt{face.pam\_panel\_ref} pointing at
the panel Rectangle, used by \texttt{get\_panel\_badge\_anchor()} to
compute the name-tag / badge attachment point. Same scale-survival
pattern as \texttt{pam\_head\_ref}. See §5.5.4.

\paragraph{5.4.5 Side views}\label{side-views}

\texttt{build\_face()} accepts a \texttt{view} parameter passed through
from \texttt{attach\_face}:

\begin{Shaded}
\begin{Highlighting}[]
\FunctionTok{\{}\DataTypeTok{"action"}\FunctionTok{:} \StringTok{"attach\_face"}\FunctionTok{,} \DataTypeTok{"who"}\FunctionTok{:} \StringTok{"thalia"}\FunctionTok{,}
 \DataTypeTok{"image"}\FunctionTok{:} \StringTok{"thalia"}\FunctionTok{,} \DataTypeTok{"view"}\FunctionTok{:} \StringTok{"lside"}\FunctionTok{,} \DataTypeTok{"scale"}\FunctionTok{:} \DecValTok{0}\ErrorTok{.}\DecValTok{35}\FunctionTok{\}}
\end{Highlighting}
\end{Shaded}

\begin{longtable}[]{@{}ll@{}}
\toprule\noalign{}
Value & Direction \\
\midrule\noalign{}
\endhead
\bottomrule\noalign{}
\endlastfoot
\texttt{"front"} (default) & Front-facing portrait. \\
\texttt{"lside"} & Profile facing LEFT (-x direction). \\
\texttt{"rside"} & Profile facing RIGHT (+x direction). \\
\end{longtable}

Use \texttt{view:\ "lside"} / \texttt{"rside"} whenever the character's
skeleton is in a side pose (e.g.~during a \texttt{walk\_to}, or while
standing in a side-facing pose). The profile face is built from a
separate code path --- wider head ellipse
(\texttt{\_SIDE\_HEAD\_RX\ =\ 0.30} vs front
\texttt{HEAD\_RX\ =\ 0.25}), single eye/brow/mouth on the face side, a
nose polygon protruding beyond the head edge, a single far-side ear
behind the head oval, and a far-side hair back panel for longer styles.
See §5.6 for the full side-view z-order.

Unknown \texttt{view} values raise \texttt{ValueError} from
\texttt{build\_face()}.

\begin{center}\rule{0.5\linewidth}{0.5pt}\end{center}

\subsubsection{5.5 Hat and hair geometry reference}\label{hat-geometry-reference}

This section documents the constants used by the hat and hair builders,
and the per-character override fields that can adjust hair geometry
without editing source. When a hat clips the hair, fails to cover a bun,
or sits at the wrong height, the fix is usually to adjust
\texttt{bun\_offset\_y} (§5.2.10) or the hair geometry overrides
(§5.5.5) rather than to change the constants --- but the constants are
listed here for completeness so anomalies can be diagnosed without
reading source.

All dimensions are in Manim units at scale=1.0.

\paragraph{5.5.1 Hat shapes}\label{hat-shapes}

\textbf{\texttt{"triangle"}} --- Triangle hat. Base at \texttt{HEAD\_RY}
(top of the head oval), tip pointing up. Defaults are equilateral; both
dimensions are overridable per character (v0.9.18).

\begin{longtable}[]{@{}
  >{\raggedright\arraybackslash}p{(\columnwidth - 4\tabcolsep) * \real{0.2200}}
  >{\raggedright\arraybackslash}p{(\columnwidth - 4\tabcolsep) * \real{0.3300}}
  >{\raggedright\arraybackslash}p{(\columnwidth - 4\tabcolsep) * \real{0.4500}}@{}}
\toprule\noalign{}
\begin{minipage}[b]{\linewidth}\raggedright
Dimension
\end{minipage} & \begin{minipage}[b]{\linewidth}\raggedright
Default
\end{minipage} & \begin{minipage}[b]{\linewidth}\raggedright
Override
\end{minipage} \\
\midrule\noalign{}
\endhead
\bottomrule\noalign{}
\endlastfoot
Base y & \texttt{HEAD\_RY\ =\ 0.28} & (anchored, not overridable) \\
Base width & \texttt{0.52} & \texttt{hat\_width} \\
Tip y above base & \texttt{0.45} (equilateral for default base) &
\texttt{hat\_height} (absolute, in Manim units) \\
\end{longtable}

When only \texttt{hat\_width} is given, the tip height auto-derives as
the equilateral value for the new base (\texttt{hat\_width\ *\ 0.866}),
so width-only changes preserve the equilateral silhouette. When only
\texttt{hat\_height} is given, the base width stays at \texttt{0.52}.
Specifying both gives a fully free triangle (e.g., a flat wide cone or a
tall narrow spike). Factor's hat. Pairs with Bevers's \texttt{flat\_top}
rectangle as a father-son shape family (square → triangle).

\textbf{\texttt{"triangle\_inv"}} --- Inverted triangle hat. Wide base
at the top, tip pointing down to just above the hairline. Defaults are
tuned to fully contain a bun; both dimensions are overridable per
character (v0.9.18).

\begin{longtable}[]{@{}
  >{\raggedright\arraybackslash}p{(\columnwidth - 4\tabcolsep) * \real{0.2800}}
  >{\raggedright\arraybackslash}p{(\columnwidth - 4\tabcolsep) * \real{0.3200}}
  >{\raggedright\arraybackslash}p{(\columnwidth - 4\tabcolsep) * \real{0.4000}}@{}}
\toprule\noalign{}
Dimension & Default & Override \\
\midrule\noalign{}
\endhead
\bottomrule\noalign{}
\endlastfoot
Tip y (anchor point) & \texttt{HEAD\_RY\ -\ 0.10\ =\ 0.18} & (anchored,
not overridable --- keeps the hat seated on the head crown) \\
Total vertical span & \texttt{0.20} (so top is at
\texttt{HEAD\_RY\ +\ 0.10\ =\ 0.38}) &
\texttt{hat\_height} (sets total span; top y = tip y +
\texttt{hat\_height}) \\
Base width (at top) & \texttt{0.80} (wide enough to fully contain a
bun) & \texttt{hat\_width} \\
\end{longtable}

The tip stays anchored at \texttt{HEAD\_RY\ -\ 0.10} regardless of
override, so the hat keeps its bun-covering relationship to the head
when only width is tweaked. Nona's hat.

\textbf{Side-view triangle hats:} The same \texttt{hat\_width} and
\texttt{hat\_height} values apply directly in side view (no front/side
scaling --- unlike \texttt{"square"} which scales front width down for
the profile). Side-view defaults are slightly different:
\texttt{triangle} default base in profile is
\texttt{2\ *\ \_SIDE\_HEAD\_RX\ +\ 0.04\ ≈\ 0.64} to match the head's
front-to-back depth; \texttt{triangle\_inv} default base in profile is
\texttt{0.76}. Specifying \texttt{hat\_width} overrides both views to
the same absolute value, so the hat looks identical from front and
side when overridden --- this is usually what you want for a
recognizable silhouette.

\textbf{Worked syntax examples:}

\textbf{Example 1 --- the original Factor case (flat wide
triangle).} The exact configuration that motivated the v0.9.18
parameterization:

\begin{Shaded}
\begin{Highlighting}[]
\FunctionTok{\{}
  \DataTypeTok{"factor"}\FunctionTok{:} \FunctionTok{\{}
    \DataTypeTok{"hat\_style"}\FunctionTok{:}  \StringTok{"triangle"}\FunctionTok{,}
    \DataTypeTok{"hat\_color"}\FunctionTok{:}  \StringTok{"\#2a6830"}\FunctionTok{,}
    \DataTypeTok{"hat\_width"}\FunctionTok{:}  \FloatTok{0.72}\FunctionTok{,}
    \DataTypeTok{"hat\_height"}\FunctionTok{:} \FloatTok{0.03}
  \FunctionTok{\}}
\FunctionTok{\}}
\end{Highlighting}
\end{Shaded}

Produces a base \texttt{0.72} wide at \texttt{HEAD\_RY}, tip
\texttt{0.03} above --- effectively a flat triangular ridge. Same shape
in front and side view.

\textbf{Example 2 --- width-only override (auto-equilateral).} A
wider equilateral triangle. The tip auto-derives at
\texttt{hat\_width\ *\ 0.866\ ≈\ 0.52} above the base:

\begin{Shaded}
\begin{Highlighting}[]
\FunctionTok{\{}
  \DataTypeTok{"factor"}\FunctionTok{:} \FunctionTok{\{}
    \DataTypeTok{"hat\_style"}\FunctionTok{:}  \StringTok{"triangle"}\FunctionTok{,}
    \DataTypeTok{"hat\_color"}\FunctionTok{:}  \StringTok{"\#2a6830"}\FunctionTok{,}
    \DataTypeTok{"hat\_width"}\FunctionTok{:}  \FloatTok{0.60}
  \FunctionTok{\}}
\FunctionTok{\}}
\end{Highlighting}
\end{Shaded}

\textbf{Example 3 --- tall narrow spike.}

\begin{Shaded}
\begin{Highlighting}[]
\FunctionTok{\{}
  \DataTypeTok{"factor"}\FunctionTok{:} \FunctionTok{\{}
    \DataTypeTok{"hat\_style"}\FunctionTok{:}  \StringTok{"triangle"}\FunctionTok{,}
    \DataTypeTok{"hat\_color"}\FunctionTok{:}  \StringTok{"\#2a6830"}\FunctionTok{,}
    \DataTypeTok{"hat\_width"}\FunctionTok{:}  \FloatTok{0.30}\FunctionTok{,}
    \DataTypeTok{"hat\_height"}\FunctionTok{:} \FloatTok{0.80}
  \FunctionTok{\}}
\FunctionTok{\}}
\end{Highlighting}
\end{Shaded}

\textbf{Example 4 --- narrower \texttt{triangle\_inv} keeping default
height.} The base narrows but the hat still seats at the same y
relative to the head:

\begin{Shaded}
\begin{Highlighting}[]
\FunctionTok{\{}
  \DataTypeTok{"nona"}\FunctionTok{:} \FunctionTok{\{}
    \DataTypeTok{"hat\_style"}\FunctionTok{:} \StringTok{"triangle\_inv"}\FunctionTok{,}
    \DataTypeTok{"hat\_color"}\FunctionTok{:} \StringTok{"\#5a4030"}\FunctionTok{,}
    \DataTypeTok{"hat\_width"}\FunctionTok{:} \FloatTok{0.60}
  \FunctionTok{\}}
\FunctionTok{\}}
\end{Highlighting}
\end{Shaded}

\textbf{Example 5 --- taller \texttt{triangle\_inv} (taller dome over a
bigger bun).} Tip stays at \texttt{HEAD\_RY\ -\ 0.10}; top extends up
to \texttt{HEAD\_RY\ -\ 0.10\ +\ 0.30\ =\ 0.48}:

\begin{Shaded}
\begin{Highlighting}[]
\FunctionTok{\{}
  \DataTypeTok{"nona"}\FunctionTok{:} \FunctionTok{\{}
    \DataTypeTok{"hat\_style"}\FunctionTok{:}  \StringTok{"triangle\_inv"}\FunctionTok{,}
    \DataTypeTok{"hat\_color"}\FunctionTok{:}  \StringTok{"\#5a4030"}\FunctionTok{,}
    \DataTypeTok{"hat\_height"}\FunctionTok{:} \FloatTok{0.30}
  \FunctionTok{\}}
\FunctionTok{\}}
\end{Highlighting}
\end{Shaded}

\textbf{Example 6 --- baseline (no overrides, current equilateral
default).} Both fields omitted; renders exactly as before v0.9.18:

\begin{Shaded}
\begin{Highlighting}[]
\FunctionTok{\{}
  \DataTypeTok{"factor"}\FunctionTok{:} \FunctionTok{\{}
    \DataTypeTok{"hat\_style"}\FunctionTok{:} \StringTok{"triangle"}\FunctionTok{,}
    \DataTypeTok{"hat\_color"}\FunctionTok{:} \StringTok{"\#2a6830"}
  \FunctionTok{\}}
\FunctionTok{\}}
\end{Highlighting}
\end{Shaded}

\paragraph{5.5.2 Module-level geometry
constants}\label{module-level-geometry-constants}

These constants live in \texttt{face\_builder.py} at module scope and
can be overridden at runtime to re-tune face geometry globally.
Per-character tuning should use FACE\_DATA fields instead; modifying the
constants affects every face.

\begin{longtable}[]{@{}
  >{\raggedright\arraybackslash}p{(\columnwidth - 4\tabcolsep) * \real{0.3333}}
  >{\raggedright\arraybackslash}p{(\columnwidth - 4\tabcolsep) * \real{0.3333}}
  >{\raggedright\arraybackslash}p{(\columnwidth - 4\tabcolsep) * \real{0.3333}}@{}}
\toprule\noalign{}
\begin{minipage}[b]{\linewidth}\raggedright
Constant
\end{minipage} & \begin{minipage}[b]{\linewidth}\raggedright
Value
\end{minipage} & \begin{minipage}[b]{\linewidth}\raggedright
Description
\end{minipage} \\
\midrule\noalign{}
\endhead
\bottomrule\noalign{}
\endlastfoot
\texttt{OUTLINE\_COLOR} & \texttt{"\#18120c"} & Default stroke color for
all elements (near-black warm). \\
\texttt{STROKE\_WIDTH} & \texttt{2.2} & Outline stroke weight (Manim
units × 100). \\
\texttt{HEAD\_RX} & \texttt{0.25} & Head oval x-radius (front view). \\
\texttt{HEAD\_RY} & \texttt{0.28} & Head oval y-radius (slightly taller
than wide). \\
\texttt{\_SIDE\_HEAD\_RX} & \texttt{0.30} & Profile head front-to-back
radius (side view). \\
\texttt{EYE\_Y} & \texttt{0.05} & Eye center y (above head center). \\
\texttt{EYE\_DX} & \texttt{0.09} & Eye center x offset from vertical
centerline. \\
\texttt{BROW\_Y} & \texttt{0.14} & Eyebrow arc centerline y. \\
\texttt{NOSE\_Y} & \texttt{-0.04} & Nose arc centerline y. \\
\texttt{MOUTH\_Y} & \texttt{-0.13} & Mouth arc/line centerline y. \\
\texttt{NECK\_W}, \texttt{NECK\_H} & \texttt{0.12}, \texttt{0.12} & Neck
rectangle dimensions. \\
\texttt{EYE\_R\_OUTER} & \texttt{0.038} & Outer eye circle radius
(sclera or alien iris). \\
\texttt{EYE\_R\_IRIS} & \texttt{0.020} & Human iris ring radius. \\
\texttt{EYE\_R\_PUPIL} & \texttt{0.010} & Pupil dark circle radius. \\
\texttt{EYE\_R\_LIGHT} & \texttt{0.005} & Catchlight specular dot
radius. \\
\texttt{PANEL\_WIDTH} (v0.9.X) & \texttt{0.32} & Width of the
\texttt{cloth\_style:\ "panel"} rectangular bib. \\
\texttt{PANEL\_TOP\_Y} (v0.9.X) & \texttt{-0.22} & Top edge of the panel
--- the symbolic ``neck\_bottom'' reference. \\
\texttt{PANEL\_BOTTOM\_Y} (v0.9.X) & \texttt{-0.60} & Bottom edge of the
default panel. \\
\texttt{PANEL\_DEFAULT\_LENGTH} (v0.9.X) & \texttt{0.38} & Derived bib
height, \texttt{PANEL\_TOP\_Y\ -\ PANEL\_BOTTOM\_Y}. \\
\texttt{PANEL\_SIDE\_OFFSET\_X} (v0.9.X) & \texttt{0.04} & Side-view
shift onto face side (multiplied by \texttt{xs}). \\
\texttt{PANEL\_BADGE\_Y\_OFFSET} (v0.9.X) & \texttt{0.08} & Badge anchor
distance below \texttt{PANEL\_TOP\_Y}. \\
\end{longtable}

Approximate bounding box of a built face at scale=1.0:

\begin{longtable}[]{@{}ll@{}}
\toprule\noalign{}
Feature & y \\
\midrule\noalign{}
\endhead
\bottomrule\noalign{}
\endlastfoot
Top of tallest updo hair & \texttt{≈\ +0.86} \\
Top of typical short cap & \texttt{≈\ +0.50} \\
Head center & \texttt{0.00} \\
Chin & \texttt{≈\ −0.28} \\
Bottom of clothing shape & \texttt{≈\ −0.45} \\
\end{longtable}

Side-view bounding box differs slightly:

\begin{longtable}[]{@{}ll@{}}
\toprule\noalign{}
Feature & y or x \\
\midrule\noalign{}
\endhead
\bottomrule\noalign{}
\endlastfoot
Top of updo & \texttt{y\ ≈\ +0.87} \\
Top of \texttt{flat\_top} & \texttt{y\ ≈\ +0.54} \\
Head top & \texttt{y\ ≈\ +0.28} \\
Nose tip (lside) & \texttt{x\ ≈\ −0.38} (rside: \texttt{x\ ≈\ +0.38}) \\
Back of head & \texttt{x\ ≈\ ±0.30} \\
Clothing bottom & \texttt{y\ ≈\ −0.45} \\
\end{longtable}

\paragraph{\texorpdfstring{5.5.3 Bun coverage by
\texttt{triangle\_inv}}{5.5.3 Bun coverage by triangle\_inv}}\label{bun-coverage-by-triangle_inv}

The default \texttt{bun\_offset\_y\ =\ 0.155} places Nona's bun so that:

\begin{longtable}[]{@{}ll@{}}
\toprule\noalign{}
Feature & y \\
\midrule\noalign{}
\endhead
\bottomrule\noalign{}
\endlastfoot
Bun center & \texttt{HEAD\_RY\ +\ 0.155\ =\ 0.435} \\
Bun top (radius \texttt{0.15}) & \texttt{0.585} \\
\texttt{triangle\_inv} base (top of hat) &
\texttt{HEAD\_RY\ +\ 0.52\ =\ 0.800} \\
Clearance & \texttt{0.215} \\
\end{longtable}

If a custom \texttt{bun\_offset\_y} pushes the bun top above
\texttt{0.65}, the \texttt{triangle\_inv} will clip the bun crown.
Either lower the bun (decrease \texttt{bun\_offset\_y}) or accept the
clip --- the hat builder does not currently expose a ``rise'' parameter
for the hat itself.

In side view the bun shifts toward the far side of the head, and the bun
y-offset is recomputed as \texttt{bun\_offset\_y\ *\ 0.71}. A bun that
just barely clears the hat in front view will have additional clearance
in side view because the side-view multiplier brings the bun lower.

\paragraph{\texorpdfstring{5.5.4 \texttt{cloth\_style:\ "panel"}
geometry
(v0.9.X)}{5.5.4 cloth\_style: "panel" geometry (v0.9.X)}}\label{cloth_style-panel-geometry-v0.9.x}

When \texttt{cloth\_style\ ==\ "panel"}, the default trapezoid torso is
replaced by a Rectangle of dimensions
\texttt{PANEL\_WIDTH\ ×\ PANEL\_DEFAULT\_LENGTH\ =\ 0.32\ ×\ 0.38}.

\begin{longtable}[]{@{}ll@{}}
\toprule\noalign{}
Dimension & Value \\
\midrule\noalign{}
\endhead
\bottomrule\noalign{}
\endlastfoot
Width & \texttt{PANEL\_WIDTH\ =\ 0.32} \\
Height & \texttt{PANEL\_DEFAULT\_LENGTH\ =\ 0.38} \\
Top edge y & \texttt{PANEL\_TOP\_Y\ =\ -0.22} \\
Bottom edge y & \texttt{PANEL\_BOTTOM\_Y\ =\ -0.60} \\
Center y &
\texttt{PANEL\_TOP\_Y\ -\ PANEL\_DEFAULT\_LENGTH\ /\ 2\ =\ -0.41} \\
Front-view center x & \texttt{0} \\
Side-view center x &
\texttt{xs\ *\ PANEL\_SIDE\_OFFSET\_X\ =\ xs\ *\ 0.04} \\
\end{longtable}

\texttt{PANEL\_TOP\_Y\ =\ -0.22} is the symbolic ``neck\_bottom''
reference --- it matches the y-coordinate of the inner top edge of the
default trapezoid (the torso's neck opening). The panel hangs from this
y down to \texttt{PANEL\_BOTTOM\_Y\ =\ -0.60}. The top edge stays fixed
at the neck anchor; changing the panel length pushes only the bottom
edge downward.

\textbf{Side view:} the same Rectangle dimensions are used; only the
x-shift differs. Far-side arm occlusion when a character is in a side
pose comes from skeleton draw order --- \texttt{face\_builder.py} does
no special handling for the far-side arm crossing the panel.

\textbf{Badge / name-tag anchor.} \texttt{build\_face()} tags the panel
Rectangle as \texttt{face.pam\_panel\_ref} (analogous to
\texttt{face.pam\_head\_ref} for the head oval --- see §5.4.4).
Consumers should retrieve the world-space anchor via the public helper:

\begin{Shaded}
\begin{Highlighting}[]
\ImportTok{from}\NormalTok{ pam.face\_builder }\ImportTok{import}\NormalTok{ get\_panel\_badge\_anchor}
\NormalTok{anchor }\OperatorTok{=}\NormalTok{ get\_panel\_badge\_anchor(face)   }\CommentTok{\# np.ndarray (x, y, z), or None}
\end{Highlighting}
\end{Shaded}

The helper returns the live world-space point at
\texttt{PANEL\_BADGE\_Y\_OFFSET\ =\ 0.08} below the panel top, computed
from \texttt{panel.get\_top()} and \texttt{panel.get\_bottom()} so the
anchor scales correctly with the face --- at scale=1.0 the offset is
\texttt{0.08}; at scale=0.35 it returns \texttt{0.028}. Returns
\texttt{None} for non-panel faces (no \texttt{pam\_panel\_ref}
attribute).

In face-local coordinates the anchor is \texttt{(0,\ -0.30,\ 0)} front
view and \texttt{(xs\ *\ 0.04,\ -0.30,\ 0)} side view, but you should
not hardcode these --- use \texttt{get\_panel\_badge\_anchor()} so the
value stays correct after \texttt{act\_attach\_face} has applied scale
and shift.

\textbf{Exclusivity.} Setting \texttt{cloth\_style:\ "panel"} skips all
the other torso decorations: blazer crease, military/uniform seam, and
shoulder epaulettes (even when \texttt{"epaulettes"} is listed in
\texttt{extras} or \texttt{cloth\_style\ ==\ "military"} would otherwise
apply --- but those two values are mutually exclusive with
\texttt{"panel"} anyway). Attach badges and name tags as separate props
using the anchor above.

\paragraph{\texorpdfstring{5.5.5 Hair geometry overrides
(v0.9.18)}{5.5.5 Hair geometry overrides (v0.9.18)}}\label{hair-geometry-overrides}

The defaults built into \texttt{\_build\_hair\_front} and
\texttt{\_build\_hair\_side\_front} are tuned for the original eleven-or-so
TNTD characters; on other faces they often read slightly too big or too
small. Three optional FACE\_DATA fields let you tune the \emph{primary
cap shape} of any \texttt{hair\_style} in place, without editing the
builder.

\begin{longtable}[]{@{}
  >{\raggedright\arraybackslash}p{(\columnwidth - 4\tabcolsep) * \real{0.2222}}
  >{\raggedright\arraybackslash}p{(\columnwidth - 4\tabcolsep) * \real{0.2222}}
  >{\raggedright\arraybackslash}p{(\columnwidth - 4\tabcolsep) * \real{0.5556}}@{}}
\toprule\noalign{}
\begin{minipage}[b]{\linewidth}\raggedright
Field
\end{minipage} & \begin{minipage}[b]{\linewidth}\raggedright
Applies to
\end{minipage} & \begin{minipage}[b]{\linewidth}\raggedright
Effect
\end{minipage} \\
\midrule\noalign{}
\endhead
\bottomrule\noalign{}
\endlastfoot
\texttt{hair\_width} & Front view only & Cap width in Manim units.
Replaces the per-style default. Smaller = thinner hair, larger = wider.
Not applied in side view because the side cap width is locked to the
head's front-to-back depth (\texttt{2\ *\ \_SIDE\_HEAD\_RX\ +\ 0.04\ ≈\ 0.64}),
which is a different physical dimension. \\
\texttt{hair\_height} & Both views & Cap height in Manim units. Replaces
the per-style default. Smaller = shorter hair, larger = taller. Cap
extends both above and below its center, so a taller cap lowers the
apparent hairline unless paired with \texttt{hair\_offset\_y}. \\
\texttt{hair\_offset\_y} & Both views & Vertical offset from
\texttt{HEAD\_RY} to the cap CENTER. Replaces the per-style default
offset. Use to raise or lower the hairline independently of height.
Positive = higher, negative = lower. \\
\end{longtable}

All three default to absent (i.e., to the per-style hardcoded value), so
omitting them preserves existing rendering exactly. Setting them
replaces the default outright (absolute override, not a multiplier ---
same semantic as the existing \texttt{hat\_width} / \texttt{hat\_height}
fields).

\textbf{Scope --- primary cap shape only.} The overrides affect the
single largest hair element (the cap ellipse, or for \texttt{flat\_top}
the cap rectangle). Secondary shapes are \emph{not} affected:

\begin{itemize}
\tightlist
\item
  \texttt{"updo"} --- the dome (\texttt{0.38\ ×\ 0.46}) and the top knot
  (\texttt{radius\ 0.10}) keep their hardcoded geometry.
\item
  \texttt{"bun"} --- the bun ball (\texttt{radius\ 0.15}) keeps its
  hardcoded geometry. Use \texttt{bun\_offset\_y} (§5.2.10) to move it
  vertically.
\item
  \texttt{"slicked"} --- the highlight arc keeps its hardcoded geometry.
\item
  \texttt{"grey\_wavy"} --- the wave highlight arc keeps its hardcoded
  geometry.
\item
  \texttt{"medium\_female"}, \texttt{"grey\_wavy"}, \texttt{"bob"} --- the
  back panels drawn behind the head oval (the longer-hair side panels)
  keep their hardcoded geometry. \texttt{hair\_back\_width} /
  \texttt{hair\_back\_height} fields may follow in a later release.
\end{itemize}

If you need to tune those secondary shapes, edit
\texttt{face\_builder.py} directly for now.

\paragraph{\texorpdfstring{5.5.6 Per-style hair defaults
(v0.9.18)}{5.5.6 Per-style hair defaults (v0.9.18)}}\label{per-style-hair-defaults}

Defaults the overrides replace, by style and view. Front-view defaults
matter for both width and height tuning; side-view defaults matter only
for height (width is locked).

\textbf{Front view --- primary cap shape:}

\begin{longtable}[]{@{}
  >{\raggedright\arraybackslash}p{(\columnwidth - 8\tabcolsep) * \real{0.2200}}
  >{\raggedright\arraybackslash}p{(\columnwidth - 8\tabcolsep) * \real{0.1500}}
  >{\raggedright\arraybackslash}p{(\columnwidth - 8\tabcolsep) * \real{0.1500}}
  >{\raggedright\arraybackslash}p{(\columnwidth - 8\tabcolsep) * \real{0.1800}}
  >{\raggedright\arraybackslash}p{(\columnwidth - 8\tabcolsep) * \real{0.2800}}@{}}
\toprule\noalign{}
\texttt{hair\_style} & Default width & Default height & Default offset y
& Shape \\
\midrule\noalign{}
\endhead
\bottomrule\noalign{}
\endlastfoot
\texttt{"flat\_top"} & \texttt{0.54} & \texttt{0.29} & \texttt{+0.06} &
Rectangle \\
\texttt{"stubble"} & \texttt{0.35} & \texttt{0.10} & \texttt{-0.01} &
Ellipse \\
\texttt{"short\_male"} & \texttt{0.52} & \texttt{0.22} & \texttt{+0.03} &
Ellipse \\
\texttt{"slicked"} & \texttt{0.54} & \texttt{0.18} & \texttt{+0.01} &
Ellipse + arc \\
\texttt{"medium\_female"} & \texttt{0.52} & \texttt{0.26} & \texttt{+0.05}
& Ellipse + back panels \\
\texttt{"grey\_wavy"} & \texttt{0.52} & \texttt{0.26} & \texttt{+0.05} &
Ellipse + back panels + wave arc \\
\texttt{"bob"} & \texttt{0.52} & \texttt{0.26} & \texttt{+0.05} &
Ellipse + shorter back panels \\
\texttt{"updo"} & \texttt{0.52} & \texttt{0.26} & \texttt{+0.05} &
Ellipse + dome + knot \\
\texttt{"bun"} & \texttt{0.50} & \texttt{0.22} & \texttt{+0.02} &
Ellipse + bun ball \\
(fallback) & \texttt{0.50} & \texttt{0.22} & \texttt{+0.02} & Ellipse \\
\end{longtable}

\textbf{Side view --- primary cap shape (height \& offset only):}

\begin{longtable}[]{@{}
  >{\raggedright\arraybackslash}p{(\columnwidth - 6\tabcolsep) * \real{0.3333}}
  >{\raggedright\arraybackslash}p{(\columnwidth - 6\tabcolsep) * \real{0.3333}}
  >{\raggedright\arraybackslash}p{(\columnwidth - 6\tabcolsep) * \real{0.3333}}@{}}
\toprule\noalign{}
\texttt{hair\_style} & Default height & Default offset y \\
\midrule\noalign{}
\endhead
\bottomrule\noalign{}
\endlastfoot
\texttt{"flat\_top"} & \texttt{0.26} & \texttt{+0.08} \\
\texttt{"stubble"} & \texttt{0.13} & \texttt{+0.01} \\
\texttt{"short\_male"} & \texttt{0.22} & \texttt{+0.03} \\
\texttt{"slicked"} & \texttt{0.17} & \texttt{+0.01} \\
\texttt{"medium\_female"} / \texttt{"grey\_wavy"} / \texttt{"bob"} &
\texttt{0.24} & \texttt{+0.04} \\
\texttt{"updo"} & \texttt{0.24} & \texttt{+0.04} \\
\texttt{"bun"} & \texttt{0.20} & \texttt{+0.02} \\
(fallback) & \texttt{0.20} & \texttt{+0.02} \\
\end{longtable}

Side-view heights differ slightly from front-view (most are
\texttt{0.02--0.04} smaller because the profile head is taller relative
to its width). The same \texttt{hair\_height} value applies to both
views; if you tune for front-view, the side view scales by the same
absolute amount, which generally reads correctly. Width is locked in
side view --- the cap always spans the head's front-to-back depth,
regardless of \texttt{hair\_width}.

\paragraph{\texorpdfstring{5.5.7 Syntax examples
(v0.9.18)}{5.5.7 Syntax examples (v0.9.18)}}\label{hair-geometry-examples}

\textbf{Example 1 --- baseline (no overrides, current default).} Factor
keeps the per-style \texttt{slicked} defaults (\texttt{0.54\ ×\ 0.18}
cap at \texttt{HEAD\_RY\ +\ 0.01}):

\begin{Shaded}
\begin{Highlighting}[]
\FunctionTok{\{}
  \DataTypeTok{"factor"}\FunctionTok{:} \FunctionTok{\{}
    \DataTypeTok{"hair"}\FunctionTok{:}       \StringTok{"\#a0a0a0"}\FunctionTok{,}
    \DataTypeTok{"hair\_style"}\FunctionTok{:} \StringTok{"slicked"}
  \FunctionTok{\}}
\FunctionTok{\}}
\end{Highlighting}
\end{Shaded}

\textbf{Example 2 --- narrower slicked hair.} Override width only, leave
height and offset alone. Useful when the default \texttt{0.54}-wide cap
overhangs a narrower head:

\begin{Shaded}
\begin{Highlighting}[]
\FunctionTok{\{}
  \DataTypeTok{"factor"}\FunctionTok{:} \FunctionTok{\{}
    \DataTypeTok{"hair"}\FunctionTok{:}        \StringTok{"\#a0a0a0"}\FunctionTok{,}
    \DataTypeTok{"hair\_style"}\FunctionTok{:}  \StringTok{"slicked"}\FunctionTok{,}
    \DataTypeTok{"hair\_width"}\FunctionTok{:}  \FloatTok{0.46}
  \FunctionTok{\}}
\FunctionTok{\}}
\end{Highlighting}
\end{Shaded}

\textbf{Example 3 --- taller cap, hairline preserved.} If you just bump
\texttt{hair\_height} the cap extends both up and down from its center,
dropping the hairline by half the height delta. Pair with
\texttt{hair\_offset\_y} to lift the cap by the same amount so the
hairline stays put. Slicked default is height \texttt{0.18} at offset
\texttt{+0.01}; adding \texttt{0.10} to both gives a cap of height
\texttt{0.28} sitting \texttt{0.05} higher than before, with the
hairline unchanged:

\begin{Shaded}
\begin{Highlighting}[]
\FunctionTok{\{}
  \DataTypeTok{"factor"}\FunctionTok{:} \FunctionTok{\{}
    \DataTypeTok{"hair"}\FunctionTok{:}           \StringTok{"\#a0a0a0"}\FunctionTok{,}
    \DataTypeTok{"hair\_style"}\FunctionTok{:}     \StringTok{"slicked"}\FunctionTok{,}
    \DataTypeTok{"hair\_height"}\FunctionTok{:}    \FloatTok{0.28}\FunctionTok{,}
    \DataTypeTok{"hair\_offset\_y"}\FunctionTok{:}  \FloatTok{0.06}
  \FunctionTok{\}}
\FunctionTok{\}}
\end{Highlighting}
\end{Shaded}

\textbf{Example 4 --- shorter cap, hairline preserved.} Mirror of
Example 3. To shrink the cap by \texttt{0.06} (from \texttt{0.18}
default to \texttt{0.12}) while keeping the hairline at the original
position, lower \texttt{hair\_offset\_y} by half the height reduction
(\texttt{0.01\ -\ 0.03\ =\ -0.02}):

\begin{Shaded}
\begin{Highlighting}[]
\FunctionTok{\{}
  \DataTypeTok{"factor"}\FunctionTok{:} \FunctionTok{\{}
    \DataTypeTok{"hair"}\FunctionTok{:}           \StringTok{"\#a0a0a0"}\FunctionTok{,}
    \DataTypeTok{"hair\_style"}\FunctionTok{:}     \StringTok{"slicked"}\FunctionTok{,}
    \DataTypeTok{"hair\_height"}\FunctionTok{:}    \FloatTok{0.12}\FunctionTok{,}
    \DataTypeTok{"hair\_offset\_y"}\FunctionTok{:} \OtherTok{-}\FloatTok{0.02}
  \FunctionTok{\}}
\FunctionTok{\}}
\end{Highlighting}
\end{Shaded}

\textbf{Example 5 --- ``barely-there'' stubble.} The
\texttt{"stubble"} style is already minimal (\texttt{0.35\ ×\ 0.10}),
but you can push it further:

\begin{Shaded}
\begin{Highlighting}[]
\FunctionTok{\{}
  \DataTypeTok{"yannos"}\FunctionTok{:} \FunctionTok{\{}
    \DataTypeTok{"hair"}\FunctionTok{:}        \StringTok{"\#403028"}\FunctionTok{,}
    \DataTypeTok{"hair\_style"}\FunctionTok{:}  \StringTok{"stubble"}\FunctionTok{,}
    \DataTypeTok{"hair\_width"}\FunctionTok{:}  \FloatTok{0.28}\FunctionTok{,}
    \DataTypeTok{"hair\_height"}\FunctionTok{:} \FloatTok{0.06}
  \FunctionTok{\}}
\FunctionTok{\}}
\end{Highlighting}
\end{Shaded}

\textbf{Example 6 --- wider \texttt{flat\_top}.} Bevers's defining
shape, made noticeably blockier:

\begin{Shaded}
\begin{Highlighting}[]
\FunctionTok{\{}
  \DataTypeTok{"bevers"}\FunctionTok{:} \FunctionTok{\{}
    \DataTypeTok{"hair"}\FunctionTok{:}        \StringTok{"\#e8c870"}\FunctionTok{,}
    \DataTypeTok{"hair\_style"}\FunctionTok{:}  \StringTok{"flat\_top"}\FunctionTok{,}
    \DataTypeTok{"hair\_width"}\FunctionTok{:}  \FloatTok{0.66}\FunctionTok{,}
    \DataTypeTok{"hair\_height"}\FunctionTok{:} \FloatTok{0.34}
  \FunctionTok{\}}
\FunctionTok{\}}
\end{Highlighting}
\end{Shaded}

\textbf{Example 7 --- tighter \texttt{bun} cap, bun ball
unaffected.} The cap shrinks; the bun ball (radius \texttt{0.15})
ignores the override because it's a secondary shape. Use
\texttt{bun\_offset\_y} to reposition the ball if needed:

\begin{Shaded}
\begin{Highlighting}[]
\FunctionTok{\{}
  \DataTypeTok{"nona"}\FunctionTok{:} \FunctionTok{\{}
    \DataTypeTok{"hair"}\FunctionTok{:}          \StringTok{"\#a8a8a8"}\FunctionTok{,}
    \DataTypeTok{"hair\_style"}\FunctionTok{:}    \StringTok{"bun"}\FunctionTok{,}
    \DataTypeTok{"hair\_width"}\FunctionTok{:}    \FloatTok{0.44}\FunctionTok{,}
    \DataTypeTok{"hair\_height"}\FunctionTok{:}   \FloatTok{0.18}\FunctionTok{,}
    \DataTypeTok{"bun\_offset\_y"}\FunctionTok{:}  \FloatTok{0.18}
  \FunctionTok{\}}
\FunctionTok{\}}
\end{Highlighting}
\end{Shaded}

\textbf{Example 8 --- \texttt{updo} with low base cap.} The dome and
knot keep their hardcoded geometry (Thalia's silhouette stays
recognizable), but the base cap can be tuned independently:

\begin{Shaded}
\begin{Highlighting}[]
\FunctionTok{\{}
  \DataTypeTok{"thalia"}\FunctionTok{:} \FunctionTok{\{}
    \DataTypeTok{"hair"}\FunctionTok{:}           \StringTok{"\#5a3080"}\FunctionTok{,}
    \DataTypeTok{"hair\_style"}\FunctionTok{:}     \StringTok{"updo"}\FunctionTok{,}
    \DataTypeTok{"hair\_height"}\FunctionTok{:}    \FloatTok{0.20}\FunctionTok{,}
    \DataTypeTok{"hair\_offset\_y"}\FunctionTok{:}  \FloatTok{0.02}
  \FunctionTok{\}}
\FunctionTok{\}}
\end{Highlighting}
\end{Shaded}

\textbf{Example 9 --- offset-only tweak.} If the hairline is the only
thing off (the cap shape is fine, it just sits too high or too low), use
\texttt{hair\_offset\_y} alone. Here \texttt{short\_male} sits
\texttt{0.04} higher than its default \texttt{+0.03}:

\begin{Shaded}
\begin{Highlighting}[]
\FunctionTok{\{}
  \DataTypeTok{"freydoon"}\FunctionTok{:} \FunctionTok{\{}
    \DataTypeTok{"hair"}\FunctionTok{:}          \StringTok{"\#2c2018"}\FunctionTok{,}
    \DataTypeTok{"hair\_style"}\FunctionTok{:}    \StringTok{"short\_male"}\FunctionTok{,}
    \DataTypeTok{"hair\_offset\_y"}\FunctionTok{:} \FloatTok{0.07}
  \FunctionTok{\}}
\FunctionTok{\}}
\end{Highlighting}
\end{Shaded}

\textbf{Example 10 --- the same override on multiple characters with
different defaults.} The override is absolute, not a multiplier --- so
the same \texttt{hair\_width:\ 0.46} produces different effects on
characters whose defaults differ. On \texttt{flat\_top} (default
\texttt{0.54}) it narrows by \texttt{0.08}; on \texttt{bun} (default
\texttt{0.50}) it narrows by \texttt{0.04}; on \texttt{stubble} (default
\texttt{0.35}) it widens by \texttt{0.11}. Always cross-reference
§5.5.6 before applying the same value across characters.

\begin{center}\rule{0.5\linewidth}{0.5pt}\end{center}

\subsubsection{5.6 Layering order
reference}\label{layering-order-reference}

Both views use a fixed z-order --- face elements are added to the VGroup
in the order below, so later elements draw on top of earlier ones.

\paragraph{5.6.1 Front view (back to
front)}\label{front-view-back-to-front}

\begin{longtable}[]{@{}
  >{\raggedright\arraybackslash}p{(\columnwidth - 4\tabcolsep) * \real{0.3333}}
  >{\raggedright\arraybackslash}p{(\columnwidth - 4\tabcolsep) * \real{0.3333}}
  >{\raggedright\arraybackslash}p{(\columnwidth - 4\tabcolsep) * \real{0.3333}}@{}}
\toprule\noalign{}
\begin{minipage}[b]{\linewidth}\raggedright
Layer
\end{minipage} & \begin{minipage}[b]{\linewidth}\raggedright
Element
\end{minipage} & \begin{minipage}[b]{\linewidth}\raggedright
Notes
\end{minipage} \\
\midrule\noalign{}
\endhead
\bottomrule\noalign{}
\endlastfoot
1 & Clothing & Trapezoid torso + any collar / epaulette detail. \\
2 & Neck & Skin-colored rectangle. \\
3 & Ears & Visible only for hair styles in
\texttt{\_HAIR\_SHOWS\_EARS}. \\
4 & Hair back panels & For \texttt{medium\_female}, \texttt{bob},
\texttt{grey\_wavy} --- drawn BEHIND head oval. \\
5 & \textbf{Head oval} & Skin-colored Ellipse. Covers inner portions of
hair panels. Tagged as \texttt{face.pam\_head\_ref} (§5.4.4). \\
6 & Hair front cap & Drawn ON TOP of head oval. \\
7 & Blush & Low-opacity ellipses, if extras includes
\texttt{"blush"}. \\
8 & Eyes & Sclera + iris + pupil + catchlight (human) or disc + pupil +
catchlight (alien). \\
9 & Eyebrows & Stroke-only arcs. \\
10 & Nose & Skin-colored arc. \\
11 & Mouth & Arc or line per \texttt{mouth} value. \\
12 & Earrings & If extras includes \texttt{"gold\_earrings"} or
\texttt{"pearl\_earrings"}. \\
13 & Hat & If \texttt{hat\_style} is set. Drawn LAST --- on top of all
hair and face. \\
\end{longtable}

\paragraph{5.6.2 Side view (back to
front)}\label{side-view-back-to-front}

\begin{longtable}[]{@{}
  >{\raggedright\arraybackslash}p{(\columnwidth - 4\tabcolsep) * \real{0.3333}}
  >{\raggedright\arraybackslash}p{(\columnwidth - 4\tabcolsep) * \real{0.3333}}
  >{\raggedright\arraybackslash}p{(\columnwidth - 4\tabcolsep) * \real{0.3333}}@{}}
\toprule\noalign{}
\begin{minipage}[b]{\linewidth}\raggedright
Layer
\end{minipage} & \begin{minipage}[b]{\linewidth}\raggedright
Element
\end{minipage} & \begin{minipage}[b]{\linewidth}\raggedright
Notes
\end{minipage} \\
\midrule\noalign{}
\endhead
\bottomrule\noalign{}
\endlastfoot
1 & Clothing & Profile trapezoid + optional far-shoulder epaulette. \\
2 & Neck & Offset slightly toward face side. \\
3 & Far-side ear & Drawn BEFORE head oval --- head oval covers the inner
half. \\
4 & Far-side hair back panel & For \texttt{medium\_female},
\texttt{bob}, \texttt{grey\_wavy} only. \\
5 & Nose polygon & Drawn BEFORE head oval --- tip protrudes beyond face
edge; head oval covers the base. \\
6 & \textbf{Head oval} & Wider profile Ellipse. Covers ear/panel/nose
bases. Tagged as \texttt{face.pam\_head\_ref}. \\
7 & Hair front cap & Drawn ON TOP of head oval. \\
8 & Blush & Single spot, near/face side only. \\
9 & Single eye & Near/face side. \\
10 & Single eyebrow & Near/face side. \\
11 & Mouth & Near/face side, shorter span than front view. \\
12 & Earring & Far-side, short-hair styles only. \\
13 & Hat & If \texttt{hat\_style} is set. Drawn LAST. \\
\end{longtable}

The position of layer 5 (nose polygon, side view) is the key trick that
makes the side-view nose work --- drawing the polygon \emph{before} the
head oval lets the head oval mask the polygon's base, leaving only the
protruding tip visible beyond the face edge.

\begin{center}\rule{0.5\linewidth}{0.5pt}\end{center}

\subsubsection{5.7 Debugging checklist --- when faces don't appear or
mis-anchor}\label{debugging-checklist-when-faces-dont-appear-or-mis-anchor}

Ordered by frequency of occurrence in past renders.

\textbf{1. Console warning:
\texttt{PAMPlayer:\ attach\_face\ —\ unknown\ face\_builder\ key\ \textquotesingle{}sidel\_smile\textquotesingle{}}.}

The variant key is not in \texttt{FACE\_DATA} at lookup time. Check, in
order:

\begin{itemize}
\item
  The base character (\texttt{"sidel"}) is declared in the
  \texttt{"faces"} block in tntd\_characters.json. Without the base,
  \texttt{register\_variants()} returns early (silent skip) and the
  variant never enters \texttt{FACE\_DATA}.
\item
  The \texttt{"expressions"} block on the cast entry for
  \texttt{"sidel"} lists \texttt{"sidel\_smile"} with the expected
  override fields.
\item
  \texttt{pam\_player.py} is copying \texttt{"expressions"} from the
  cast spec into the live cast dict:

\begin{Shaded}
\begin{Highlighting}[]
\CommentTok{"expressions"}\NormalTok{: spec.get(}\StringTok{"expressions"}\NormalTok{, \{\}),}
\end{Highlighting}
\end{Shaded}

  If this line is absent from the cast handler,
  \texttt{act\_attach\_face} receives an empty expressions dict and
  \texttt{register\_variants()} registers nothing.
\item
  \texttt{face\_builder.py} is v0.9.17 or later (has
  \texttt{register\_variants()}).
\item
  Variant key spelling matches exactly --- \texttt{FACE\_DATA} keys are
  case-sensitive. \texttt{"Sidel\_Smile"} ≠ \texttt{"sidel\_smile"}.
\end{itemize}

\textbf{2. Face appears but is positioned wrong --- head dot visible
behind/below the face, or face floats above the body.}

\begin{itemize}
\tightlist
\item
  The face has tall hair (bun, updo) or a hat, and
  \texttt{pam\_head\_ref} is not being honored. This means either (a)
  the face was built from a bitmap/SVG asset (no \texttt{pam\_head\_ref}
  attribute), or (b) \texttt{act\_attach\_face} has fallen through to
  the legacy \texttt{move\_to()} path. Verify which branch fires in
  \texttt{act\_attach\_face}.
\item
  The character has been moved (\texttt{walk\_to},
  \texttt{group\_translate}, etc.) and the head-dot follower updater has
  lagged or been removed. Re-attach the face after motion completes.
\end{itemize}

\textbf{3. Hat clips the bun.}

\begin{itemize}
\tightlist
\item
  Decrease \texttt{bun\_offset\_y} (e.g.~\texttt{0.155\ →\ 0.140}) or
  switch from \texttt{triangle} to \texttt{triangle\_inv}. See §5.5.3
  for the clearance arithmetic.
\end{itemize}

\textbf{4. Expression swap produces a visible size jump.}

\begin{itemize}
\tightlist
\item
  \texttt{scale} differs between the two \texttt{attach\_face} calls.
  Make them identical. See §5.4.2.
\end{itemize}

\textbf{5. Hat disappears on expression swap.}

\begin{itemize}
\tightlist
\item
  The expression variant was declared in \texttt{"expressions"} (so it
  should inherit the base face's hat), but at some point a separate
  FACE\_DATA entry without \texttt{hat\_style} got loaded under the
  variant's key. Check the \texttt{"faces"} block --- only the base
  character should be declared there; variants must come through
  \texttt{"expressions"}, not through duplicate face entries.
\end{itemize}

\textbf{6. Side-view face has front-view geometry.}

\begin{itemize}
\tightlist
\item
  \texttt{view} parameter missing from the \texttt{attach\_face} step.
  Default is \texttt{"front"}. Pass \texttt{"view":\ "lside"} or
  \texttt{"view":\ "rside"} explicitly when the figure is in a side
  pose.
\end{itemize}

\textbf{7. Earrings missing in front view despite
\texttt{"gold\_earrings"} in extras.}

\begin{itemize}
\tightlist
\item
  The character has a long hair style (\texttt{medium\_female},
  \texttt{grey\_wavy}, or \texttt{bob}) which suppresses earring
  rendering. Either switch to a hair style in
  \texttt{\_HAIR\_SHOWS\_EARS} (\texttt{short\_male}, \texttt{stubble},
  \texttt{slicked}, \texttt{flat\_top}, \texttt{bun}, \texttt{updo}), or
  accept the suppression --- earrings will still appear at the far-side
  ear in side view if the hair style is in \texttt{\_HAIR\_SHOWS\_EARS}.
\end{itemize}

\textbf{8. Epaulettes set in \texttt{extras} (or
\texttt{cloth\_style:\ "military"}) but not rendering.}

\begin{itemize}
\tightlist
\item
  Character has \texttt{cloth\_style:\ "panel"}, which silently drops
  shoulder epaulettes --- the panel replaces the trapezoid, so the
  shoulder coordinates the epaulette builder uses
  (\texttt{x\ =\ ±0.30,\ y\ =\ -0.26}) fall outside the panel's
  geometry. This is intentional exclusivity (§5.5.4). Move rank insignia
  to a badge prop attached at \texttt{get\_panel\_badge\_anchor(face)}.
\end{itemize}

\textbf{9. Brow color barely visible against hair.}

\begin{itemize}
\tightlist
\item
  The default brow is hair darkened by \texttt{0.82}, which is too
  subtle for dark base hair. Set \texttt{brow\_color} explicitly to a
  darker tone (e.g.~Thalia's \texttt{"\#3a1858"} against hair
  \texttt{"\#5a3080"}).
\end{itemize}

\textbf{10. Wry mouth looks neutral on a side-pose beat.}

\begin{itemize}
\tightlist
\item
  This is expected behavior --- \texttt{"wry"} maps to
  \texttt{"neutral"} in side view because the asymmetry does not read
  from the side. See §5.2.6. If the wry read is critical for the beat,
  either rework the blocking so the character is front-facing, or accept
  the neutral substitution.
\end{itemize}

\textbf{11. List all registered face keys for verification.}

After a render, dump the FACE\_DATA registry to see what's actually
registered (including lazily-registered variants):

\begin{Shaded}
\begin{Highlighting}[]
\ImportTok{from}\NormalTok{ pam.face\_builder }\ImportTok{import}\NormalTok{ list\_characters}
\NormalTok{list\_characters()}
\end{Highlighting}
\end{Shaded}

Run this immediately after a scene's first \texttt{attach\_face} step to
verify both base characters and expression variants made it into
\texttt{FACE\_DATA}.

\begin{center}\rule{0.5\linewidth}{0.5pt}\end{center}

\emph{End of Section 5.}


\subsection{6. Wardrobe overlays, badges, and name
tags}\label{wardrobe-overlays-badges-and-name-tags}

This section covers decorative mobjects that attach to an already-live
figure. These are not pose data, prop registry entries, or skeleton
joints. They are small Manim mobjects with updaters that follow an
anchor point on the character every frame.

The authoring pattern mirrors \texttt{attach\_face}: attach once after
\texttt{fade\_in}, let the updater track walks and morphs, and either
detach explicitly or allow \texttt{fade\_out} to clean up the attachment
bundle.

\subsubsection{6.1 Common lifecycle rules}\label{common-lifecycle-rules}

All attached wardrobe overlays share the same lifecycle conventions:

\begin{longtable}[]{@{}
  >{\raggedright\arraybackslash}p{(\columnwidth - 2\tabcolsep) * \real{0.5000}}
  >{\raggedright\arraybackslash}p{(\columnwidth - 2\tabcolsep) * \real{0.5000}}@{}}
\toprule\noalign{}
\begin{minipage}[b]{\linewidth}\raggedright
Rule
\end{minipage} & \begin{minipage}[b]{\linewidth}\raggedright
Consequence
\end{minipage} \\
\midrule\noalign{}
\endhead
\bottomrule\noalign{}
\endlastfoot
Idempotent attach & Re-attaching the same kind of item first removes the
previous one. \\
Decorative only & They do not change pose dictionaries, joint positions,
hit detection, or action targeting. \\
Updater-driven & They follow live joint locations through
\texttt{walk\_to}, \texttt{run\_to}, \texttt{morph}, \texttt{turn}, and
scale changes. \\
Explicit detach available & Each attach action has a matching
\texttt{detach\_*} action. \\
Fade-out cleanup & \texttt{fade\_out} removes faces, torso icons,
gloves, shoes, name tags, and dog harnesses in the same concurrent fade
bundle as the body. \\
Joint-dependent no-op & If the figure lacks the needed joints, the
handler prints a diagnostic or skips cleanly. \\
\end{longtable}

These overlays are best used for readable character coding: uniforms,
delivery gloves, shoes, badges, harness labels, or temporary disguise
markers.

\subsubsection{\texorpdfstring{6.2 \texttt{attach\_torso\_icon} /
\texttt{detach\_torso\_icon}}{6.2 attach\_torso\_icon / detach\_torso\_icon}}\label{attach_torso_icon-detach_torso_icon}

\texttt{attach\_torso\_icon} anchors a small icon to the figure's torso
midpoint. In the attached source, the JSON-facing preset factory
currently ships one preset: \texttt{name\_tag}. The low-level
\texttt{HumanGraph.attach\_torso\_icon(icon,\ scene)} method can attach
any pre-built VMobject or VGroup, but the JSON action exposes only named
presets.

\begin{longtable}[]{@{}llrll@{}}
\toprule\noalign{}
Key & Type & Required & Default & Meaning \\
\midrule\noalign{}
\endhead
\bottomrule\noalign{}
\endlastfoot
\texttt{who} & string & yes & - & Character key. \\
\texttt{preset} & string & yes & - & Currently \texttt{"name\_tag"}. \\
\texttt{text} & string & no & \texttt{"NAME"} & Text for the
\texttt{name\_tag} preset. \\
\texttt{fill} & hex & no & \texttt{"\#f0e4b8"} & Plate fill color. \\
\texttt{stroke} & hex & no & \texttt{"\#3a2a14"} & Plate outline and
text color. \\
\texttt{scale} & float & no & \texttt{1.0} & Scale applied to the built
icon. \\
\end{longtable}

\begin{Shaded}
\begin{Highlighting}[]
\FunctionTok{\{}\DataTypeTok{"action"}\FunctionTok{:} \StringTok{"attach\_torso\_icon"}\FunctionTok{,} \DataTypeTok{"who"}\FunctionTok{:} \StringTok{"yannos"}\FunctionTok{,}
 \DataTypeTok{"preset"}\FunctionTok{:} \StringTok{"name\_tag"}\FunctionTok{,} \DataTypeTok{"text"}\FunctionTok{:} \StringTok{"Y. YANNOS"}\FunctionTok{,}
 \DataTypeTok{"fill"}\FunctionTok{:} \StringTok{"\#f0e4b8"}\FunctionTok{,} \DataTypeTok{"stroke"}\FunctionTok{:} \StringTok{"\#3a2a14"}\FunctionTok{,} \DataTypeTok{"scale"}\FunctionTok{:} \FloatTok{1.0}\FunctionTok{\}}
\end{Highlighting}
\end{Shaded}

\begin{Shaded}
\begin{Highlighting}[]
\FunctionTok{\{}\DataTypeTok{"action"}\FunctionTok{:} \StringTok{"detach\_torso\_icon"}\FunctionTok{,} \DataTypeTok{"who"}\FunctionTok{:} \StringTok{"yannos"}\FunctionTok{\}}
\end{Highlighting}
\end{Shaded}

The torso anchor uses the midpoint between \texttt{lshoulder} and
\texttt{lhip}, with the x-coordinate forced to the figure centerline.
This keeps the icon centered on both ordinary humans and split-torso
alien builds.

\subsubsection{\texorpdfstring{6.3 \texttt{attach\_gloves} /
\texttt{detach\_gloves}}{6.3 attach\_gloves / detach\_gloves}}\label{attach_gloves-detach_gloves}

\texttt{attach\_gloves} creates one flat-cartoon mitten ellipse at each
wrist. It is intentionally low-detail: a colored mass with a dark
outline, not articulated fingers.

\begin{longtable}[]{@{}
  >{\raggedright\arraybackslash}p{(\columnwidth - 8\tabcolsep) * \real{0.1875}}
  >{\raggedright\arraybackslash}p{(\columnwidth - 8\tabcolsep) * \real{0.1875}}
  >{\raggedleft\arraybackslash}p{(\columnwidth - 8\tabcolsep) * \real{0.2500}}
  >{\raggedright\arraybackslash}p{(\columnwidth - 8\tabcolsep) * \real{0.1875}}
  >{\raggedright\arraybackslash}p{(\columnwidth - 8\tabcolsep) * \real{0.1875}}@{}}
\toprule\noalign{}
\begin{minipage}[b]{\linewidth}\raggedright
Key
\end{minipage} & \begin{minipage}[b]{\linewidth}\raggedright
Type
\end{minipage} & \begin{minipage}[b]{\linewidth}\raggedleft
Required
\end{minipage} & \begin{minipage}[b]{\linewidth}\raggedright
Default
\end{minipage} & \begin{minipage}[b]{\linewidth}\raggedright
Meaning
\end{minipage} \\
\midrule\noalign{}
\endhead
\bottomrule\noalign{}
\endlastfoot
\texttt{who} & string & yes & - & Character key. \\
\texttt{color} & hex & yes & - & Glove fill color. \\
\texttt{size} & float & no & \texttt{0.22\ *\ fig.\_scale\_sy} & Glove
width in world units. \\
\texttt{stroke} & hex & no & \texttt{"\#18120c"} & Outline color. \\
\end{longtable}

\begin{Shaded}
\begin{Highlighting}[]
\FunctionTok{\{}\DataTypeTok{"action"}\FunctionTok{:} \StringTok{"attach\_gloves"}\FunctionTok{,} \DataTypeTok{"who"}\FunctionTok{:} \StringTok{"bevers"}\FunctionTok{,} \DataTypeTok{"color"}\FunctionTok{:} \StringTok{"\#cc3333"}\FunctionTok{\}}
\FunctionTok{\{}\DataTypeTok{"action"}\FunctionTok{:} \StringTok{"attach\_gloves"}\FunctionTok{,} \DataTypeTok{"who"}\FunctionTok{:} \StringTok{"freydoon"}\FunctionTok{,} \DataTypeTok{"color"}\FunctionTok{:} \StringTok{"\#3a2a14"}\FunctionTok{,} \DataTypeTok{"size"}\FunctionTok{:} \DecValTok{0}\ErrorTok{.}\DecValTok{18}\FunctionTok{\}}
\FunctionTok{\{}\DataTypeTok{"action"}\FunctionTok{:} \StringTok{"detach\_gloves"}\FunctionTok{,} \DataTypeTok{"who"}\FunctionTok{:} \StringTok{"bevers"}\FunctionTok{\}}
\end{Highlighting}
\end{Shaded}

The method requires \texttt{lwrist} and \texttt{rwrist}. Quadrupeds and
non-humanoid figures do not have the required wrist joints, so the
action is not useful for dogs or the Governor.

\subsubsection{\texorpdfstring{6.4 \texttt{attach\_shoes} /
\texttt{detach\_shoes}}{6.4 attach\_shoes / detach\_shoes}}\label{attach_shoes-detach_shoes}

\texttt{attach\_shoes} creates one elongated ellipse at each ankle. The
shoe updater tracks both ankle position and the figure's facing
direction. When \texttt{walk\_to} or \texttt{run\_to} changes
\texttt{fig.facing}, the shoes mirror so their long axis still points in
the direction of travel.

\begin{longtable}[]{@{}
  >{\raggedright\arraybackslash}p{(\columnwidth - 8\tabcolsep) * \real{0.1875}}
  >{\raggedright\arraybackslash}p{(\columnwidth - 8\tabcolsep) * \real{0.1875}}
  >{\raggedleft\arraybackslash}p{(\columnwidth - 8\tabcolsep) * \real{0.2500}}
  >{\raggedright\arraybackslash}p{(\columnwidth - 8\tabcolsep) * \real{0.1875}}
  >{\raggedright\arraybackslash}p{(\columnwidth - 8\tabcolsep) * \real{0.1875}}@{}}
\toprule\noalign{}
\begin{minipage}[b]{\linewidth}\raggedright
Key
\end{minipage} & \begin{minipage}[b]{\linewidth}\raggedright
Type
\end{minipage} & \begin{minipage}[b]{\linewidth}\raggedleft
Required
\end{minipage} & \begin{minipage}[b]{\linewidth}\raggedright
Default
\end{minipage} & \begin{minipage}[b]{\linewidth}\raggedright
Meaning
\end{minipage} \\
\midrule\noalign{}
\endhead
\bottomrule\noalign{}
\endlastfoot
\texttt{who} & string & yes & - & Character key. \\
\texttt{color} & hex & yes & - & Shoe fill color. \\
\texttt{size} & float & no & \texttt{0.32\ *\ fig.\_scale\_sy} & Shoe
length in world units. \\
\texttt{stroke} & hex & no & \texttt{"\#18120c"} & Outline color. \\
\end{longtable}

\begin{Shaded}
\begin{Highlighting}[]
\FunctionTok{\{}\DataTypeTok{"action"}\FunctionTok{:} \StringTok{"attach\_shoes"}\FunctionTok{,} \DataTypeTok{"who"}\FunctionTok{:} \StringTok{"bevers"}\FunctionTok{,} \DataTypeTok{"color"}\FunctionTok{:} \StringTok{"\#1a1a1a"}\FunctionTok{\}}
\FunctionTok{\{}\DataTypeTok{"action"}\FunctionTok{:} \StringTok{"attach\_shoes"}\FunctionTok{,} \DataTypeTok{"who"}\FunctionTok{:} \StringTok{"thalia"}\FunctionTok{,} \DataTypeTok{"color"}\FunctionTok{:} \StringTok{"\#3a2a14"}\FunctionTok{,} \DataTypeTok{"size"}\FunctionTok{:} \DecValTok{0}\ErrorTok{.}\DecValTok{30}\FunctionTok{\}}
\FunctionTok{\{}\DataTypeTok{"action"}\FunctionTok{:} \StringTok{"detach\_shoes"}\FunctionTok{,} \DataTypeTok{"who"}\FunctionTok{:} \StringTok{"thalia"}\FunctionTok{\}}
\end{Highlighting}
\end{Shaded}

Shoe height is about 40\% of the requested length, giving a readable 3:1
flat shoe silhouette.

\subsubsection{\texorpdfstring{6.5 \texttt{attach\_name\_tag} /
\texttt{detach\_name\_tag}}{6.5 attach\_name\_tag / detach\_name\_tag}}\label{attach_name_tag-detach_name_tag}

\texttt{attach\_name\_tag} attaches a text-only label at a semantic
badge anchor. It is different from the \texttt{attach\_torso\_icon}
\texttt{name\_tag} preset: the torso icon is a rectangular plate
centered on the torso; \texttt{attach\_name\_tag} is a free text mobject
attached to either a dog harness or a panel-clothed face.

Anchor resolution order:

\begin{enumerate}
\def\labelenumi{\arabic{enumi}.}
\tightlist
\item
  \texttt{fig.get\_harness\_nametag\_anchor()} for a DogGraph with a
  harness.
\item
  \texttt{face\_builder.get\_panel\_badge\_anchor(fig.head\_face)} for a
  human/alien face with \texttt{cloth\_style:\ "panel"}.
\item
  No anchor available: diagnostic and no-op.
\end{enumerate}

\begin{longtable}[]{@{}llrll@{}}
\toprule\noalign{}
Key & Type & Required & Default & Meaning \\
\midrule\noalign{}
\endhead
\bottomrule\noalign{}
\endlastfoot
\texttt{who} & string & yes & - & Character key. \\
\texttt{text} & string & yes & - & Label text; embedded
\texttt{\textbackslash{}n} is supported. \\
\texttt{font\_size} & float & no & \texttt{18} & Manim text font
size. \\
\texttt{color} & hex & no & \texttt{"\#ffffff"} & Text color. \\
\texttt{dy} & float & no & \texttt{0.0} & Extra vertical nudge in world
units. \\
\end{longtable}

\begin{Shaded}
\begin{Highlighting}[]
\FunctionTok{\{}\DataTypeTok{"action"}\FunctionTok{:} \StringTok{"attach\_name\_tag"}\FunctionTok{,} \DataTypeTok{"who"}\FunctionTok{:} \StringTok{"chekov"}\FunctionTok{,} \DataTypeTok{"text"}\FunctionTok{:} \StringTok{"CHEKOV"}\FunctionTok{\}}
\end{Highlighting}
\end{Shaded}

\begin{Shaded}
\begin{Highlighting}[]
\FunctionTok{\{}\DataTypeTok{"action"}\FunctionTok{:} \StringTok{"attach\_face"}\FunctionTok{,} \DataTypeTok{"who"}\FunctionTok{:} \StringTok{"sidel"}\FunctionTok{,} \DataTypeTok{"image"}\FunctionTok{:} \StringTok{"sidel"}\FunctionTok{,} \DataTypeTok{"scale"}\FunctionTok{:} \DecValTok{0}\ErrorTok{.}\DecValTok{35}\FunctionTok{\}}
\FunctionTok{\{}\DataTypeTok{"action"}\FunctionTok{:} \StringTok{"attach\_name\_tag"}\FunctionTok{,} \DataTypeTok{"who"}\FunctionTok{:} \StringTok{"sidel"}\FunctionTok{,}
 \DataTypeTok{"text"}\FunctionTok{:} \StringTok{"CMDR}\CharTok{\textbackslash{}n}\StringTok{SIDEL"}\FunctionTok{,} \DataTypeTok{"font\_size"}\FunctionTok{:} \DecValTok{14}\FunctionTok{,} \DataTypeTok{"color"}\FunctionTok{:} \StringTok{"\#c8a020"}\FunctionTok{,} \DataTypeTok{"dy"}\FunctionTok{:} \DecValTok{0}\ErrorTok{.02}\FunctionTok{\}}
\end{Highlighting}
\end{Shaded}

For panel-clothed humans and aliens, \texttt{attach\_face} must happen
before \texttt{attach\_name\_tag}, because the panel badge anchor lives
on the attached face VGroup as \texttt{face.pam\_panel\_ref}. For dogs,
the harness must be present in the dog style before fade-in.

\subsubsection{6.6 Dog harnesses and service
labels}\label{dog-harnesses-and-service-labels}

The attached \texttt{DogGraph} supports an optional triangular harness
icon. It is controlled by style keys on the dog figure, not by a
separate attach action.

\begin{longtable}[]{@{}
  >{\raggedright\arraybackslash}p{(\columnwidth - 6\tabcolsep) * \real{0.2500}}
  >{\raggedright\arraybackslash}p{(\columnwidth - 6\tabcolsep) * \real{0.2500}}
  >{\raggedright\arraybackslash}p{(\columnwidth - 6\tabcolsep) * \real{0.2500}}
  >{\raggedright\arraybackslash}p{(\columnwidth - 6\tabcolsep) * \real{0.2500}}@{}}
\toprule\noalign{}
\begin{minipage}[b]{\linewidth}\raggedright
Style key
\end{minipage} & \begin{minipage}[b]{\linewidth}\raggedright
Type
\end{minipage} & \begin{minipage}[b]{\linewidth}\raggedright
Default
\end{minipage} & \begin{minipage}[b]{\linewidth}\raggedright
Meaning
\end{minipage} \\
\midrule\noalign{}
\endhead
\bottomrule\noalign{}
\endlastfoot
\texttt{harness\_style} & string & omitted & Enables the harness.
\texttt{"standard"} is implemented; \texttt{"service"} is reserved. \\
\texttt{harness\_color} & hex & \texttt{edge\_color} & Fill color for
the triangular harness. \\
\end{longtable}

\begin{Shaded}
\begin{Highlighting}[]
\FunctionTok{\{}
  \DataTypeTok{"action"}\FunctionTok{:} \StringTok{"cast"}\FunctionTok{,}
  \DataTypeTok{"characters"}\FunctionTok{:} \FunctionTok{\{}
    \DataTypeTok{"chekov"}\FunctionTok{:} \FunctionTok{\{}
      \DataTypeTok{"figure\_type"}\FunctionTok{:} \StringTok{"dog"}\FunctionTok{,}
      \DataTypeTok{"offset"}\FunctionTok{:} \OtherTok{[}\DecValTok{{-}2}\OtherTok{,} \DecValTok{0}\OtherTok{,} \DecValTok{0}\OtherTok{]}\FunctionTok{,}
      \DataTypeTok{"style"}\FunctionTok{:} \FunctionTok{\{}
        \DataTypeTok{"harness\_style"}\FunctionTok{:} \StringTok{"standard"}\FunctionTok{,}
        \DataTypeTok{"harness\_color"}\FunctionTok{:} \StringTok{"\#2d6a4f"}
      \FunctionTok{\}}
    \FunctionTok{\}}
  \FunctionTok{\}}
\FunctionTok{\}}
\end{Highlighting}
\end{Shaded}

The harness is a rigid right-triangle badge anchored to
\texttt{spine\_mid}. It tracks the dog through trot cycles and flips its
head-side/tail-side geometry when the dog faces left. \texttt{fade\_in}
draws it after the dog skeleton so it sits visually above the body;
\texttt{fade\_out} includes it in the cleanup bundle.

\subsection{7. Figure style additions}\label{figure-style-additions}

This section documents figure-level style fields added in the attached
\texttt{figure.py}. These live in the character's \texttt{style} dict,
not in the \texttt{faces} block.

\subsubsection{7.1 Clothed limbs}\label{clothed-limbs}

\texttt{style{[}"clothed\_limbs"{]}} replaces selected arm and leg edges
with filled polygon bands. The rest of PAM still treats those bands as
edges: they live in \texttt{self.lines}, belong to \texttt{edge\_group},
and animate through \texttt{become(new\_band)} whenever the pose
changes.

\begin{Shaded}
\begin{Highlighting}[]
\FunctionTok{\{}
  \DataTypeTok{"action"}\FunctionTok{:} \StringTok{"cast"}\FunctionTok{,}
  \DataTypeTok{"characters"}\FunctionTok{:} \FunctionTok{\{}
    \DataTypeTok{"driver"}\FunctionTok{:} \FunctionTok{\{}
      \DataTypeTok{"build"}\FunctionTok{:} \StringTok{"broad"}\FunctionTok{,}
      \DataTypeTok{"offset"}\FunctionTok{:} \OtherTok{[}\DecValTok{{-}2}\OtherTok{,} \DecValTok{0}\OtherTok{,} \DecValTok{0}\OtherTok{]}\FunctionTok{,}
      \DataTypeTok{"style"}\FunctionTok{:} \FunctionTok{\{}
        \DataTypeTok{"edge\_color"}\FunctionTok{:} \StringTok{"\#304050"}\FunctionTok{,}
        \DataTypeTok{"clothed\_limbs"}\FunctionTok{:} \FunctionTok{\{}
          \DataTypeTok{"mode"}\FunctionTok{:} \StringTok{"tapered"}\FunctionTok{,}
          \DataTypeTok{"arm\_proximal\_width"}\FunctionTok{:} \DecValTok{0}\ErrorTok{.}\DecValTok{12}\FunctionTok{,}
          \DataTypeTok{"arm\_distal\_width"}\FunctionTok{:} \DecValTok{0}\ErrorTok{.07}\FunctionTok{,}
          \DataTypeTok{"leg\_proximal\_width"}\FunctionTok{:} \DecValTok{0}\ErrorTok{.}\DecValTok{16}\FunctionTok{,}
          \DataTypeTok{"leg\_distal\_width"}\FunctionTok{:} \DecValTok{0}\ErrorTok{.}\DecValTok{10}\FunctionTok{,}
          \DataTypeTok{"fill\_color"}\FunctionTok{:} \StringTok{"\#304050"}\FunctionTok{,}
          \DataTypeTok{"stroke\_color"}\FunctionTok{:} \StringTok{"\#18120c"}\FunctionTok{,}
          \DataTypeTok{"stroke\_width"}\FunctionTok{:} \FloatTok{1.8}
        \FunctionTok{\}}
      \FunctionTok{\}}
    \FunctionTok{\}}
  \FunctionTok{\}}
\FunctionTok{\}}
\end{Highlighting}
\end{Shaded}

\begin{longtable}[]{@{}
  >{\raggedright\arraybackslash}p{(\columnwidth - 6\tabcolsep) * \real{0.2500}}
  >{\raggedright\arraybackslash}p{(\columnwidth - 6\tabcolsep) * \real{0.2500}}
  >{\raggedright\arraybackslash}p{(\columnwidth - 6\tabcolsep) * \real{0.2500}}
  >{\raggedright\arraybackslash}p{(\columnwidth - 6\tabcolsep) * \real{0.2500}}@{}}
\toprule\noalign{}
\begin{minipage}[b]{\linewidth}\raggedright
Key
\end{minipage} & \begin{minipage}[b]{\linewidth}\raggedright
Values / type
\end{minipage} & \begin{minipage}[b]{\linewidth}\raggedright
Default
\end{minipage} & \begin{minipage}[b]{\linewidth}\raggedright
Meaning
\end{minipage} \\
\midrule\noalign{}
\endhead
\bottomrule\noalign{}
\endlastfoot
\texttt{mode} & \texttt{"uniform"} or \texttt{"tapered"} &
\texttt{"uniform"} & In uniform mode, each limb segment uses one width;
in tapered mode, proximal joints are thicker than distal joints. \\
\texttt{arm\_proximal\_width} & float & \texttt{0.10} & Full band width
near shoulder. \\
\texttt{arm\_distal\_width} & float & \texttt{0.06} & Full band width
near wrist. \\
\texttt{leg\_proximal\_width} & float & \texttt{0.14} & Full band width
near hip. \\
\texttt{leg\_distal\_width} & float & \texttt{0.09} & Full band width
near ankle. \\
\texttt{fill\_color} & hex & \texttt{edge\_color} & Band fill color. \\
\texttt{stroke\_color} & hex & \texttt{"\#18120c"} & Band outline
color. \\
\texttt{stroke\_width} & float & \texttt{1.8} & Band outline width. \\
\end{longtable}

The taper is continuous across elbows and knees: the width at an elbow
or knee is the arithmetic mean of proximal and distal widths. Coincident
endpoints produce invisible placeholder polygons, matching the existing
convention for degenerate lines.

\subsubsection{7.2 Character speech-bubble
colors}\label{character-speech-bubble-colors}

Human and alien figures support two style-level bubble keys:

\begin{longtable}[]{@{}
  >{\raggedright\arraybackslash}p{(\columnwidth - 4\tabcolsep) * \real{0.3333}}
  >{\raggedright\arraybackslash}p{(\columnwidth - 4\tabcolsep) * \real{0.3333}}
  >{\raggedright\arraybackslash}p{(\columnwidth - 4\tabcolsep) * \real{0.3333}}@{}}
\toprule\noalign{}
\begin{minipage}[b]{\linewidth}\raggedright
Style key
\end{minipage} & \begin{minipage}[b]{\linewidth}\raggedright
Default
\end{minipage} & \begin{minipage}[b]{\linewidth}\raggedright
Meaning
\end{minipage} \\
\midrule\noalign{}
\endhead
\bottomrule\noalign{}
\endlastfoot
\texttt{bubble\_color} & \texttt{None}, then fallback to
\texttt{edge\_color} & Drives the saturated border and the pale
tinted-white bubble fill. \\
\texttt{bubble\_text\_color} & \texttt{None}, then fallback to the dark
head color & Overrides text color when the automatic dark text
clashes. \\
\end{longtable}

\begin{Shaded}
\begin{Highlighting}[]
\FunctionTok{\{}
  \DataTypeTok{"style"}\FunctionTok{:} \FunctionTok{\{}
    \DataTypeTok{"edge\_color"}\FunctionTok{:} \StringTok{"\#cc3333"}\FunctionTok{,}
    \DataTypeTok{"bubble\_color"}\FunctionTok{:} \StringTok{"\#cc3333"}\FunctionTok{,}
    \DataTypeTok{"bubble\_text\_color"}\FunctionTok{:} \StringTok{"\#301010"}
  \FunctionTok{\}}
\FunctionTok{\}}
\end{Highlighting}
\end{Shaded}

At runtime, the implementation blends \texttt{bubble\_color} about 10\%
toward the source color and 90\% toward white for the fill. The result
is still visibly associated with the character, but pale enough for dark
text.

\subsubsection{7.3 Facing direction}\label{facing-direction}

\texttt{HumanGraph} now tracks \texttt{self.facing} as \texttt{"right"}
or \texttt{"left"}. \texttt{walk\_to} and \texttt{run\_to} update this
field from the sign of horizontal motion. This matters for attached
shoes, side-view face selection, and any later accessory whose
silhouette needs to mirror with travel direction.

\begin{Shaded}
\begin{Highlighting}[]
\OtherTok{[}
  \FunctionTok{\{}\DataTypeTok{"action"}\FunctionTok{:} \StringTok{"attach\_shoes"}\FunctionTok{,} \DataTypeTok{"who"}\FunctionTok{:} \StringTok{"bevers"}\FunctionTok{,} \DataTypeTok{"color"}\FunctionTok{:} \StringTok{"\#1a1a1a"}\FunctionTok{\}}\OtherTok{,}
  \FunctionTok{\{}\DataTypeTok{"action"}\FunctionTok{:} \StringTok{"walk\_to"}\FunctionTok{,} \DataTypeTok{"who"}\FunctionTok{:} \StringTok{"bevers"}\FunctionTok{,} \DataTypeTok{"x"}\FunctionTok{:} \FloatTok{3.0}\FunctionTok{\}}\OtherTok{,}
  \FunctionTok{\{}\DataTypeTok{"action"}\FunctionTok{:} \StringTok{"walk\_to"}\FunctionTok{,} \DataTypeTok{"who"}\FunctionTok{:} \StringTok{"bevers"}\FunctionTok{,} \DataTypeTok{"x"}\FunctionTok{:} \FloatTok{{-}2.0}\FunctionTok{\}}
\OtherTok{]}
\end{Highlighting}
\end{Shaded}

After the second walk, the shoe ellipses have flipped to preserve the
direction-of-travel read.


\chapter{Prop Properties Reference}\label{prop-properties-reference}

\textbf{PAM v0.9.13 --- props reference}

Companion document to \texttt{CHARACTER\_PROPERTIES.md}. Complete
reference for everything that can be created, modified, or done with a
prop in a PAM screenplay. Grouped by lifecycle: creation-time fields,
prop metadata, registry behavior, action targets, and meta concerns.
Heavy on syntax examples --- like the \texttt{to\_x}/\texttt{x} bite on
the character side, the \texttt{tv\_monitor} ``missing \texttt{type}
defaults to \texttt{hat}'' bite came from sparse examples.

\begin{center}\rule{0.5\linewidth}{0.5pt}\end{center}

\section{1. Creation-time fields}\label{creation-time-fields}

A prop is defined by an entry in either the manifest \texttt{props}
block (for props present at scene start) or a \texttt{spawn\_prop}
action (for props that appear mid-scene). Both share the same field set.

\subsection{\texorpdfstring{1.1 Required: \texttt{type},
\texttt{x}}{1.1 Required: type, x}}\label{required-type-x}

Every prop needs a \texttt{type} (the builder key) and an \texttt{x}
(world-space horizontal position, in Manim units; screen center is
\texttt{x:\ 0}).

\begin{Shaded}
\begin{Highlighting}[]
\FunctionTok{\{}\DataTypeTok{"type"}\FunctionTok{:} \StringTok{"chair"}\FunctionTok{,} \DataTypeTok{"x"}\FunctionTok{:} \FloatTok{{-}2.0}\FunctionTok{\}}
\end{Highlighting}
\end{Shaded}

In a manifest \texttt{props} block, each prop is keyed by a
\textbf{registry name} (unique, used in later actions):

\begin{Shaded}
\begin{Highlighting}[]
\FunctionTok{\{}\DataTypeTok{"action"}\FunctionTok{:} \StringTok{"props"}\FunctionTok{,} \DataTypeTok{"items"}\FunctionTok{:} \FunctionTok{\{}
  \DataTypeTok{"alice\_chair"}\FunctionTok{:} \FunctionTok{\{}\DataTypeTok{"type"}\FunctionTok{:} \StringTok{"chair"}\FunctionTok{,} \DataTypeTok{"x"}\FunctionTok{:} \FloatTok{{-}2.0}\FunctionTok{\}}
\FunctionTok{\}\}}
\end{Highlighting}
\end{Shaded}

In \texttt{spawn\_prop}, the registry name is given as \texttt{prop}:

\begin{Shaded}
\begin{Highlighting}[]
\FunctionTok{\{}\DataTypeTok{"action"}\FunctionTok{:} \StringTok{"spawn\_prop"}\FunctionTok{,} \DataTypeTok{"prop"}\FunctionTok{:} \StringTok{"alice\_chair"}\FunctionTok{,}
 \DataTypeTok{"type"}\FunctionTok{:} \StringTok{"chair"}\FunctionTok{,} \DataTypeTok{"x"}\FunctionTok{:} \FloatTok{{-}2.0}\FunctionTok{\}}
\end{Highlighting}
\end{Shaded}

\textbf{Gotcha:} if \texttt{spawn\_prop} omits \texttt{type},
\texttt{pam\_player.py} (line 1732 area) falls back to \texttt{"hat"}.
Always include \texttt{type} even when you've also set it in the
manifest. See §5.4.

\subsection{\texorpdfstring{1.2 Common positional: \texttt{y},
\texttt{color}, \texttt{accent},
\texttt{label}}{1.2 Common positional: y, color, accent, label}}\label{common-positional-y-color-accent-label}

\begin{longtable}[]{@{}
  >{\raggedright\arraybackslash}p{(\columnwidth - 6\tabcolsep) * \real{0.2500}}
  >{\raggedright\arraybackslash}p{(\columnwidth - 6\tabcolsep) * \real{0.2500}}
  >{\raggedright\arraybackslash}p{(\columnwidth - 6\tabcolsep) * \real{0.2500}}
  >{\raggedright\arraybackslash}p{(\columnwidth - 6\tabcolsep) * \real{0.2500}}@{}}
\toprule\noalign{}
\begin{minipage}[b]{\linewidth}\raggedright
Field
\end{minipage} & \begin{minipage}[b]{\linewidth}\raggedright
Type
\end{minipage} & \begin{minipage}[b]{\linewidth}\raggedright
Default
\end{minipage} & \begin{minipage}[b]{\linewidth}\raggedright
Notes
\end{minipage} \\
\midrule\noalign{}
\endhead
\bottomrule\noalign{}
\endlastfoot
\texttt{y} & float & varies by type & World-space vertical. Furniture
defaults to floor (-2.5ish); carried/wall props need explicit
\texttt{y}. \\
\texttt{color} & hex string & per-type & Primary fill / stroke color.
Builders style accents automatically. \\
\texttt{accent} & hex string & per-type & Secondary color (dodecahedron
faces, monitor frames). \\
\texttt{label} & string & none & Optional text label drawn on the prop.
Used heavily for ``Elevator'', ``PolymerCorp'', etc. \\
\end{longtable}

\begin{Shaded}
\begin{Highlighting}[]
\FunctionTok{\{}\DataTypeTok{"alice\_chair"}\FunctionTok{:} \FunctionTok{\{}\DataTypeTok{"type"}\FunctionTok{:} \StringTok{"chair"}\FunctionTok{,} \DataTypeTok{"x"}\FunctionTok{:} \FloatTok{{-}2.0}\FunctionTok{,}
                 \DataTypeTok{"color"}\FunctionTok{:} \StringTok{"\#ff9999"}\FunctionTok{,} \DataTypeTok{"label"}\FunctionTok{:} \StringTok{"A"}\FunctionTok{\}\}}
\end{Highlighting}
\end{Shaded}

\subsection{\texorpdfstring{1.3 Style variants: \texttt{style},
\texttt{size}, \texttt{animate},
\texttt{closed}}{1.3 Style variants: style, size, animate, closed}}\label{style-variants-style-size-animate-closed}

A handful of builders take additional shape/state parameters.

\textbf{\texttt{style}} --- selects a builder variant. Most useful on
\texttt{phone}:

\begin{Shaded}
\begin{Highlighting}[]
\FunctionTok{\{}\DataTypeTok{"phone1"}\FunctionTok{:} \FunctionTok{\{}\DataTypeTok{"type"}\FunctionTok{:} \StringTok{"phone"}\FunctionTok{,} \DataTypeTok{"x"}\FunctionTok{:} \FloatTok{5.3}\FunctionTok{,} \DataTypeTok{"y"}\FunctionTok{:} \FloatTok{{-}1.7}\FunctionTok{,}
            \DataTypeTok{"style"}\FunctionTok{:} \StringTok{"landline\_flat"}\FunctionTok{\}\}}
\end{Highlighting}
\end{Shaded}

\begin{longtable}[]{@{}
  >{\raggedright\arraybackslash}p{(\columnwidth - 2\tabcolsep) * \real{0.5000}}
  >{\raggedright\arraybackslash}p{(\columnwidth - 2\tabcolsep) * \real{0.5000}}@{}}
\toprule\noalign{}
\begin{minipage}[b]{\linewidth}\raggedright
Builder
\end{minipage} & \begin{minipage}[b]{\linewidth}\raggedright
Allowed \texttt{style} values
\end{minipage} \\
\midrule\noalign{}
\endhead
\bottomrule\noalign{}
\endlastfoot
\texttt{phone} & \texttt{cellphone} (default), \texttt{landline\_flat},
\texttt{landline\_wall} \\
\texttt{cell\_phone} / \texttt{smartphone} & aliases for \texttt{phone}
with \texttt{style="cellphone"} \\
\end{longtable}

\textbf{\texttt{size}} --- used by \texttt{floral\_arrangement}:

\begin{longtable}[]{@{}ll@{}}
\toprule\noalign{}
Value & Use \\
\midrule\noalign{}
\endhead
\bottomrule\noalign{}
\endlastfoot
\texttt{"carry"} (default) & Small bouquet held at chest height \\
\texttt{"large"} & Set-down arrangement on desk/table surface \\
\end{longtable}

\begin{Shaded}
\begin{Highlighting}[]
\FunctionTok{\{}\DataTypeTok{"bouquet"}\FunctionTok{:} \FunctionTok{\{}\DataTypeTok{"type"}\FunctionTok{:} \StringTok{"floral\_arrangement"}\FunctionTok{,} \DataTypeTok{"size"}\FunctionTok{:} \StringTok{"carry"}\FunctionTok{,}
             \DataTypeTok{"parent"}\FunctionTok{:} \StringTok{"chava"}\FunctionTok{,} \DataTypeTok{"attach"}\FunctionTok{:} \StringTok{"rhand"}\FunctionTok{\}\}}
\end{Highlighting}
\end{Shaded}

\textbf{\texttt{animate}} --- currently used only by
\texttt{dodecahedron}:

\begin{longtable}[]{@{}ll@{}}
\toprule\noalign{}
Value & Effect \\
\midrule\noalign{}
\endhead
\bottomrule\noalign{}
\endlastfoot
(omitted) & Static \\
\texttt{"spin"} & Continuous rotation \\
\end{longtable}

\textbf{\texttt{closed}} --- added in v0.9.13 to \texttt{build\_laptop}:

\begin{longtable}[]{@{}ll@{}}
\toprule\noalign{}
Value & Effect \\
\midrule\noalign{}
\endhead
\bottomrule\noalign{}
\endlastfoot
\texttt{false} (default) & Lid open, screen visible \\
\texttt{true} & Lid closed flat (storage / pre-reveal state) \\
\end{longtable}

\subsection{\texorpdfstring{1.4 Scene-graph: \texttt{parent},
\texttt{attach},
\texttt{prop\_registry}}{1.4 Scene-graph: parent, attach, prop\_registry}}\label{scene-graph-parent-attach-prop_registry}

Any prop can be positioned relative to another prop or a character joint
instead of in world space. Three fields together do this:

\begin{longtable}[]{@{}
  >{\raggedright\arraybackslash}p{(\columnwidth - 2\tabcolsep) * \real{0.5000}}
  >{\raggedright\arraybackslash}p{(\columnwidth - 2\tabcolsep) * \real{0.5000}}@{}}
\toprule\noalign{}
\begin{minipage}[b]{\linewidth}\raggedright
Field
\end{minipage} & \begin{minipage}[b]{\linewidth}\raggedright
Meaning
\end{minipage} \\
\midrule\noalign{}
\endhead
\bottomrule\noalign{}
\endlastfoot
\texttt{parent} & Registry name of the parent prop (or character). \\
\texttt{attach} & Named attachment point on the parent (see §2.4). \\
\texttt{prop\_registry} & The live registry dict --- required only when
building a child prop programmatically. \\
\end{longtable}

\begin{Shaded}
\begin{Highlighting}[]
\FunctionTok{\{}\DataTypeTok{"action"}\FunctionTok{:} \StringTok{"spawn\_prop"}\FunctionTok{,} \DataTypeTok{"prop"}\FunctionTok{:} \StringTok{"monitor"}\FunctionTok{,}
 \DataTypeTok{"type"}\FunctionTok{:} \StringTok{"dodecahedron"}\FunctionTok{,} \DataTypeTok{"radius"}\FunctionTok{:} \DecValTok{0}\ErrorTok{.}\DecValTok{12}\FunctionTok{,}
 \DataTypeTok{"parent"}\FunctionTok{:} \StringTok{"main\_desk"}\FunctionTok{,} \DataTypeTok{"attach"}\FunctionTok{:} \StringTok{"surface"}\FunctionTok{,}
 \DataTypeTok{"attrs"}\FunctionTok{:} \FunctionTok{\{}\DataTypeTok{"inclination"}\FunctionTok{:} \DecValTok{10}\FunctionTok{\}\}}
\end{Highlighting}
\end{Shaded}

A child prop snaps its origin to the parent's named attachment point;
any explicit \texttt{x}/\texttt{y} is then applied as a local offset.

For character-parent attachment (e.g.~carrying), see §4.2 and §4.6.

\subsection{\texorpdfstring{1.5 Visual attrs dict: \texttt{scale},
\texttt{inclination}}{1.5 Visual attrs dict: scale, inclination}}\label{visual-attrs-dict-scale-inclination}

The \texttt{attrs} dict carries scale and rotation overrides that apply
after the prop is built but before it's placed:

\begin{Shaded}
\begin{Highlighting}[]
\FunctionTok{\{}\DataTypeTok{"mon1"}\FunctionTok{:} \FunctionTok{\{}\DataTypeTok{"type"}\FunctionTok{:} \StringTok{"dodecahedron"}\FunctionTok{,} \DataTypeTok{"radius"}\FunctionTok{:} \DecValTok{0}\ErrorTok{.}\DecValTok{12}\FunctionTok{,}
          \DataTypeTok{"attrs"}\FunctionTok{:} \FunctionTok{\{}\DataTypeTok{"scale"}\FunctionTok{:} \DecValTok{0}\ErrorTok{.}\DecValTok{8}\FunctionTok{,} \DataTypeTok{"inclination"}\FunctionTok{:} \DecValTok{15}\FunctionTok{\}\}\}}
\end{Highlighting}
\end{Shaded}

\begin{longtable}[]{@{}llll@{}}
\toprule\noalign{}
Key & Type & Units & Default \\
\midrule\noalign{}
\endhead
\bottomrule\noalign{}
\endlastfoot
\texttt{scale} & float & multiplier & 1.0 \\
\texttt{inclination} & float & degrees & 0 \\
\end{longtable}

\subsection{1.6 Per-type parameter
tables}\label{per-type-parameter-tables}

A summary of the type-specific extras most often used. See §6 for the
full prop catalog and §1.4 of \texttt{CHARACTER\_PROPERTIES.md} analog
notes on parameter discoverability.

\textbf{\texttt{desk}}

\begin{longtable}[]{@{}llll@{}}
\toprule\noalign{}
Param & Type & Default & Notes \\
\midrule\noalign{}
\endhead
\bottomrule\noalign{}
\endlastfoot
\texttt{width} & float & 3.0 & Desk top width \\
\texttt{monitor} & bool & false & Adds a small dodecahedron monitor \\
\texttt{monitor\_color} & hex & accent & Color of the added monitor \\
\end{longtable}

\textbf{\texttt{dodecahedron}}

\begin{longtable}[]{@{}lll@{}}
\toprule\noalign{}
Param & Default & Notes \\
\midrule\noalign{}
\endhead
\bottomrule\noalign{}
\endlastfoot
\texttt{radius} & 0.4 & Geometric size \\
\texttt{accent} & per-type & Face accent color \\
\texttt{animate} & (none) & \texttt{"spin"} for rotation \\
\end{longtable}

\textbf{\texttt{tv\_monitor}}

\begin{longtable}[]{@{}lll@{}}
\toprule\noalign{}
Param & Default & Notes \\
\midrule\noalign{}
\endhead
\bottomrule\noalign{}
\endlastfoot
\texttt{screen\_text} & ``\,'' & On-screen text \\
\texttt{screen\_color} & dark & Background fill \\
\texttt{width} & 2.2 & Frame width \\
\end{longtable}

\textbf{\texttt{avatar\_pod}}

\begin{longtable}[]{@{}lll@{}}
\toprule\noalign{}
Param & Default & Notes \\
\midrule\noalign{}
\endhead
\bottomrule\noalign{}
\endlastfoot
\texttt{occupied} & false & Use occupied color palette \\
\end{longtable}

\textbf{\texttt{elevator}}

\begin{longtable}[]{@{}lll@{}}
\toprule\noalign{}
Param & Default & Notes \\
\midrule\noalign{}
\endhead
\bottomrule\noalign{}
\endlastfoot
\texttt{label} & ``Elevator'' & Sign text \\
\texttt{capacity\_label} & ``Max Capacity: 4'' & Secondary sign \\
\end{longtable}

\textbf{\texttt{backpack}} (v0.9.7)

\begin{longtable}[]{@{}lll@{}}
\toprule\noalign{}
Param & Default & Notes \\
\midrule\noalign{}
\endhead
\bottomrule\noalign{}
\endlastfoot
\texttt{accent} & teal & Edge color of the Fano-plane graph \\
\end{longtable}

\begin{center}\rule{0.5\linewidth}{0.5pt}\end{center}

\section{2. Prop metadata}\label{prop-metadata}

Every prop VGroup carries metadata in three forms, all set by the
\texttt{\_attach\_pam\_attrs} helper in \texttt{props\_core.py}.

\subsection{2.1 Flat attributes
(backward-compat)}\label{flat-attributes-backward-compat}

The original metadata layer, preserved so older code in
\texttt{pam\_player.py} keeps working unchanged.

\begin{longtable}[]{@{}
  >{\raggedright\arraybackslash}p{(\columnwidth - 4\tabcolsep) * \real{0.3333}}
  >{\raggedright\arraybackslash}p{(\columnwidth - 4\tabcolsep) * \real{0.3333}}
  >{\raggedright\arraybackslash}p{(\columnwidth - 4\tabcolsep) * \real{0.3333}}@{}}
\toprule\noalign{}
\begin{minipage}[b]{\linewidth}\raggedright
Attribute
\end{minipage} & \begin{minipage}[b]{\linewidth}\raggedright
Type
\end{minipage} & \begin{minipage}[b]{\linewidth}\raggedright
Meaning
\end{minipage} \\
\midrule\noalign{}
\endhead
\bottomrule\noalign{}
\endlastfoot
\texttt{.pam\_name} & str & Registry name (unique key used in screenplay
actions) \\
\texttt{.pam\_type} & str & Type string (\texttt{"chair"},
\texttt{"desk"}, etc.) \\
\texttt{.pam\_x} & float & World x-coordinate (center of the prop) \\
\texttt{.pam\_y} & float & World y-coordinate \\
\texttt{.pam\_surface\_y} & float & Y of the top surface (desk top,
chair seat, etc.) \\
\end{longtable}

\begin{Shaded}
\begin{Highlighting}[]
\NormalTok{prop }\OperatorTok{=}\NormalTok{ registry[}\StringTok{"main\_desk"}\NormalTok{]}
\BuiltInTok{print}\NormalTok{(prop.pam\_x, prop.pam\_y, prop.pam\_surface\_y)}
\end{Highlighting}
\end{Shaded}

Action handlers in \texttt{pam\_player.py} rely on these heavily ---
anything that needs to know where a prop is reads \texttt{.pam\_x} /
\texttt{.pam\_y} directly.

\subsection{\texorpdfstring{2.2 The \texttt{.pam\_node} scene-graph
dict}{2.2 The .pam\_node scene-graph dict}}\label{the-.pam_node-scene-graph-dict}

Added in the scene-graph pass. Carries the relational metadata.

\begin{Shaded}
\begin{Highlighting}[]
\NormalTok{prop.pam\_node }\OperatorTok{=}\NormalTok{ \{}
    \StringTok{"name"}\NormalTok{:   }\StringTok{"monitor"}\NormalTok{,}
    \StringTok{"kind"}\NormalTok{:   }\StringTok{"dodecahedron"}\NormalTok{,}
    \StringTok{"parent"}\NormalTok{: }\StringTok{"main\_desk"}\NormalTok{,     }\CommentTok{\# or None for world coords}
    \StringTok{"attach"}\NormalTok{: }\StringTok{"surface"}\NormalTok{,       }\CommentTok{\# name of attachment point on parent}
    \StringTok{"attrs"}\NormalTok{:  \{}\StringTok{"scale"}\NormalTok{: }\FloatTok{1.0}\NormalTok{, }\StringTok{"inclination"}\NormalTok{: }\DecValTok{10}\NormalTok{\},}
    \StringTok{"x"}\NormalTok{:      }\FloatTok{0.0}\NormalTok{,}
    \StringTok{"y"}\NormalTok{:      }\FloatTok{0.06}\NormalTok{,}
\NormalTok{\}}
\end{Highlighting}
\end{Shaded}

This is the canonical place to read parent/attach relationships. The
flat attrs (§2.1) are the legacy mirror --- \texttt{pam\_node} is
authoritative when they disagree.

\subsection{\texorpdfstring{2.3 The \texttt{.pam\_attachments}
dict}{2.3 The .pam\_attachments dict}}\label{the-.pam_attachments-dict}

A dict of named 3-D positions (numpy arrays) on the prop geometry where
child props or characters may connect.

\begin{Shaded}
\begin{Highlighting}[]
\NormalTok{desk.pam\_attachments}
\CommentTok{\# \{}
\CommentTok{\#   "surface":    [0.0, {-}0.4, 0.0],}
\CommentTok{\#   "left{-}edge":  [{-}1.5, {-}0.4, 0.0],}
\CommentTok{\#   "right{-}edge": [ 1.5, {-}0.4, 0.0],}
\CommentTok{\#   "floor":      [0.0, {-}2.5, 0.0],}
\CommentTok{\# \}}
\end{Highlighting}
\end{Shaded}

Used by \texttt{resolve\_position(node,\ registry)} to compute the world
position of any child prop given its parent and attach point.

\subsection{2.4 Per-type attachment-point
catalog}\label{per-type-attachment-point-catalog}

A reference for which attach points each prop kind exposes.

\begin{longtable}[]{@{}
  >{\raggedright\arraybackslash}p{(\columnwidth - 2\tabcolsep) * \real{0.5000}}
  >{\raggedright\arraybackslash}p{(\columnwidth - 2\tabcolsep) * \real{0.5000}}@{}}
\toprule\noalign{}
\begin{minipage}[b]{\linewidth}\raggedright
Prop type
\end{minipage} & \begin{minipage}[b]{\linewidth}\raggedright
Attachment points
\end{minipage} \\
\midrule\noalign{}
\endhead
\bottomrule\noalign{}
\endlastfoot
\texttt{chair} & \texttt{surface} (seat), \texttt{back-top},
\texttt{floor}, \texttt{left-edge}, \texttt{right-edge} \\
\texttt{desk} & \texttt{surface}, \texttt{left-edge},
\texttt{right-edge}, \texttt{floor} \\
\texttt{door} / \texttt{pocket\_door} & \texttt{surface},
\texttt{handle}, \texttt{threshold}, \texttt{floor}, \texttt{left-edge},
\texttt{right-edge} \\
\texttt{elevator} & \texttt{surface}, \texttt{floor},
\texttt{threshold}, \texttt{centre}, \texttt{left-edge},
\texttt{right-edge}, \texttt{button\_panel} \\
\texttt{avatar\_pod} & \texttt{surface} (top of lid when closed),
\texttt{floor}, \texttt{centre}, \texttt{lid\_hinge} \\
\texttt{dodecahedron} & \texttt{surface} (top of polygon),
\texttt{centre}, \texttt{floor} \\
\texttt{tv\_monitor} & \texttt{surface} (wall-mounted top),
\texttt{centre}, \texttt{screen\_center} \\
\texttt{desk\_lamp} & \texttt{bulb}, \texttt{base}, \texttt{arm\_top}
(bug target) \\
\texttt{phone} (cell) & \texttt{screen}, \texttt{top} \\
\texttt{phone} (landline\_flat) & \texttt{cradle}, \texttt{receiver},
\texttt{dial\_pad} \\
\texttt{backpack} & \texttt{top\_strap}, \texttt{bottom},
\texttt{interior} \\
\texttt{flower} / \texttt{floral\_arrangement} & \texttt{stem\_base},
\texttt{bloom\_center} \\
\end{longtable}

A child prop with \texttt{attach} set to a name not in this dict raises
a \texttt{KeyError} at build time; check the builder source if
uncertain.

\begin{center}\rule{0.5\linewidth}{0.5pt}\end{center}

\section{3. Registry behavior}\label{registry-behavior-1}

The prop registry is a dict-like object on the player, maintaining which
props are currently live in the scene and where to find them.

\subsection{3.1 The prop registry}\label{the-prop-registry-1}

\begin{longtable}[]{@{}
  >{\raggedright\arraybackslash}p{(\columnwidth - 4\tabcolsep) * \real{0.3333}}
  >{\raggedright\arraybackslash}p{(\columnwidth - 4\tabcolsep) * \real{0.3333}}
  >{\raggedright\arraybackslash}p{(\columnwidth - 4\tabcolsep) * \real{0.3333}}@{}}
\toprule\noalign{}
\begin{minipage}[b]{\linewidth}\raggedright
Method
\end{minipage} & \begin{minipage}[b]{\linewidth}\raggedright
Returns
\end{minipage} & \begin{minipage}[b]{\linewidth}\raggedright
Notes
\end{minipage} \\
\midrule\noalign{}
\endhead
\bottomrule\noalign{}
\endlastfoot
\texttt{registry.get(name)} & the prop VGroup, or \texttt{None} & Safe
lookup \\
\texttt{registry.get\_raw(name)} & the prop VGroup or raises & Strict
--- for code paths that must have the prop \\
\texttt{registry.add(name,\ prop)} & None & Used by \texttt{spawn\_prop}
internals \\
\texttt{registry.world\_pos(name,\ attach)} & numpy array & Resolved
world position of a named attachment point \\
\end{longtable}

Most action handlers in \texttt{pam\_player.py} go through \texttt{get}
and fall through quietly if a prop isn't found --- see §3.3 for why this
matters.

\subsection{\texorpdfstring{3.2 \texttt{spawn\_prop} --- mid-scene
appearance}{3.2 spawn\_prop --- mid-scene appearance}}\label{spawn_prop-mid-scene-appearance}

Adds a new prop to the registry and fades it in.

\begin{Shaded}
\begin{Highlighting}[]
\FunctionTok{\{}\DataTypeTok{"action"}\FunctionTok{:} \StringTok{"spawn\_prop"}\FunctionTok{,} \DataTypeTok{"prop"}\FunctionTok{:} \StringTok{"bouquet"}\FunctionTok{,}
 \DataTypeTok{"type"}\FunctionTok{:} \StringTok{"floral\_arrangement"}\FunctionTok{,} \DataTypeTok{"size"}\FunctionTok{:} \StringTok{"carry"}\FunctionTok{,}
 \DataTypeTok{"parent"}\FunctionTok{:} \StringTok{"chava"}\FunctionTok{,} \DataTypeTok{"attach"}\FunctionTok{:} \StringTok{"rhand"}\FunctionTok{,}
 \DataTypeTok{"color"}\FunctionTok{:} \StringTok{"\#ffcc88"}\FunctionTok{\}}
\end{Highlighting}
\end{Shaded}

Required fields: \texttt{prop} (registry name), \texttt{type}. All other
fields match the manifest \texttt{props} block (§1).

\textbf{Authoring tip:} \texttt{spawn\_prop} is the right action
whenever the prop's existence is itself a story beat (a phone appears in
a pocket, a folder materializes on the desk). For props that are simply
present from the beginning of the scene, use the manifest \texttt{props}
block.

\subsection{\texorpdfstring{3.3 \texttt{remove\_prop} --- despawn and
symmetric-clear}{3.3 remove\_prop --- despawn and symmetric-clear}}\label{remove_prop-despawn-and-symmetric-clear}

Removes a prop from the registry and fades it out.

\begin{Shaded}
\begin{Highlighting}[]
\FunctionTok{\{}\DataTypeTok{"action"}\FunctionTok{:} \StringTok{"remove\_prop"}\FunctionTok{,} \DataTypeTok{"prop"}\FunctionTok{:} \StringTok{"bouquet"}\FunctionTok{\}}
\end{Highlighting}
\end{Shaded}

\textbf{Symmetric-clear gotcha:} if a prop has been
\texttt{parent}-attached as a child of another prop (or character),
\texttt{remove\_prop} clears it from the prop registry but does
\textbf{not} automatically clear the parent's attachment-point
reservation. This can leave a stale entry that affects later
\texttt{spawn\_prop} calls reusing the same attach point. v1.0.0 will
add a \texttt{clear\_attachments=True} helper; until then, the safest
pattern is:

\begin{Shaded}
\begin{Highlighting}[]
\FunctionTok{\{}\DataTypeTok{"action"}\FunctionTok{:} \StringTok{"remove\_prop"}\FunctionTok{,} \DataTypeTok{"prop"}\FunctionTok{:} \StringTok{"old\_monitor"}\FunctionTok{\}}\ErrorTok{,}
\FunctionTok{\{}\DataTypeTok{"action"}\FunctionTok{:} \StringTok{"spawn\_prop"}\FunctionTok{,} \DataTypeTok{"prop"}\FunctionTok{:} \StringTok{"new\_monitor"}\FunctionTok{,}
 \DataTypeTok{"type"}\FunctionTok{:} \StringTok{"dodecahedron"}\FunctionTok{,}
 \DataTypeTok{"parent"}\FunctionTok{:} \StringTok{"main\_desk"}\FunctionTok{,} \DataTypeTok{"attach"}\FunctionTok{:} \StringTok{"surface"}\FunctionTok{\}}
\end{Highlighting}
\end{Shaded}

\ldots where the new spawn explicitly overwrites the attach.

\subsection{\texorpdfstring{3.4 The \texttt{hidden}
flag}{3.4 The hidden flag}}\label{the-hidden-flag-1}

A prop registered with \texttt{hidden:\ true} is added to the registry
at full opacity 0 --- present but invisible. Used for pre-positioning
props that will appear later via \texttt{spawn\_prop} or a custom reveal
animation.

\begin{Shaded}
\begin{Highlighting}[]
\FunctionTok{\{}\DataTypeTok{"laptop"}\FunctionTok{:} \FunctionTok{\{}\DataTypeTok{"type"}\FunctionTok{:} \StringTok{"laptop"}\FunctionTok{,} \DataTypeTok{"x"}\FunctionTok{:} \FloatTok{{-}1.5}\FunctionTok{,} \DataTypeTok{"y"}\FunctionTok{:} \DecValTok{0}\ErrorTok{.}\DecValTok{8}\FunctionTok{,}
            \DataTypeTok{"closed"}\FunctionTok{:} \KeywordTok{true}\FunctionTok{,} \DataTypeTok{"hidden"}\FunctionTok{:} \KeywordTok{true}\FunctionTok{\}\}}
\end{Highlighting}
\end{Shaded}

Later in the scene:

\begin{Shaded}
\begin{Highlighting}[]
\FunctionTok{\{}\DataTypeTok{"action"}\FunctionTok{:} \StringTok{"reveal\_laptop"}\FunctionTok{,} \DataTypeTok{"prop"}\FunctionTok{:} \StringTok{"laptop"}\FunctionTok{,} \DataTypeTok{"from"}\FunctionTok{:} \StringTok{"backpack"}\FunctionTok{\}}
\end{Highlighting}
\end{Shaded}

\textbf{Scope note:} \texttt{hidden} is a \emph{prop} concept. An early
v1.0.0 draft of \texttt{CHARACTER\_PROPERTIES.md} referenced a
character-side \texttt{hidden} flag, but that implementation did not
land --- characters use \texttt{fade\_in} / \texttt{fade\_out} and the
cast-block manifest controls visibility. Props are the only entities
where \texttt{hidden:\ true} is meaningful at construction.

\subsection{\texorpdfstring{3.5 Manifest \texttt{props} block vs
\texttt{spawn\_prop}}{3.5 Manifest props block vs spawn\_prop}}\label{manifest-props-block-vs-spawn_prop}

\begin{longtable}[]{@{}
  >{\raggedright\arraybackslash}p{(\columnwidth - 4\tabcolsep) * \real{0.3333}}
  >{\raggedright\arraybackslash}p{(\columnwidth - 4\tabcolsep) * \real{0.3333}}
  >{\raggedright\arraybackslash}p{(\columnwidth - 4\tabcolsep) * \real{0.3333}}@{}}
\toprule\noalign{}
\begin{minipage}[b]{\linewidth}\raggedright
Question
\end{minipage} & \begin{minipage}[b]{\linewidth}\raggedright
Manifest \texttt{props}
\end{minipage} & \begin{minipage}[b]{\linewidth}\raggedright
\texttt{spawn\_prop}
\end{minipage} \\
\midrule\noalign{}
\endhead
\bottomrule\noalign{}
\endlastfoot
When does it run? & At scene start, before any other action & At its
position in the action sequence \\
Visible immediately? & Yes (opacity 1) unless \texttt{hidden:\ true} &
Fades in with \texttt{rt} (default 0.4s) \\
Use for? & Furniture and props present from the start & Story-driven
appearances \\
Audit-friendly? & Yes --- single block, easy to scan & Distributed
through scene \\
\end{longtable}

A useful convention: list every prop in the manifest with appropriate
\texttt{hidden:\ true} flags, then use \texttt{spawn\_prop} only for
animation timing of reveals. This keeps the at-a-glance prop manifest
complete.

\begin{center}\rule{0.5\linewidth}{0.5pt}\end{center}

\section{4. Action targets}\label{action-targets-1}

The set of JSON actions that operate on props. Grouped by interaction
type. Every action here resolves its target via the prop registry
(§3.1).

\subsection{4.1 Pointing and reaching}\label{pointing-and-reaching}

\textbf{\texttt{reach\_for}} --- extend one arm toward a prop, hold,
then retract.

\begin{Shaded}
\begin{Highlighting}[]
\FunctionTok{\{}\DataTypeTok{"action"}\FunctionTok{:} \StringTok{"reach\_for"}\FunctionTok{,} \DataTypeTok{"who"}\FunctionTok{:} \StringTok{"sidel"}\FunctionTok{,} \DataTypeTok{"target"}\FunctionTok{:} \StringTok{"button\_panel"}\FunctionTok{,}
 \DataTypeTok{"arm"}\FunctionTok{:} \StringTok{"r"}\FunctionTok{,} \DataTypeTok{"hold"}\FunctionTok{:} \DecValTok{0}\ErrorTok{.}\DecValTok{6}\FunctionTok{\}}
\end{Highlighting}
\end{Shaded}

\begin{longtable}[]{@{}llll@{}}
\toprule\noalign{}
Key & Type & Default & Notes \\
\midrule\noalign{}
\endhead
\bottomrule\noalign{}
\endlastfoot
\texttt{who} & str & required & Character key \\
\texttt{target} & str & required & Prop registry name \\
\texttt{arm} & \texttt{"r"} / \texttt{"l"} & \texttt{"r"} & Which arm to
extend \\
\texttt{hold} & float & 0.4 & Seconds at full extension \\
\texttt{offset\_x} & float & 0 & Adjust reach landing point \\
\texttt{offset\_y} & float & 0 & Adjust reach landing point \\
\end{longtable}

\textbf{\texttt{punch\_button}} --- sharp jab and retract.

\begin{Shaded}
\begin{Highlighting}[]
\FunctionTok{\{}\DataTypeTok{"action"}\FunctionTok{:} \StringTok{"punch\_button"}\FunctionTok{,} \DataTypeTok{"who"}\FunctionTok{:} \StringTok{"sidel"}\FunctionTok{,} \DataTypeTok{"target"}\FunctionTok{:} \StringTok{"button\_panel"}\FunctionTok{\}}
\end{Highlighting}
\end{Shaded}

Same key set as \texttt{reach\_for} but with snappier timing (shorter
\texttt{hold}, faster extend/retract). v0.9.5 patch added
\texttt{offset\_x}/\texttt{offset\_y} for fine targeting within a
multi-button panel; an earlier bug where the arm failed to retract was
fixed by deep-copying the rest pose before the jab so \texttt{morph\_to}
couldn't overwrite it.

\textbf{\texttt{peel\_from\_hand}} (v0.9.5) --- open-palm gesture that
slides a tiny prop off the palm, leaving the prop in
\texttt{fig.\_held\_prop}, ready for \texttt{stick\_to}.

\begin{Shaded}
\begin{Highlighting}[]
\FunctionTok{\{}\DataTypeTok{"action"}\FunctionTok{:} \StringTok{"peel\_from\_hand"}\FunctionTok{,} \DataTypeTok{"who"}\FunctionTok{:} \StringTok{"sidel"}\FunctionTok{,} \DataTypeTok{"prop"}\FunctionTok{:} \StringTok{"bug"}\FunctionTok{,} \DataTypeTok{"arm"}\FunctionTok{:} \StringTok{"r"}\FunctionTok{\}}
\end{Highlighting}
\end{Shaded}

Best paired with \texttt{FRAMING=insert} --- the prop is tiny in wide
shots.

\subsection{4.2 Carrying}\label{carrying}

\textbf{\texttt{grab}} --- decisive reach + retract with the prop now in
the character's hand. Faster than \texttt{pick\_up}.

\begin{Shaded}
\begin{Highlighting}[]
\FunctionTok{\{}\DataTypeTok{"action"}\FunctionTok{:} \StringTok{"grab"}\FunctionTok{,} \DataTypeTok{"who"}\FunctionTok{:} \StringTok{"nona"}\FunctionTok{,} \DataTypeTok{"prop"}\FunctionTok{:} \StringTok{"folder"}\FunctionTok{,} \DataTypeTok{"arm"}\FunctionTok{:} \StringTok{"r"}\FunctionTok{\}}
\end{Highlighting}
\end{Shaded}

Sets \texttt{fig.\_held\_prop\ =\ prop} and reparents the prop to the
relevant hand joint. Subsequent \texttt{walk\_to} actions will drag the
held prop along (see §4.6).

\textbf{\texttt{place\_on}} --- set a carried prop onto a target
surface.

\begin{Shaded}
\begin{Highlighting}[]
\FunctionTok{\{}\DataTypeTok{"action"}\FunctionTok{:} \StringTok{"place\_on"}\FunctionTok{,} \DataTypeTok{"who"}\FunctionTok{:} \StringTok{"nona"}\FunctionTok{,}
 \DataTypeTok{"prop"}\FunctionTok{:} \StringTok{"bouquet"}\FunctionTok{,} \DataTypeTok{"target"}\FunctionTok{:} \StringTok{"main\_desk"}\FunctionTok{\}}
\end{Highlighting}
\end{Shaded}

Reparents the prop to the target's \texttt{surface} attachment point and
clears \texttt{fig.\_held\_prop}.

\textbf{\texttt{move\_aside}} --- push a prop laterally without picking
it up.

\begin{Shaded}
\begin{Highlighting}[]
\FunctionTok{\{}\DataTypeTok{"action"}\FunctionTok{:} \StringTok{"move\_aside"}\FunctionTok{,} \DataTypeTok{"who"}\FunctionTok{:} \StringTok{"sidel"}\FunctionTok{,}
 \DataTypeTok{"prop"}\FunctionTok{:} \StringTok{"desk\_lamp"}\FunctionTok{,} \DataTypeTok{"direction"}\FunctionTok{:} \StringTok{"left"}\FunctionTok{,} \DataTypeTok{"distance"}\FunctionTok{:} \DecValTok{0}\ErrorTok{.}\DecValTok{4}\FunctionTok{\}}
\end{Highlighting}
\end{Shaded}

\begin{longtable}[]{@{}lll@{}}
\toprule\noalign{}
Key & Type & Default \\
\midrule\noalign{}
\endhead
\bottomrule\noalign{}
\endlastfoot
\texttt{direction} & \texttt{"left"} / \texttt{"right"} & required \\
\texttt{distance} & float & 0.3 \\
\texttt{arm} & \texttt{"r"} / \texttt{"l"} & inferred from
\texttt{direction} \\
\end{longtable}

\textbf{\texttt{stick\_to}} --- attach a tiny prop to a target prop's
surface.

\begin{Shaded}
\begin{Highlighting}[]
\FunctionTok{\{}\DataTypeTok{"action"}\FunctionTok{:} \StringTok{"stick\_to"}\FunctionTok{,} \DataTypeTok{"who"}\FunctionTok{:} \StringTok{"sidel"}\FunctionTok{,}
 \DataTypeTok{"prop"}\FunctionTok{:} \StringTok{"bug"}\FunctionTok{,} \DataTypeTok{"target"}\FunctionTok{:} \StringTok{"desk\_lamp"}\FunctionTok{\}}
\end{Highlighting}
\end{Shaded}

Reparents the prop and consumes \texttt{fig.\_held\_prop}.

\textbf{\texttt{snap\_photo}} --- point a phone at a target and flash.

\begin{Shaded}
\begin{Highlighting}[]
\FunctionTok{\{}\DataTypeTok{"action"}\FunctionTok{:} \StringTok{"snap\_photo"}\FunctionTok{,} \DataTypeTok{"who"}\FunctionTok{:} \StringTok{"nona"}\FunctionTok{,} \DataTypeTok{"target"}\FunctionTok{:} \StringTok{"governor"}\FunctionTok{\}}
\end{Highlighting}
\end{Shaded}

Combines \texttt{reach\_for} + \texttt{flash} (§4.4) with phone-specific
framing.

\subsection{4.3 Stateful props:
open/close}\label{stateful-props-openclose}

Some prop types carry their own animation methods, bound to the VGroup
at build time. Action handlers dispatch by method name.

\textbf{\texttt{open\_doors} / \texttt{close\_doors}} --- used by both
\texttt{pocket\_door} and \texttt{elevator}.

\begin{Shaded}
\begin{Highlighting}[]
\FunctionTok{\{}\DataTypeTok{"action"}\FunctionTok{:} \StringTok{"open\_doors"}\FunctionTok{,} \DataTypeTok{"prop"}\FunctionTok{:} \StringTok{"elev\_1"}\FunctionTok{\}}
\FunctionTok{\{}\DataTypeTok{"action"}\FunctionTok{:} \StringTok{"close\_doors"}\FunctionTok{,} \DataTypeTok{"prop"}\FunctionTok{:} \StringTok{"elev\_1"}\FunctionTok{,} \DataTypeTok{"run\_time"}\FunctionTok{:} \DecValTok{0}\ErrorTok{.}\DecValTok{6}\FunctionTok{\}}
\end{Highlighting}
\end{Shaded}

\begin{longtable}[]{@{}lll@{}}
\toprule\noalign{}
Key & Type & Default \\
\midrule\noalign{}
\endhead
\bottomrule\noalign{}
\endlastfoot
\texttt{prop} & str & required \\
\texttt{run\_time} & float & 0.6 \\
\end{longtable}

The prop's own \texttt{is\_open} flag prevents redundant animations.

\textbf{\texttt{open\_lid} / \texttt{close\_lid}} --- used by
\texttt{avatar\_pod}.

\begin{Shaded}
\begin{Highlighting}[]
\FunctionTok{\{}\DataTypeTok{"action"}\FunctionTok{:} \StringTok{"open\_lid"}\FunctionTok{,} \DataTypeTok{"prop"}\FunctionTok{:} \StringTok{"pod\_1"}\FunctionTok{\}}
\FunctionTok{\{}\DataTypeTok{"action"}\FunctionTok{:} \StringTok{"close\_lid"}\FunctionTok{,} \DataTypeTok{"prop"}\FunctionTok{:} \StringTok{"pod\_1"}\FunctionTok{\}}
\end{Highlighting}
\end{Shaded}

The lid hinges open along a configured axis and swaps in the
``occupied'' color palette when closed over a character.

\textbf{Generic dispatch.} Since \texttt{pocket\_door} and
\texttt{elevator} share method names but are different prop types, the
action handler dispatches by \texttt{getattr(prop,\ "open\_doors")}
rather than checking \texttt{pam\_type}. This means \textbf{any future
prop with an \texttt{.open\_doors} method automatically inherits the
action} --- which is both a feature and a small footgun (see §5.5).

\subsection{\texorpdfstring{4.4 \texttt{flash} --- temporary
recolor}{4.4 flash --- temporary recolor}}\label{flash-temporary-recolor}

Recolor a prop, character, or both for a brief duration, then restore.

\begin{Shaded}
\begin{Highlighting}[]
\FunctionTok{\{}\DataTypeTok{"action"}\FunctionTok{:} \StringTok{"flash"}\FunctionTok{,} \DataTypeTok{"prop"}\FunctionTok{:} \StringTok{"monitor"}\FunctionTok{,}
 \DataTypeTok{"color"}\FunctionTok{:} \StringTok{"\#44aaff"}\FunctionTok{,} \DataTypeTok{"duration"}\FunctionTok{:} \DecValTok{0}\ErrorTok{.}\DecValTok{3}\FunctionTok{\}}
\end{Highlighting}
\end{Shaded}

\begin{longtable}[]{@{}llll@{}}
\toprule\noalign{}
Key & Type & Default & Notes \\
\midrule\noalign{}
\endhead
\bottomrule\noalign{}
\endlastfoot
\texttt{prop} & str & (one of) & Target prop \\
\texttt{who} & str & (one of) & Target character \\
\texttt{color} & hex & \texttt{"\#44aaff"} & Flash color \\
\texttt{duration} & float & 0.3 & Total time in flashed state \\
\end{longtable}

At least one of \texttt{prop} or \texttt{who} is required; both may be
set to flash both simultaneously. Used for camera flashes
(\texttt{snap\_photo}), highlight beats, and avatar-transfer effects
(\texttt{highlight\_edges} was an earlier specialized variant;
\texttt{flash} superseded it).

\subsection{\texorpdfstring{4.5 \texttt{prop\_say} --- speech from
non-character
props}{4.5 prop\_say --- speech from non-character props}}\label{prop_say-speech-from-non-character-props}

Make a prop speak via a bubble anchored to the prop position.

\begin{Shaded}
\begin{Highlighting}[]
\FunctionTok{\{}\DataTypeTok{"action"}\FunctionTok{:} \StringTok{"prop\_say"}\FunctionTok{,} \DataTypeTok{"prop"}\FunctionTok{:} \StringTok{"phone1"}\FunctionTok{,}
 \DataTypeTok{"text"}\FunctionTok{:} \StringTok{"Hello, is anyone there?"}\FunctionTok{,}
 \DataTypeTok{"side"}\FunctionTok{:} \StringTok{"right"}\FunctionTok{,} \DataTypeTok{"hold"}\FunctionTok{:} \FloatTok{1.4}\FunctionTok{\}}
\end{Highlighting}
\end{Shaded}

\begin{longtable}[]{@{}lll@{}}
\toprule\noalign{}
Key & Type & Default \\
\midrule\noalign{}
\endhead
\bottomrule\noalign{}
\endlastfoot
\texttt{prop} & str & required \\
\texttt{text} & str & required \\
\texttt{side} & \texttt{"left"} / \texttt{"right"} & \texttt{"right"} \\
\texttt{hold} & float & 1.4 \\
\texttt{font\_size} & int & 18 \\
\texttt{rt\_in} & float & 0.35 \\
\texttt{rt\_out} & float & 0.25 \\
\texttt{persist} & bool & false \\
\texttt{duration} & float & none \\
\end{longtable}

The bubble appears above the prop with a tail pointing at it. Used
heavily for phones, TVs, and the GovernorGraph autonomous prop-char.

\textbf{Persistent variant.} Setting \texttt{persist:\ true} or a
\texttt{duration:\ t} leaves the bubble on screen until explicitly
cleared by \texttt{clear\_bubble} or \texttt{clear\_all\_bubbles}. Watch
for the focus\_reset gotcha --- see §5.2.

\subsection{\texorpdfstring{4.6 Walk-following via
\texttt{\_drag\_attached\_props}}{4.6 Walk-following via \_drag\_attached\_props}}\label{walk-following-via-_drag_attached_props}

When a character holds a prop (set by \texttt{grab}, \texttt{pick\_up},
or manifest-level \texttt{parent:\ \textless{}character\textgreater{}}),
\texttt{walk\_to} / \texttt{run\_to} / \texttt{rush\_to} automatically
drag the prop along.

The mechanism is \texttt{\_drag\_attached\_props(fig)}, called inside
the locomotion handlers (v0.9.7+). It iterates the global prop registry,
finds every prop whose \texttt{pam\_node{[}"parent"{]}} matches the
character's key, and applies the same per-frame translation that the
character is undergoing.

\textbf{Carry positions.} Held props snap to the character's hand joint
(\texttt{rhand} or \texttt{lhand}) by default. For carry styles that
need a different anchor (chest, hip, head), set the \texttt{attach}
explicitly when spawning:

\begin{Shaded}
\begin{Highlighting}[]
\FunctionTok{\{}\DataTypeTok{"action"}\FunctionTok{:} \StringTok{"spawn\_prop"}\FunctionTok{,} \DataTypeTok{"prop"}\FunctionTok{:} \StringTok{"bouquet"}\FunctionTok{,}
 \DataTypeTok{"type"}\FunctionTok{:} \StringTok{"floral\_arrangement"}\FunctionTok{,}
 \DataTypeTok{"parent"}\FunctionTok{:} \StringTok{"chava"}\FunctionTok{,} \DataTypeTok{"attach"}\FunctionTok{:} \StringTok{"rhand"}\FunctionTok{\}}
\end{Highlighting}
\end{Shaded}

\textbf{Known gotcha} --- \texttt{say}/\texttt{prop\_say} with
\texttt{persist} was rejected during walk-and-talk design because speech
bubbles anchor to world-space coordinates at the moment of creation, not
to the figure's head joint during movement. A walking character with a
persistent bubble leaves the bubble behind. For walk-and-talk dialogue,
use short non-persisted \texttt{say} calls per beat or design the
dialogue around standing/bench rest points.

\subsection{4.7 Prop-characters: DogGraph,
GovernorGraph}\label{prop-characters-doggraph-governorgraph}

Some entities are simultaneously \textbf{characters} (have figures, can
pose and speak) and \textbf{props} (can be spawned, removed, parented).
They're dual-registered in both the cast registry and the prop registry.
See also \texttt{CHARACTER\_PROPERTIES.md} §4.1 for the
dual-registration mechanics.

\textbf{DogGraph.} Built via \texttt{build\_dog(...)} in
\texttt{characters.py} but also exposed as a prop type for spawning
convenience. Has a \texttt{facing} parameter and supports
\texttt{trot\_to}, \texttt{prop\_say}, and the usual locomotion verbs.
Chekov (the robot dog) is the canonical example.

\begin{Shaded}
\begin{Highlighting}[]
\FunctionTok{\{}\DataTypeTok{"action"}\FunctionTok{:} \StringTok{"spawn\_prop"}\FunctionTok{,} \DataTypeTok{"prop"}\FunctionTok{:} \StringTok{"chekov"}\FunctionTok{,}
 \DataTypeTok{"type"}\FunctionTok{:} \StringTok{"dog"}\FunctionTok{,} \DataTypeTok{"x"}\FunctionTok{:} \FloatTok{{-}3.0}\FunctionTok{,} \DataTypeTok{"facing"}\FunctionTok{:} \StringTok{"right"}\FunctionTok{\}}
\FunctionTok{\{}\DataTypeTok{"action"}\FunctionTok{:} \StringTok{"trot\_to"}\FunctionTok{,} \DataTypeTok{"who"}\FunctionTok{:} \StringTok{"chekov"}\FunctionTok{,} \DataTypeTok{"x"}\FunctionTok{:} \FloatTok{2.0}\FunctionTok{\}}
\FunctionTok{\{}\DataTypeTok{"action"}\FunctionTok{:} \StringTok{"prop\_say"}\FunctionTok{,} \DataTypeTok{"prop"}\FunctionTok{:} \StringTok{"chekov"}\FunctionTok{,} \DataTypeTok{"text"}\FunctionTok{:} \StringTok{"Beep boop!"}\FunctionTok{\}}
\end{Highlighting}
\end{Shaded}

\texttt{trot\_to} uses \texttt{step{[}"x"{]}} (note: \texttt{x}, not
\texttt{to\_x}) for the destination --- a known inconsistency with
\texttt{walk\_to} flagged for normalization in v1.0.0.

\textbf{GovernorGraph.} A dodecahedral autonomous prop-char that pulses
gold/red and rotates on its own. Designed to be a background presence
(``the Governor of Venus'') that doesn't need explicit
\texttt{morph}/\texttt{pose} actions. Built via
\texttt{build\_governor(...)}.

\begin{Shaded}
\begin{Highlighting}[]
\FunctionTok{\{}\DataTypeTok{"action"}\FunctionTok{:} \StringTok{"spawn\_prop"}\FunctionTok{,} \DataTypeTok{"prop"}\FunctionTok{:} \StringTok{"governor"}\FunctionTok{,}
 \DataTypeTok{"type"}\FunctionTok{:} \StringTok{"governor"}\FunctionTok{,} \DataTypeTok{"x"}\FunctionTok{:} \DecValTok{0}\ErrorTok{.}\DecValTok{0}\FunctionTok{,} \DataTypeTok{"y"}\FunctionTok{:} \FloatTok{1.5}\FunctionTok{\}}
\FunctionTok{\{}\DataTypeTok{"action"}\FunctionTok{:} \StringTok{"prop\_say"}\FunctionTok{,} \DataTypeTok{"prop"}\FunctionTok{:} \StringTok{"governor"}\FunctionTok{,}
 \DataTypeTok{"text"}\FunctionTok{:} \StringTok{"PLEASE WAIT..."}\FunctionTok{,} \DataTypeTok{"persist"}\FunctionTok{:} \KeywordTok{true}\FunctionTok{\}}
\end{Highlighting}
\end{Shaded}

GovernorGraph's rotation runs continuously; it does not need a per-frame
action to keep moving.

\subsection{\texorpdfstring{4.8 The \texttt{tv\_monitor} /
\texttt{news\_anchor}
pattern}{4.8 The tv\_monitor / news\_anchor pattern}}\label{the-tv_monitor-news_anchor-pattern}

A stateful display prop pattern for when a prop's contents change over
time but its position does not.

\textbf{The pattern.} A \texttt{tv\_monitor} prop holds
\texttt{screen\_text} and \texttt{screen\_color}. To ``change
channels,'' the scene removes the prop and respawns it with new params
at the same position.

\begin{Shaded}
\begin{Highlighting}[]
\FunctionTok{\{}\DataTypeTok{"action"}\FunctionTok{:} \StringTok{"remove\_prop"}\FunctionTok{,} \DataTypeTok{"prop"}\FunctionTok{:} \StringTok{"monitor"}\FunctionTok{\}}\ErrorTok{,}
\FunctionTok{\{}\DataTypeTok{"action"}\FunctionTok{:} \StringTok{"spawn\_prop"}\FunctionTok{,} \DataTypeTok{"prop"}\FunctionTok{:} \StringTok{"monitor"}\FunctionTok{,}
 \DataTypeTok{"type"}\FunctionTok{:} \StringTok{"tv\_monitor"}\FunctionTok{,} \DataTypeTok{"x"}\FunctionTok{:} \FloatTok{4.0}\FunctionTok{,} \DataTypeTok{"y"}\FunctionTok{:} \FloatTok{1.8}\FunctionTok{,}
 \DataTypeTok{"screen\_text"}\FunctionTok{:} \StringTok{"BREAKING NEWS"}\FunctionTok{,} \DataTypeTok{"screen\_color"}\FunctionTok{:} \StringTok{"\#cc2222"}\FunctionTok{\}}
\end{Highlighting}
\end{Shaded}

\textbf{The footgun.} A \texttt{spawn\_prop} block that omits
\texttt{type} defaults to \texttt{"hat"} (see §5.4). When iterating on
tv\_monitor scenes, it's easy to copy a \texttt{spawn\_prop} block and
forget the \texttt{type}. The fix is to always include \texttt{type}
even on the re-spawn --- see the canonical example above.

A v1.0.0 candidate is a \texttt{set\_prop\_state} action that mutates
the prop in-place instead of remove+respawn, with \texttt{tv\_monitor}
as the first client. Not yet scoped.

\begin{center}\rule{0.5\linewidth}{0.5pt}\end{center}

\section{5. Meta and gotchas}\label{meta-and-gotchas}

\subsection{5.1 Z-order for props}\label{z-order-for-props}

Manim's compositing follows the order in which mobjects are added to the
scene. For props this typically gives:

\begin{enumerate}
\def\labelenumi{\arabic{enumi}.}
\tightlist
\item
  Backdrop / building props (lowest)
\item
  Floor-level furniture (desk, chair, bench, sofa)
\item
  Wall-mounted props (tv\_monitor, signs)
\item
  Characters
\item
  Hand-held props attached to characters
\item
  Mid-scene spawned props (often above characters intentionally)
\item
  Speech bubbles (highest, set by \texttt{bring\_to\_front})
\end{enumerate}

\textbf{Footgun.} A \texttt{focus\_reset} re-plays characters into the
scene at restored opacity, which can push them above props/bubbles added
earlier in the scene. See §5.2.

PAM does not currently expose explicit z-order overrides for props. If a
prop must sit above a character (e.g.~a foreground bouquet during a wide
shot), the practical solution today is to \texttt{remove\_prop} + later
\texttt{spawn\_prop} after the character is staged, so the prop is added
last.

\subsection{\texorpdfstring{5.2 Speech-bubble interactions:
\texttt{clear\_bubble},
\texttt{focus\_reset}}{5.2 Speech-bubble interactions: clear\_bubble, focus\_reset}}\label{speech-bubble-interactions-clear_bubble-focus_reset}

Persistent bubbles (\texttt{say} or \texttt{prop\_say} with
\texttt{persist:\ true} or \texttt{duration:\ t}) survive across actions
until explicitly cleared.

\textbf{The \texttt{focus\_reset} gotcha.} When a \texttt{focus\_reset}
runs after a persistent bubble is on screen, the reset restores
characters to full opacity and z-order, which can paint character
mobjects on top of the bubble --- visually overlaying it. Two
workarounds:

\begin{Shaded}
\begin{Highlighting}[]
\FunctionTok{\{}\DataTypeTok{"action"}\FunctionTok{:} \StringTok{"clear\_bubble"}\FunctionTok{,} \DataTypeTok{"who"}\FunctionTok{:} \StringTok{"thalia"}\FunctionTok{\}}\ErrorTok{,}
\FunctionTok{\{}\DataTypeTok{"action"}\FunctionTok{:} \StringTok{"focus\_reset"}\FunctionTok{\}}
\end{Highlighting}
\end{Shaded}

\ldots or:

\begin{Shaded}
\begin{Highlighting}[]
\FunctionTok{\{}\DataTypeTok{"action"}\FunctionTok{:} \StringTok{"say"}\FunctionTok{,} \DataTypeTok{"who"}\FunctionTok{:} \StringTok{"thalia"}\FunctionTok{,} \DataTypeTok{"text"}\FunctionTok{:} \StringTok{"..."}\FunctionTok{,}
 \DataTypeTok{"duration"}\FunctionTok{:} \FloatTok{2.5}\FunctionTok{\}}\ErrorTok{,}
\FunctionTok{\{}\DataTypeTok{"action"}\FunctionTok{:} \StringTok{"wait"}\FunctionTok{,} \DataTypeTok{"t"}\FunctionTok{:} \FloatTok{3.0}\FunctionTok{\}}\ErrorTok{,}
\FunctionTok{\{}\DataTypeTok{"action"}\FunctionTok{:} \StringTok{"focus\_reset"}\FunctionTok{\}}
\end{Highlighting}
\end{Shaded}

\ldots where \texttt{duration} is tuned to expire before the reset.

\textbf{Action reference:}

\begin{longtable}[]{@{}
  >{\raggedright\arraybackslash}p{(\columnwidth - 2\tabcolsep) * \real{0.5000}}
  >{\raggedright\arraybackslash}p{(\columnwidth - 2\tabcolsep) * \real{0.5000}}@{}}
\toprule\noalign{}
\begin{minipage}[b]{\linewidth}\raggedright
Action
\end{minipage} & \begin{minipage}[b]{\linewidth}\raggedright
Effect
\end{minipage} \\
\midrule\noalign{}
\endhead
\bottomrule\noalign{}
\endlastfoot
\texttt{clear\_bubble} & Dismiss one persistent bubble. Sub-key:
\texttt{who} or \texttt{prop}. Optional \texttt{rt}. \\
\texttt{clear\_prop\_bubble} & v0.9.9 alias for \texttt{clear\_bubble}.
Identical semantics. \\
\texttt{clear\_all\_bubbles} & Dismiss every persistent bubble. Optional
uniform \texttt{rt}. \\
\end{longtable}

\subsection{5.3 Bubble world-space anchoring during
walks}\label{bubble-world-space-anchoring-during-walks}

Speech bubbles attach to a fixed world-space position at the moment of
creation. They do \textbf{not} track the character or prop during
subsequent motion. Implications:

\begin{itemize}
\tightlist
\item
  \texttt{say} with \texttt{persist:\ true} followed by
  \texttt{walk\_to} leaves the bubble behind at the character's starting
  position.
\item
  \texttt{prop\_say} followed by reparenting / moving the prop leaves
  the bubble behind at the prop's starting position.
\end{itemize}

\textbf{Authoring pattern for walk-and-talk.} Break dialogue into short
non-persisted \texttt{say} beats per step of the walk, or design the
dialogue to occur during a standing rest point (treadmill-style:
walk-and-talk -\textgreater{} bench rest -\textgreater{} group-translate
reset).

\subsection{\texorpdfstring{5.4 The \texttt{tv\_monitor}
\texttt{type}-field
requirement}{5.4 The tv\_monitor type-field requirement}}\label{the-tv_monitor-type-field-requirement}

When \texttt{spawn\_prop} is used to swap a stateful prop (typically
\texttt{tv\_monitor}), the block must include \texttt{type}. Omitting it
triggers a fallback in \texttt{pam\_player.py} (around line 1732) that
defaults to \texttt{"hat"}, producing a silent visual bug --- the new
``channel'' appears as a small hat at the wall position.

\begin{Shaded}
\begin{Highlighting}[]
\ErrorTok{//} \ErrorTok{WRONG} \ErrorTok{—} \ErrorTok{defaults} \ErrorTok{to} \ErrorTok{hat}
\FunctionTok{\{}\DataTypeTok{"action"}\FunctionTok{:} \StringTok{"spawn\_prop"}\FunctionTok{,} \DataTypeTok{"prop"}\FunctionTok{:} \StringTok{"monitor"}\FunctionTok{,}
 \DataTypeTok{"x"}\FunctionTok{:} \FloatTok{4.0}\FunctionTok{,} \DataTypeTok{"y"}\FunctionTok{:} \FloatTok{1.8}\FunctionTok{,} \DataTypeTok{"screen\_text"}\FunctionTok{:} \StringTok{"BREAKING"}\FunctionTok{\}}

\ErrorTok{//} \ErrorTok{RIGHT}
\FunctionTok{\{}\DataTypeTok{"action"}\FunctionTok{:} \StringTok{"spawn\_prop"}\FunctionTok{,} \DataTypeTok{"prop"}\FunctionTok{:} \StringTok{"monitor"}\FunctionTok{,}
 \DataTypeTok{"type"}\FunctionTok{:} \StringTok{"tv\_monitor"}\FunctionTok{,}
 \DataTypeTok{"x"}\FunctionTok{:} \FloatTok{4.0}\FunctionTok{,} \DataTypeTok{"y"}\FunctionTok{:} \FloatTok{1.8}\FunctionTok{,} \DataTypeTok{"screen\_text"}\FunctionTok{:} \StringTok{"BREAKING"}\FunctionTok{\}}
\end{Highlighting}
\end{Shaded}

Also worth checking the \texttt{prop} key matches the \textbf{manifest
registry name}, not the type. If the manifest has
\texttt{"monitor":\ \{"type":\ "tv\_monitor",\ ...\}}, all later
references must say \texttt{"prop":\ "monitor"} --- not
\texttt{"prop":\ "tv\_monitor"}.

\subsection{5.5 Stateful prop method-name
overloading}\label{stateful-prop-method-name-overloading}

Because \texttt{open\_doors} / \texttt{close\_doors} are dispatched by
\texttt{getattr} on the prop VGroup (§4.3), any prop type that exposes
those methods inherits the action. This is mostly fine ---
\texttt{pocket\_door} and \texttt{elevator} work identically from the
screenplay's point of view --- but means an authoring error like
specifying a \texttt{pocket\_door} action on an \texttt{avatar\_pod}
(which has \texttt{open\_lid} instead) produces a silent no-op rather
than an error.

The unification candidate flagged for v1.0.0+ is a \texttt{StatefulProp}
mixin with a generic \texttt{set\_state(scene,\ name)} method. Each
stateful prop would declare its states (\texttt{open}, \texttt{closed},
\texttt{partial}) and the per-state visual config; one JSON action
\texttt{\{"action":\ "prop\_set\_state",\ "prop":\ "elev\_1",\ "state":\ "open"\}}
covers all four current stateful prop types. Not yet scoped; useful past
five or six stateful props.

\subsection{5.6 Back-burner}\label{back-burner}

Items relevant to props, drawn from \texttt{BACK\_BURNER.md} and the
cleanup-pass surface:

\begin{itemize}
\tightlist
\item
  \textbf{\texttt{bench} prop missing from
  \texttt{props\_furniture.py}.} Currently flagged for the walk-and-talk
  -\textgreater{} bench rest design. Build pattern follows
  \texttt{chair} with widened seat and no back.
\item
  \textbf{\texttt{props\_furniture.py} refactor.} Duplicated
  stateful-prop logic across \texttt{tv\_monitor}, \texttt{avatar\_pod},
  \texttt{pocket\_door}, and \texttt{elevator} is ripe for the
  \texttt{StatefulProp} mixin in §5.5.
\item
  \textbf{\texttt{to\_x} / \texttt{x} normalization.} \texttt{run\_to}
  uses \texttt{step{[}"x"{]}}, \texttt{walk\_to} uses
  \texttt{step{[}"to\_x"{]}}, \texttt{trot\_to} uses
  \texttt{step{[}"x"{]}}. Inconsistent across locomotion handlers that
  drive walk-following. Cross-referenced in
  \texttt{CHARACTER\_PROPERTIES.md} §4.6.
\item
  \textbf{Elevator door animation approach.} Implemented in an earlier
  scene; the specific approach (slide direction, timing) is worth
  documenting in a \texttt{pocket\_door} / \texttt{elevator} section of
  the pam\_player manual when that's redone.
\item
  \textbf{Carried-prop visibility in pan-up shots.} When a character
  carries flowers (or any chest-attached prop) and the camera pans up,
  the prop may be partially out of frame depending on \texttt{attach}
  choice. Worth experimenting with carry-point alternatives.
\item
  \textbf{Prop gallery.} Analogous to \texttt{character\_gallery.py}.
  Useful for visual regression and documentation. Not yet scoped.
\end{itemize}

\begin{center}\rule{0.5\linewidth}{0.5pt}\end{center}

\section{6. Prop type catalog}\label{prop-type-catalog}

Full registry of prop types as of v0.9.13. Aliases route to the same
builder. Asterisk (*) = stateful (carries animation methods).

\subsection{\texorpdfstring{6.1 Furniture
(\texttt{props\_furniture.py})}{6.1 Furniture (props\_furniture.py)}}\label{furniture-props_furniture.py}

\begin{longtable}[]{@{}
  >{\raggedright\arraybackslash}p{(\columnwidth - 6\tabcolsep) * \real{0.2500}}
  >{\raggedright\arraybackslash}p{(\columnwidth - 6\tabcolsep) * \real{0.2500}}
  >{\raggedright\arraybackslash}p{(\columnwidth - 6\tabcolsep) * \real{0.2500}}
  >{\raggedright\arraybackslash}p{(\columnwidth - 6\tabcolsep) * \real{0.2500}}@{}}
\toprule\noalign{}
\begin{minipage}[b]{\linewidth}\raggedright
Type
\end{minipage} & \begin{minipage}[b]{\linewidth}\raggedright
Aliases
\end{minipage} & \begin{minipage}[b]{\linewidth}\raggedright
Builder
\end{minipage} & \begin{minipage}[b]{\linewidth}\raggedright
Notes
\end{minipage} \\
\midrule\noalign{}
\endhead
\bottomrule\noalign{}
\endlastfoot
\texttt{chair} & --- & \texttt{build\_chair} & \texttt{surface},
\texttt{back-top}, \texttt{floor} attach points \\
\texttt{desk} & \texttt{table}, \texttt{console}, \texttt{computer},
\texttt{workstation}, \texttt{terminal} & \texttt{build\_desk} &
Optional \texttt{monitor=true} \\
\texttt{door} & --- & \texttt{build\_door} & Static panel \\
\texttt{pocket\_door}* & --- & \texttt{build\_pocket\_door} &
\texttt{open\_doors} / \texttt{close\_doors} \\
\texttt{elevator}* & --- & \texttt{build\_elevator} &
\texttt{open\_doors} / \texttt{close\_doors} / \texttt{partial\_close};
exposes \texttt{button\_panel} sub-prop \\
\texttt{desk\_lamp} & \texttt{lamp} & \texttt{build\_desk\_lamp} &
L-shaped, bug target \\
\texttt{bench} & --- & \emph{(missing --- back-burner)} & --- \\
\end{longtable}

\subsection{\texorpdfstring{6.2 Carried
(\texttt{props\_carried.py})}{6.2 Carried (props\_carried.py)}}\label{carried-props_carried.py}

\begin{longtable}[]{@{}
  >{\raggedright\arraybackslash}p{(\columnwidth - 6\tabcolsep) * \real{0.2500}}
  >{\raggedright\arraybackslash}p{(\columnwidth - 6\tabcolsep) * \real{0.2500}}
  >{\raggedright\arraybackslash}p{(\columnwidth - 6\tabcolsep) * \real{0.2500}}
  >{\raggedright\arraybackslash}p{(\columnwidth - 6\tabcolsep) * \real{0.2500}}@{}}
\toprule\noalign{}
\begin{minipage}[b]{\linewidth}\raggedright
Type
\end{minipage} & \begin{minipage}[b]{\linewidth}\raggedright
Aliases
\end{minipage} & \begin{minipage}[b]{\linewidth}\raggedright
Builder
\end{minipage} & \begin{minipage}[b]{\linewidth}\raggedright
Notes
\end{minipage} \\
\midrule\noalign{}
\endhead
\bottomrule\noalign{}
\endlastfoot
\texttt{hat} & --- & \texttt{build\_hat} & Default \texttt{spawn\_prop}
fallback type (gotcha) \\
\texttt{briefcase} & --- & \texttt{build\_briefcase} & Carried at
side \\
\texttt{folder} & --- & \texttt{build\_folder} & Flat, carried at
chest \\
\texttt{phone} & --- & \texttt{build\_phone} & \texttt{style}:
\texttt{cellphone} /\allowbreak{} \texttt{landline\_flat} /\allowbreak{} \texttt{landline\_wall}\\
\texttt{cell\_phone} & \texttt{smartphone} &
\texttt{build\_phone(style="cellphone")} & Aliases \\
\texttt{audio\_video\_bug} & \texttt{bug} &
\texttt{build\_audio\_video\_bug} & Tiny dot prop \\
\texttt{backpack}* & --- & \texttt{build\_backpack} (v0.9.7) &
Fano-plane graph; \texttt{reveal\_laptop(scene,\ laptop)} \\
\texttt{laptop} & --- & \texttt{build\_laptop} (v0.9.7) &
\texttt{closed} param (v0.9.13) \\
\end{longtable}

\subsection{\texorpdfstring{6.3 Flora
(\texttt{props\_flora.py})}{6.3 Flora (props\_flora.py)}}\label{flora-props_flora.py}

\begin{longtable}[]{@{}
  >{\raggedright\arraybackslash}p{(\columnwidth - 6\tabcolsep) * \real{0.2500}}
  >{\raggedright\arraybackslash}p{(\columnwidth - 6\tabcolsep) * \real{0.2500}}
  >{\raggedright\arraybackslash}p{(\columnwidth - 6\tabcolsep) * \real{0.2500}}
  >{\raggedright\arraybackslash}p{(\columnwidth - 6\tabcolsep) * \real{0.2500}}@{}}
\toprule\noalign{}
\begin{minipage}[b]{\linewidth}\raggedright
Type
\end{minipage} & \begin{minipage}[b]{\linewidth}\raggedright
Aliases
\end{minipage} & \begin{minipage}[b]{\linewidth}\raggedright
Builder
\end{minipage} & \begin{minipage}[b]{\linewidth}\raggedright
Notes
\end{minipage} \\
\midrule\noalign{}
\endhead
\bottomrule\noalign{}
\endlastfoot
\texttt{flower} & --- & \texttt{build\_flower} & Single stem \\
\texttt{floral\_arrangement} & \texttt{bouquet} &
\texttt{build\_floral\_arrangement} & \texttt{size}:
\texttt{carry}/\texttt{large} \\
\end{longtable}

\subsection{\texorpdfstring{6.4 Environment
(\texttt{props\_environment.py})}{6.4 Environment (props\_environment.py)}}\label{environment-props_environment.py}

\begin{longtable}[]{@{}
  >{\raggedright\arraybackslash}p{(\columnwidth - 6\tabcolsep) * \real{0.2500}}
  >{\raggedright\arraybackslash}p{(\columnwidth - 6\tabcolsep) * \real{0.2500}}
  >{\raggedright\arraybackslash}p{(\columnwidth - 6\tabcolsep) * \real{0.2500}}
  >{\raggedright\arraybackslash}p{(\columnwidth - 6\tabcolsep) * \real{0.2500}}@{}}
\toprule\noalign{}
\begin{minipage}[b]{\linewidth}\raggedright
Type
\end{minipage} & \begin{minipage}[b]{\linewidth}\raggedright
Aliases
\end{minipage} & \begin{minipage}[b]{\linewidth}\raggedright
Builder
\end{minipage} & \begin{minipage}[b]{\linewidth}\raggedright
Notes
\end{minipage} \\
\midrule\noalign{}
\endhead
\bottomrule\noalign{}
\endlastfoot
\texttt{dodecahedron} & --- & \texttt{build\_dodecahedron} &
\texttt{radius}, \texttt{accent}, \texttt{animate="spin"} \\
\texttt{building} & --- & \texttt{build\_building} &
\texttt{pam\_height} for pan-up camera moves \\
\texttt{sun} & --- & \texttt{build\_sun} & Background \\
\texttt{moon} & --- & \texttt{build\_moon} & Background \\
\texttt{solar\_panel} & --- & \texttt{build\_solar\_panel} &
Background \\
\texttt{wire} & --- & \texttt{build\_wire} & Background \\
\texttt{backdrop} & --- & \texttt{build\_backdrop} & Full-frame
background \\
\texttt{tv\_monitor}* & \texttt{monitor} & \texttt{build\_tv\_monitor} &
\texttt{screen\_text}, \texttt{screen\_color} \\
\texttt{avatar\_pod}* & \texttt{pod} & \texttt{build\_avatar\_pod} &
\texttt{open\_lid} / \texttt{close\_lid}; \texttt{occupied} palette \\
\texttt{governor} & --- & \texttt{build\_governor} (in characters.py ---
dual-registered) & Autonomous rotation, pulses gold/red \\
\end{longtable}

\subsection{\texorpdfstring{6.5 Accessories
(\texttt{props\_accessories.py})}{6.5 Accessories (props\_accessories.py)}}\label{accessories-props_accessories.py}

These were originally scoped as character accessories but per v0.9.5
design landed as CAPTION text rather than built props. Listed here for
completeness; the build functions exist but are not registered for
\texttt{spawn\_prop}.

\begin{longtable}[]{@{}
  >{\raggedright\arraybackslash}p{(\columnwidth - 2\tabcolsep) * \real{0.5000}}
  >{\raggedright\arraybackslash}p{(\columnwidth - 2\tabcolsep) * \real{0.5000}}@{}}
\toprule\noalign{}
\begin{minipage}[b]{\linewidth}\raggedright
Type
\end{minipage} & \begin{minipage}[b]{\linewidth}\raggedright
Notes
\end{minipage} \\
\midrule\noalign{}
\endhead
\bottomrule\noalign{}
\endlastfoot
\texttt{name\_tag} & Caption-only. E.g. \emph{``Chava name\_tag: Eve
Smith''} \\
\texttt{delivery\_cap} & Caption-only \\
\texttt{cheap\_suit} & Caption-only \\
\texttt{silver\_hair} & Builder exists (\texttt{build\_silver\_hair})
--- accessible via \texttt{spawn\_prop} for Mrs Bosch and similar.
Render-check on \texttt{narrow}/\texttt{female} figures recommended. \\
\end{longtable}

\begin{center}\rule{0.5\linewidth}{0.5pt}\end{center}

\emph{End of PROP\_PROPERTIES.md v0.9.13 draft. Companion to
\texttt{CHARACTER\_PROPERTIES.md}. Revision pass planned after the
\texttt{props\_furniture.py} stateful-prop refactor lands.}

\chapter{References and Source
Documents}\label{references-and-source-documents}

\section{Source Markdown files}\label{source-markdown-files}

The book was assembled from the following source Markdown files, each
folded into the single-file LaTeX output:

\begin{itemize}
\tightlist
\item
  \texttt{README-pam0914.md}
\item
  \texttt{pam\_player-manual\_v0914.md}
\item
  \texttt{CHARACTER\_PROPERTIES\_section\_1.md}
\item
  \texttt{CHARACTER\_PROPERTIES\_section\_2.md}
\item
  \texttt{CHARACTER\_PROPERTIES\_section\_3.md}
\item
  \texttt{CHARACTER\_PROPERTIES\_section\_4.md}
\item
  \texttt{CHARACTER\_PROPERTIES\_section\_4b.md}
\item
  \texttt{CHARACTER\_PROPERTIES\_section\_4c.md}
\item
  \texttt{face\_builder.py}
\item
  \texttt{actions.py}
\item
  \texttt{figure.py}
\item
  \texttt{pam\_player.py}
\item
  \texttt{PROP\_PROPERTIES-pam0914.md}
\end{itemize}

\section{External links}\label{external-links}

These links appeared in the source Markdown and are preserved as URL
references in the body text where relevant.

\begin{itemize}
\tightlist
\item
  Manim Community: \url{https://www.manim.community/}
\item
  screenplain: \url{https://pypi.org/project/screenplain/}
\end{itemize}

\backmatter
\end{document}
